% !TEX program = xelatex
\documentclass[12pt,a4paper,fontset=fandol]{ctexbook}

% ==================== 页面 ====================
\usepackage[top=2.5cm,bottom=2.5cm,left=2.5cm,right=2.5cm]{geometry}

% ==================== 字体 ====================
\usepackage{fontspec}
\usepackage{iffont}
% CJK 钉 fandol（两侧都取自 TeX 树里的同一份 FandolSong-Regular.otf）；
% mono 钉 DejaVu Sans Mono：CI 由 fontconfig 命中同名，本机没有该系统字体，
% 就用 MiKTeX dejavu 包里的同一版本 TTF（kpathsea 按文件名命中）。
\IfFontExistsTF{DejaVu Sans Mono}
  {\setmonofont{DejaVu Sans Mono}[Scale=MatchLowercase]}
  {\setmonofont[Extension=.ttf,
      UprightFont=DejaVuSansMono,
      BoldFont=DejaVuSansMono-Bold,
      ItalicFont=DejaVuSansMono-Oblique,
      BoldItalicFont=DejaVuSansMono-BoldOblique]
     {DejaVuSansMono}[Scale=MatchLowercase]}

% ==================== 颜色 ====================
\usepackage[dvipsnames,svgnames,x11names]{xcolor}
\definecolor{codebg}{HTML}{F5F5F5}
\definecolor{codeframe}{HTML}{E0E0E0}
\definecolor{cppkeyword}{HTML}{1A1AFF}
\definecolor{cppcomment}{HTML}{2E8B2E}
\definecolor{cppstring}{HTML}{B22222}
\definecolor{linenumgray}{HTML}{999999}
\definecolor{toolheadbg}{HTML}{2D3E50}
\definecolor{yamlkey}{HTML}{7A3E9D}
\definecolor{diffadd}{HTML}{1A7F37}
\definecolor{diffdel}{HTML}{B62324}

% ==================== listings ====================
\usepackage{listings}

% ---- 所有代码块共用的排版（先定义，各语言 style 组合它）----
\lstdefinestyle{codebase}{
  basicstyle=\small\ttfamily,
  keywordstyle=\color{cppkeyword}\bfseries,
  commentstyle=\color{cppcomment}\itshape,
  stringstyle=\color{cppstring},
  numberstyle=\tiny\color{linenumgray},
  numbers=left,
  frame=single,
  rulecolor=\color{codeframe},
  framesep=6pt,
  breaklines=true,
  tabsize=4,
  columns=flexible,
  keepspaces=true,
  backgroundcolor=\color{codebg},
  showstringspaces=false,
  escapeinside={(*@}{@*)},
}

\lstdefinestyle{cppstyle}{
  style=codebase,
  language=C++,
  morekeywords={constexpr,consteval,constinit,decltype,auto,
    concept,requires,co_await,co_yield,co_return,
    noexcept,override,final,nullptr,static_assert,
    template,typename,namespace,using,mutable,volatile,
    explicit,operator,friend,inline,virtual,
    thread_local,alignas,alignof},
  deletekeywords={int,char,bool,float,double,long,short,unsigned,signed,void,wchar_t},
  emph={size_t,int32_t,uint32_t,int64_t,uint64_t,ssize_t,
    string,vector,map,set,unique_ptr,shared_ptr,optional,variant,any,
    span,string_view,FILE,wchar_t},
  emphstyle=\color{teal!80!black},
}
\lstset{style=cppstyle}

% ---- C：UDF、原生插件与 C ABI 边界示例 ----
\lstdefinestyle{cstyle}{
  style=codebase,
  language=C,
  morekeywords={bool,FILE,size_t,restrict,inline},
  emph={size_t,int32_t,uint32_t,int64_t,uint64_t,real,ThreadData,
    cell_t,face_t,vectorField,scalarField},
  emphstyle=\color{teal!80!black},
}

% ---- YAML：求解器配置、引擎 .ini/.yaml、CI 描述文件 ----
\lstdefinestyle{yamlstyle}{
  style=codebase,
  morecomment=[l]{\#},
  morestring=[b]",
  morestring=[d]',
  alsoletter={-,:},
  morekeywords=[1]{Version,Checks,CheckOptions,HeaderFilterRegex,FormatStyle,
    User,WarningsAsErrors,CompilationDatabase,ExtraArgs,Add,Remove,
    Diagnostics,ClangTidy,UnusedIncludes,MissingIncludes,index,background-index,
    clang-tidy,InlayHints,Style,BasedOnStyle,IndentWidth,TabWidth,UseTab,
    ColumnLimit,BreakBeforeBraces,AlignAfterOpenBracket,AllowShortIfStatementsOnASingleLine,
    IncludeBlocks,SortIncludes,SeparateDefinitionBlocks,SpacesInAngles,
    Language,If,Else,Disable,Key,Name,files,exclude,repo,rev,hooks,id,entry,
    language_types,on,push_pull_request,steps,uses,runs-on,defaults},
  keywordstyle=[1]\color{yamlkey}\bfseries,
}

% ---- JSON：compile_commands.json / settings.json / CMakePresets.json ----
\lstdefinestyle{jsonstyle}{
  style=codebase,
  morestring=[b]",
  alsoletter={-},
  morekeywords={directory,file,output,arguments,command,version,
    cmakeMinimumRequired,name,generator,binaryDir,sourceDir,caches,
    toolchainFile,C_CXX_FLAGS,defineConstants,environments,
    editorMode,rulerColumn,tabSize,insertSpaces,files,settings},
  keywordstyle=\color{yamlkey}\bfseries,
}

% ---- Shell：命令行与工具调用 ----
\lstdefinelanguage{ShellLD}{
  morekeywords={if,then,else,elif,fi,for,in,do,done,while,case,esac,
    function,return,local,export,source,echo,cd,set,exec,trap,read,
    printf,test,exit,shift,eval,alias},
  sensitive=true,
  morecomment=[l]{\#},
  morestring=[b]",
  morestring=[d]',
}
\lstdefinestyle{shellstyle}{
  style=codebase,
  language=ShellLD,
  keywordstyle=\color{cppkeyword}\bfseries,
  emph={clang-format,clang-tidy,clangd,clang,clang++,cppcheck,iwyu,
    compile-commands,grep,sed,awk,xargs,jq,python3,pip,pip3,pwsh,bash,git,
    cmake,ninja,make,ccache,lldb,gdb,code,nvim,vim,emacs,tar,mkdir,touch,ls,cat,
    pre-commit,actions,gh,yamllint,wc,diff,find,sort,uniq,tee,rm},
  emphstyle=\color{teal!80!black}\bfseries,
}

% ---- CMake ----
\lstdefinelanguage{CMakeLD}{
  morekeywords={cmake_minimum_required,project,add_executable,
    add_library,target_link_libraries,target_include_directories,
    target_compile_definitions,target_compile_options,set,if,elseif,
    else,endif,find_package,option,message,install,file,
    target_sources,add_subdirectory,enable_language,set_target_properties,
    FetchContent_Declare,FetchContent_MakeAvailable,list,string},
  sensitive=false,
  morecomment=[l]{\#},
  morestring=[b]",
}
\lstdefinestyle{cmakestyle}{
  style=codebase,
  language=CMakeLD,
  keywordstyle=\color{cppkeyword}\bfseries,
}

% ---- Fortran：Abaqus/LS-DYNA 一类求解器的用户子程序 ----
\lstdefinestyle{fortranstyle}{
  style=codebase,
  language=[90]Fortran,
  morekeywords={subroutine,implicit,none,double,precision,intent,
    dimension,parameter,allocate,deallocate,use,module,endmodule},
  keywordstyle=\color{cppkeyword}\bfseries,
  sensitive=false,
}

% ---- Python：PyAnsys、引擎编辑器脚本与后处理 ----
\lstdefinestyle{pythonstyle}{
  style=codebase,
  language=Python,
  morekeywords={self,lambda,async,await,yield,with,as,pass,del,
    global,nonlocal,is,in,not,and,or,True,False,None},
  keywordstyle=\color{cppkeyword}\bfseries,
  emph={np,pl,anp,ansys,mechanical,fluent,pymechanical,plt,numpy,
    matplotlib,json,os,sys,pathlib},
  emphstyle=\color{teal!80!black},
}

% ---- 统一 diff：补丁式改动对比 ----
\lstdefinestyle{diffstyle}{
  style=codebase,
  numbers=none,
  frame=none,
  backgroundcolor=\color{white},
  morecomment=[l][\color{diffdel}]{-},
  morecomment=[l][\color{diffadd}]{+},
  morecomment=[l][\color{blue!70!black}\bfseries]{@@},
}

% ==================== tcolorbox ====================
\usepackage[most]{tcolorbox}

\newtcolorbox{guideline}[1][]{
  enhanced,breakable,
  colback=blue!5!white,colframe=blue!75!black,
  fonttitle=\bfseries,title=要点,#1,
  boxed title style={colback=blue!75!black},
  attach boxed title to top left={yshift=-2mm,xshift=2mm},
  arc=2mm,boxrule=0.8mm,
}
\newtcolorbox{compare}[1][]{
  enhanced,breakable,
  colback=green!3!white,colframe=green!50!black,
  fonttitle=\bfseries,title=对比,#1,
  boxed title style={colback=green!50!black},
  attach boxed title to top left={yshift=-2mm,xshift=2mm},
  arc=2mm,boxrule=0.8mm,
}
\newtcolorbox{warning}[1][]{
  enhanced,breakable,
  colback=red!5!white,colframe=red!75!black,
  fonttitle=\bfseries,title=常见陷阱,#1,
  boxed title style={colback=red!75!black},
  attach boxed title to top left={yshift=-2mm,xshift=2mm},
  arc=2mm,boxrule=0.8mm,
}
\newtcolorbox{note}[1][]{
  enhanced,breakable,
  colback=orange!5!white,colframe=orange!80!black,
  fonttitle=\bfseries,title=注意,#1,
  boxed title style={colback=orange!80!black},
  attach boxed title to top left={yshift=-2mm,xshift=2mm},
  arc=2mm,boxrule=0.8mm,
}
\newtcolorbox{bestpractice}[1][]{
  enhanced,breakable,
  colback=purple!5!white,colframe=purple!60!black,
  fonttitle=\bfseries,title=最佳实践,#1,
  boxed title style={colback=purple!60!black},
  attach boxed title to top left={yshift=-2mm,xshift=2mm},
  arc=2mm,boxrule=0.8mm,
}

% ==================== 章节标题 ====================
\usepackage{titlesec}
\usepackage{fancyhdr}
\usepackage{graphicx}
\usepackage{amssymb}
\usepackage{amsmath}
\usepackage{tabularx}
\usepackage{booktabs}
\usepackage{longtable}
\usepackage{array}
\usepackage{multirow}
\usepackage{enumitem}
\usepackage{seqsplit}
\usepackage{float}

\titleformat{\chapter}[display]
  {\normalfont\huge\bfseries\color{toolheadbg}}
  {\CTEXthechapter}{20pt}{\Huge}
\titleformat{\section}
  {\normalfont\Large\bfseries\color{blue!70!black}}
  {\thesection}{1em}{}
\titleformat{\subsection}
  {\normalfont\large\bfseries\color{blue!60!black}}
  {\thesubsection}{1em}{}

\pagestyle{fancy}
\fancyhf{}
\fancyhead[LE,RO]{\thepage}
\fancyhead[RE]{\leftmark}
\fancyhead[LO]{\rightmark}
\setlength{\headheight}{14.5pt}

\setlist[itemize]{leftmargin=2em,itemsep=2pt}
\setlist[enumerate]{leftmargin=2em,itemsep=2pt}
\emergencystretch=6em

% ==================== 超链接 ====================
\usepackage{hyperref}
\hypersetup{
  colorlinks=true,
  linkcolor=blue!70!black,
  urlcolor=blue!60!black,
  bookmarks=true,
  pdfauthor={C++ Reference},
  pdftitle={C/C++ 仿真与游戏引擎脚本参考手册：Ansys 等非实时求解器与 Unreal、Unity、Godot 的扩展层},
}

% ==================== 自定义命令 ====================
\newcommand{\cpp}[1]{C++#1}
\newcommand{\Fn}[1]{\texttt{#1()}}
\newcommand{\Tp}[1]{\texttt{#1}}
\newcommand{\Opt}[1]{\texttt{-\/-#1}}
\newcommand{\Bk}{\discretionary{}{}{}}
\newcommand{\Lib}[1]{\textbf{\texttt{#1}}}
\newcommand{\Lang}[1]{\textsf{#1}}
\newcommand{\CL}{\Lang{C}}
\newcommand{\CPP}{\Lang{C++}}
\newcommand{\PY}{\Lang{Python}}
% 域标签：本书讲"两侧对比"时统一用这两个写法，避免同义多形
\newcommand{\DomCAE}{\Lang{CAE}}
\newcommand{\DomEngine}{\Lang{Engine}}

% ====================================================================
% book 双页模式默认 \flushbottom，会靠拉伸行距把每页撑满；本仓库框体与长表多，
% 撑页产生的 Underfull \vbox 是排版噪声而非缺陷，这里统一改用小底部对齐。
\raggedbottom

\begin{document}

\frontmatter

\begin{titlepage}
\centering
\vspace*{4cm}
{\Huge\bfseries C/C++ 仿真与游戏引擎脚本参考手册\par}
\vspace{1cm}
{\LARGE 非实时求解器扩展与引擎脚本层\par}
\vspace{0.4cm}
{\Large Ansys、Abaqus、OpenFOAM 与 Unreal、Unity、Godot 的 C/C++ 接口\par}
\vfill
{\large \today\par}
\end{titlepage}

\tableofcontents

\mainmatter

\cleardoublepage

% ====================================================================
% 章节正文由 frag_chNN.tex 合并插入到下面这行之前（merge_frags.py）

% ============================================================
% ---- frag_ch01.tex ----
\part{非实时仿真软件：求解器与它的扩展层}

\chapter{分界与成因：为什么批处理仿真落在 C/C++}
\label{ch:ch01-boundary}

这一册讲两样东西：一是求解器，二是贴在求解器外面的那层脚本。两者用的语言、验收标准、失败时的表现都不一样，混在一起谈很容易讲糊。所以第一章先把分界画清楚，再解释为什么批处理仿真的内核最终都落在 \CL/{\Bk}\CPP{}/Fortran 上，最后把用户能够介入的三种形态摆出来——后面所有章节都是这三种形态的具体化。

\section{两种``仿真''的分界}
\label{sec:ch01-divide}

日常说的``仿真``至少盖着两类完全不同的工程活动。第一类是 \DomCAE{} 意义上的批处理仿真：结构有限元、流体求解器、电磁全波、热网络、多体动力学。它的产物是一份精度可交代的结果，一次算例从几十分钟到几天都算正常，规模以百万自由度计。第二类是 \DomEngine{} 意义上的实时仿真：游戏引擎、编辑器预览、数字孪生的可视化端。它的产物是每帧的画面或每帧的物理步进，帧预算是硬的（60 帧意味着每帧 16.7 毫秒，90 帧意味着 11.1 毫秒），精度可以让位于流畅度。

两者的差别不在``用的是不是同一个学科''，而在时间从哪里来、失败怎么定义。一张表说完：

\begin{table}[htbp]
\centering
\caption{两个域的同名不同义：批处理求解与实时步进}
\label{tab:ch01-divide}
\begin{tabular}{>{\raggedright\arraybackslash}p{0.14\textwidth}
                >{\raggedright\arraybackslash}p{0.34\textwidth}
                >{\raggedright\arraybackslash}p{0.32\textwidth}}
\toprule
\textbf{维度} & \textbf{\DomCAE{}：非实时批处理} & \textbf{\DomEngine{}：实时步进} \\
\midrule
目标 & 精度、收敛、可复现；结果要能签字 & 帧率、手感、可交互；画面要不断流 \\
时间约束 & 只有排队截止时间；单次求解可跑几小时 & 每帧固定预算，超时立刻可见 \\
误差处理 & 网格无关性验证、残差判据、时间步收敛研究 & 允许降级：减步数、换近似模型、跳细节 \\
失败定义 & 发散、不收敛、结果不可信，``慢''也是失败 & 掉帧、卡顿、输入延迟；``慢''是致命失败 \\
扩展层形态 & 宏语言、编译进内核的用户子程序、批处理脚本 & 脚本语言、反射元数据、插件模块、热重载 \\
\bottomrule
\end{tabular}
\end{table}

这张表决定了全书的写法：\DomCAE{} 侧的扩展层是``把用户意图翻译成求解器能执行的指令''，
而 \DomEngine{} 侧的扩展层是``把用户代码塞进一个已经在跑的循环''。前者可以慢，
因为它只在建模期执行；后者必须快，因为它每个渲染帧都要执行若干次。

\section{为什么求解器内核是 \CL/\CPP{}/Fortran}
\label{sec:ch01-kernel}

求解器不是``算得多''才需要原生语言，而是它的访存模式决定了必须用原生语言。

\textbf{其一，网格与稀疏矩阵的内存布局是设计出来的，不是默认的。} 一套工业级网格通常
把节点坐标存成 SoA（$\Bk$ 每分量一个连续数组），单元连接表按``单元 $\times$ 节点数''
拉平成一维 \Tp{int} 数组；刚度矩阵用 CSR 三数组（行指针、列索引、值）存储，因为非零元
位置在装配前就能由连接表统计出来。这些布局的共同点是：内层循环遍历的是连续地址，
迭代顺序固定，索引可以提前算。脚本语言的对象模型把每一次取元素都变成一次哈希查找加
类型检查，缓存局部性无从谈起。

\textbf{其二，热点集中在两处：单元装配与线性求解。} 装配是``取出单元刚度、scatter-add
进全局矩阵''；线性求解是稀疏矩阵-向量积加若干归约。两者的算术强度都不高，性能几乎全部
来自能否向量化、能否避免原子操作竞争、能否让行数据落在同一条缓存行上。

\begin{lstlisting}[language=C++, style=cppstyle,
  caption={装配热点的通用形状：CSR 定位加 scatter-add},
  label={lst:ch01-csr-assembly}]
// 三数组稀疏存储：非零位置由连接表统计得到，装配期只写值
struct Csr {
  int n;                             // 方程数 = 节点数 x 自由度
  std::vector<int> rowPtr, colIdx;   // rowPtr 长度 n+1
  std::vector<double> val, rhs;      // val 与非零元一一对应
};

// (i,j) 在行 i 内的槽位；一行只扫几个元素，热在 L1 里
int slotOf(const Csr& A, int i, int j) {
  for (int k = A.rowPtr[i]; k < A.rowPtr[i + 1]; ++k)
    if (A.colIdx[k] == j) return k;
  return -1;                         // 结构零：本例不应发生
}

void assemble(const Mesh& m, Csr& A) {
  std::fill(A.val.begin(), A.val.end(), 0.0);
  std::fill(A.rhs.begin(), A.rhs.end(), 0.0);
  for (int e = 0; e < m.nElem; ++e) {        // 外层：单元，天然可并行
    const double k = m.ea[e];                // 一维杆：EA/L
    const int i = m.conn[2 * e], j = m.conn[2 * e + 1];
    const int idx[2] = {i, j};
    const double ke[2][2] = {{k, -k}, {-k, k}};
    for (int a = 0; a < 2; ++a) {
      for (int b = 0; b < 2; ++b)            // 内层：4 次读改写
        A.val[slotOf(A, idx[a], idx[b])] += ke[a][b];
      A.rhs[idx[a]] += m.fe[2 * e + a];
    }
  }
}
\end{lstlisting}

清单~\ref{lst:ch01-csr-assembly} 的形状对一维杆、四面体、砖单元都成立，只是 \Tp{npe}
的规模和内层次数变化。它同时暴露了并行的两个真问题：外层``按单元''并行时同一个全局
槽位会被两个线程同时写；``按行着色''（graph coloring）把互不共享非零元的单元分到同一轮，
才能安全地无锁累加。这就是为什么求解器里到处能看到颜色数组。

\textbf{其三，SIMD 与并行的落地手段只有原生语言够得着。} 自动向量化要求内层循环无
控制流分叉、无指针别名（这也是 \Tp{restrict} 的意义）；线程级并行要求能精确控制临界区
粒度；内存分配要与求解器的分配器一致，否则百万次小分配会把堆拖垮。

\textbf{其四，内核常常不是自己写的，而是别人编译好的库。} 直接调 BLAS/LAPACK、
MUMPS/UMFPACK/PETSc 这一类库时，跨过的是语言边界，不是运行时边界。这里的坑是约定：
Fortran 的名字修饰（尾随下划线、部分平台加额外后缀）、列主序、LP64 与 ILP64 的索引宽度
不一致会让内核读越界。因此 ``\CL/{\Bk}\CPP{}/Fortran'' 三者共存不是历史包袱，而是 ABI
现实——线性代数层保留 \CPP{} 与 Fortran 实现的混合，外面包一层窄 \CL{} 接口。

这类边界值得单独列一遍，因为每一处约定都是编译期查不出来的：

\begin{table}[htbp]
\centering
\caption{内核常用的四类外部库及其跨语言边界约定}
\label{tab:ch01-abi}
\begin{tabular}{>{\raggedright\arraybackslash}p{0.20\textwidth}
                >{\raggedright\arraybackslash}p{0.24\textwidth}
                >{\raggedright\arraybackslash}p{0.34\textwidth}}
\toprule
\textbf{库族} & \textbf{承担什么} & \textbf{必须对齐的约定} \\
\midrule
BLAS、LAPACK & 稠密小矩阵、特征值、直接因子化 & 列主序、Fortran 名字修饰、整型宽度（LP64 与 ILP64 不可混链） \\
稀疏直接解 & 大规模稀疏因式分解、多右端回代 & 符号分析与数值分析分两阶段调用，索引基（0 或 1）要问清 \\
迭代法与预处理 & 共轭梯度、Krylov 子空间、多重网格 & 矩阵与向量的并行分布方式属于对象本身，不能跨分布套用 \\
MPI & 进程级并行、域分解后的通信 & 通信子与求解器内部的是同一份，自定义扩展里要谨慎插入集合调用 \\
\bottomrule
\end{tabular}
\end{table}

\section{用户能介入的三种形态}
\label{sec:ch01-three-forms}

用户代码进入求解器的路径只有三条，区别在``什么时候把代码变成机器指令''。

\textbf{形态一：脚本宏语言。} \Lib{APDL}、\Lib{Lua}、\Lib{Python}、各引擎的编辑器
脚本都属于这类。求解器里嵌一个解释器或维持一个命令通道，用户写的是``指令序列''而不是
``程序''。适用条件：任务由建模参数、工况枚举、结果提取这些粗粒度步骤组成，每一步的
计算都在求解器内部完成。代价：一次调用一次往返，循环若写在脚本层就会退化成成百上千次
往返；类型与单位全靠人守；没有编译期检查，错误往往到求解阶段才炸。

\textbf{形态二：把用户代码编译进内核。} 用户材料、用户单元、源项与边界剖面
（\DomCAE{} 侧的 UDF、UMAT 一类），以及求解器自带的 ``C++ 扩展''。用户提供符合固定
入口形状的函数，用指定工具链编成目标文件，由求解器在求解前链接或装入。适用条件：用户
逻辑会被调用很多次（每个积分点、每个面、每个时间步），或者需要访问内核的数据结构。
代价：工具链与求解器构建必须口径一致，ABI 不对就是崩溃；调试要跨进程边界；版本升级
常需重编。

\textbf{形态三：运行时插件。} \DomEngine{} 侧的主力形态：\CPP{} 动态库（DLL、SO、
DYLIB）加一份元数据描述。主程序在启动或热重载时按符号名取入口。适用条件：需要在不
重启进程的前提下换实现；需要把第三方库带进项目。代价：必须自己维护一个稳定 ABI；
加载顺序、符号可见性、静态初始化顺序都会咬人。

\begin{table}[htbp]
\centering
\caption{三种介入形态的适用条件与代价}
\label{tab:ch01-forms}
\begin{tabular}{>{\raggedright\arraybackslash}p{0.14\textwidth}
                >{\raggedright\arraybackslash}p{0.20\textwidth}
                >{\raggedright\arraybackslash}p{0.22\textwidth}
                >{\raggedright\arraybackslash}p{0.23\textwidth}}
\toprule
 & \textbf{脚本宏} & \textbf{编译进内核} & \textbf{运行时插件} \\
\midrule
绑定时机 & 运行期逐条解释 & 求解前静态链接或整份装入 & 首次用到时按符号解析 \\
能力上限 & 只能组合已有命令 & 与内核同速、可改本构 & 可替换子系统、可热重载 \\
粒度 & 命令、工况级 & 积分点、单元、面级 & 模块、系统级 \\
主要代价 & IPC 往返、无类型检查 & 工具链与 ABI 口径、重编 & ABI 稳定性、初始化顺序 \\
\bottomrule
\end{tabular}
\end{table}

形态三的最小可用形状，本质上是一份 ``\CL{} 结构体 + 版本号'' 的合同。它同时也是
\DomEngine{} 侧各引擎插件入口的抽象共同点，值得在这里先立起来：

\begin{lstlisting}[language=C, style=cstyle,
  caption={运行时插件的极简 ABI：入口符号只有一个，返回一份合同},
  label={lst:ch01-plugin-abi}]
/* 稳定的是这个结构体的布局与版本号，不是内部类型 */
#define SIM_PLUGIN_ABI 1

typedef struct SimPluginDesc {
  int abi;                    /* 必须等于 SIM_PLUGIN_ABI */
  const char* name;
  void* (*open)(const char* cfg);          /* 解析配置、分配句柄 */
  int   (*step)(void* h, double t, double dt); /* 每个时间步一次 */
  void  (*close)(void* h);
} SimPluginDesc;

#if defined(_WIN32)
#define SIM_EXPORT __declspec(dllexport)
#else
#define SIM_EXPORT __attribute__((visibility("default")))
#endif

SIM_EXPORT SimPluginDesc* sim_plugin_desc(void) {
  static SimPluginDesc d = {SIM_PLUGIN_ABI, "my_bc",
                            my_open, my_step, my_close};
  return &d;
}
\end{lstlisting}

\section{本册的阅读路线与统一口径}
\label{sec:ch01-roadmap}

全书按``同一件事在两侧各长什么样''来组织，来回对照而不是分两半各讲一遍。路线：
第 2 章起把 \DomCAE{} 侧的脚本层（\Lib{APDL}、参数化、批处理）讲清形状；随后进入原生
扩展层，即用户材料/单元与 \DomCAE{} 的 C/C++ 钩子；中段转到 \DomEngine{}，看脚本语言
如何与原生层交接、插件与反射元数据如何配合；末段回到共同话题——构建、性能测量、
可复现性、调试两侧混合代码。

统一口径有三条，全书各章都按这套说法写：

\begin{itemize}
\item \textbf{层名固定三个}：内核层（求解器或引擎本体）、脚本层（宏与解释器）、
  原生扩展层（编译进来的用户代码与插件）。不混用``插件''指脚本。
\item \textbf{域标签固定两个}：\DomCAE{} 指批处理仿真域，\DomEngine{} 指实时引擎域。
  跨域对比一律用这两个标签。
\item \textbf{事实纪律}：产品接口只写有把握的名称；不确定的参数个数与拼写一律改成
  描述调用形状，并注明以官方文档为准。版本口径写``以某一代发布为参照''，不锁小版本号。
\end{itemize}

跨域对照时最容易出的歧义是同一个词指两侧的不同东西，所以本册的层名固定按下表使用：

\begin{table}[htbp]
\centering
\caption{本册统一用语在两个域里的对应物}
\label{tab:ch01-terms}
\begin{tabular}{>{\raggedright\arraybackslash}p{0.16\textwidth}
                >{\raggedright\arraybackslash}p{0.40\textwidth}
                >{\raggedright\arraybackslash}p{0.24\textwidth}}
\toprule
\textbf{本册用语} & \textbf{\DomCAE{} 里的对应物} & \textbf{\DomEngine{} 里的对应物} \\
\midrule
内核层 & 求解器进程：装配、迭代、域分解 & 引擎运行时：渲染循环、物理步进 \\
脚本层 & 宏语言与宿主脚本通道 & 编辑器脚本、组件脚本 \\
原生扩展层 & 用户材料、单元、剖面与源项钩子 & \CPP{} 模块与插件动态库 \\
参数表 & 全局参数与表数组 & 序列化字段与反射元数据 \\
装入 & 求解前整份装入或解释执行 & 启动期或热重载期解析符号 \\
结果 & 结果集文件与后处理量 & 帧、统计叠加层、采样快照 \\
\bottomrule
\end{tabular}
\end{table}

\begin{note}[title={同一句``循环''在两侧的成本差三个数量级}]
\DomCAE{} 里在脚本层写 2000 次求解循环是可以接受的，因为每次循环体内部是几百万次
浮点操作；\DomEngine{} 里在脚本层写 2000 次每帧调用的循环一定会掉帧，因为每次循环体
只有几微秒的工作量，而一次脚本调用的固定开销就有微秒级。写代码前先问``这个循环外层
是求解循环还是渲染循环''。
\end{note}

\section{贯穿例子：一条最小的求解线}
\label{sec:ch01-thread-example}

后面所有章节都需要挂在同一条线上，否则例子会散。这条线是：网格 $\rightarrow$ 装配
$\rightarrow$ 求解 $\rightarrow$ 后处理。这里用一维杆的静力问题把它写完——四步齐全、
能直接编译运行，又小到能一眼看完。它不是玩具：脚本层的参数、扩展层的钩子、结果文件的
内容，都能从这四个环节里找到位置。

\begin{lstlisting}[language=C++, style=cppstyle,
  caption={贯穿例子：网格、装配、共轭梯度、后处理四步都在一个文件里},
  label={lst:ch01-mini1d}]
// mini1d.cpp：L=1m 的等截面杆，分布载荷与端载，左端固定
#include <algorithm>
#include <cmath>
#include <cstdio>
#include <vector>

int main() {
  const int n = 2001;                       // 节点数
  const double L = 1.0, E = 2.1e11, A = 1e-4;
  const double h = L / (n - 1), q = 1e6;    // q：单位长度上的力
  std::vector<double> x(n), conn(2 * (n - 1));
  for (int i = 0; i < n; ++i) {
    x[i] = i * h;
    if (i + 1 < n) { conn[2 * i] = i; conn[2 * i + 1] = i + 1; }
  }

  // 1) 装配：三对角，只存对角 d 与下次对角 l
  std::vector<double> d(n, 0.0), l(n - 1, 0.0), f(n, 0.0);
  const double k = E * A / h;
  for (int e = 0; e < n - 1; ++e) {
    l[e] = -k; d[e] += k; d[e + 1] += k;
    f[e] += q * h * 0.5; f[e + 1] += q * h * 0.5;   // 等效力
  }
  f[n - 1] += 2000.0;                        // 端载，N

  // 2) 边界条件：左端固定 -> 该行改成单位矩阵
  d[0] = 1.0; l[0] = 0.0; f[0] = 0.0;

  auto apply = [&](const std::vector<double>& p) {
    std::vector<double> y(n);
    for (int i = 0; i < n; ++i) {
      y[i] = d[i] * p[i];
      if (i > 0) y[i] += l[i - 1] * p[i - 1];
      if (i + 1 < n) y[i] += l[i] * p[i + 1];
    }
    return y;
  };

  // 3) 共轭梯度：矩阵-向量积是唯一的热点
  std::vector<double> u(n, 0.0), r = f, p = f, g(n);
  double rr = 0.0;
  for (int i = 0; i < n; ++i) rr += r[i] * r[i];
  const double r0 = std::sqrt(rr);
  int it = 0;
  for (; it < 2000 && rr > 1e-24 * r0 * r0; ++it) {
    g = apply(p);
    double pg = 0.0;
    for (int i = 0; i < n; ++i) pg += p[i] * g[i];
    const double alpha = rr / pg;            // 精确线搜索
    double nrr = 0.0;
    for (int i = 0; i < n; ++i) {
      u[i] += alpha * p[i];
      r[i] -= alpha * g[i];
      nrr += r[i] * r[i];
    }
    const double beta = nrr / rr;
    for (int i = 0; i < n; ++i) p[i] = r[i] + beta * p[i];
    rr = nrr;
  }

  // 4) 后处理：应变、应力，只算不存全量
  double epsMax = 0.0;
  for (int i = 0; i + 1 < n; ++i)
    epsMax = std::max(epsMax, std::abs((u[i + 1] - u[i]) / h));
  std::printf("iter=%d  u_tip=%.6e m  sigma_max=%.6e Pa\n",
              it, u[n - 1], E * epsMax);
  return 0;
}
\end{lstlisting}

清单~\ref{lst:ch01-mini1d} 的求解段是标准共轭梯度，但为了讲清形状写得刻意直白：
\Tp{apply} 每次都返回一个新向量（一次堆分配），残差与方向的更新分在三段循环里做，
也没有预处理。这三点恰好是性能章里第一批要动的手术：把临时缓冲提到循环外、把两次归约
融进一次遍历、再用一个不完全 Cholesky 或对角缩放把条件数压下来。改完之后迭代次数与
每次迭代的成本同时下降，这正是``热点在访存不在算术''的直观证据。

这条线上四处都可以被``外部意图''改写，对应关系如下：网格与参数进入第 1 段之前
（脚本层、批处理宏的参数入口）；单元刚度与等效力是原生扩展层的第一个钩子（用户材料、
用户载荷）；求解器选择与预处理是配置层的事；第 4 段的后处理是``取回哪些量、取回多少''
的策略问题，在真实产品里直接决定网络与磁盘成本。

\begin{bestpractice}[title={扩展代码要沿这条线定位，不要贴着功能命名}]
写任何一段扩展代码之前先回答三个问题：它改写四步中的哪一步、它的调用频次是每次求解
一次还是每个积分点一次、它失败时是``结果不可信''还是``进程崩''。三问能确定它该放脚本层
还是原生扩展层：调用频次是微秒级的必须编译进内核，是分钟级的留在脚本层更好维护；
一旦答案落在``结果不可信''，就不要用脚本层的静默错误传播去赌，改为在扩展层里做显式
校验并让它直接失败退出。
\end{bestpractice}

\textbf{小结}：批处理仿真的内核之所以是原生语言，因为它卖的是精度与规模而不是响应；
用户的介入之所以分层，因为``什么时候绑定''决定了能力上限与代价。下一章先把
\DomCAE{} 侧的脚本层与原生扩展层讲出形状，再逐点深入。

% ============================================================
% ---- frag_ch02.tex ----
\chapter{Ansys 生态的脚本层次：APDL、Workbench 与 C/C++ 可介入点}
\label{ch:ch02-ansys-layers}

这一族产品是 \DomCAE{} 侧脚本层最完整的样本：一层项目管理，一层宏语言，一层编译进去的用户代码。本章不列命令手册，只把三层的形状、边界与交接点讲清楚——哪些事只能在某一层做，越层去做要付什么代价。名称与参数以对应版本的官方命令手册为准；本章只写有把握的入口形状。

\section{三层结构}
\label{sec:ch02-three-layers}

\textbf{项目层}是 Workbench 这一类图形外壳：它管理工程文件、把几何、网格、求解、结果串成节点图，并负责把上游的改动传播给下游。它不解释载荷的物理含义，含义由各求解器解释。

\textbf{脚本层}由两半组成。一半是命令流：APDL 一类语言，把前处理、求解、后处理写成一个 ``命令 + 参数'' 序列。另一半是宿主脚本：把一段脚本挂到项目流水线的某个节点上，在图形的生命周期里插入用户动作。

\textbf{原生扩展层}是把用户代码编译进求解器：用户材料与本构、用户单元、边界剖面与源项钩子、自定义结果量。

\begin{table}[htbp]
\centering
\caption{``我要做这件事''应该落在哪一层}
\label{tab:ch02-layer-choice}
\begin{tabular}{>{\raggedright\arraybackslash}p{0.36\textwidth}
                >{\raggedright\arraybackslash}p{0.15\textwidth}
                >{\raggedright\arraybackslash}p{0.28\textwidth}}
\toprule
\textbf{要做的事} & \textbf{用哪层} & \textbf{为什么不是别层} \\
\midrule
同一模型跑 20 个参数组合 & 脚本层宏 & 项目层没有循环原语，扩展层不该管排队 \\
从 CAD 改名后重建几何并更新网格 & 项目层 & 只有它持有节点间的传播关系 \\
提取若干节点的时程曲线写成表格 & 脚本层 & 扩展层没有面向用户的文件接口 \\
非线性本构、与应力历史相关 & 原生扩展层 & 脚本层无法在每个积分点插入逻辑 \\
给入口面挂一个抛物线速度剖面 & 原生扩展层 & 每步、每个面都要执行一次 \\
加一套自己的单位换算与校验 & 哪层都不合适 & 这是外部预处理；落进求解器必然越界 \\
\bottomrule
\end{tabular}
\end{table}

\begin{guideline}[title={三层之间只有向下的交接是稳定的}]
项目层向下交给脚本层的是参数、文件路径与更新顺序；脚本层向下交给原生扩展层的是命令里引用的常数与装入的产物。反向只传结果文件的位置与少数标量。让扩展层回调脚本层、让脚本层去改项目层的节点图，都超出该层的合同：这类用法在文档里通常找不到，行为随版本变化，而且不受任何升级兼容性承诺保护。
\end{guideline}

\section{APDL 的参数化与批处理}
\label{sec:ch02-apdl}

APDL 的价值全在 ``参数化'' 三个字：所有命令的位置参数都可以是变量，于是建模、求解、后处理三段被拉成一个可重入的函数体。几个要点：

\begin{itemize}
\item \textbf{宏文件就是一次调用。} 一个文本文件里写命令流，外层用 ``读入并执行'' 一类命令把它当成一条语句调用；参数通过固定的形参名在宏体内引用。具体形参个数与命名规则以手册为准，形状是：定义时列出形参槽位，调用时按位置传值。
\item \textbf{``星号''命令族是控制流与 I/O。} 循环、分支、表数组、消息输出、文件读写这一类能力由它们提供，与建模、求解命令分工明确。表数组（TABLE 型 \Tp{DIM}）是参数化扫掠最常用的落地形式，因为它能被后处理直接插值。
\item \textbf{批处理是默认工作方式。} 不带图形界面时，输入是一份命令流文件，输出是一份文本日志加若干结果文件；``执行成功'' 与 ``求解成功'' 是两件事，日志必须自己扫。
\end{itemize}

参数传递的真实约束比语法更值得记。第一，宏内可见的变量属于求解器的全局参数表，所谓 ``局部'' 靠命名前缀维持而不是靠作用域，嵌套调用与递归共享同一张表。第二，数组与表也是全局对象，跨宏传值通常是 ``被调宏写出、主宏读回''，因此宏之间的接口只存在于命名约定里，签名不会被检查。第三，单位与量级不会随参数一起传，``这个数是 mm 还是 m'' 只能靠命名、注释与自检兜住。

表数组的类型选哪个也不是风格问题：ARRAY 是纯数组，只适合存中间量；TABLE 带插值与外推语义，扫掠结果与随时间变化的载荷曲线一般落在这里，因为求解器与后处理都能直接消费它；PARM 是标量参数，跨宏传单个量时比数组更容易对齐命名。

\begin{table}[htbp]
\centering
\caption{``星号''命令族按能力分组的形状（命令名以手册为准，语义长期稳定）}
\label{tab:ch02-starfamily}
\begin{tabular}{>{\raggedright\arraybackslash}p{0.14\textwidth}
                >{\raggedright\arraybackslash}p{0.33\textwidth}
                >{\raggedright\arraybackslash}p{0.33\textwidth}}
\toprule
\textbf{能力} & \textbf{形状} & \textbf{批处理里的角色} \\
\midrule
循环 & \Tp{*DO} 到 \Tp{*ENDDO}，计数变量是参数 & 扫掠的主干；每次迭代是一段完整求解 \\
分支 & \Tp{*IF} \Tp{...*THEN} 配 \Tp{*ELSEIF} 与 \Tp{*ENDIF}，比较符固定 & 收敛失败时改路径，而不是中止整批 \\
数组与表 & \Tp{*DIM} 声明 ARRAY、TABLE、PARM 三种 & TABLE 是扫掠结果的落点，能被插值 \\
参数读写 & \Tp{*SET} 赋值，\Tp{*STATUS} 列出当前参数表 & 宏与外层之间唯一的可观察面 \\
结果取标量 & \Tp{*GET}：参数名、实体类型、编号、标签、分量 & 把 GB 级结果压成一个数 \\
文本输出 & \Tp{*CFOPEN} 与 \Tp{*CFCLOSE} 夹文件，\Tp{*VWRITE} 配格式行 & 归档与 CI 断言的输入；格式行必须紧跟 \\
消息与控制 & \Tp{*MSG} 分级提示，读入外部命令流，\Tp{/EXIT} 收尾 & 日志可定位性的来源 \\
\bottomrule
\end{tabular}
\end{table}

\begin{lstlisting}[language=C, style=cstyle,
  caption={参数化扫掠宏：循环、表数组、后处理取量与结果落地},
  label={lst:ch02-sweep}]
! sweep.mac —— 一次调用跑完 ns 个厚度，结果写进表再落文件
! 本文件是 APDL 宏（不是 C）：! 之后为注释，逗号分隔实参
*SET, tmin, 8.0            ! 单位口径全程自洽：mm / N / MPa
*SET, tmax, 12.0
*SET, ns, 5
*DIM, out, TABLE, ns, 3    ! 表：厚度、UX 峰值、UY 峰值
/PREP7
ET, 1, 185                 ! 单元类型编号，含义见单元库
MP, EX, 1, 210000
MP, PRXY, 1, 0.3
! 几何、网格、约束与载荷的命令流由建模段给出，这里只留参数入口
*DO, i, 1, ns
  *SET, tk, tmin + (tmax - tmin) * (i - 1) / (ns - 1)
  ! --- 用 tk 重铺厚度方向网格，然后求解 ---
  /SOLU
  ANTYPE, STATIC
  SOLVE
  /POST1
  SET, LAST
  ! *GET 的形状：参数名, 实体类型, 编号, 标签, 分量
  *GET, ux, NODE, , UX, MAX
  *GET, uy, NODE, , UY, MAX
  *SET, out(i, 1), tk
  *SET, out(i, 2), ux
  *SET, out(i, 3), uy
  FINISH
*ENDDO
! 格式行必须紧跟 *VWRITE，FORTRAN 风格
*CFOPEN, sweep, csv
*VWRITE, out(1,1), out(1,2), out(1,3)
(F10.3,2E18.7)
*CFCLOSE
\end{lstlisting}

清单~\ref{lst:ch02-sweep} 里的 ``重铺厚度方向网格'' 一行是本章故意留空的：真实宏里它是十几条前处理命令，长度取决于建模方式，而不是控制结构本身。把控制结构写对、建模段从图形界面转录，是这类宏最常见的生产用法。分支命令族还有一处常被忽略的价值：把 ``本步不收敛就退回上一步、缩小步长重试'' 这类判断写成 \Tp{*IF} 加 \Tp{*MSG}，比让整批作业在第三次失败时崩掉要省事得多。

宏被当成函数调用时的接口形状如下。命令名、槽位个数与形参命名规则各代版本改动过，这里只给调用形状，落地前查手册确认：

\begin{lstlisting}[language=C, style=cstyle,
  caption={把宏当函数用：主文件绑定实参，被调宏用槽位名取值},
  label={lst:ch02-macro-call}]
! driver.inp —— 主文件：参数、表、然后把宏当一次调用
*SET, t0, 9.5              ! 厚度下界
*SET, dt, 0.5              ! 厚度步长
*DIM, res, TABLE, 4, 3     ! 由被调宏写出、主文件读回
MACRO, run_case, case.mac, t0, dt, res
run_case                   ! 一次调用跑完建模、求解、取量三段
! --- case.mac 的头部：实参在宏体内以固定槽位名出现 ---
! *SET, tk, PAR1           ! 取第 1 个实参（槽位名以手册为准）
! *SET, step, PAR2
! 建模与求解段同清单 (参数化扫掠宏)，结果写进 res
*STATUS                    ! 调用结束后列出参数表，核对 res 已被填好
\end{lstlisting}

``全局参数表加位置槽位'' 是双刃剑：好处是宏之间不需要任何声明就能传表，坏处是接口没有签名、拼错名字不会报错、两个宏用了同名参数就互相覆盖。工程上唯一可靠的缓解是给参数名加宏前缀，并在每个宏的末尾用 \Tp{*STATUS} 的输出做一次落盘核对。

批处理调用与 CI 里的判定：

\begin{lstlisting}[language=ShellLD, style=shellstyle,
  caption={CI 中一次批处理求解：日志与结果文件分别处理},
  label={lst:ch02-batch}]
# 可执行名带版本号（ansys<NNN> 一类的写法），以安装文档为准
BIN=ansys231
JOB=sweep
mkdir -p run
"$BIN" -b -np 4 -j "$JOB" -dir run \
  -i sweep.mac -o run/solver.log
# 1) 进程退出码不代表求解成功：必须扫日志
n=$(grep -c '\*\*\* ERROR' run/solver.log)
echo "solver errors: $n"
[ "$n" -eq 0 ] || exit 1
# 2) 数值健全性：宏自己写出的小文件，比拉回结果集便宜得多
awk -F, 'NR>1 && $2>1e-2 {bad++} END{exit bad?1:0}' run/sweep.csv
# 3) 只归档小产物；大结果集留在计算节点的暂存区
cp run/sweep.csv run/solver.log artifacts/
\end{lstlisting}

\begin{note}[title={命令行选项的口径要按版本重新确认}]
批处理模式、并行规模、工作目录、任务名这些选项在不同代版本里改过名，部分开关还随平台变化。清单~\ref{lst:ch02-batch} 写的是常见形状，落地前用安装目录里的帮助输出核对一次，并把核对结果写进 CI 的注释——否则换一台机器就不是同一条命令了。
\end{note}

\section{脚本驱动：Python 通道}
\label{sec:ch02-python}

Python 侧的封装（\Lib{PyMAPDL}、\Lib{PyMechanical} 这一类）看起来像 ``在 Python 里做仿真''，底下做的事其实很窄：启动或连接一个求解器进程，把方法调用还原成一条命令串发给它，把返回的文本或数组取回来。较新的版本走 gRPC 一类的进程间通道，更早的走别的中间件；具体传输与序列化方式以官方文档为准。

\begin{lstlisting}[language=Python, style=pythonstyle,
  caption={Python 通道：一次调用一次往返，取结果是整份数组拷贝},
  label={lst:ch02-python}]
# 名称与参数以 PyMAPDL 官方文档为准，这里只给出调用形状
from ansys.mapdl.core import launch_mapdl
import numpy as np

mapdl = launch_mapdl(ncore=2, cwd=r"C:\tmp\ch02")  # 起一个本地求解器进程
mapdl.run("/PREP7")            # 一次 run = 一次序列化 + 一次往返
mapdl.run("ET,1,185")
mapdl.run("MP,EX,1,210000")
# 建模、网格、边界段省略：与 APDL 命令流逐行同形
mapdl.run("/SOLU")
mapdl.run("ANTYPE,STATIC")
mapdl.run("SOLVE")
mapdl.run("/POST1")
mapdl.run("SET,LAST")

# 标量：把 *GET 的返回值取回 Python
ux_max = mapdl.get("UX", "NODE", "", "MAX")        # 形状以手册为准
# 全场数组：一次往返换回 n_node 个浮点，走的是数据通道
disp = mapdl.post_processing.nodal_displacement("X")

# 反面写法：逐节点循环 = n_node 次往返，比上一条慢三个量级
# for nid in range(1, disp.size + 1):
#     mapdl.run("NODE,%d" % nid)                   # 千万不要这样
print(ux_max, disp.max(), np.__version__)
mapdl.exit()
\end{lstlisting}

它快不起来的原因是结构性的，三条：

\begin{enumerate}
\item \textbf{IPC 往返是固定成本。} 每个 ``送一条命令'' 的调用要付序列化、进程切换、求解器主循环响应的代价，与命令本身的计算量无关。循环放在 Python 层就意味着把这笔固定成本乘上循环次数。
\item \textbf{结果是拷贝出来的。} 求解器把结果写在自己的二进制结果集里；Python 侧要用就得读文件、解码、再复制成宿主语言的数组。节点数上百万时，一次全场取回就是几十到几百 MB。
\item \textbf{两侧状态要手工同步。} Python 对象不知道求解器里的模型是否已被改写，缓存下来的数组随时可能是旧值；会话断开重连后，参数表与工作目录的可见性也可能变。
\end{enumerate}

由此得到一条实用的分工线：控制流（扫掠哪些参数、跑哪些工况、失败怎么重试）留在 Python；数据归约（取极值、按区域求和、抽若干条曲线）要么写成宏让求解器自己算完只回传标量，要么整份读入后在 \PY{} 侧用数组运算一次算完。中间态——逐节点、逐步取数——几乎总是错的。

另一条通道长在项目层而不是命令层：宿主脚本对象拿到的是项目与模型的句柄，写的是 ``遍历树、改属性、触发更新'' 这一类动作，而不是逐条求解命令。它的成本结构与 APDL 宏不同，但结论同一条：一次属性写入或一次更新触发同样是一次跨层往返，把循环留在脚本侧就要乘上次数。项目层的对象名与方法表在各代版本里改动频繁，本章只给调用形状，具体名称以官方脚本指南为准。

\begin{compare}[title={同样一个循环，放在两个域的脚本层里命运不同}]
\DomCAE{} 侧的循环写在宏或宿主脚本里：一次迭代触发的是几百万次浮点操作，往返成本可以忽略，2000 次求解的调度用 \Tp{*DO} 还是 \PY{} 的 \Tp{for} 都一样顺。
\DomEngine{} 侧的循环写在宿主脚本里：每帧每个对象回调一次，往返成本与循环体本身的工作量同量级，于是同样的逻辑要么挪进原生层，要么改用脚本层提供的批量化接口。
两边共同的判据只有一条：数一数循环体里除 ``发起一次调用'' 之外还剩多少计算。剩得少就不该留在这层。
\end{compare}

\section{C/C++ 可介入点清单}
\label{sec:ch02-cpp-hooks}

原生扩展层的共同形状：求解器在一个确定的时机调用你提供的一个确定入口。要谈的从来不是 ``API 有什么函数''，而是四个问题：入口在哪个阶段被调用、需要什么工具链口径、并行时调用多少次、失败时会怎样。

\textbf{入口形状。} 除少数带完整头文件与宏体系的接口外，本族大多数用户例程的入口是一个 C 函数：第一个实参是求解器给的线程或上下文数据指针，其后是单元或面索引、进出变量（应变、应力、状态变量数组）、时间步与材料常数。用户材料与本构类入口的进出通常是 ``本步应变增量进、应力量与雅可比出''；具体名称、实参顺序与数组约定以官方扩展手册为准，不要凭其它产品的记忆来写。

\textbf{编译期口径。} 这一层没有 ABI 兼容余地：编译器的选择、目标平台位宽、运行时库的链接方式（动态或静态）、调用约定，都必须与求解器构建一致。求解器会指定可接受的编译器族与版本范围，越界的表现通常是装入阶段报错，或者更糟——运行中越界写坏内存。

\textbf{装入时机两种。} 一是求解前整份装入：求解器在开始迭代前解析符号并驻留，失败立刻可见，运行期成本最低；二是每次调用解释或即时编译的 ``快改'' 模式，省去重编但也丢掉优化与部分类型检查。两种时机决定了调试手段完全不同：前者能挂原生调试器，后者只能在钩子里打印。

\textbf{并行调用形状。} 分布式并行下每个进程持有一本子域，入口函数在每个进程的每个单元（或每个面、每个积分点）上被调用；共享内存并行下同一个函数会被多个线程同时调用。由此推出三条硬规矩：不得有可写的静态或全局状态；不得在入口里申请大块的进程级资源；不得假定调用顺序。

\begin{lstlisting}[language=C, style=cstyle,
  caption={钩子入口的两个典型形状：面剖面与体源项（Fluent UDF 风格）},
  label={lst:ch02-udf-shape}]
/* 这一段是 UDF 风格示例：入口宏的实参表由求解器头文件约定 */
#include "udf.h"

/* 形状一：边界剖面。入口按"线程指针 + 边界编号"给出，
   内部遍历属于本进程的那部分面，写入待求未知量 */
DEFINE_PROFILE(inlet_axial_velocity, thread, index)
{
    real x[ND_ND];                 /* 质心坐标，维度由求解器决定 */
    real r, r0 = 0.01;
    face_t f;

    begin_f_loop(f, thread)        /* 并行下只覆盖本地面集合 */
    {
        F_CENTROID(x, f, thread);
        r = SQRT(x[1] * x[1] + x[2] * x[2]);
        F_U(f, thread) = 20.0 * (1.0 - (r / r0) * (r / r0));
    }
    end_f_loop(f, thread)
}

/* 形状二：体源项。入口按"单元 + 线程 + 梯度槽 + 方程号"给出，
   返回源项值，同时把对方程未知量的解析偏导写进梯度槽 */
DEFINE_SOURCE(porosity_drag, c, thread, dS, eqn)
{
    const real K = 1.0e5;          /* 示意用的多孔阻力系数 */
    dS[eqn] = -K;                  /* 解析导数利于线性化收敛 */
    return -K * C_U(c, thread);
}
\end{lstlisting}

清单~\ref{lst:ch02-udf-shape} 值得对照着看的是两条入口的差别：剖面钩子返回的是赋值，源项钩子既要返回值又要填导数槽。这种 ``多通道出参'' 是本层所有入口的共同特征，也是它不能被脚本层替掉的原因——脚本层没有表达偏导数的地方，而非线性求解的收敛质量恰恰取决于这个偏导。

\begin{table}[htbp]
\centering
\caption{C/C++ 介入点的调用时机、装入阶段与代价}
\label{tab:ch02-hooks}
\begin{tabular}{>{\raggedright\arraybackslash}p{0.20\textwidth}
                >{\raggedright\arraybackslash}p{0.17\textwidth}
                >{\raggedright\arraybackslash}p{0.20\textwidth}
                >{\raggedright\arraybackslash}p{0.21\textwidth}}
\toprule
\textbf{介入点} & \textbf{调用时机} & \textbf{装入阶段} & \textbf{代价} \\
\midrule
用户材料、本构 & 每个积分点、每次迭代 & 求解前整份装入 & 收敛性由你负责，雅可比必须给对 \\
用户单元、自定义自由度 & 装配阶段与结果阶段 & 求解前装入并注册 & 要自己保证一致性与数值稳定 \\
自定义载荷、边界剖面 & 每个时间步或每个子步 & 装入或解释执行 & 频率高，禁全局状态与文件 I/O \\
自定义结果量、后处理函数 & 结果写盘或读回时 & 后处理阶段 & 影响结果集大小与取回成本 \\
几何与网格导入适配 & 求解器启动之前的转换阶段 & 项目层触发 & 产品不一定开放该接口 \\
求解流程与迭代控制 & 每个子步的迭代循环 & 多为配置而非代码 & 越权改动易破坏收敛判据 \\
\bottomrule
\end{tabular}
\end{table}

多数产品给了自己的编译入口（面板按钮、生成的 makefile、随安装提供的模板脚本）。不必依赖那些生成物也能把这层讲清楚：需要的是一个 MODULE 目标，一套与求解器构建一致的编译口径，外加一个 ABI 版本宏。下面这段是等价形状，头文件目录与运行时库取值必须换成你所用版本实际要求的那一套。

\begin{lstlisting}[language=CMakeLD, style=cmakestyle,
  caption={把用户例程编成可装入模块：口径与求解器构建对齐},
  label={lst:ch02-build}]
cmake_minimum_required(VERSION 3.24)
project(userext C)
add_library(usermat MODULE usermat.c)     # MODULE：主程序按符号取入口
target_compile_definitions(usermat PRIVATE USEREXT_ABI=1)
set_target_properties(usermat PROPERTIES
    PREFIX ""                              # 文件名由求解器的约定决定
    C_STANDARD 11)
if(WIN32)
  # 运行时库必须与求解器构建同款，否则堆跨模块边界
  set(CMAKE_MSVC_RUNTIME_LIBRARY "MultiThreaded$<$<CONFIG:Debug>:Debug>DLL")
else()
  target_compile_options(usermat PRIVATE -fPIC -O2)
  target_link_options(usermat PRIVATE "-shared")
endif()
# 头文件与库目录来自求解器安装：路径变量名以官方扩展文档为准
target_include_directories(usermat PRIVATE "$ENV{SOLVER_INCLUDE}")
\end{lstlisting}

\section{工程约束}
\label{sec:ch02-constraints}

\textbf{单位制。} 求解器内核里没有单位这个概念，它只处理数字：输入是量值，输出也是量值。所谓 ``求解器的单位制'' 其实是某一层为图形显示或导入导出做的换算，一旦换算散落在脚本层、扩展层与材料库里，就出现 mm 配 Pa 而结果差六个数量级这类事故。可行做法是：整个工程固定一套自洽单位并在命令流文件头写明，任何跨层传值（尤其是扩展层的材料常数与钩子返回值）都当作带单位的量来审计，只在最外层做显示换算。

\textbf{许可证与版本。} 三处必须同代：许可证服务端能力、求解器可执行文件、扩展层编译产物。CI 镜像里求解器版本与集群节点不一致时，常见现象是作业排到队首却被拒；而编译产物跨代复用更隐蔽——符号能解析、行为变了。做法是把版本与编译器写成产物元数据一起归档，作业脚本显式声明所需版本，许可证占用量（并行规模直接影响核数）在排产时一并计入。

\textbf{大规模结果的取回。} 完整结果集（每个时间步、全场分量、若干单元量）很容易到几十 GB，从计算节点搬运的成本会盖过计算本身。三条策略按优先级：

\begin{itemize}
\item 归约推回求解侧：宏或后处理函数算完只写标量与小表，本地拿到的是 KB 级文件；清单~\ref{lst:ch02-sweep} 末尾那张 CSV 就是这个策略的最小形式。
\item 必须取整份时取子集：指定时间步区间、指定节点集合、指定分量，并分块读；一次 ``所有步所有量'' 的读法在交互会话里几乎从不正确。
\item 远端解算、本地只看小产物：取回路径写进日志与产物目录，而不是靠终端会话的历史，否则复现时没人知道那批数字是从哪一次运行来的。
\end{itemize}

\begin{table}[htbp]
\centering
\caption{四类工程约束的症状与对策}
\label{tab:ch02-constraints}
\begin{tabular}{>{\raggedright\arraybackslash}p{0.14\textwidth}
                >{\raggedright\arraybackslash}p{0.33\textwidth}
                >{\raggedright\arraybackslash}p{0.33\textwidth}}
\toprule
\textbf{约束} & \textbf{典型症状} & \textbf{对策} \\
\midrule
单位制 & 结果量级差六个数量级；应力像压强又像应力；日志一切正常 & 工程级固定一套自洽单位写在文件头；跨层传值做量级断言；只在显示层换算 \\
许可证 & 作业排到队首被拒；交互会话占着核数不放；批处理排队超时 & 批处理与交互分队列；显式声明所需核数；脚本末尾保证释放 \\
版本口径 & 同一条命令流在两台机器上一个收敛一个不收敛；扩展产物装入失败 & 记录求解器版本与编译器到产物元数据；作业脚本声明版本；禁止跨代复用二进制 \\
结果规模 & 一次扫掠产出几百 GB；本地拉回结果比求解还慢；磁盘先满 & 归约推到求解侧；只取需要的步与分量并分块读；大结果留在暂存区 \\
\bottomrule
\end{tabular}
\end{table}

\begin{bestpractice}[title={宏与扩展层之间留一份 ``自检契约''}]
在命令流文件头把单位口径、求解器版本、所需环境变量与产物名写成注释块，并在宏的第一个 \Tp{*IF} 里对关键参数做量级断言（厚度落在哪几个数量级、载荷是否为正）。原生扩展层同理：入口函数第一次被调用时检查传入的单元体积、时间步长与材料常数是否落在物理合理区间，越界就显式报错并中止，而不是继续返回一个数量级错了的应力。两层的成本都是几行代码，换来的是 ``错误在产生它的那一层暴露''。
\end{bestpractice}

\begin{warning}[title={把脚本层当求解器层用：最常见的三类翻车}]
第一类是粒度错配：在宏或 \PY{} 里对节点逐个循环、每个节点发一条命令，把微秒级的一次往返放大成几万次。归约要么交给求解器，要么整份取回后用数组运算处理。第二类是单位错配：脚本层按 mm 建模，材料常数从手册里抄来是 m 口径，而求解器不会报任何错——它只是算完。单位约定必须写死在文件头并由扩展层自检（量值不落在物理合理区间就断言失败）。第三类是把扩展层的责任拉回脚本层：非线性本构、依赖历史的量、需要解析导数的源项，在脚本层用 ``分段近似加循环调用'' 硬做，结果既不可信也不可收敛。判据很简单：只要你的逻辑需要每个积分点执行一次，它就不属于脚本层。
\end{warning}

% ============================================================
% ---- frag_ch03.tex ----
\chapter{Fluent UDF：C 宏库、DEFINE\_* 族与两种加载方式}
\label{ch:ch03-fluent-udf}

本章讲 \Lib{Fluent} 的用户定义函数（\Tp{UDF}）：用户把 \CL{} 源码交给求解器，
求解器把它编进自己的进程里，在每个迭代的固定时刻回调。这是 \DomCAE{} 侧最典型的
``把用户代码搬进求解器内部''形态，也是本册后面讲引擎侧扩展时的对照组。

\begin{guideline}[title={读本章前先记住三件事}]
\begin{itemize}
  \item UDF 的语言是 \CL{}，不是脚本：它跑在求解器进程里，链接的是求解器自己的
        运行时，因此一次越界写坏的是整个算例，而不是一条被丢弃的消息。
  \item UDF 的入口不是函数签名而是宏：\Tp{DEFINE\_*} 展开成求解器认得的符号名，
        宏名决定回调时机，宏体决定算法。
  \item 加载方式决定性能与可用功能：解释加载免编译但慢，编译加载要配套编译器
        但能进热点循环。
\end{itemize}
具体宏名、参数个数与可用功能集合都按版本演进，本章只给稳定口径与调用形状，
落笔前请对照所使用版本的官方 UDF 手册。
\end{guideline}

\section{UDF 在求解循环里的位置}
\label{sec:ch03-udf-loop}

一个典型的分离式求解器迭代是这样的：读网格、初始化场量、然后反复执行
``边界与源项赋值 $\rightarrow$ 组装并求解输运方程 $\rightarrow$ 更新梯度与物性
$\rightarrow$ 判敛''。UDF 不是插在这个循环外面的调用者，而是被这个循环回调的一方：

\begin{itemize}
  \item \textbf{初始化时刻}：求解器做一次全域单元/面遍历的钩子，用来把初值、
        历史变量、区域标记写进求解器自己的存储。
  \item \textbf{每迭代一次}：全局量汇总、监测值计算、按总体平均量修正边界，
        这一类钩子每次迭代被调用一次，实参是求解域句柄。
  \item \textbf{遍历时刻}：边界剖面与源项这类钩子在求解器组装方程时被逐点调用，
        调用次数等于该边界上的面数或域内的单元数，一次调用只处理一个面或一个单元。
  \item \textbf{迭代之后}：执行钩子（\Tp{DEFINE\_EXECUTE\_AT\_END} 这一族，
        在迭代结束处运行），适合做后处理统计、
        写文件、更新外部状态。
\end{itemize}

由``逐点调用''这一形状可以直接推出三条写法纪律：

\begin{enumerate}
  \item \textbf{无状态}：钩子之间不能靠函数里的 \Tp{static} 变量传数据。
        并行时每个 rank 有自己的进程副本，\Tp{static} 只在本 rank 可见，
        且求解器可能重排、重复或跳过调用（重启、迭代数改小、判敛提前）。
        需要跨迭代保存的东西，写进求解器的存储：历史变量、用户定义的
        材料参数槽位、或者边界上的自定义场变量。
  \item \textbf{可重入}：同一个钩子会被不同线程并发调用（求解器的
        并行计算线程或 OpenMP 分块），因此钩子里不许有共享可变态、
        不许有阻塞同步、不许依赖调用顺序。
  \item \textbf{局部性}：一次调用只允许碰当前单元/面及其邻居的数据。
        需要邻居数据时应使用求解器提供的插值或梯度接口，而不是自己按
        索引猜内存布局。
\end{enumerate}

\section{DEFINE\_* 宏族的分工}
\label{sec:ch03-define-family}

\Tp{DEFINE\_PROFILE}、\Tp{DEFINE\_SOURCE}、\Tp{DEFINE\_ADJUST} 这一组宏的共同
结构是：第一个实参是用户起的函数名（宏会把它变成求解器查找的符号），
后面的实参由宏名决定。最常见的三种形状是

\Tp{DEFINE\_PROFILE(名字, thread, position)} —— \Tp{thread} 是边界线程指针，
\Tp{position} 是求解器传进来的槽位号，写回时原样传给
\Tp{F\_PROFILE(f, thread, position)}；

\Tp{DEFINE\_SOURCE(名字, c, t, dS, eqn)} —— \Tp{c} 是单元索引，\Tp{t} 是单元
线程，\Tp{dS} 是源项对因变量的导数数组（求解器用它进雅可比），
\Tp{eqn} 是当前方程号；返回值就是这个单元的源项数值；

\Tp{DEFINE\_ADJUST(名字, domain)} —— 每迭代一次，实参是求解域句柄，
在里面自己发起遍历。

物性类钩子（密度、黏度、比热一类）用的是``单元进、标量出''的形状：
一次调用算一个单元的物性值，返回值即物性，没有输出实参。
初始化与执行钩子则是``域句柄进、无返回值''的形状，遍历由用户自己写。
再往边界方向还有壁面温度、入口湍流量、辐射与组分等专用钩子，
它们的分工同样是``指定被改的量 + 指定遍历对象''，
具体宏名以官方文档为准，不要把本章的写法当成完整清单。

遍历对象由两套循环宏承担，形状对称：

\begin{itemize}
  \item 面遍历：\Tp{begin\_f\_loop(f, thread)} \ldots \Tp{end\_f\_loop(f, thread)}，
        用于边界剖面类钩子。
  \item 单元遍历：\Tp{begin\_c\_loop(c, thread)} \ldots \Tp{end\_c\_loop(c, thread)}，
        用于源项与物性类钩子（这两个宏在新版本里还分裂出若干带不同
        语义的变体，用途差别在``是否包含鬼单元''，选哪个以官方说明为准）。
  \item 域级遍历：先按线程迭代（把域内所有同类线程走一遍），再在线程内
        做单元或面遍历。求解器提供的域级迭代宏名本章只描述形状，
        落笔时请查手册。
\end{itemize}

数据访问器按前缀分族，这是最容易记的口径：\Tp{F\_*} 取面上的量
（速度 \Fn{F\_U}、压力 \Fn{F\_P}、面积向量 \Fn{F\_AREA}、面心 \Fn{F\_CENTROID}），
\Tp{C\_*} 取单元里的量（\Fn{C\_U}、\Fn{C\_T}、体积 \Fn{C\_V}、质心
\Fn{C\_CENTROID}），\Tp{N\_*} 取节点量，\Tp{THREAD\_*} 取线程元数据。
大部分读访问器同时可写，但``能写''不等于``该写''：直接改求解器存储里的
单元值会绕过它的一致性更新，见 \S\ref{sec:ch03-cost}。

\begin{table}[htbp]
\centering
\caption{我要改什么 $\rightarrow$ 用哪类钩子 $\rightarrow$ 遍历的是面还是单元}
\label{tab:ch03-hook-map}
\begin{tabular}{>{\raggedright\arraybackslash}p{0.22\textwidth} >{\raggedright\arraybackslash}p{0.23\textwidth} >{\raggedright\arraybackslash}p{0.09\textwidth} >{\raggedright\arraybackslash}p{0.20\textwidth}}
\toprule
想改的东西 & 钩子族（形状） & 遍历 & 调用时机 \\
\midrule
边界上的入流剖面、壁面值 & \Tp{DEFINE\_{\Bk}PROFILE} 一类 & 面 & 组装方程前，逐面 \\
方程里的源项、阻尼项 & \Tp{DEFINE\_SOURCE} 一类 & 单元 & 组装方程时，逐单元 \\
物性（密度、黏度、比热） & \Tp{DEFINE\_PROPERTY} 一类 & 单元 & 需要物性时，逐单元 \\
按总体量修正参数、监测 & \Tp{DEFINE\_{\Bk}ADJUST} 一类 & 域级遍历 & 每迭代一次 \\
初值、历史变量清零 & \Tp{DEFINE\_{\Bk}INIT} 一类 & 全域遍历 & 初始化时刻，仅一次 \\
算完一步后写文件、统计 & 执行钩子（迭代末调用） & 域级遍历 & 迭代结束处 \\
\bottomrule
\end{tabular}
\end{table}

\section{最小可用例子：入口剖面与源项}
\label{sec:ch03-examples}

清单 \ref{lst:ch03-udf-profile} 是一个可以直接编译运行的入口速度剖面：
在圆管入口上给一条抛物线，中心 10 m/s，边界层内有平滑削峰。

\begin{lstlisting}[language=C, style=cstyle, caption={UDF：入口速度剖面（面遍历 + F\_ 访问器）}, label={lst:ch03-udf-profile}]
/* inlet\_profile.c —— 编译加载与解释加载都可用 */
#include "udf.h"            /* 求解器的宏与访问器都在这里 */

DEFINE_PROFILE(inlet_x_vel, thread, position)
{
    real x[ND_ND];          /* ND_ND 随 2D/3D 求解器展开 */
    real y;
    face_t f;               /* 面索引：只是一个句柄整数 */

    begin_f_loop(f, thread) /* thread 已绑定到该边界 */
    {
        F_CENTROID(x, f, thread);   /* 输出：面心坐标 */
        y = x[1];                   /* 径向坐标 */
        if (y > 0.05)               /* 近壁削峰，避免负值 */
            y = 0.05;
        F_PROFILE(f, thread, position) =
            10.0 * (1.0 - (y / 0.05) * (y / 0.05));
    }
    end_f_loop(f, thread)
}
\end{lstlisting}

实参逐个说清：\Tp{inlet\_x\_vel} 是宏生成的函数名，加载后在边界条件面板里
按这个名字选它；\Tp{thread} 是求解器传来的边界线程指针，它同时代表
``这个边界''和``这个边界的数据在哪里''，所有 \Tp{F\_*} 访问都要把它带上；
\Tp{position} 是求解器分配的剖面槽位号（同一边界可以同时挂速度剖面和
温度剖面），写值时必须原样传回 \Fn{F\_PROFILE}，否则写到别人的槽里。
\Tp{real} 是求解器的浮点类型，随单精度/双精度构建而变，不要写死 \Tp{double}。

清单 \ref{lst:ch03-udf-source} 给源项类钩子的约定：返回值就是源项本身，
而 \Tp{dS} 数组里放源项对因变量的导数。

\begin{lstlisting}[language=C, style=cstyle, caption={UDF：线性热源项（单元遍历 + 导数入雅可比）}, label={lst:ch03-udf-source}]
/* heat\_source.c —— 能量方程上的 S = A - B * T */
#include "udf.h"

#define A\_COEF  5.0e3       /* 常数项 W/m^3 */
#define B\_COEF  2.0e2       /* 线性项 W/(m^3 K) */

DEFINE_SOURCE(qlinear, c, t, dS, eqn)
{
    real source;

    /* eqn 由面板指定的方程决定，本例只服务能量方程 */
    dS[eqn] = -B_COEF;      /* dS/dT：负号利于隐式稳定 */
    source  = A_COEF - B_COEF * C_T(c, t);
    return source;          /* 返回值即该单元的源项 */
}

/* 同一个文件里再挂一个每迭代一次的钩子：只演示形状 */
DEFINE_ADJUST(max_T_probe, domain)
{
    /* 实参 domain 是求解域句柄：在这里自己发起域级遍历，
       把汇总值写进求解器的参数槽，供其它钩子读取。
       逐点钩子里不要做这种汇总。 */
    Message("adjust hook called once per iteration\n");
}
\end{lstlisting}

\begin{note}[title={源项一定要给导数}]
\Tp{dS[eqn]} 不是可选项。给成零等于告诉求解器``这一项与因变量无关''，
隐式组装里就少了这条对角修正，强非线性源项会直接表现为振荡或步长塌缩。
导数取解析值；实在推不出来时，取一个偏保守的常数斜率也比填零好。
\end{note}

清单 \ref{lst:ch03-property} 是物性钩子的形状：没有输出实参，返回一个标量。

\begin{lstlisting}[language=C, style=cstyle, caption={UDF：表观黏度（物性钩子，单元进标量出）}, label={lst:ch03-property}]
/* viscosity.c —— 幂律表观黏度 mu = k * shear^(n-1) */
#include "udf.h"

#define K\_CONS  0.55
#define N\_PW   0.42
#define MU\_MIN  1.0e-4
#define MU\_MAX  5.0e2

DEFINE_PROPERTY(mu\_apparent, c, t)
{
    real shear;

    shear = C_STRAIN_RATE_MAG(c, t);   /* 单元剪切率量级 */
    if (shear < 1.0e-6)
        shear = 1.0e-6;                /* 防 0 的负幂次 */
    if (shear > 1.0e4)
        shear = 1.0e4;
    return MAX(MU_MIN, MIN(MU_MAX,
           K_CONS * pow(shear, N_PW - 1.0)));
}
\end{lstlisting}

\Fn{C\_STRAIN\_RATE\_MAG} 这类由求解器算好的场量访问器，在部分版本或
部分求解设置下需要额外激活对应的存储；具体可用性以官方文档为准，
遇到取不到值的编译或加载错误时，先看该量在当前设置里是否被计算。

\section{解释加载与编译加载}
\label{sec:ch03-load-modes}

两种加载方式在``什么时候把源码变成机器码''这一点上分道：

\begin{itemize}
  \item \textbf{解释加载}：求解器用内置的 C 解释器直接吃源码，运行时机是
        每次调用时解释执行。不需要外部编译器（早期版本会借一个编译器
        做语法预处理，但产物不是本地机器码），改完源码要显式重新解释。
        限制包括：可用的宏与库函数是子集、调用开销按次数计费、
        复杂 UDF 常被拒绝解释。
  \item \textbf{编译加载}：求解器在自己的工作目录下搭一套 \Tp{libudf}
        源树与 makefile，调用外部 \CL{} 编译器把源码编成共享库
        （Linux 下 \Tp{libudf.so}、Windows 下 \Tp{libudf.dll}），
        运行时由求解器进程动态装载。需要受支持的编译器与对应
        SDK（Windows 上要 MSVC 一类，Linux 上通常是 gcc/Intel 系列），
        编译器的版本、位数与运行库必须和求解器构建一致，这一步
        是新手最常见的失败点。
\end{itemize}

\begin{lstlisting}[language=ShellLD, style=shellstyle, caption={编译加载的典型动作序列（TUI 路径与产物检查）}, label={lst:ch03-libudf}]
# 源码放到求解器工作目录下的 libudf 源树里
mkdir -p libudf/src
cp inlet_profile.c heat_source.c libudf/src/

# 起求解器：3 维双精度、8 个并行 rank
fluent 3ddp -t8 -cnf=./host_list

# TUI 路径：define -> user-defined -> compiled-functions
#   compile  逐个填源文件、编译选项、链接选项
#   load     装载共享库（并行时每个 rank 各自装载同一份库）
#   unload   卸载后才能重新编译，否则文件被占用

# 产物按维度/精度分目录，检查它是否真的存在
ls -l libudf/3ddp/libudf.so 2>/dev/null
ls -l libudf/src makefile_udf 2>/dev/null
\end{lstlisting}

并行下的装载语义要写清楚：共享库不是``主进程装一次然后共享内存''，
而是每个 rank 在自己的进程地址空间里各装一份。于是
\begin{itemize}
  \item 每个 rank 的库里跑的是同一份代码、各自的状态，跨 rank 的
        汇总必须显式用求解器的归约接口；
  \item 库文件必须放在所有 rank 都能读到的路径上，共享文件系统上的
        相对路径最容易出问题；
  \item 装载失败往往只在某个 rank 上报，症状是``部分核崩、部分核等''，
        先找日志里第一个报错的 rank。
\end{itemize}
调试输出同理：\Tp{printf} 走的是那个进程的标准输出，并行时多数
计算进程的 stdout 不一定被收集，看上去就像``我的 print 没执行''。
更稳的做法是用求解器自己的日志宏（\Tp{Message} 一类，会进求解器日志；
主控进程专用的变体只从 host 输出一次），并给每行日志带上
迭代号与边界名，否则日志会被大量重复淹没。求解器捕获到 UDF 出错时
会打印带函数名的警告并把该钩子停用，看到 \Tp{caught error thrown by UDF}
这类字样先回去看函数里有没有越界与除零。
各日志宏的确切名字与行为按版本略有差异，以官方文档为准。

\begin{compare}[title={解释加载与编译加载的分工}]
\begin{itemize}
  \item 起步与探路：解释加载。省去编译器对齐的坑，适合验证剖面形状、
        源项符号这类逻辑问题。
  \item 进热点循环：编译加载。单元级与面级的逐点钩子（物性、源项、
        剖面）在解释模式下通常慢一到两个数量级，迭代时间肉眼可见上涨。
  \item 需要完整宏库、外部库调用、并行归约、mesh motion 一类
        高级钩子：只能编译加载。
  \item 交付与复现：把 UDF 源码、\Tp{libudf} 目录与编译选项一起归档；
        解释模式的源码虽轻，但换版本后能不能解释不保证。
\end{itemize}
\end{compare}

\section{代价：热点循环、解释惩罚与内部缓存}
\label{sec:ch03-cost}

UDF 的开销不体现在代码行数，而体现在被调用次数。一个挂在入口上的
剖面钩子每次迭代被调用的次数等于入口面数；一个物性钩子的次数等于
全域单元数，还可能在同一迭代里因不同的组装阶段被重复调用。因此：

\begin{itemize}
  \item 逐点钩子里禁止出现超越函数之外的重活：堆分配、字符串格式化、
        三角/对数的重复计算、把同一个常数每次迭代重新算一遍。
        把不变量提出来：用编译期常量、初始化钩子里算一次、
        或存在求解器的参数槽里。
  \item 访问器看着像数组下标，实际常常是一串指针解引用与位移计算。
        热点里应当一次取多个量、局部变量算完再写回，而不是在同一条
        表达式里反复调访问器。
  \item 解释模式的惩罚按调用次数线性放大，这就是``单点测试没问题、
        上大网格后时间翻倍''的典型原因。
\end{itemize}

还要考虑与求解器内部缓存的相互作用。求解器会缓存面上的重构值、
梯度与限制器结果，并在自己的时序里决定何时重算。UDF 若绕过钩子
直接改存储（某些 \Tp{C\_*} 访问器是可写的），改动可能落在缓存之后，
导致``这一迭代用的是旧梯度''；反过来，把一个昂贵钩子的结果缓存到
自定义存储里，就要自己负责失效时机（网格改变、并行分区改变、
重启都可能让它过期）。哪一类量允许在哪个钩子里写、写完是否要求
重算，属于版本相关的约定，以官方文档为准；本章的立场是：
默认只改你挂的那个钩子负责的槽位，其他存储只读。

\begin{warning}[title={在 UDF 里做全局归约或阻塞 IO}]
最贵的一类事故不是算得慢，而是把并行打散：
\begin{itemize}
  \item 在每迭代的域级钩子里做全局求和，等于给每个迭代插入一次集合
        通信；单元级钩子里更不该有任何通信。
  \item 在逐点钩子里打开文件写一行日志，几万个 rank 各开各的句柄，
        共享存储上会退化成串行；文件锁还会让快 rank 等慢 rank。
  \item 用 \Tp{static} 变量``缓存''上次调用的结果，在并行下每个 rank
        各自一份，于是同一个算例里出现两套历史。
  \item 需要汇总值：把归约放进每迭代一次的钩子，交给求解器的归约
        接口做，并把结果写成求解器参数；需要输出：交给执行钩子，
        只在主控进程写，且按迭代间隔抽稀。
\end{itemize}
\end{warning}

\subsection*{本章小结}
UDF 的接口是``时机 + 遍历对象 + 访问器前缀''三件事的组合：先由
\Tp{DEFINE\_*} 宏名确定钩子何时被调、遍历面还是单元，再由
\Fn{F\_PROFILE} 或 \Tp{dS[eqn]} 这类约定确定把结果写到哪。
选加载方式的标准只有一条——钩子是否在热点里。
下一章把同一套形态搬到 \DomCAE{} 的结构求解器侧，那里的扩展面是材料、
单元与结果，代价结构却完全类似。

% ============================================================
% ---- frag_ch04.tex ----
\chapter{Mechanical 的用户扩展：材料、自定义单元与结果后处理}
\label{ch:ch04-mechanical-ext}

上一章把 \DomCAE{} 侧流体求解器的扩展面写成了``宏名决定时机''；本章在结构求解器
（\Lib{Mechanical} 这一族）上看同一个形态的另一面：宏换成例程，
遍历面/单元换成遍历积分点，写回剖面换成写回应力与切线矩阵。
产品接口名在不同版本间确实变过，因此本章凡不确定的地方一律只给\emph{调用形状}
并注明以官方文档为准；能给出完整代码的部分，是自己就能编译验证的数值内核。

\section{三条扩展路线与取舍}
\label{sec:ch04-routes}

结构侧的用户扩展按``改动离求解器内核有多近''排成一条线：

\begin{itemize}
  \item \textbf{脚本层参数化}：几何、网格、材料常数、边界条件、求解设置、
        结果提取都由外部脚本驱动。求解器二进制不动，扩展物是脚本与数据。
        \PY{}、批处理命令文件、图形界面里可记录的命令流都属于这一层。
  \item \textbf{用户材料（本构积分）}：把``给定应变增量，求应力与切线''这件事
        交给你写的例程。求解器仍在自己的循环里调你，但调用粒度细到
        每个积分点、每个增量步，于是这一层与上一章的流体钩子有同一条纪律：
        无状态、可重入。
  \item \textbf{外部数据与结果的双向通道}：把外部场（温度、载荷谱、
        另一台机器的结果）映射成求解器的载荷或初值，再把结果取出去。
        它改的是输入输出，不是算法。
  \item 边缘的第四层是\textbf{自定义单元与附加自由度}：需要注册自由度类型、
        提供形函数与单元刚度/内力组装规则，接口面最大、文档最薄，
        通常只在现有单元族确实表达不了你的物理时才碰。
\end{itemize}

\begin{table}[htbp]
\centering
\caption{选哪条路线：改动位置、前置成本与主要风险}
\label{tab:ch04-route-choice}
\begin{tabular}{>{\raggedright\arraybackslash}p{0.15\textwidth} >{\raggedright\arraybackslash}p{0.25\textwidth} >{\raggedright\arraybackslash}p{0.19\textwidth} >{\raggedright\arraybackslash}p{0.20\textwidth}}
\toprule
路线 & 适合做什么 & 前置成本 & 主要风险 \\
\midrule
脚本层参数化 & 批量算例、扫参、统一建模规范、结果批导 & 会写脚本即可，不碰编译 &
建模语义与手工操作不一致 \\
用户材料例程 & 新本构、损伤、粘塑性、试验标定的模型 & 编译器、例程模板、回归框架 &
切线写错导致不收敛或假收敛 \\
外部数据通道 & 热—力耦合、载荷谱、其他求解器结果映射 & 数据格式与映射设置 &
网格不一致时插值误差 \\
自定义单元 & 新自由度、特殊连接、降阶单元 & 最深，需接口文档逐条对齐 &
与求解器的单元机制互相假设，难定位 \\
\bottomrule
\end{tabular}
\end{table}

三句话的取舍：能用脚本表达的不要进内核；只有``应力—应变关系不在内置库里''
才写用户材料；只有``连自由度都不够''才写自定义单元。
自定义单元的形状与用户材料同构（例程按单元被调用，输入节点坐标、
节点自由度值与材料参数，输出刚度阵与内力向量，再额外提供 DOF 注册与形函数），
但可用的入口和命名要逐版本核对官方扩展文档，本章不编造签名。

\section{用户材料例程的调用形状}
\label{sec:ch04-umat-shape}

把这一族的入口抽象出来，它就是一个例程，由求解器在每个增量步、
每个积分点调用一次，形如：

\begin{itemize}
  \item \textbf{输入}：应变增量（或总应变，取决于版本约定）、应力历史、
        温度与时间增量、材料常数数组、方向余弦（各向异性时）、
        旧的历史变量。
  \item \textbf{输出}：更新后的应力、切线刚度矩阵（进整体刚度 assembly）、
        更新后的历史变量、可选的应变能与塑性耗散能。
  \item \textbf{约定}：应力/应变按 Voigt 六分量顺序给定；
        剪切分量到底是张量剪应变还是工程剪应变、以及切线矩阵剪切行的
        系数是否要求乘二，各家产品不同，这一条必须以官方文档为准，
        写反了的症状是``刚度不对称或收敛极慢''。
\end{itemize}

清单 \ref{lst:ch04-umat-kernel} 给一个可以直接编译、可以脱离求解器跑的
本构内核：小应变 $J_2$ 弹塑性、线性各向同性硬化、径向返回（return mapping），
并同时返回一致切线。它就是上面那个入口去掉外壳之后的数值部分。

\begin{lstlisting}[language=Fortran, style=fortranstyle, caption={用户材料内核骨架：J2 径向返回 + 一致切线（\Lang{Fortran} 自由格式）}, label={lst:ch04-umat-kernel}]
! umat_kernel.f90 —— 本构内核：单积分点一次调用
! 约定：分量顺序 11,22,33,12,23,13；剪切用张量剪应变
! 历史变量只有一个：statev(1) 存等效塑性应变
subroutine umat_kernel(stress, jtang, ddestr, statev, nhisv)
  implicit none
  integer,          intent(in)                :: nhisv
  double precision, intent(in)                :: ddestr(6)   ! 应变增量
  double precision, intent(inout)             :: stress(6)   ! 应力
  double precision, intent(out)               :: jtang(6, 6) ! 一致切线
  double precision, intent(inout)             :: statev(nhisv)
  double precision, parameter :: e0 = 2.1d11  ! 杨氏模量
  double precision, parameter :: v0 = 0.3d0   ! 泊松比
  double precision, parameter :: sy = 2.5d8   ! 初始屈服应力
  double precision, parameter :: hh = 1.0d9   ! 硬化模量
  double precision :: lam, mu, bulk, vol, pmean, strn, seqv
  double precision :: rq, ftr, dg, aa, hp, denom, eeq
  double precision :: dev(6), nvec(6), mvec(6), idv(6, 6)
  integer :: i, j

  mu    = 0.5d0 * e0 / (1.0d0 + v0)
  lam   = e0 * v0 / ((1.0d0 + v0) * (1.0d0 - 2.0d0 * v0))
  bulk  = lam + 2.0d0 * mu / 3.0d0
  do i = 1, 6
    mvec(i) = 0.0d0
    if (i .le. 3) mvec(i) = 1.0d0
  end do

  ! ---- 1) 弹性预测 ----
  vol = ddestr(1) + ddestr(2) + ddestr(3)
  do i = 1, 6
    stress(i) = stress(i) + 2.0d0 * mu * ddestr(i)
    if (i .le. 3) stress(i) = stress(i) + lam * vol
  end do
  pmean = (stress(1) + stress(2) + stress(3)) / 3.0d0
  do i = 1, 3
    dev(i) = stress(i) - pmean
  end do
  do i = 4, 6
    dev(i) = stress(i)
  end do
  strn = dsqrt(dev(1)**2 + dev(2)**2 + dev(3)**2 &
       + 2.0d0 * (dev(4)**2 + dev(5)**2 + dev(6)**2))
  nvec = dev
  if (strn > 0.0d0) nvec = dev / strn
  seqv = dsqrt(1.5d0) * strn     ! von Mises 等效应力
  rq   = sy + hh * statev(1)     ! 当前屈服半径
  ftr  = seqv - rq

  ! ---- 2) 返回映射与切线系数 ----
  aa = 1.0d0
  hp = 0.0d0
  eeq = 0.0d0
  if (ftr > 0.0d0) then
    denom = 2.0d0 * mu * dsqrt(1.5d0) &
            + hh * dsqrt(2.0d0 / 3.0d0)
    dg = ftr / denom             ! 塑性乘子增量
    aa = 1.0d0 - 2.0d0 * mu * dg / strn
    hp = dsqrt(2.0d0 / 3.0d0) * hh / denom
    eeq = dsqrt(2.0d0 / 3.0d0) * dg
    do i = 1, 6
      stress(i) = aa * dev(i)    ! 偏量按比例缩回屈服面
    end do
    do i = 1, 3
      stress(i) = stress(i) + pmean ! 平均应力不变（塑性不可压缩）
    end do
    statev(1) = statev(1) + eeq
  end if

  ! ---- 3) 一致切线：bulk*mm + 2mu*aa*(Idev - nn) + 2mu*hp*nn ----
  do i = 1, 6
    do j = 1, 6
      idv(i, j) = 0.0d0
      if (i == j) idv(i, j) = 1.0d0
      if (i <= 3 .and. j <= 3) idv(i, j) = idv(i, j) - 1.0d0 / 3.0d0
      jtang(i, j) = bulk * mvec(i) * mvec(j) &
        + 2.0d0 * mu * aa * (idv(i, j) - nvec(i) * nvec(j)) &
        + 2.0d0 * mu * hp * nvec(i) * nvec(j)
    end do
  end do
end subroutine umat_kernel
\end{lstlisting}

读这段代码的三个要点：\Tp{statev} 是本构唯一允许的``记忆''载体，
它的尺寸由你在接口里声明的历史变量个数决定；\Tp{aa} 是弹性卸载比例，
纯弹性时退化为 1；切线在塑性分支里被拆成体变项、偏量投影项与
流动方向项三块，这正是径向返回一致切线的标准形状。
材料常数在这里被写死成参数，实战中它们必须来自求解器传来的材料数据数组，
否则参数化研究（\S\ref{sec:ch04-routes} 的第一条路线）与你的例程会脱节。

\section{为什么必须返回切线：一致切线与数值微分}
\label{sec:ch04-tangent}

Newton 型求解在每一步要求 $\Delta F = -K_T^{-1} R$，其中 $K_T$ 由各个积分点的
切线模量组装。你的本构只给出应力而不给模量时，有两条退路，都不便宜：

\begin{itemize}
  \item \textbf{用连续弹性刚度}（即所谓 \Tp{DRAG\_IN} / 不更新刚度）：
        迭代方向不再是牛顿方向，收敛从二次降到线性甚至停滞，
        步数常常翻几倍。
  \item \textbf{数值微分切线}：对六个应变分量各施加一个扰动，
        重新跑一遍你的本构积分，取差分。六个分量扰动六次，
        每次调用都要完整积分一遍，热点成本大约按分量数放大；
        同时扰动步长带来截断与舍入误差，取小了两头失真、
        取大了切线偏软，症状是``勉强收敛但对步长敏感''。
\end{itemize}

一致切线（algorithmic/consistent tangent）指的是：把``你这一步实际做了
什么''线性化，而不是把``真实材料的连续切线''抄过来。上例里塑性分支的
$2\mu\,h'\,n\otimes n$ 这一项就来自返回映射本身；用了它，残差的线性化
与你更新的应力一致，Newton 才保住二次收敛率。写错的代价很有欺骗性：
只要矩阵偏保守（偏硬），算例照样收敛，只是荷载—位移曲线偏刚，
最后被当成``材料本来就这样''。

\begin{note}[title={切线自检的三条硬指标}]
\begin{itemize}
  \item 对称性：$D_{ij} = D_{ji}$，各向同性关联塑性的切线必然对称；
        不对称多半是 Voigt 顺序或剪切系数搞反了。
  \item 弹性退化：应变增量足够小时必须退回弹性刚度阵。
  \item 方向性：单轴拉伸塑性段的模量应落在弹性模量与硬化模量之间，
        且随应变单调；出现超过弹性模量的值就是算错了。
\end{itemize}
\end{note}

\section{历史变量、状态数组与并行}
\label{sec:ch04-statev}

状态数组的约定形状是：你在例程的声明/模板里给出历史变量的\emph{个数}，
求解器按``每个积分点一份''的私有存储去分配，把指针递给你；
每个增量步调用你时，先带来上一步结束时的值，你负责写回新值。
因此：

\begin{itemize}
  \item 尺寸要在编译期或接口声明里定死（数组必须能越界检查），
        历史变量个数改变意味着必须重新编译并同步所有已有算例，
        旧结果文件里的状态维数与新例程不一致时会静默读错或报错。
  \item 存储位置归求解器管：重启、子模型、结果读回、网格迁移
        （remap）都要靠它把每个积分点的状态搬到新位置。
        自己维护外部数组的话，这些机制全部失效。
  \item 状态变量的顺序就是协议：把温度、屈服面、背应力、损伤量
        按固定顺序编号并写进文档，改顺序等于改文件格式版本。
\end{itemize}

并行是当前结构求解器把同一个例程在多个线程上并发跑：域分解后每个
rank 有自己的一套积分点，同一 rank 内部还可能有计算线程按单元分块。
于是与上一章流体侧完全同构的纪律成立——\textbf{禁止共享可变态}：
不要企图用模块级的全局数组缓存``上次算的结果''，它的可见范围是
本进程/本线程，不同 rank 上的同名变量各自漂移；
不要做全局归约，不要让一次积分点调用等待别的调用；
不要在这里打印文件。要跨步保存的，一律走状态数组。

\begin{warning}[title={状态与并行的四类典型事故}]
\begin{itemize}
  \item 例程内部用 \Tp{save} 或模块全局量做缓存：分区改变后结果变化，
        同一算例换个核数跑不出同一答案。
  \item 历史变量个数改了但没重编译：新维数与旧库不一致，
        症状是状态值看起来是随机数。
  \item 增量步内让本构``自己再迭代一次''并把中间量写回状态数组的
        固定位置：出错时写入半截状态，重启后材料点直接坏掉。
        要么整个更新成功后才写回，要么在局部变量里算完。
  \item 在例程里调用带阻塞或文件 IO 的外部库：把一个逐点热点
        变成同步点，核数越多越慢。
\end{itemize}
\end{warning}

\section{结果侧扩展：自定义后处理量与批量导出}
\label{sec:ch04-results}

结果侧有两件事可做：定义求解器里没有现成按钮的量，以及把结果搬到外部。

\begin{enumerate}
  \item \textbf{自定义后处理量}：最省的形式是给一个表达式（用已有结果项
        组合，例如某个应力分量差、塑性应变比、某种等效值），
        由求解器在自己内部遍历求值；这属于结果层的``钩子''，
        能避免把整场数据搬出来再算。
  \item \textbf{脚本层批量导出}：在脚本里请求``给定节点集/单元集 +
        结果项 + 时间步''的数组，落盘。调用形状各家相似，
        名字与参数以官方文档为准。优点是能做选择与聚合，
        缺点是脚本进程与求解器进程要来回搬数据，全模型遍历时会很慢。
  \item \textbf{直接读结果文件}：结果文件是二进制、按步/按分区组织的，
        直接读省掉一次导出，但格式属实现细节，跨版本不保证兼容，
        而且并行写出的结果可能分散在多份分区文件里。
        除非要做第三方软件的插件式读取，一般不推荐。
  \item \textbf{用户材料自己送出的量}：应变能、塑性耗散这类
        由积分点累加的量，通常靠例程的输出参数交给求解器存场，
        然后作为普通结果项取用；历史变量本身也常被当作结果输出。
\end{enumerate}

大规模结果只取子集，四条常用招法，按有效性排序：

\begin{itemize}
  \item 先选集合再取值：用节点集/单元集/Named Selection 把范围钉死，
        导出接口只接受集合，不要``取全部然后在客户端筛''。
  \item 抽时间步：几百个增量步里只输出关心的若干步（写结果时设间隔），
        这比读的时候筛更省，因为省的是 IO 与磁盘。
  \item 抽分量：只请求需要的张量分量而不是整个张量，
        体积按分量数线性下降。
  \item 在求解器内部做归约：路径、截面力、监测点的时程都在内部算完
        只送标量，让搬运量降到时间步数量级。
\end{itemize}

清单 \ref{lst:ch04-post} 是``取出去之后''的典型一段：读入脚本层导出的
长表 CSV，按节点聚合，只保留关心的量并写回一张小表。它不依赖任何
产品 API，因此可以原样进回归流水线。

\begin{lstlisting}[language=Python, style=pythonstyle, caption={结果后处理：从导出的长表算自定义量并只留子集}, label={lst:ch04-post}]
"""后处理：输入是脚本层导出的长表（step,node,eps_eq,ux,uy,uz）。

输出：每个关心的节点，加载路径 sqrt(ux^2+uy^2+uz^2) 随步的变化，
以及增量位移幅值最大的若干节点。纯标准库，可直接进回归流水线。
"""
import csv
import math
from collections import defaultdict

WATCH = {101, 205, 317}          # 关心的节点号（来自 Named Selection）


def load_rows(path):
    by_node = defaultdict(list)
    with open(path, newline="", encoding="utf-8") as f:
        for row in csv.DictReader(f):
            node = int(row["node"])
            if node not in WATCH:      # 先筛后算，别把全场搬进内存
                continue
            path_len = math.sqrt(sum(float(row[k]) ** 2
                                     for k in ("ux", "uy", "uz")))
            by_node[node].append((int(row["step"]),
                                  float(row["eps_eq"]), path_len))
    return by_node


def report(by_node, out_path):
    with open(out_path, "w", newline="", encoding="utf-8") as f:
        w = csv.writer(f)
        w.writerow(["node", "step", "eps_eq", "path", "dpath"])
        for node, rows in sorted(by_node.items()):
            rows.sort()
            prev = None
            for step, eps, path_len in rows:
                d = 0.0 if prev is None else path_len - prev
                w.writerow([node, step, f"{eps:.6g}",
                            f"{path_len:.6g}", f"{d:.6g}"])
                prev = path_len


if __name__ == "__main__":
    import sys
    report(load_rows(sys.argv[1]), sys.argv[2])
\end{lstlisting}

\section{验证与回归：从单元素到解析解}
\label{sec:ch04-verify}

用户材料必须先在没有求解器的环境里被验通，再进算例。三级测试：

\begin{enumerate}
  \item \textbf{单元素 / 单积分点测试}：自己构造应变历史，直接调内核，
        检查应力、状态量与切线。这是最便宜也最有效的一层，
        能抓住绝大多数实现错误。清单 \ref{lst:ch04-single-elem} 是骨架。
  \item \textbf{均匀化测试}：一个单元网格（一层高斯点）施加周期/均匀应变，
        期望``全场应力相同、荷载—位移斜率等于切线模量预测值''。
        它检查例程与求解器的装配是否对齐，也就是 Voigt 顺序与
        剪切系数那一类约定错误——单元素纯数值测试抓不住这类问题。
  \item \textbf{解析解对照}：梁的纯弯、厚壁筒的内压（拉梅解）、
        半空间压痕这类有闭式解的模型，用同一套本构跑一次，
        并把残差阈值写进回归脚本而不是肉眼看图。
\end{enumerate}

\begin{lstlisting}[language=C, style=cstyle, caption={单元素回归测试骨架：链接 \Lang{Fortran} 内核做三条断言}, label={lst:ch04-single-elem}]
/* test_umat.c —— 与求解器无关，只做积分点层面的检查
 * 注意：Fortran 符号名修饰（下划线/大小写）随工具链而定，
 *       链接前先确认实际的导出符号。 */
#include <stdio.h>
#include <math.h>

extern void umat_kernel(double s[6], double d[6][6],
                        const double de[6], double sv[1],
                        int nhisv);

static int fails = 0;

static void check(const char *tag, int ok)
{
    if (!ok) { printf("FAIL %s\n", tag); fails++; }
}

/* 弹性刚度：与解析的线弹性阵逐元素比 */
static void expect_elastic(double d[6][6], double e, double v)
{
    double lam = e * v / ((1.0 + v) * (1.0 - 2.0 * v));
    double mu = 0.5 * e / (1.0 + v);
    check("D11", fabs(d[0][0] - (lam + 2 * mu)) < 1e-6);
    check("D12", fabs(d[0][1] - lam) < 1e-6);
    check("D44", fabs(d[3][3] - 2.0 * mu) < 1e-6);
}

int main(void)
{
    double s[6] = {0}, d[6][6];
    double sv[1] = {0};
    double de[6] = {2.0e-4, 0, 0, 0, 0, 0};  /* 单轴小步：纯弹性 */

    umat_kernel(s, d, de, sv, 1);
    expect_elastic(d, 2.1e11, 0.3);
    check("s22_eq_s33", fabs(s[1] - s[2]) < 1e-6);
    check("no_plastic", sv[0] == 0.0);

    de[0] = 5.0e-3;                          /* 大增量：进入塑性 */
    s[0] = 0; s[1] = 0; s[2] = 0;
    umat_kernel(s, d, de, sv, 1);
    check("yielded", sv[0] > 0.0);
    for (int i = 0; i < 6; i++)
        for (int j = 0; j < 6; j++)
            check("tangent_symmetric", fabs(d[i][j] - d[j][i]) < 1e-6);
    printf("fails=%d ep=%g\n", fails, sv[0]);
    return fails != 0;
}
\end{lstlisting}

能量检查是第二道闸：一个增量步里，外力功应等于应变能增量加塑性耗散
（$\sigma : \Delta\varepsilon^{p}$ 的累加），两者之差按步长应当单调收敛；
如果你的例程顺手返回了应变能/耗散能，就把这两个量也拿去做同一条曲线。
刚度对称性用上面这种逐元素比较即可，代价是一次调用；不要等求解器
报``矩阵不对称''再回头查，那时定位成本高得多。

\begin{bestpractice}[title={用户材料扩展的交付清单}]
\begin{itemize}
  \item 例程源码、构建脚本、单元素测试、均匀化测试、解析解算例
        五件同仓；改本构必跑前两件，发版必跑后两件。
  \item 回归的判据是数字：把最大残差、能量差、切线对称性上限
        写进脚本返回值，不要只留一张云图。
  \item 历史变量的含义与顺序写进代码注释和文档同一处，
        并给出``旧算例在新例程下是否可重启''的结论。
  \item 记录求解器版本、编译器版本与优化选项；本构这类数值代码
        在极端优化下可能因浮点重排改变收敛行为，出问题先降优化复现。
  \item 本构里的常数一律从材料数据传入；写死一个常数，
        就等于给未来的参数化研究埋了一颗雷。
\end{itemize}
\end{bestpractice}

\subsection*{本章小结}
结构侧的扩展面比剖面钩子更硬：它要求你不仅知道被调用的时机，
还要给出线性化的结果（切线），并为数值一致性负责。
与上一章相同的是三条纪律——无状态、可重入、只碰自己负责的存储；
与上一章不同的是代价的重心：流体侧慢在遍历次数，结构侧错在切线方向。

% ============================================================
% ---- frag_ch05.tex ----
\chapter{Abaqus 与 LS-DYNA：用户子程序的调用约定与 C 混编}
\label{ch:ch05-abaqus-usub}

\DomCAE{} 里两个以结构求解闻名的产品——\Lib{Abaqus} 与 \Lib{LS-DYNA}——共享同一套扩展哲学：
求解器的核心不动，用户只在\textbf{材料本构}与若干挂钩点上提供一段代码，由求解器在积分点层面
反复调用。这段代码的母语是 Fortran（历史包袱，也是数值社区的通用语），而工程上真正想写的
东西往往是 \CPP{} 的材料模型、状态机、损伤累积或外部耦合。于是这一章的重心不在本构理论，
而在\textbf{调用约定}：谁在什么粒度上调用你、你必须往回写什么、以及跨越
Fortran/{\Bk}\CL{} ABI 时会被咬到的五处地方。

签名口径先说清楚：下面只在 \Tp{UMAT} 这类形状稳定的入口上写形参名，其余一律按
\textbf{调用形状}描述（入口是一段 Fortran 子程序、第几个实参是数组首地址、返回什么方向），
具体形参顺序、完整清单与新增参数以官方文档和你所用版本的用户子程序条目为准。

\section{隐式与显式的分界：切线刚度到底从哪来}
\label{sec:ch05-implicit-explicit}

用户子程序生态里最重要的一个分界，不是产品，而是求解器类型。

隐式（\Lib{Abaqus}/Standard 一类）求解的是非线性方程组的 Newton 迭代：每步、每次迭代
要求解 $\partial R/\partial u$ 型的整体刚度，单元刚度由材料层给出的
\textbf{应力—应变增量导数}拼装而来。这就意味着材料子程序必须回答一个二阶问题：
``给定这一增量应变，应力增量是多少，而且这个映射对我的迭代格式来说导数是多少''。
如果 \Tp{DDSDDE} 只是连续体（率型）本构矩阵，而应力更新却用了一个显式返回映射算法，
残差的线性化就和实际更新不一致，Newton 失去二次收敛，表现为迭代次数膨胀、
步长被反复削小、最后在某个奇怪的角度上发散。所谓\textbf{一致切线刚度}
（consistent tangent），就是把返回映射这一步也微分进去得到的矩阵。

显式（\Lib{Abaqus}/Explicit、\Lib{LS-DYNA} 的显式序列）不做整体线性化：
它是中心差分的时间推进，节点力由单元应力直接装配，稳定性由临界时间步
（网格特征长度除以材料波速）决定。因此用户入口只需要给出\textbf{应力}
（以及内能/密度一类守恒量），完全不需要 Jacobian——线性化在这个格式里没有消费者。
代价换了个方向：单步必须小到应变增量落在本构曲线的局部直线段内，
算法一致性变成``增量足够小''这一外部保证；反过来，写子程序的门槛低很多，
调试时也没有``收敛历史''这种中间信息可看，只能靠能量曲线判断哪里错了。

\begin{compare}[title={一致切线、连续体切线与干脆不给切线}]
三者都不是免费的。\textbf{一致切线}让隐式迭代次数最少且对步长不敏感，
推导量大（非关联流动、各向异性损伤时矩阵还不对称，装配与求解都要走非对称路径）；
\textbf{连续体切线}推导容易、对极端硬化软化更``钝''因而有时更稳，但迭代次数随步长
显著增加，Newton 退化为线性收敛甚至振荡；\textbf{不给切线}只在显式里合法。
一个实用折中：隐式里可以先用连续体切线打通流程，再逐段替换为一致切线，
用迭代次数曲线量化差别，而不是先花两周推导。
\end{compare}

\section{UMAT 的调用形状：按积分点、应力回来、切线出去}
\label{sec:ch05-umat-shape}

\Tp{UMAT} 的调用粒度是\textbf{单个积分点单次增量}：求解器为每个单元中的每个积分点、
每个增量步、每次 Newton 迭代都调一次。这意味着子程序内部不能有跨调用的循环假设，
也不能假定调用顺序与单元编号顺序一致。形状稳定的部分只有一句：
进来的\textbf{上步末应力 + 本步应变增量 + 历史变量}，出去的
\textbf{新应力 + 新历史变量 + 切线矩阵}。骨架如下。

\begin{lstlisting}[language={[95]Fortran}, style=fortranstyle, caption={UMAT 骨架：常引用形参与逐参数注释
（形参完整清单与顺序以官方用户子程序条目为准）}, label={lst:ch05-umat-skeleton}]
      SUBROUTINE UMAT(STRESS,STATEV,DDSDDE,
     1 SSE,SPD,SCD,STRAN,STRANINC,TIME,DTIME,
     2 CMNAME,MATERIAL,NPROP,NDI,NSH,NTENS,NSTATV,PROPS,
     3 TEMP,NOEL,NPT,LAYER,KSPT,KSTEP,KINC,NTIME)
      IMPLICIT NONE
!      --- 方向标记：IN 只读，INOUT 必须在原值上累加，OUT 由你写 ---
      INTEGER NTENS, NSTATV, NPROP, MATERIAL
      INTEGER NOEL, NPT, LAYER, KSPT
      INTEGER NDI, NSH, KSTEP, KINC, NTIME
      DOUBLE PRECISION STRESS(NTENS)     ! INOUT 应力向量，IN 是上步末值
      DOUBLE PRECISION STATEV(NSTATV)    ! INOUT 历史变量，尺寸来自材料卡声明
      DOUBLE PRECISION DDSDDE(NTENS,NTENS)! OUT 列主序，把应力增量映到应变增量
      DOUBLE PRECISION STRAN(NTENS)      ! IN  总应变（相对本步起始构型）
      DOUBLE PRECISION STRANINC(NTENS)   ! IN  本步应变增量，DDSDDE 乘的是它
      DOUBLE PRECISION TIME(2)           ! IN  TIME(1) 总时间, TIME(2) 步内进度
      DOUBLE PRECISION DTIME, TEMP       ! IN  增量时间步长, 当前温度
      DOUBLE PRECISION PROPS(NPROP)      ! IN  材料卡上给的常数数组
      DOUBLE PRECISION SSE, SPD, SCD     ! OUT 弹性能/塑性耗散/蠕变耗散
      CHARACTER*8 CMNAME                 ! IN  材料名，用于调试打印
      INTEGER I, J
      DOUBLE PRECISION EQPL, TRIAL, DGAM
!      1) 从历史变量取回上步状态：UMAT 是无状态的，一切靠 STATEV
      EQPL = STATEV(1)
!      2) 用 DDSDDE 做试弹性更新（这里只演示形状，不是可用本构）
      DO J = 1, NTENS
        DO I = 1, NTENS
          STRESS(I) = STRESS(I) + DDSDDE(I,J) * STRANINC(J)
        END DO
      END DO
!      3) 返回映射：把应力拉回屈服面，同时累加等效塑性应变
      TRIAL = 1.0D0
      DGAM  = 0.0D0
      EQPL  = EQPL + DGAM
      STATEV(1) = EQPL
!      4) 三个能量量必须自洽，否则全局能量平衡会露馅
      SSE = 0.5D0 * TRIAL
      SPD = 0.0D0
      SCD = 0.0D0
      RETURN
      END SUBROUTINE UMAT
\end{lstlisting}

\begin{table}[htbp]
\centering
\caption{\Tp{UMAT} 常引用参数的方向与语义（顺序不保证，以官方为准）}
\label{tab:ch05-umat-args}
\begin{tabular}{>{\raggedright\arraybackslash}p{0.19\textwidth}>{\raggedright\arraybackslash}p{0.09\textwidth}>{\raggedright\arraybackslash}p{0.50\textwidth}}
\toprule
参数 & 方向 & 语义与常见误用 \\
\midrule
\Tp{STRESS} & in/out & 进来是上步末的客观应力，必须在其上累加增量后写回；直接覆盖成``全量本构''结果会破坏多步历史。 \\
\Tp{STATEV} & in/out & 历史状态向量，长度由材料卡的变量数决定；每个积分点私有，是唯一合法的``记忆''。 \\
\Tp{DDSDDE} & out & \Tp{NTENS} 乘 \Tp{NTENS} 的切线矩阵，列主序；隐式收敛速度基本由它决定。 \\
\Tp{STRANINC} & in & 本步应变增量；\Tp{STRAN} 是总应变，增量式本构几乎只用前者。 \\
\Tp{TIME}、\Tp{DTIME} & in & 全局时间与步内归一化进度；蠕变/粘弹按 \Tp{DTIME} 推进内部时间。 \\
\Tp{NDI}、\Tp{NSH}、\Tp{NTENS} & in & 正分量数、剪切分量数、向量总长；工程剪应变与张量剪应变的因子 2 差在这里暴露。 \\
\Tp{NOEL}、\Tp{NPT}、\Tp{LAYER}、\Tp{KSPT} & in & 定位当前积分点，只适合打印与断言，不适合做索引。 \\
\Tp{SSE}、\Tp{SPD}、\Tp{SCD} & out & 弹性能、塑性耗散、蠕变耗散，是能量平衡检查的原料。 \\
\bottomrule
\end{tabular}
\end{table}

两个必须记住的细节。其一，\Tp{DDSDDE} 与 \Tp{STRESS} 处于同一客观应力量测
（\Lib{Abaqus} 的 UMAT 约定是 Jaumann 率框架），大变形下把弹性矩阵当成``固定参考系''
里的常量是最常见的错误来源。其二，剪切分量的应变约定可能是工程值（因子 2），
写各向异性时若不在 \Tp{DDSDDE} 的对应行列上补偿，刚度矩阵就会不对称，
而求解器默认走对称求解器——症状是收敛极慢而不是报错。

\Lib{Abaqus}/Explicit 一侧的对应入口常被称为 \Tp{VUMAT}，调用形状换成
\textbf{一次一块}：实参里给出本增量要更新的积分点个数，应力与状态变量以二维块
传入，你必须循环这个块；不需要 \Tp{DDSDDE}，取而代之的是能量/质量更新类返回值。
\Lib{LS-DYNA} 的用户材料入口同样是 Fortran 子程序族，按关键字选择接口类型，
返回偏应力率、内能/密度演化一类的量，具体入口名与形参以官方用户材料手册为准；
它的一个实际差别是历史状态要按求解器提供的存储机制显式登记，
而不是像 \Tp{STATEV} 那样由关键字统一给尺寸。

\subsection{加载方式：共享对象与重链可执行文件}
\label{sec:ch05-loading}

\Lib{Abaqus} 走\textbf{运行时加载}：作业参数里指定用户子程序目标或共享库，
求解器进程启动后按符号名找到入口，改了材料只需重编子程序、重跑作业。
\Lib{LS-DYNA} 的传统路径是\textbf{把用户材料的 Fortran 目标文件与求解器源码一起
重新链接}出一个定制可执行文件（串行与并行 MPP 版本还要各出一份），
因为 MPI 版里子程序要在每个 rank 上存在。两者对混编的影响见 \S\ref{sec:ch05-cpp-shell}
末尾：符号名与运行时库必须和求解器编译链一致。具体工具链、开关与并行版本要求以官方为准。

\section{与 \CL{}/\CPP{} 混编：ABI 的五个咬合点}
\label{sec:ch05-cpp-shell}

把本构写成 \CPP{} 的理由通常有三个：已有代码库、需要模板化的状态机、
要用外部库（特殊函数、优化器、自动微分）。求解器不会因此变成 \CPP{}，
所以必须造一层壳。五个咬合点见表 \ref{tab:ch05-abi}：
\begin{table}[htbp]
\centering
\caption{Fortran 与 \CL{} 边界的约定差异与统一做法}
\label{tab:ch05-abi}
\begin{tabular}{>{\raggedright\arraybackslash}p{0.16\textwidth}>{\raggedright\arraybackslash}p{0.20\textwidth}>{\raggedright\arraybackslash}p{0.22\textwidth}>{\raggedright\arraybackslash}p{0.22\textwidth}}
\toprule
事项 & Fortran 侧 & \CL{}/{\Bk}\CPP{} 侧 & 统一做法 \\
\midrule
名字修饰 & 外部符号常带下划线或全小写 & \CPP{} 带类型编码 & C++ 核外面加带 \CL{} 链接规范的扁平接口（\Tp{extern\textquotedbl C\textquotedbl}），禁用重载与命名空间 \\
传参方式 & 默认按地址（隐式引用） & 默认按值 & 接口里全部写指针，标量也传地址 \\
数组存储 & 列主序 & 行主序 & 只传首元素指针，用维度参数换算下标，别在两侧都按自然序索引 \\
实数种类 & 默认双精度 & \Tp{double} & 统一 \Tp{double}，禁用长双精度；整型固定 32/64 位 \\
字符与派生类型 & 字符串带隐藏长度；派生类型布局依编译器 & \Tp{char*}；\Tp{struct} 布局可指定 & 边界上只走裸数组，字符串转成定长字符数组或干脆不传 \\
\bottomrule
\end{tabular}
\end{table}

做法是把材料核写成纯粹的 \CPP{}（不依赖求解器、不知道自己在哪个积分点），
外面压一层 \CL{} 扁平接口，再让 Fortran 壳按官方签名调用它。

\begin{lstlisting}[style=cppstyle, caption={材料核（\CPP{}）加 \CL{} 扁平接口：
被 Fortran 壳调用的那一侧}, label={lst:ch05-cpp-kernel}]
// kernel.cpp —— 材料核：输入应变增量，输出应力、状态与切线
// 约定：所有实参都是地址；ddsdde 是 Fortran 列主序的 NTENS*NTENS
#include <cstddef>
#include <cmath>

namespace mat {

struct Core {
  double E, nu, sy, H;             // 由 PROPS 前几项填入
  // 只有一个历史量（等效塑性应变），便于检查序列化与恢复
  static double qmetric(const double* s, int n) {
    double j2 = 0.0;
    for (int i = 0; i < n; ++i) j2 += s[i] * s[i];
    return std::sqrt(1.5 * j2);          // 等效应力的粗化版
  }

  // 应力在传入值上累加；ddsdde 就地覆写为一致切线
  void step(double* stress, double* statev, double* ddsdde,
            const double* dstran, int n, double dt) {
    const double lam = E * nu / ((1.0 + nu) * (1.0 - 2.0 * nu));
    const double mu  = 0.5 * E / (1.0 + nu);
    for (int i = 0; i < n; ++i)
      for (int j = 0; j < n; ++j)
        ddsdde[static_cast<std::size_t>(j) * n + i] =   // (i,j)
          (i == j ? lam + 2.0 * mu * (i < 3 ? 1.0 : 0.5) : lam);
    for (int i = 0; i < n; ++i) {                       // 试弹性
      double acc = 0.0;
      for (int j = 0; j < n; ++j)
        acc += ddsdde[static_cast<std::size_t>(j) * n + i] * dstran[j];
      stress[i] += acc;
    }
    const double q = qmetric(stress, n);                // 屈服检查
    if (q > sy && H > 0.0) statev[0] += (q - sy) / H;
    (void)dt;
  }
};

}  // namespace mat

extern "C" {   // 唯一暴露给 Fortran 的符号，参数全是指针

void mat_step(double* stress, double* statev, double* ddsdde,
              const double* dstran, const double* props,
              const double* dtime, const int* ntens,
              int* ierr) {
  using namespace mat;
  if (stress == nullptr || ddsdde == nullptr) { *ierr = 1; return; }
  Core c{props[0], props[1], props[2], props[3]};
  c.step(stress, statev, ddsdde, dstran, *ntens, *dtime);
  *ierr = 0;      // 非零返回给壳，由壳转成求解器的作业级错误
}

}  // extern "C"
\end{lstlisting}

\begin{lstlisting}[language={[95]Fortran}, style=fortranstyle, caption={Fortran 壳：显式 interface 加
\Tp{BIND(C)} 声明，把官方形参原样转交给 \Tp{mat\_step}}, label={lst:ch05-fortran-shell}]
!      壳的全部价值：把 Fortran ABI 的形参按裸指针转交
      SUBROUTINE UMAT(STRESS,STATEV,DDSDDE,STRANINC,
     1 PROPS,DTIME,NTENS,NSTATV,NOEL,NPT)
      IMPLICIT NONE
!      显式接口：BIND(C) 关掉名字修饰，NAME 指定链接名
      INTERFACE
        SUBROUTINE MAT_STEP(STRESS,STATEV,DDSDDE,DSTRAN,
     1       PROPS,DTIME,NTENS,IERR) BIND(C, NAME='mat_step')
          IMPORT
          DOUBLE PRECISION STRESS(*), STATEV(*), DDSDDE(*)
          DOUBLE PRECISION DSTRAN(*), PROPS(*), DTIME
          INTEGER NTENS, IERR
        END SUBROUTINE MAT_STEP
      END INTERFACE
      DOUBLE PRECISION STRESS(*), STATEV(*), DDSDDE(*)
      DOUBLE PRECISION STRANINC(*), PROPS(*), DTIME
      INTEGER NTENS, NSTATV, NOEL, NPT, IERR
      CALL MAT_STEP(STRESS,STATEV,DDSDDE,STRANINC,
     1              PROPS,DTIME,NTENS,IERR)
      IF (IERR .NE. 0) THEN
        WRITE(6,*) 'UMAT 核失败, 单元', NOEL, '点', NPT
!        以官方提供的中止手段结束作业，不要 return 一个半成品应力
      END IF
      RETURN
      END SUBROUTINE UMAT
\end{lstlisting}

工程上还有两条经验。第一，\textbf{不要试图让求解器直接调 \CPP{} 成员函数}：
修饰名对不上，而且 \Tp{this} 指针会多占一个实参槽位，症状是``看起来跑起来了
但数值全错''。第二，混编的运行时库风险不对称：如果求解器与你的库用了不同的
编译器大版本，堆、\Tp{errno}、OpenMP 运行时可能各自一份——\Lib{LS-DYNA} 那条
``重链可执行文件''的路径至少把编译链钉死成一套，反而更少这类问题；
\Lib{Abaqus} 的动态加载则要显式对齐工具链。工具链认证清单以官方兼容性说明为准。

\section{状态数组与并行：每积分点私有意味着什么}
\label{sec:ch05-statev-parallel}

历史变量（\Tp{STATEV} 及 \Lib{LS-DYNA} 侧的等价存储）的语义有三条，
每一条都会在你踩坑后显得理所当然。

\textbf{一、尺寸是声明出来的，不是算出来的}。必须在材料卡上显式声明状态变量个数，
声明少了是越界写（症状千奇百怪），声明多了只是浪费内存与结果文件体积。
尺寸应在整个作业生命周期固定，中途改变历史变量含义会让结果无法与旧算例比对。

\textbf{二、私有性}。每个积分点一份，求解器保证不会让两个积分点共享同一段
\Tp{STATEV} 存储；但也\textbf{绝不}保证你在堆上、在模块变量里、在静态数组里的东西
是私有的。子程序应当是无状态纯函数：输入是形参，输出是形参。

\textbf{三、并行下的可变态是禁区的根因}。

\begin{table}[htbp]
\centering
\caption{三种并行形态对用户子程序状态的要求}
\label{tab:ch05-parallel}
\begin{tabular}{>{\raggedright\arraybackslash}p{0.18\textwidth}>{\raggedright\arraybackslash}p{0.24\textwidth}>{\raggedright\arraybackslash}p{0.36\textwidth}}
\toprule
并行形态 & 积分点归属 & 对可变态的约束 \\
\midrule
串行 & 顺序遍历，顺序可依赖 & 仍然别写模块级可变态，否则换并行立刻失配 \\
线程共享内存 & 同一时间窗内多核并发调同一入口 & 静态变量、共享缓存、\Tp{save} 的属性都必须加锁或干脆去掉；锁会让标量本构比原实现更慢 \\
区域分解（MPI） & 每个 rank 只看到本地子域，边界积分点可能易主 & 全局聚合量（总损伤、累积用量）不能靠子程序内部累加，必须走求解器的输出通道 \\
\bottomrule
\end{tabular}
\end{table}

还有一个更隐蔽的：同一台机器上，线程数变化会改变调用交错顺序。如果你的本构
依赖``上一次调用的应力''来差分（比如为了做粘性的隐式积分），而 \Tp{STRESS}
已经被求解器更新过，那么单线程正确、多线程结果随核数漂移。正确做法是把
需要跨步的任何量放进 \Tp{STATEV}，让求解器替你搬运。

检查点/续算是同一约束的另一个方向：求解器重启时会把每积分点的历史变量从结果文件
恢复，再按当时的并行拓扑重新分配积分点。若你在 \CPP{} 侧维护了自己的查找表
（例如按 \Tp{NOEL} 索引的缓存），重启后拓扑变了、\Tp{NOEL} 与积分点的对应关系变了，
缓存就指错了地方；同时 \CPP{} 侧的表也不会被写进结果文件，重启后是空的。
结论：\textbf{能被重启恢复的唯一状态，是你交给求解器的状态数组}。

\section{调试与验证：先让一个单元对得上解析解}
\label{sec:ch05-verify}

用户子程序的调试没有断点友好环境（求解器是黑盒进程、积分点调用次数以百万计），
可行的流程是四步，按顺序做，不要跳。

\textbf{单单元回归}。建一个只有一两个积分点的模型（一个八节点单元、
位移或应变控制边界），把本构在解析上能写成闭式（closed form）的量算出来，与子程序输出逐点比对。
这一步能抓住 ABI 错位、剪应变因子 2、主序换算这三类最贵的 bug。

\textbf{能量平衡}。把求解器输出的总内能、\Tp{SSE}/{\Bk}\Tp{SPD} 累计耗散、外力功
放在一起画曲线。显式算例还要看沙漏能（hourglass energy）与人工体积粘性能。
能量对不上时，先怀疑状态变量的演化而不是本构公式。

\textbf{应变路径对照}。单轴拉伸、纯剪切、双轴等比、加载-卸载-反向加载四条路径，
每条都应与解析解或独立实现的标量积分器（\PY{} 写一个 20 行的
返回映射，作为``参考实现''）一致。把这条做成脚本化的回归，
每次改本构就跑，比任何调试器都值。

\textbf{并行一致性}。同一算例分别用串行、2 线程、8 线程、2 与 8 个 MPI rank 跑，
比较关键测点与总能量到打印精度。任何随核数变化的结果都是状态泄漏。

\begin{warning}[title={把历史变量当缓存用，结果随核数变化}]
最常见的事故形态是这样的：为了让``损伤演化更平滑''，在模块里放了一个全局
数组或静态变量当缓存，跨积分点、跨增量共享一些中间量；或者把 \Tp{STATEV} 的
某几位挪用成``这一帧的临时数据''（下一步会重写）。单线程串行跑得很好，
上多核之后结果开始依赖核数，区域分解还让边界附近的单元先错，
最后表现为损伤带出现在错误的地方——因为\textbf{错误的顺序}就是错误的输入。
这类问题在提交前几乎不可能被肉眼发现，只能靠上面的并行一致性回归。
规矩只有两条：状态一律走求解器给的历史数组；入口函数不保留任何跨调用可变态。
若确实需要跨步历史信息（例如为应变率做记忆），把它编码进历史变量的具体槽位，
并在文档里写死每个槽位的含义与单位，避免后续版本挪用时静默错位。
\end{warning}

\begin{note}[title={作业失败时先看第一条错误而不是最后一条}]
子程序里 \Tp{WRITE(6,*)} 的输出散落在求解器日志中，百万积分点会把它刷爆。
约定：只在 \Tp{NOEL}/{\Bk}\Tp{NPT} 命中某个环境变量或硬编码值时打印；
错误路径上先把状态写成一个可识别的哨兵值（例如把应力置为 0 并返回错误码），
再让求解器中止，避免留下``部分更新''的应力与状态。日志文件的具体名称
与错误码约定以官方文档为准。
\end{note}

% ============================================================
% ---- frag_ch06.tex ----
\chapter{OpenFOAM：C++ 求解器、场代数与自定义边界条件}
\label{ch:ch06-openfoam}

与前一章正相反：\Lib{Abaqus}/{\Bk}\Lib{LS-DYNA} 的用户扩展是一段被求解器调用的 Fortran，
而 \Lib{OpenFOAM} 的扩展\textbf{就是直接写 \CPP{}}——没有 ABI 边界、没有符号修饰，
也没有``官方只给挂钩点''这层隔离。这一章用它来说明另一端的取舍：当你把几何、拓扑、场
这套抽象完全暴露给使用者时，换来的是任意改方程的自由度，付出的是编译链、模板报错、
运行期字典这三笔持续的利息。

类名口径先钉住：\Tp{geometric{\Bk}Field}、\Tp{vol{\Bk}Scalar{\Bk}Field}、\Tp{dimensioned{\Bk}Scalar}、
\Tp{fv{\Bk}Vector{\Bk}Matrix}、\Tp{mixed{\Bk}Fv{\Bk}Patch{\Bk}Field} 这一族名称与角色是稳定的；
模板参数个数、构造函数的具体形参清单、注册宏的名字随发行版变化，
下文按调用形状描述，以你所用版本的官方文档与头文件为准。

\section{case 与 code 的分界}
\label{sec:ch06-case-vs-code}

一个算例（case）是纯数据：网格、初边值、离散与控制参数。代码侧是求解器可执行文件
与库。分界如下表。

\begin{table}[htbp]
\centering
\caption{case 侧与 code 侧各管什么}
\label{tab:ch06-case-code}
\begin{tabular}{>{\raggedright\arraybackslash}p{0.14\textwidth}>{\raggedright\arraybackslash}p{0.33\textwidth}>{\raggedright\arraybackslash}p{0.33\textwidth}}
\toprule
& case 侧 & code 侧 \\
\midrule
内容 & 网格（\Tp{constant/polyMesh}）、场初值与边界类型（\Tp{0/U}、\Tp{0/p}）、字典（\Tp{system} 与 \Tp{constant}） & 求解器 \Tp{main} 循环、离散算子、物性模型、边界条件类 \\
改动成本 & 改文本、直接重跑 & 编译、链接、再重跑 \\
出错时机 & 运行期（字典解析、量纲检查） & 编译期为主，加运行期注册 \\
\bottomrule
\end{tabular}
\end{table}

最小 case 的形状与运行方式如下（路径与命令是惯例形式，字段名以官方模板为准）。

\begin{lstlisting}[style=shellstyle, caption={最小 case 的目录形状与运行命令形状},
label={lst:ch06-case-tree}]
# 一个不可压缩稳态算例的最小目录
simpleCavity/
  0/
    U                        # 速度：internalField 加各 patch 的边界类型
    p                        # 压力：量纲写在文件头，边界类型逐 patch 声明
  constant/
    polyMesh/                # blockMesh 生成，或 gmsh/snappy 转进来
    transportProperties      # 物性，例如 nu
    turbulenceProperties      # laminar 或 RAS/LES 模型选择
  system/
    controlDict              # 起止时间、步长、写间隔、libs 与 functions
    blockMeshDict            # 几何（用 blockMesh 时）
    fvSchemes                # 每个方程每项用哪个离散格式
    fvSolution               # 每个方程用哪个线性/非线性求解器
    decomposeParDict         # 区域分解参数

# 运行形状：先网格，再求解，最后并行与合并
blockMesh
simpleFoam
decomposePar
mpirun -np 4 simpleFoam -parallel > log.simpleFoam 2>&1
reconstructPar -latestTime
# 自定义库的加载声明写在 system/controlDict 里
# libs ("libmyBC.so");   然后 0/U 的某个 patch 才能写 type myBC
\end{lstlisting}

``改一个方程=写一个求解器''这句话是这套架构的直接后果：方程的装配语句
（哪个项、用 \Tp{fvm} 还是 \Tp{fvc}、怎么松弛、怎么耦合）就写在求解器
\Tp{main} 的 \Tp{while} 循环里，字典只能改系数与格式，不能改方程结构。
所以官方给出的不是``一个可配置的全能求解器''，而是几十个工作循环各异的求解器
（密度基、可压、多相、共轭传热……），用户复制一个最接近的、改那几行装配代码、
起自己的名字。这条路看起来笨，但它保证了每个求解器的循环是显式可读的，
代价是求解器之间的代码漂移（同一件事在十个求解器里有十种写法）。

\section{场、网格与单位：抽象层到底给了你什么}
\label{sec:ch06-fields}

场的类型名基本是``位置 + 拓扑 + 张量阶''的组合：单元中心上的标量场
\Tp{vol{\Bk}Scalar{\Bk}Field}、单元中心上的矢量场 \Tp{vol{\Bk}Vector{\Bk}Field}、面上的矢量场
\Tp{surface{\Bk}Vector{\Bk}Field}。它们都是 \Tp{geometric{\Bk}Field} 家族的特化：一个内部场
（\Tp{internal{\Bk}Field}，按单元/面编号连续存放）加一组边界场
（\Tp{boundary{\Bk}Field}，每个 patch 一个子数组），再挂一份网格引用与一份量纲。
有三条性质对代码可见性影响最大：

\textbf{一、并行分布是透明的，但索引不是。}每个 rank 只持有自己的那段子域网格，
\Tp{mesh.nCells()} 是本地单元数；内部场的下标是本地编号，跨 rank 的邻居
以 ghost 单元（\Tp{processor} 类型的 patch）出现。所以任何``按全局编号思考''
的写法（例如用单元号做哈希键、按 \Tp{cell} 顺序累加统计量）在串行下正确、
并行下错，而且不会报错。全局聚合要走 \Tp{gSum}、\Tp{reduce} 这一类通信封装。

\textbf{二、量纲是运行期检查的。}\Tp{dimensionedScalar} 带一个 \Tp{dimensionSet}
（七个基本量：质量、长度、时间、温度、物质量、电流、发光强度），
加减与赋值时比对量纲，不匹配就抛运行期异常。这让``给速度场加了个压力的量''
这类错误必炸，但也意味着单位不写在类型里、而写在字典里——字典里
把 \Tp{nu} 从 \Tp{m2/s} 写成别的，代码不会拦你。

\textbf{三、场的读写是带注册表的。}场构造时向网格/场存储登记，析构时注销，
写结果文件由 \Tp{write()} 触发；这就是为什么 \Tp{controlDict} 里
\Tp{writeControl} 与 \Tp{writeInterval} 改一下就能改输出密度，
也是为什么临时场不该在每步构造（会污染输出列表并带来分配开销）。

\begin{lstlisting}[style=cppstyle, caption={构造场、算源项、写回：注释里点明索引语义},
label={lst:ch06-field-code}]
// createFields.H —— 求解器每步之前的一次性构造
volScalarField p
(
    IOobject::MUST_READ,   // 从 0/p 读初值与边界类型
    "p",
    mesh
);
volVectorField U(IOobject::MUST_READ, "U", mesh);

// 物性：量纲 (0 2 -1 0 0 0 0) 即 m2/s，运行期比对
dimensionedScalar nu("nu", dimensionSet(0, 2, -1, 0, 0, 0, 0), 1.0e-6);

// 源项：按单元中心索引。cellI 是"本地"编号，
// 并行下只覆盖本 rank 的 owned 单元，不含 ghost
dimensionedScalar beta
(
    "beta", dimensionSet(0, 0, -1, 0, 0, 0, 0), 1.0   // 阻尼系数，量纲 1/s
);
forAll(U, cellI)
{
    U.internalField()[cellI] -= beta.value() * U[cellI];
}

// 面上的通量场：phi 建在 Faces 上，索引语义是面号，不是面所属单元
surfaceScalarField phi
(
    IOobject::NO_READ,     // 不读文件，由 U 插值出来
    IOobject::NO_WRITE,
    fvc::interpolate(U) & mesh.Sf()   // Sf 是有向面矢量（面积x法向）
);

Info<< "kinetic energy = "
    << magSqr(fvc::average(U)).max() << nl;   // 归约后才能打印标量
p.write();                   // 显式写出一次（受 writeControl 影响）
\end{lstlisting}

其中 \Tp{fvc::} 与 \Tp{fvm::} 的区别下一节展开：前者把表达式直接算成场，
后者装配成矩阵。清单里 \Tp{internalField}、\Tp{faceCells} 等访问器名与
返回是否 \Tp{tmp} 包装，都以所用版本的头文件为准。

\section{离散、方程装配与自定义边界条件}
\label{sec:ch06-equation}

\Tp{fvm::laplacian(nu, U)} 这类写法做的是``把算子变成矩阵''：它返回一个
\Tp{fv{\Bk}Scalar{\Bk}Matrix} 或 \Tp{fv{\Bk}Vector{\Bk}Matrix}，里面是稀疏系数（对角与非对角）
加右端源项。与之配对的是 \Tp{fvc::}——比如压力梯度 \Tp{-fvc::grad(p)}
在动量方程里作为显式右端项出现。用哪个不只是风格问题：
放进 \Tp{fvm} 意味着进入隐式矩阵（更稳、可松弛），放进 \Tp{fvc}
意味着显式（便宜、受步长限制）。

边界条件进入方程的位置在矩阵装配的最后一步：矩阵对每个 patch 询问其
\Tp{fvPatchField}，由该补丁场的类型决定如何修改系数——
\Tp{fixedValue} 把该面系数写成极大对角并置右端项为目标值，
\Tp{zeroGradient} 把面值取为相邻内部单元值，对称面与周期面各自有不同处理。
因此自定义边界条件不是回调，而是一个真正的类，它同时拥有
``面值的更新规则''和``对矩阵的贡献规则''。

\begin{table}[htbp]
\centering
\caption{自定义边界条件的四个必备点}
\label{tab:ch06-bc-four}
\begin{tabular}{>{\raggedright\arraybackslash}p{0.17\textwidth}>{\raggedright\arraybackslash}p{0.24\textwidth}>{\raggedright\arraybackslash}p{0.39\textwidth}}
\toprule
点 & 常见基类选择 & 要做什么 / 漏掉的后果 \\
\midrule
构造与析构 & \Tp{mixedFvPatchField} 族（值与梯度都可插手的混合类） & 四类构造：从字典、从补丁加内部场、拷贝、克隆；缺一个就在某种重建/映射路径上崩 \\
\Tp{updateCoeffs} & 同上 & 每个时间步把系数（如目标值、传热系数）算新一遍；漏写则边界用旧值，稳得可疑 \\
\Tp{evaluate} & 同上或 \Tp{fixedValue} 派生的补丁场 & 显式给出面值；非耦合求解里必须能被调用到，否则矩阵拿到未初始化面场 \\
读写与注册 & 同上 & 定义字典可读的参数与 \Tp{read}/{\Bk}\Tp{write} 流操作符，并注册进运行时选择表；否则字典里 \Tp{type} 写不出这个类 \\
\bottomrule
\end{tabular}
\end{table}

\begin{lstlisting}[style=cppstyle, caption={一个混合类边界条件：外壳与
\Tp{updateCoeffs} 的核心（注册宏名以版本为准）}, label={lst:ch06-bc-code}]
// pDrivenFvPatchVectorField.H —— 混合类边界条件的骨架
class pDrivenFvPatchVectorField
:   public mixedFvPatchField<vector>
{
    // 字典可读参数：构造时从表里取，read 里再覆盖一次
    dimensionedScalar pExt_;     // 外部压力
    dimensionedScalar kSoft_;    // 标定系数

public:
    // 四个构造：从补丁加字典、从字典加内部场、拷贝、克隆
    // 基类形参清单随版本变化，这里只按调用形状写
    pDrivenFvPatchVectorField(const fvPatch&, const dictionary&);
    pDrivenFvPatchVectorField(const fvPatch&, const dictionary&,
                             const vectorField&, const word& n);

    // 注册：运行时选择表宏把类型名与构造函数绑定，宏名以版本为准
    // 混合类必须回答三件事：追哪个值、追哪个梯度、各占多少权重
    virtual tmp<vectorField> refValue() const;
    virtual tmp<vectorField> refGrad() const;
    virtual tmp<scalarField> valueFraction() const;

    virtual void updateCoeffs();          // 每步刷新系数
    virtual void write(Ostream&) const;   // 写回结果里的字典
};
\end{lstlisting}

\begin{lstlisting}[style=cppstyle, caption={\Tp{updateCoeffs} 的实现形状：
把面速度按外部压力与标定系数拉回目标值}, label={lst:ch06-bc-impl}]
// pDrivenFvPatchVectorField.C —— 系数更新是所有逻辑的落点
void pDrivenFvPatchVectorField::updateCoeffs()
{
    if (this->updated())     // 已更新就返回：一步之内可能被要多次
    {
        return;
    }

    // size() 是本补丁上的面数；faceCells 把面号换成本地单元号
    const vectorField& Uc = internalField();
    const labelList& owner = patch().faceCells();
    vectorField target(this->size(), vector::zero);

    forAll(*this, faceI)
    {
        const vector n = patch().normals()[faceI];
        // 相邻内部单元速度在法向上的投影
        const scalar un = Uc[owner[faceI]] & n;
        const scalar drive = un + kSoft_.value() * pExt_.value();
        target[faceI] = n * drive;
    }

    refValue() = target;
    refGrad() = vectorField(this->size(), vector::zero);
    valueFraction() = tmp<scalarField>(   // 权重 1 表示完全追值
        new scalarField(this->size(), 1.0)
    );

    mixedFvPatchField<vector>::updateCoeffs();  // 置 updated 标志
}
\end{lstlisting}

三个必须留心的形状问题。其一，\Tp{size()} 是补丁上的\textbf{面}数，不是单元数，
和 \Tp{internalField} 混用时靠 \Tp{faceCells} 换算。其二，
\Tp{updated} 标志是幂等保护：一个求解器迭代里 \Tp{updateCoeffs} 与
\Tp{evaluate} 都可能被多次调用，不做检查会把系数连乘两遍。其三，
不可压求解里压力补丁还参与 \Tp{adjustPhi}（维持连续性方程的总量平衡），
自定义速度边界时若不同时想清楚压力边界，残差会表现为整体压力漂移。

\section{编译与并行下的加载}
\label{sec:ch06-build}

\Lib{OpenFOAM} 的构建封装习惯叫 \Tp{wmake}：它把``项目目录变量、依赖库清单、
每个源文件的编译命令''这套东西收进两个文本文件，你只写清单。

\begin{lstlisting}[style=cmakestyle, caption={构建清单形状：\Tp{Make/files} 与
\Tp{Make/options} 两个文本文件}, label={lst:ch06-build}]
# Make/files —— 编译哪些源、产出叫什么
pDrivenFvPatchVectorField.C
LIB = $(FOAM_USER_LIBBIN)/libmyBC

# Make/options —— 头文件搜索路径与要链的库
EXE_INC = \
    -I$(LIB_SRC)/finiteVolume/lnInclude \
    -I$(LIB_SRC)/meshTools/lnInclude \
    -I$(LIB_SRC)/transportModels/incompressible/singlePhase
EXE_LIBS = \
    -lfiniteVolume \
    -lmeshTools

# 求解器侧把库加进去，再重编求解器
EXE_INC  += -I$(FOAM_USER_LIBBIN)
EXE_LIBS += -lmyBC

# 常用命令形状
wclean ; wmake            # 编库；libso 变体用于动态库
wmake -debug              # 带调试符号的编译
# 项目根与用户库目录来自环境：WM_PROJECT_DIR, FOAM_USER_LIBBIN
\end{lstlisting}

\Tp{lnInclude} 这一层是这套构建的第二件事：它把一个库对外可见的头文件
软链接到一个目录，求解器只需包含那一个路径。好处是包含面窄、增量重编快；
坏处是新增类忘了被收集时，报错是``找不到头文件''而不是``你没有导出它''。

并行（MPI）下的加载语义值得单独讲：\textbf{每个 rank 是独立进程}，
各自按 \Tp{controlDict} 里的 \Tp{libs} 加载同一份动态库。因此库文件路径必须在
所有节点上一致可见（共享文件系统或分发脚本同步），且所有 rank 必须加载\textbf{同一版本}
的 \CPP{} 运行时；只更新登录节点上的库、计算节点仍是旧库，症状是各 rank
构造出的补丁类型不一致，报错点常常离真正原因很远。

与 \Tp{cmake} 生态的摩擦点有三处：清单格式不通用（\Tp{Make/files} 不是
\Tp{CMakeLists.txt}）、IDE 拿不到 \Tp{compile\_commands} 数据库因而跳转与
补全半残、以及依赖表达方式是``库名加相对路径''而非目标（target）。
可行的折中是：保持用 \Tp{wmake} 真编译，另外手工维护一份
\Tp{CMakeLists.txt} 只做索引与静态分析用（不参与产出）。
较新的若干发行版已把顶层构建向 \Tp{cmake} 迁移，用户库的惯用路径仍以
\Tp{wmake} 为主，具体以所用发行版的文档为准。

\section{代价：模板、字典与调试}
\label{sec:ch06-cost}

把抽象直接暴露给 \CPP{} 使用者，收上来的账单有四张。

\textbf{编译时间}。场类型是深层模板嵌套（几何场叠场叠并行派发），
一个新算例文件的编译时间主要由模板实例化决定；把 \Tp{.C} 拆细、
显式限制实例化的类型集合是唯一稳妥的缓解。

\textbf{报错可读性}。表达式模板（\Tp{fvc::} 那一串嵌套表达式）把整条算式
压成一个惰性求值对象，好处是不产生临时场，坏处是一个量纲不匹配会给出
跨十几个模板层的错误文本。工程做法是把长式子拆成带类型的中间量
（\Tp{tmp<fvVectorMatrix>} 一类），宁可多一次求值，换来错误定位在行号上。

\textbf{字典与类型的耦合}。类型名在字典里是字符串，编译器不参与检查：
拼错 \Tp{type}、少写一个必填系数，都是运行期才炸。缓解顺序是先用官方样例的
字典片段做底、只改数值，再逐步改类型；可复用的边界条件参数要写进文档，
并在 \Tp{read} 里对缺省值给警告。

\begin{table}[htbp]
\centering
\caption{五个扩展层次与各自的门槛}
\label{tab:ch06-layers}
\begin{tabular}{>{\raggedright\arraybackslash}p{0.18\textwidth}>{\raggedright\arraybackslash}p{0.20\textwidth}>{\raggedright\arraybackslash}p{0.42\textwidth}}
\toprule
层次 & 要不要编 \CPP{} & 适用判断 \\
\midrule
改字典（\Tp{fvSchemes}、\Tp{fvSolution}） & 不要 & 换离散格式、换线性求解器、调松弛；这是最便宜的一层，先试 \\
\Tp{functions}、\Tp{fvOptions} 条目 & 通常不要 & 加源项、加监测量、加后处理；改的是``问题''不是``方程'' \\
自定义边界条件库 & 要，一个类 & 边界规律变了（壁面函数、入流剖面、共轭接口）；方程不变 \\
写/改求解器 & 要，一个可执行文件 & 方程结构变了（新增 transported 变量、改耦合顺序）；这才是``改方程''的正当位置 \\
改离散/场核心库 & 要，且要懂模板 & 新算子、新网格类型；影响所有下游算例，成本最高 \\
\bottomrule
\end{tabular}
\end{table}

\begin{warning}[title={边界条件写了但字典不生效，或旧系数被静默复用}]
几个高频形态：把 \Tp{updateCoeffs} 写成每步都算、却不检查 \Tp{updated} 标志，
于是同一系数在一个迭代里被叠加多次；只实现了拷贝构造忘了克隆入口，
串行正常、一旦网格重划分或读 restart 就报找不到类型；
在 \Tp{evaluate} 之外直接改面场却没有置 \Tp{updated}，矩阵拿到的还是上一时刻的值——
结果是残差照常下降、物理量单调不对，最难查。另外，补丁场的参数字典少写一项时，
若构造函数给了默认值就不会有任何提示，所以必填参数应该在解析失败时主动报错。
\end{warning}

\begin{bestpractice}[title={什么时候改方程，什么时候换求解器}]
判断线只有一条：\textbf{需要改的是方程结构，还是求解顺序}。
若是新增一个输运方程、把耦合改成全耦合、把时间离散改成二阶隐式，
这些是``改方程''，可以在现有求解器上动那几行装配语句，
但要把残差定义与线性求解器设置一起改；若是耦合框架本身不对
（不可压 SIMPLE 去算强可压燃烧），那不是改方程，是选错求解器，
应该换一个工作循环本来就正确的求解器，再往上加物理。
中间地带（多一个源项、多一个监测量、边界规律变了）请留在
\Tp{fvOptions} 与边界条件这两层，别动求解器——
你从 \Tp{simpleFoam} 复制出的那份 \Tp{myFoam}，一年后就是你的维护负担：
上游修了收敛、改了默认松弛，而你的分支不会自动收到。
\end{bestpractice}

% ============================================================
% ---- frag_ch07.tex ----
\chapter{其余扩展面：COMSOL 与 Simulink 的 S-Function，以及开源求解器}
\label{ch:ch07-other-extensions}

前几章讲的都是某一家求解器怎么接你的代码。本章把镜头拉远：\DomCAE{} 侧的扩展面其实只有四种基本形状，差别集中在\textbf{语言}、\textbf{编译时机}、\textbf{运行时机}这三条轴上。把三轴想清楚，换一个平台时就不用重学一遍世界观，只需要重查一遍入口函数的参数表。本章依次看模型树驱动的多物理场平台（以 COMSOL 一类为参照）、控制侧的 Simulink S-Function、开源求解器特有的``改源码还是写插件''这条分岔，最后给一份选平台时的提问清单。具体 API 名只写有把握的部分，其余一律描述调用形状并注明以官方文档为准。

% ------------------------------------------------------------------
\section{四种扩展入口的谱系}
\label{sec:ch07-entry-lineage}

按``代码什么时候变成机器码''把见过的入口排一下，得到一条谱系：

\begin{itemize}
\item \textbf{纯脚本宏。}语言是产品自带的脚本方言或宿主脚本（\PY、Java 形态的对象接口、专有 Basic 一类）。没有独立编译步骤，运行前解释或即时编译。能看到的东西只有平台主动暴露的对象树：模型、节点、参数、结果算子；看不到单元内部的数组。
\item \textbf{表达式级。}你在 GUI 的属性框里写一个数学式，平台把它解析成一棵求值树。编译时机很特别：既不是保存时也不是运行期，而是\emph{模型装配时}——有限元组装矩阵的那一刻，这棵树被嵌进每个积分点的循环里。语言只支持函数调用与算符，不支持分支与副作用。
\item \textbf{编译型用户函数。}你在外部用 \CL{}、\CPP{} 或 Fortran 写一段代码，编成动态库；求解器启动或首次用到时加载它，之后在求解循环里反复调用。这一档第一次拿到原生指针和原生调试器。
\item \textbf{源码级。}你自己编译整个求解器。内部量全可见，代价是每次改动都要过一遍全量构建，并从此背上一个需要持续合并的分支。
\end{itemize}

三条轴经常被混淆，值得单独点破。\textbf{编译时机}决定反馈周期：表达式级写错通常在解析阶段就报，且能指出是哪个属性框；编译型写错，报错落在链接器或运行期的符号解析；源码级写错，报错在编译期但要重建整个求解器。\textbf{运行时机}决定性能预算：同一个``给边界指定通量''的需求，走表达式每个装配步求值一次；走用户函数则可能被残差评估调用几十次（非线性迭代乘上线性迭代），再乘并行分块就是每秒万次量级。在函数体里放文件 I/O、内存分配或打印，单次调用看不出，规模化跑就知道。\textbf{可见范围}决定你能做什么：表达式只能读当前求值点的自由变量；用户函数能读平台传进来的数组，但看不到邻居单元，除非平台把梯度或重构后的量一并给你；源码级什么都能看到。

\begin{table}[htbp]
\centering
{\small
\begin{tabular}{%
 >{\raggedright\arraybackslash}p{0.12\textwidth}
 >{\raggedright\arraybackslash}p{0.11\textwidth}
 >{\raggedright\arraybackslash}p{0.18\textwidth}
 >{\raggedright\arraybackslash}p{0.21\textwidth}
 >{\raggedright\arraybackslash}p{0.21\textwidth}}
\toprule
入口 & 语言 & 编译与加载时机 & 能否访问内部量 & 调试手段 \\
\midrule
脚本宏 & 脚本方言、\PY & 无独立编译，运行前解释 & 只能到暴露的对象与参数 & 脚本内打印、宿主调试器 \\
表达式级 & 内置式语言 & 模型装配时编成求值树 & 当前点的自由度、坐标、时间 & 先算标量再画场、解析报错定位 \\
编译型用户函数 & \CL、\CPP、Fortran & 外部编成库，运行时加载 & 由钩子给的索引与数组决定 & 原生调试器、单元测试、条件断点 \\
源码级 & 与求解器同栈 & 全量构建，随求解器一起发布 & 全部 & 完整调试器、剖析工具、回归测试 \\
\bottomrule
\end{tabular}
}
\caption{四种扩展入口在语言、编译时机、内部量可见性与调试手段上的对照}
\label{tab:ch07-lineage}
\end{table}

表~\ref{tab:ch07-lineage} 里最有用的一列是``能否访问内部量''。配套的经验规律是：\emph{访问深度每加一层，可移植性掉一半}。表达式级写的内容换版本基本还能用；用户函数换大版本要重编、偶尔要改签名；源码级分支在两次升级之间可能就废掉。所以选档位的正确姿势不是``哪个强''，而是``哪一档的可见范围刚好够用''——刚好够用的那一档，升级成本与调试难度都最低。反过来，把一行本可以写成表达式的逻辑塞进用户函数，换来的是编译链路、ABI 依赖和一条新的失败模式。

% ------------------------------------------------------------------
\section{多物理场平台：模型树、表达式与外部函数}
\label{sec:ch07-comsol-platform}

以 COMSOL 一类多物理场平台为参照（下面的形状描述适用于同类``模型树驱动''的平台，产品名与节点拼写以官方文档为准）。这类平台的中心对象是\textbf{模型树}：几何、材料、物理场接口、网格、研究步骤、结果，每层节点都带一批可填表达式的属性框。

\subsection*{模型树与参数化}

参数（Parameters 一类节点）是这类平台最便宜的扩展入口：把几何尺寸、物性、时间尺度写成命名参数，研究步骤里做扫描，模型就能像函数一样被调用。参数化到一定规模后你会想要外部驱动——平台通常提供一个从宿主脚本构造模型、设置参数、提交求解、取回结果的编程接口（Java 形态的对象树，另有从 \PY 或 MATLAB 一侧驱动的接口；早年的 LiveLink 系列与现在的编程接口是同一件事的不同名字）。工程上的判断标准很硬：\textbf{凡是第二次跑要改三个数的模型，都该有一条脚本化的构造路径}，否则算例一多，没人记得每个差在哪。

\subsection*{内置表达式能到哪里}

表达式级入口能覆盖的场景比直觉大得多：非线性物性、随位置变化的源项、依赖其他场的耦合项、方向相关的边界条件，都能写成式子。此外通常还能声明\emph{自定义算子}（把一个积分、求平均或滤波操作起个名字，之后像函数一样在别的表达式里调用），以及把附加的常微分或代数方程挂在几何点、边或域上——这正是``用状态变量做控制器或集总模型''的入口。内置函数库覆盖常见数学、单位换算与特殊函数；单位系统是这一档的隐藏约束：写进属性框的量要能化成正确的 SI 组合，否则求解器不会报错，只会给出一条量纲自洽但物理错误的曲线。表达式级的天花板有三块：

\begin{itemize}
\item \textbf{无过程。}循环、查表后分支、带状态的迭代只能靠已有算子拼；拼不出来就得换入口。
\item \textbf{无历史。}表达式在当前时间步被求值，看不到上一时刻的自定义中间量，除非你把它登记成状态变量或附加的微分代数方程。
\item \textbf{无外部数据流。}要读自己的求解结果、调用外部程序、走网络，表达式做不到。
\end{itemize}

\subsection*{需要 \CL{} 与 \CPP{} 时走哪条路}

常见两条，都属``编译型用户函数''档。第一条是\textbf{外部函数}（External function 一类节点）：你提供库文件路径与函数名，平台在需要点值时调用它。调用形状大致是实参为一批求值点坐标加已有场值（打包成数组），返回一个或多个值。不同版本在这里差别明显——有的直接按符号名解析 C 导出函数，有的要求中间加一层宿主语言包装（历史上有走 Java 与 JNI 的做法）。这类入口一般要求按 C 链接导出，且不允许语言级异常穿越边界。第二条是\textbf{实时编译}形态：在 GUI 里贴一段代码，平台调用本地工具链即时编成临时库。好处是反馈快、便于试验；代价是编译依赖必须装在计算节点上，批处理与集群环境下要确认工具链路径与许可。应用侧还有一条：把模型连同方法脚本编成一个可分发的应用（Application Builder 一类的形态），方法代码是宿主语言形态、编译发生在打包应用时，属于``脚本档但可交付''。

\begin{lstlisting}[language=C++, style=cppstyle,
  caption={外部编译型用户函数：通用的调用形状（签名由平台决定）},
  label={lst:ch07-extfun-shape}]
/* 形状示例：入口是导出符号，实参是打包好的数组。
   确切参数顺序与列打包方式以官方文档为准。 */
#ifdef __cplusplus
extern "C" {
#endif

/* 求值入口：n 个积分点，给坐标与已有场值，写回同名长度的 out */
void flux_at_face(const double *xyz, const double *var,
                  int n, double *out)
{
    for (int i = 0; i < n; ++i) {
        const double *p = xyz + 3 * i;  /* 坐标，勿依赖全局网格 */
        const double  t = var[2 * i];   /* 温度，列顺序由平台约定 */
        out[i] = k0 * (1.0 + a * (t - t0)) * p[2];
    }
}

/* 若入口允许交导数：dout 按同样的打包顺序给出，形如
   dout[i * nvar + j] = 该点对第 j 个变量的偏导。
   不给则平台改用差分扰动，成本与风险见正文。 */

#ifdef __cplusplus
}
#endif
\end{lstlisting}

\begin{note}[title={外部函数必须回答``导数从哪来''}]
隐式求解要的是残差对自由变量的雅可比。你交一个点值函数上去，平台只有两种选择：接受你再交一个求导函数（部分入口允许），或者对它做差分扰动。差分意味着这个点上每步多算若干次；如果扰动跨过了你库里的分支或不连续点，线性化就是错的，表现为残差震荡或线搜索失败。写分段本构时，优先把折点做成连续可导的形式（软转折、光滑逼近），比调求解器容差有效得多。
\end{note}

\subsection*{耦合场数量对入口选择的约束}

物理场越多，表达式级越占便宜：加一个耦合项就是在已有属性框里多写一项，不需要新编译单元，也不引入新的失败模式。反过来，用户函数在耦合场景里成本上升得很快——参数个数与场顺序变了要改签名；每个入口被调用的次数随方程数线性上升；雅可比缺失的代价从``一个场多算几次''变成``整块耦合矩阵多算很多次''。由此得到一张可操作的对照：

\begin{table}[htbp]
\centering
{\small
\begin{tabular}{%
 >{\raggedright\arraybackslash}p{0.20\textwidth}
 >{\raggedright\arraybackslash}p{0.17\textwidth}
 >{\raggedright\arraybackslash}p{0.23\textwidth}
 >{\raggedright\arraybackslash}p{0.21\textwidth}}
\toprule
需求形状 & 首选入口 & 该档的失败信号 & 换档依据 \\
\midrule
物性随温度或坐标变化 & 表达式 & 式子长到没人敢改 & 拆成命名参数与中间变量 \\
需要查表加插值、含历史 & 自定义算子、状态变量 & 拼不出过程 & 外部函数，把它做成纯函数 \\
边界值来自外部程序或测量 & 外部函数 & 每步阻塞在 I/O & 换成预取的插值表 \\
新的场、新的离散或求解策略 & 源码级 & 钩子拿不到内部量 & 这一档本来就在源码里 \\
\bottomrule
\end{tabular}
}
\caption{按需求形状选入口：失败信号出现时就是换档的时刻}
\label{tab:ch07-pick-entry}
\end{table}

表~\ref{tab:ch07-pick-entry} 的第三列比第四列重要。多数项目的做法是：先停在便宜的档位，等失败信号真的出现再升级，而不是``为了灵活''一上来就写编译型函数。

% ------------------------------------------------------------------
\section{控制侧的扩展：Simulink S-Function}
\label{sec:ch07-sfunction}

前面几家都在物理场这一侧，S-Function 是控制侧的对应物：在 \Lib{Simulink} 的求解循环里插一个自定义块。它的价值与代价来自同一件事——\textbf{你和求解器共享同一条时间轴}。

\subsection*{它在求解循环里的位置}

\textbf{定步长}时，块按自己的采样时间被命中，命中即``读输入、算输出、按需推进状态''，顺序由块图上的连通关系与采样率决定，可预期。\textbf{变步长}时情况完全不同：积分器为了收敛会把一步退回去重来，同一个采样点内你的块可能被反复求值。因此块必须满足三条纪律：

\begin{itemize}
\item \textbf{可重入、无副作用。}输出只依赖当前输入与已登记的状态；不要写文件、不要动外部资源、不要依赖上一次被调用的顺序。
\item \textbf{状态交给平台保管。}连续状态与离散状态要通过尺寸声明交出去，求解器才有退回与重置的余地。自己藏在静态变量里的状态，退步时不会被还原，结果随容差与最大步长漂移，且极难复现。
\item \textbf{直接馈通要显式声明。}输出是否依赖同一时刻的输入，决定块能不能处在代数环里。声明错一次，轻则报环，重则静默给出一条错误的时间响应。
\end{itemize}

多速率模型里还要当心速率转换：不同采样率的块之间需要显式的速率过渡，否则平台会拒绝编译模型或在仿真时报采样率冲突；把``上一步的值''跨速率读，语义上等于读到不确定的时刻。

\subsection*{端口、状态与回调的形状}

一个典型的 Level-2 \CL{} MEX S-Function 由这几个回调组成（签名与可选方法宏的确切拼写以官方文档为准，这里只给形状）：\Fn{mdlInitializeSizes} 声明参数个数、端口数与宽度、状态数、采样时间组数与工作向量槽位；\Fn{mdlInitializeSampleTimes} 给每组采样时间与偏移；\Fn{mdlOutputs} 每个求值点都被调；\Fn{mdlUpdate} 只在离散命中时推进状态；\Fn{mdlDerivatives} 提供连续状态导数；\Fn{mdlTerminate} 释放你挂在状态槽里的对象指针。端口数据以指针给出，对话框参数从平台提供的参数对象里取标量。

\begin{lstlisting}[language=C, style=cstyle,
  caption={Level-2 C MEX S-Function 最小骨架：只列回调与调用形状},
  label={lst:ch07-sfun-min}]
#define S_FUNCTION_NAME  gain1_sfun   /* 必须与源文件名一致 */
#define S_FUNCTION_LEVEL 2
#include "simstruc.h"

static void mdlInitializeSizes(SimStruct *S)
{
    ssSetNumSFcnParams(S, 2);             /* 对话框参数个数 */
    if (ssGetNumSFcnParams(S) != ssGetSFcnParamsCount(S)) {
        return;                    /* 不符，平台已报错 */
    }
    ssSetNumContStates(S, 0);
    ssSetNumDiscStates(S, 0);
    ssSetNumSampleTimes(S, 1);
    if (!ssSetNumInputPorts(S, 1)) { return; }
    ssSetInputPortWidth(S, 0, 1);
    ssSetInputPortRequiredContiguous(S, 0, true);
    if (!ssSetNumOutputPorts(S, 1)) { return; }
    ssSetOutputPortWidth(S, 0, 1);
    ssSetOptions(S, SS_OPTION_EXCEPTION_FREE_CODE);
}

static void mdlInitializeSampleTimes(SimStruct *S)
{
    ssSetSampleTime(S, 0, INHERITED_SAMPLE_TIME);
    ssSetOffsetTime(S, 0, SS_NOCMD_OFFSET_TIME);
}

static void mdlOutputs(SimStruct *S, int_T tid)
{
    const real_T *u;
    real_T       *y;
    const real_T  k = mxGetScalar(ssGetSFcnParams(S, 0));
    u = (const real_T *)ssGetInputPortSignal(S, 0);
    y = (real_T *)ssGetOutputPortSignal(S, 0);
    y[0] = k * u[0];              /* 纯计算，不留副作用 */
}

/* 有离散状态时实现 mdlUpdate，有连续状态时实现 mdlDerivatives，
   并声明对应的可选方法宏；状态放进平台登记的工作向量槽位。
   调试用 mdlCheckParameters 校验对话框参数的类型与尺寸。 */

static void mdlTerminate(SimStruct *S) { }

#ifdef MATLAB_MEX_FILE
#include "simulink.c"
#else
#include "cg_sfun.h"
#endif
\end{lstlisting}

\subsection*{\CL{} 写的与解释语言写的差别}

同一需求有三档入口，差别全在编译时机与可交付性：

\begin{itemize}
\item \textbf{解释型（宿主脚本写的 S-Function 或脚本函数块）。}改代码不用重编，调试就在宿主环境里断点，起步最快。代价是每次调用要跨一次宿主边界，在快积分或高频采样下这一步开销很可观。
\item \textbf{\CL{} MEX。}先编成二进制，运行期是原生调用；状态与采样率语义要自己按规矩写对，出错时表现为求解器深处的崩溃或静默的错误数值。
\item \textbf{图形化的块构造器。}由向导生成 \CL{} 骨架后编译，改逻辑要在向导里点。适合一次性的小算法，不适合版本管理与评审。
\end{itemize}

三档的取舍见表~\ref{tab:ch07-sfun-tier}：它其实是本章开头那条谱系在控制侧的重演，
连``访问深度换可移植性''的规律都一样。

\begin{table}[htbp]
\centering
{\small
\begin{tabular}{%
 >{\raggedright\arraybackslash}p{0.16\textwidth}
 >{\raggedright\arraybackslash}p{0.16\textwidth}
 >{\raggedright\arraybackslash}p{0.21\textwidth}
 >{\raggedright\arraybackslash}p{0.28\textwidth}}
\toprule
档位 & 编译时机 & 运行期成本 & 适合与不适合 \\
\midrule
解释型 & 无，运行前解释 & 每次调用跨宿主边界 & 算法验证适合；高频控制环不适合 \\
\CL{} MEX & 事先编成二进制 & 原生调用，可内联 & 量产与代码生成适合；起步门槛高 \\
块构造器 & 生成骨架后编译 & 与 MEX 同档 & 单人小算法适合；评审与 diff 不适合 \\
\bottomrule
\end{tabular}
}
\caption{S-Function 三档入口的编译时机、运行期成本与适用面}
\label{tab:ch07-sfun-tier}
\end{table}

要量产（嵌入式代码生成、硬件在环）时，通常只有 \CL{} MEX 这一档能跟着走到生成代码；这一条常反过来决定项目该用哪一档，值得在早期确认。与 \DomCAE{} 侧的连接点：多物理场平台与控制系统联合仿真的接口，本质上就是把一侧的求解器包成另一侧的一个块，接口形状仍是``端口进、端口出、状态自己保管''；字段与端口的映射细节以官方文档为准。

% ------------------------------------------------------------------
\section{开源求解器的两条路：改源码与写插件库}
\label{sec:ch07-ossolver}

开源求解器（以 \Lib{OpenFOAM} 系列、\Lib{deal.II}、\Lib{FEniCS}、\Lib{SU2} 一类为参照）把源码级这一档默认摊在你面前，于是编译型与源码级之间多出一条中间路：\textbf{不改核心，编一个求解器运行时加载的库}。

\subsection*{路径 A：插件或函数库}

形状上就是本章前面的``编译型用户函数''，只不过用的是求解器自己的 \CPP{} 类型。典型做法：实现一个继承自某基类的类，按名字注册，之后在控制文件里按名调用（OpenFOAM 的 function object 正是这一形状：每个时间步被调一次，拿到当前场对象的引用，自己决定写什么）。升级成本集中在两处：一是\textbf{签名漂移}，基类的虚函数换名字或换参数，插件就编不过；二是 \textbf{ABI 对齐}，同一版本不同构建选项（C++ 标准、运行库、断言开关）可能连符号都对不上。控制办法是把插件放进独立仓库、用求解器自带的构建脚本编译、在持续集成里对每个要支持的求解器版本各编一遍，并把``编不过''当作升级信号而不是噪声。

\subsection*{路径 B：改核心源码}

只有当你改的是\textbf{策略}而不是\textbf{量}时才值得：新的离散格式、新的预条件器或线性求解器组合、新的时间推进、新的装配顺序。代价是从此背一个分支，上游每动一次你要重看一次差异；收益是不受任何钩子可见范围的限制。判断标准很朴素：如果改动可以写成``在两个已有函数之间插一段''，那它应该做成插件；如果改动是``把这两个函数的调用顺序换掉''，那才是源码级。分支的维护成本还与版本节奏成正比：给自己定一条硬规矩——每次上游合并只做一件语义上的事，并把本地提交压成一小批可复述的补丁，否则半年后没人分得清哪些改动是必须的。

\subsection*{许可传染：给判断方法，不给结论}

GPL 类许可在这条路上是真问题，因为``你的代码链接进一个 GPL 程序再分发''正落在它管得最严的地方。这里只给判断方法：

\begin{enumerate}
\item 读\textbf{核心}的 LICENSE 与 README 原文，注意是否附带例外条款。有的项目在 GPL 之外给一段库例外，有的整块是 LGPL 或 BSD 类，差别巨大。以仓库文本为准，不要用流传的版本印象。
\item 分清\textbf{分发}与\textbf{内部使用}。只在内部跑，义务范围与把二进制连同你的模块交给客户完全不同。
\item 问三个技术问题：静态还是动态链接？你的产物是否随核心一起发布？你的模块是被解释执行的脚本，还是编进同一进程的原生库？
\item 把上面三条的\textbf{答案写成文档}放进仓库，让下一个接手的人不必重新猜。
\item 触及对外发布、双重许可或商用授权时，把判断交给懂知识产权的角色；工程师负责提供准确的技术事实。
\end{enumerate}

\begin{lstlisting}[style=cmakestyle,
  caption={把用户函数编成运行时加载的模块：与核心解耦的构建形状},
  label={lst:ch07-userlib-cmake}]
# 形状示例：核心用哪个构建系统不影响这几行的思路
add_library(myfun MODULE myfun.cpp)
target_include_directories(myfun PRIVATE "${SOLVER_INC}")
target_compile_features(myfun PRIVATE cxx_std_17)

# 运行时按名加载的库，通常不要 lib 前缀
set_target_properties(myfun PROPERTIES PREFIX "" SUFFIX ".so")

if(NOT CMAKE_BUILD_TYPE)
  set(CMAKE_BUILD_TYPE Release)
endif()
target_compile_options(myfun PRIVATE -O2 -fPIC)

# 关键一条：插件与核心的编译选项要对齐（语言标准、运行库、
# 断言开关），否则符号查得到也会在运行期崩，且崩点不在你的代码里。
\end{lstlisting}

\begin{bestpractice}[title={插件路线的四条纪律}]
一、把插件目录与核心版本一一对应，构建产物命名里带核心版本号。二、入口函数做成纯函数：同样的输入给同样的输出，不开文件、不加全局可变状态、不打时间戳。三、任何分配都要有对应的释放路径，因为库加载一次、调用十万次。四、给插件配一组不依赖求解器的单元测试：把平台传进来的数组形状在测试里照样伪造一份，就能在提交前抓越界与列顺序错误。
\end{bestpractice}

% ------------------------------------------------------------------
\section{选平台时该问的八个问题}
\label{sec:ch07-eight-questions}

上面这些差别最后都会落到选型上。这八个问题按``能不能自动化、能不能扩展、能不能带出门''的顺序问，比读功能表有效：

\begin{enumerate}
\item \textbf{能不能离线跑？}没有图形界面、没有网络许可服务器时能不能提交一个算例。这决定它能否进持续集成，也决定你半夜排队时能不能干活。
\item \textbf{批处理接口长什么样？}入口是命令行加模型归档文件，还是必须走宿主脚本；作业能否排队、能否只跑一部分网格或若干个时间步、失败后能否只重跑受影响的那一段。
\item \textbf{许可证与并行怎么绑？}核数、GPU、多机是否分别计费；并行效率的提升要不要额外授权。这决定``多快好省''里前两项能不能同享。
\item \textbf{结果导出格式是谁的？}有开放的中间格式（网格、场数据、时间序列）才能把后处理接给自己写的工具；只有私有归档格式就被绑定，这一环节在 Part I 的收束章里单独展开。
\item \textbf{用户函数的编译口径？}平台是否提供匹配的头文件与运行库、是否要求专用构建脚本、集群计算节点上有没有工具链；``在我机器上能编''在集群上不成立。
\item \textbf{内部量能读到哪一层？}只有节点值，还是能读到梯度、雅可比、网格拓扑、上一迭代值。这一条决定表~\ref{tab:ch07-lineage} 里哪一档够用。
\item \textbf{并行下你的代码被怎么调？}每个进程各加载一份还是只在主进程？钩子拿到的是局部索引还是全局索引？跨进程的汇总要不要自己写？
\item \textbf{版本升级的历史成本？}同一模型跨大版本能否打开、用户函数要改多少、有没有可读的迁移说明、社区里升级踩坑的密度如何。
\end{enumerate}

\begin{guideline}[title={本章要点：先定三轴，再查参数表}]
换平台的成本不在语法，而在三轴的重排。写代码前先定位自己在谱系的哪一档，然后确认这段逻辑被调用的\textbf{频率}是否配得上它的\textbf{可见范围}。把三件事处理掉——不连续的本构要给出导数来源、跨进程的全局假设要显式汇总、高频钩子里不做输入输出——扩展面就不会在夜里咬你。
\end{guideline}

\begin{warning}[title={这一章的四个典型陷阱}]
\begin{itemize}
\item 在高频钩子里做文件写入或加锁：症状是``随机变慢''而不是崩溃，最难定位。
\item 把状态藏进静态变量而不是平台登记的状态槽：变步长退步重算时不被还原，结果随容差漂移且无法复现。
\item 外部函数只给点值不给导数：求解器差分扰动撞上分段点，线性化失效，常被误读成``网格太粗''。
\item 用印象判断许可：以仓库里的 LICENSE 原文为准，并把结论连同技术事实一起写进仓库文档。
\end{itemize}
\end{warning}

% ============================================================
% ---- frag_ch08.tex ----
\chapter{求解器工程共性：网格 I/O、区域分解、检查点与后处理}
\label{ch:ch08-solver-engineering}

这是 Part I 的收束章。前面几章按产品逐个看扩展面，本章换一个视角：把\DomCAE{} 侧反复出现的四个工程环节抽象出来——\textbf{网格怎么进内存}、\textbf{并行怎么切}、\textbf{状态怎么存}、\textbf{结果怎么出去}。这四件事在各家求解器里的形状高度相似，也正是下一章之后与 \DomEngine{} 侧接轨的地方：引擎同样要处理资源装载、状态序列化与剖析。凡本章讲到的机制，在第 \ref{sec:ch08-part2-bridge} 节都会指一次它在实时侧的对应物。

% ------------------------------------------------------------------
\section{网格数据的表示与读写代价}
\label{sec:ch08-mesh-representation}

一个网格在内存里最少是两块数据：\textbf{节点坐标表}与\textbf{单元连接表}。连接表通常每行先给一个类型码（四面体、六面体、楔形、金字塔，或更高阶的形状），再给一串节点索引。围绕这两块还有三样常被低估的东西：

\begin{itemize}
\item \textbf{属性分组。}单元按区域分组（材料号、单元块、part），面按边界分组（边界名、patch、boundary id）。分组决定装配循环的边界：求解器``对每个组做不同的事''，你的用户函数被挂在哪一组上、以什么顺序被遍历，都由它决定。一个真实网格常有几十到几百个面分组，名字是字符串，于是读取时要做一次字典化，把名字换成整数编号。
\item \textbf{高阶单元。}自由度不止在顶点：边上、面上、体内都有节点，连接表要能容纳可变长度的节点列表。曲边网格还带几何映射信息（把参考单元映到物理单元的那个映射的系数），丢了它，``高阶''就退化成折面近似。
\item \textbf{坐标系与约定。}同一份坐标可能带单位（毫米还是米）、带参考系（旋转域要一个角速度或参考位置）、带维度约定（二维轴对称的第三维是角度还是被忽略的厚度）。这些是网格文件与物理场接口之间的\textbf{隐式契约}，跨工具时最容易出错。
\end{itemize}

\subsection*{读一次网格要付的代价}

从磁盘到能被装配循环使用，一份网格一般要过四道手：\textbf{读入}（一次大块顺序读）、\textbf{解析}（文本格式在此处最贵，行数乘以浮点转换，往往比求解本身还慢）、\textbf{类型与索引转换}（文件里的 1 基索引变 0 基，32 位全局编号升 64 位，浮点精度对齐）、\textbf{重排与分区}（为了访存连续或为了并行分块，把节点与单元重新编号）。每道手都是一次完整拷贝，粗算下来是``原始字节数 $\times$ 若干份''的峰值内存。工程上有三条通用对策：能存二进制就别存文本；把重排做成\textbf{稳定且可复现}的置换（同一输入两次分区得到同一编号，否则调试时对不上号）；在流水线里尽早确定唯一的所有者，别让三份副本同时活着。

\begin{lstlisting}[language=C++, style=cppstyle,
  caption={最小网格容器：同一份数据的 AoS 与 SoA 两种布局},
  label={lst:ch08-mesh-layout}]
#include <cstdint>
#include <vector>

// 单元类型码：与文件格式的约定一致，不要自己发明编号
enum CellKind : uint8_t { TET4 = 1, HEX8 = 2,
                          PRISM6 = 3, TET10 = 4 };

// ---- AoS：一个单元一个结构体，读文件时顺手 ----
struct CellAos {
    CellKind kind;
    uint8_t  n;          // 节点数：高阶单元可变长
    uint32_t nid[8];     // 连接表，顺序必须与类型码约定一致
    int32_t  group;      // 所属单元组（材料、part）
};

// ---- SoA：一列一列摆，装配与并行时省带宽 ----
struct MeshSoa {
    std::vector<double>   px, py, pz;    // 坐标各占一列
    std::vector<CellKind> kind;
    std::vector<uint8_t>  nn;
    std::vector<uint32_t> conn;          // 扁平拼接，配 offsets
    std::vector<uint32_t> offsets;  // 第 c 个单元的起点
    std::vector<int32_t>  cellgroup;
};

// 读文件时写 AoS 更省心；跑内核时想要 SoA。
// 关键决策：转换只做一次，并把它做成可复现的置换，
// 而不是每次启动都现场重排——否则两次运行的单元编号对不上。
void aos_to_soa(const std::vector<CellAos>& cells,
                const std::vector<double>& xyz, MeshSoa& out)
{
    out.offsets.resize(cells.size() + 1);
    out.kind.resize(cells.size());
    out.nn.resize(cells.size());
    out.cellgroup.resize(cells.size());
    out.conn.reserve(cells.size() * 4);
    uint32_t acc = 0;
    for (size_t c = 0; c < cells.size(); ++c) {
        out.offsets[c] = acc;
        out.kind[c] = cells[c].kind;
        out.nn[c] = cells[c].n;
        out.cellgroup[c] = cells[c].group;
        for (uint8_t j = 0; j < cells[c].n; ++j) {
            out.conn.push_back(cells[c].nid[j]);
            ++acc;
        }
    }
    out.offsets[cells.size()] = acc;
    out.px.resize(xyz.size() / 3);
    for (size_t i = 0; i < out.px.size(); ++i) {
        out.px[i] = xyz[3 * i];          // 交错存到列存：一次转换
        out.py.push_back(xyz[3 * i + 1]);
        out.pz.push_back(xyz[3 * i + 2]);
    }
}
\end{lstlisting}

两种布局的取舍见表~\ref{tab:ch08-layout}，判断依据是``谁在热循环里读哪一列''。

\begin{table}[htbp]
\centering
{\small
\begin{tabular}{%
 >{\raggedright\arraybackslash}p{0.17\textwidth}
 >{\raggedright\arraybackslash}p{0.20\textwidth}
 >{\raggedright\arraybackslash}p{0.20\textwidth}
 >{\raggedright\arraybackslash}p{0.26\textwidth}}
\toprule
维度 & AoS 布局 & SoA 布局 & 什么时候选谁 \\
\midrule
装载代码 & 与文件格式一一对应，写起来短 & 要维护偏移表，容易错位 & 解析阶段偏 AoS \\
内核访存 & 拉进整行，含用不到的字段 & 只拉需要的列，易向量化 & 遍历单元的内核偏 SoA \\
可变长连接 & 结构体里放定长数组，超限要另存 & 扁平数组加偏移，天然变长 & 高阶单元多则 SoA 更省事 \\
共享与拷贝 & 一次整体复制即可传给插件 & 多列要一起传，遗漏一列就错 & 接口设计上 SoA 要写检查 \\
\bottomrule
\end{tabular}
}
\caption{网格容器的两种布局：装载阶段与内核阶段的代价对比}
\label{tab:ch08-layout}
\end{table}

% ------------------------------------------------------------------
\section{区域分解与并行下的用户函数}
\label{sec:ch08-decomposition}

\subsection*{几何分解与图分割}

把网格分给若干计算进程，两条路。一条是\textbf{几何分解}：按坐标范围切块（规则地切、或按八叉树递归二分）。它快、不需要全局信息、可就地增量重建，但切面会沿坐标平面对齐，遇到细长或不规则区域时负载很不平衡。另一条是\textbf{图分割}：把网格当图（单元为顶点、共面为边），交给 \Lib{METIS}、\Lib{Scotch} 一类的分割器，同时优化接口大小与块间不平衡度。它质量高但很贵——分割一次的时间可能与求解若干步同量级，所以``每次网格变化就重新分区''的做法往往不划算；动态自适应里通常配合``分区粗调 + 单元迁移细调''。

\subsection*{halo 与 ghost 单元}

分块之后，每个进程只持有自己的单元，但多数格式要邻居的值。于是沿分块边界复制出一圈\textbf{ghost（或 halo）单元}：只读副本，每步结束时按边界面同步。要复制几层由数值格式的模板宽度决定——一阶模板一圈就够，高阶单元或大梯度重构可能要两圈。由此得到并行效率的根源限制：\textbf{同步的消息量按界面面积增长，而计算量按体积增长}，所以核数上去后通信占比必然上升；网格很碎、单元数不多时，切到一定核数就没有意义了。几何上狭长的域（沿一个方向极长）界面占比高，扩展性天生差，这不是调参数能救的。

\subsection*{MPI 与共享内存混合}

现代跑法多是两层：进程间用消息传递，进程内用线程或向量化。工程细节集中在三处。一是\textbf{绑核}：每个进程占的物理核要与线程数一致，否则内存带宽被抢占，表现为``加了核反而变慢''。二是\textbf{归约顺序}：浮点加法不满足结合律，求和顺序一变末位就变；线程数或分块数一改，结果就逐位不同。要跨核数复现，就得固定归约树或用确定的分块求和。三是\textbf{分区依赖}：分区结果依赖 rank 数，因此``同一模型 8 核与 32 核结果差多少''是一个必须知道答案的问题——通常按物理量给容差，而不是指望逐位一致。

\subsection*{你的代码在并行下被怎么调}

这一段直接对应上一章的编译型用户函数档，形状是固定的几条：

\begin{itemize}
\item \textbf{每 rank 各加载一份库。}你的代码在每个进程里都有一份独立副本，静态变量、缓存、计数器都是每份一个，会各自漂移。别把它们当全局状态用。
\item \textbf{拿到的是局部索引。}钩子给的单元或面编号通常是本进程内的局部号，可能附带一个全局 id；不要假设全局 id 连续、也不要假设它在本进程的数组里能直接下标。
\item \textbf{看不到邻居的全部。}需要``这个单元周围''的信息时，只有 halo 里那圈同步过的值可用；跨分块求全局极值、总量或平均，必须走显式归约，不能靠本地扫描。
\item \textbf{线程可能并发调用同一入口。}同一进程内多线程遍历分组时，你的函数会被并发调用；任何非可重入的东西（静态缓冲、\texttt{strtok} 形态的解析、共享文件句柄）都会在这里炸。
\item \textbf{写文件要选边。}要么由一个进程序列化，要么每进程写一片留待合并（见第 \ref{sec:ch08-postprocess} 节）；两边都写同一路径的``追加''在并行文件系统上不保证顺序。
\end{itemize}

\begin{lstlisting}[language=C++, style=cppstyle,
  caption={halo 同步与按 rank 加载的用户函数：调用形状要点清单},
  label={lst:ch08-halo-shape}]
// 每个 rank 的求解循环骨架（伪代码，函数名按平台）
for (int step = 0; step < nsteps; ++step) {
    exchange_halo(field);        // 沿边界面同步 ghost 值
    // 本 rank 只遍历自己的单元；邻居可能在别的 rank 上
    for (uint32_t c = 0; c < local_ncells; ++c) {
        // 入口只给局部索引与已同步的模板值
        // 不要在这里问"全局最大压力是多少"
        flux[c] = user_flux(centroid(c), stencil(c));
    }
    double f = local_force_sum(flux);
    all_reduce_sum(f);          // 要总量就显式归约，别本地扫
    if (rank == 0) write_monitor(step, f);  // 或每 rank 一片
}

// 三条并行纪律：
// 1. 库每 rank 一份，静态计数器会各自漂移——要总量就通信。
// 2. 索引是局部的；全局 id 只用于交叉核对，不能直接下标。
// 3. 内核可能被多线程并发调用：入口必须可重入、不写共享缓冲。
\end{lstlisting}

\begin{note}[title={``能跑通单核''不等于``能跑通并行''}]
最常见的三类只在并行下暴露的问题：把局部索引当全局索引使用（单核下两者恰好一致）；在用户函数里累加一个全局量（单核下正确，多核下变成每进程各算一份）；输出时不排序就写监视表（单核顺序稳定，多核下交错）。值得为这三类各写一条最小回归算例，用两核跑一遍并检查守恒闭合。
\end{note}

% ------------------------------------------------------------------
\section{检查点与恢复}
\label{sec:ch08-checkpoint}

\subsection*{为什么要写历史，而不只写场值}

``把当前速度场和压力场存下来就够重启了''几乎总是不成立。要重启得和连续跑一致，就得把\textbf{下一步右端项会读到的东西}全部存下来，而它们往往不止是当前场：

\begin{itemize}
\item \textbf{多步格式的历史。}二阶向后差分需要上一乃至上两步的场；变步长还要每次历史的时间间隔与当时的格式阶数，否则重启后的第一步阶数悄悄降了。
\item \textbf{求解器的热启动。}线性求解器的不完全分解因子、Krylov 子空间、非线性迭代的阻尼或线搜索历史；丢了它们，重启后前几步迭代次数会飙升，看起来像``重启后变慢''。
\item \textbf{内部变量。}塑性与损伤的历史量、湍流模型量、任何带自身输运方程的标量——它们不是``结果''，但决定下一步的物性。
\item \textbf{离散与自适应状态。}当前分区、网格自适应树与历史误差指示、 moving mesh 的当前位置与速度。
\item \textbf{离散事件的账本。}粒子的速度与接触历史、探针的滑动窗口、监视量的累加器、随机数发生器的内部状态、当前时间步长与 Courant 爬升进度。
\end{itemize}

\subsection*{版本与字段顺序带来的坑}

位置式的二进制转储最省事也最脆：加一个场、删一个场、调换两行顺序，旧文件就整体错位，而且不会报错——它读出一份``看起来合理''的垃圾。成对的对策是：

\begin{itemize}
\item 用\textbf{自描述}的容器（按名字与形状取数据，而不是按偏移量），或至少在文件头写清单：魔数、主版本、字段名、形状、元素类型与字节序。
\item \textbf{按名字映射}读取，允许新增字段缺省、允许顺序变化；主版本不匹配就明确拒绝，不要``尽量读一下''。
\item 索引与编号的\textbf{宽度}要显式（32 位还是 64 位单元号），浮点精度要显式；固定长度字符串要注意填充。
\item 加一条测试：\textbf{用上一个发布版写的文件，当前版本要能读}。这条测试是唯一能证明兼容矩阵存在的方式。
\end{itemize}

\begin{lstlisting}[style=jsonstyle,
  caption={检查点清单（manifest）骨架：让读者按名字取字段并校验形状},
  label={lst:ch08-ckpt-manifest}]
{
  "format": "ckpt", "magic": "SIMCKPT", "major": 3, "minor": 0,
  "solver": { "name": "incomp-nse", "commit": "abc1234" },
  "clock": { "time": 1.25e-01, "step": 2500, "dt": 5.0e-05,
             "dt_prev": 5.0e-05, "scheme_order": 2 },
  "partition": { "ranks": 32, "decomp": "metis",
                 "halo_layers": 1 },
  "fields": [
    { "name": "velocity", "items": 3, "n_cells": 4194304,
      "dtype": "float64", "unit": "m/s" },
    { "name": "pressure", "items": 1, "n_cells": 4194304,
      "dtype": "float64", "unit": "Pa" },
    { "name": "turb_k", "items": 1, "n_cells": 4194304,
      "dtype": "float64", "unit": "m2/s2", "since": 2 }
  ],
  "history": [ { "name": "velocity_n1", "items": 3,
                 "dtype": "float64" } ],
  "rng": { "algorithm": "pcg64", "seeds": "rank-local" },
  "files": [ "rank-%05d.ckpt" ]
}
\end{lstlisting}

\subsection*{恢复之后要校验什么}

重启不是``能跑''就算对。四类校验值得写进回归：\textbf{连续性}——重启点前后若干步的关键监视量（阻力系数、进出口压差、总能量）不应出现台阶，出现台阶说明少了某个历史量；\textbf{单调性与有界性}——损伤、体积分数、湍流量这类物理上有界的东西，重启后不应越界或被清零；\textbf{守恒闭合}——把``体积内积分的变化''与``通过边界的净通量加源项''对账，重启前后两条曲线的差应在容差内；\textbf{可复现性}——同一分区配置下从检查点跑一步，应与连续跑的对应一步逐位一致；换核数则只按容比。最后一个技巧最省钱：\textbf{在同一个时间步上做``连续跑''与``重启跑''的对拍}，差值直接指出你漏存了哪个量。

% ------------------------------------------------------------------
\section{结果导出与后处理管线}
\label{sec:ch08-postprocess}

\subsection*{三类输出，三种问题}

\textbf{整场快照}回答``当时空间上怎么分布''，代价是体积：单元数乘场数乘字节数再乘输出次数，很容易上太字节。\textbf{探针}（固定点或沿路径取值）回答``某处随时间怎么变''，体积随探针数而非网格数增长，但一旦网格自适应或分区变了，``同一个点''就不一定是同一个点——用坐标加最近单元定位，别用编号。\textbf{派生量}（积分、阻力升力、通量、残差历史）是最该优先做扎实的一类：体积最小、可直接对比、能当校验用。工程上常犯的错是反过来：为了保险把整场存得极密，结果下游读不动，而派生量的定义却因为没代码而说不清。

\begin{table}[htbp]
\centering
{\small
\begin{tabular}{%
 >{\raggedright\arraybackslash}p{0.16\textwidth}
 >{\raggedright\arraybackslash}p{0.20\textwidth}
 >{\raggedright\arraybackslash}p{0.22\textwidth}
 >{\raggedright\arraybackslash}p{0.23\textwidth}}
\toprule
输出形态 & 体积随什么增长 & 事后能回答什么 & 落盘建议 \\
\midrule
整场快照 & 单元数乘场数乘次数 & 空间分布、涡结构、切片 & 按物理意义定间隔，写时降采样 \\
探针时间序列 & 探针数乘次数 & 定点演化、频谱、反馈 & 高频可取，务必记录探针坐标 \\
派生量与积分 & 量的个数 & 载荷、守恒闭合、稳定性判据 & 最该做扎实，每步都写也行 \\
\bottomrule
\end{tabular}
}
\caption{三类输出在体积、可回答的问题与落盘节奏上的对比}
\label{tab:ch08-output}
\end{table}

\subsection*{把大结果交给下游}

并行下每 rank 写一片再留一个全局索引，是让写盘不阻塞模拟的关键：各进程只碰自己的文件，元数据在收尾时合并或写成自描述容器（HDF 类、ADIOS 类都在这个用途上）。三个实用手法：\textbf{写出时就切片}——只留几个常用截面与几个常用时刻，体积能降一到两个数量级，代价是``当初没想到要看的那个面''从此没有；\textbf{分块压缩}——按块压缩仍可随机读取，整文件压缩则退化成顺序访问；\textbf{时间轴单独存}——不规则输出时刻要显式写时刻数组，别假设``等间隔''，否则后处理按索引取时刻会整体错位。

\subsection*{与 \DomEngine{} 侧的交界}

后处理管线的下游往往是渲染端，而渲染端的约束完全不同：它要的是\textbf{固定大小、可直接上传、按帧消费}的资源。对应关系很直白：体数据快照要变成纹理（三通道打包速度分量、单通道打包标量，尺寸取 2 的幂或对齐到设备约束，必要时降采样）；面数据要变成带顶点属性的网格；探针与派生量走另一条路——它们是遥测数据，进日志与界面，不进显存。这条边界的工程含义：导出格式要在\emph{写出时}就考虑``能否被 GPU 直接映射''，否则中间必有一个转换步骤，而那一步通常既慢又没人维护。

\begin{lstlisting}[style=pythonstyle,
  caption={后处理脚本的数据流形状：读清单、取字段、算派生量、出图},
  label={lst:ch08-post-py}]
import glob, json
import numpy as np

def load_step(manifest):
    m = json.load(open(manifest))
    parts = sorted(glob.glob('rank-*.ckpt'))   # 每 rank 一片
    fields, cells = {}, 0
    for spec in m['fields']:
        dt = spec['dtype']
        acc = [np.fromfile(p, dtype=dt) for p in parts]
        fields[spec['name']] = np.concatenate(acc)
        cells = fields[spec['name']].size // spec['items']
    return m, cells, fields

def derived(cells, f):
    vel = f['velocity'].reshape(cells, 3)
    p = f['pressure']
    return {'speed': np.linalg.norm(vel, axis=1),
            'cp': (p - p.min()) / (p.max() - p.min() + 1e-12),
            'outflow': vel.sum(axis=0)}        # 粗略守恒指示

def probe(series, xyz):          # 按坐标取，不用单元编号
    t = series['times']
    v = np.interp(np.asarray(xyz), series['probe_xyz'],
                  series['values'])
    return list(zip(t, v))

def main():
    m, cells, f = load_step('ckpt-0003.json')
    d = derived(cells, f)
    np.save('speed_step2500.npy', d['speed'])  # 交给渲染端
    print(cells, d['outflow'])

if __name__ == '__main__':
    main()
\end{lstlisting}

\begin{warning}[title={导出与后处理的四个陷阱}]
\begin{itemize}
\item 只存场值不存时刻数组与分区信息，重启后时间轴与几何对不上，``同一个点''换了位置。
\item 派生量没有代码定义：报表里的升力系数没人能复算。定义必须与算例同仓库、同版本。
\item 大结果一次性写完：进程中途失败时文件不完整且不可读。分段写、最后合并，或加一个收尾才写的索引。
\item 后处理里用编号做键：网格自适应或重新分区后编号全变，只应把坐标当键。
\end{itemize}
\end{warning}

% ------------------------------------------------------------------
\section{通往 Part II：确定性时间步、序列化与剖析，以及 ch09--ch16 的分组}
\label{sec:ch08-part2-bridge}

三个机制在求解器与引擎两侧都出现，只是预算不同。想清楚这三个，后半本书读起来会顺得多。

\textbf{确定性时间步。}本章的检查点纪律（整型步计数、显式存时间步长、不依赖挂钟、可播种的随机数）在 \DomEngine{} 侧是同一套东西的另一份实现：帧序号、固定逻辑步长与可变渲染步长分离、随机数带种子进存档。差别只在预算——求解器容忍每秒几十步，引擎要求每秒六十步且抖动必须小于人眼阈。两侧共同的硬规矩是\textbf{逻辑状态不许读挂钟}，一旦读了，回放、检查点与联机同步全部失效。

\textbf{序列化格式。}检查点、模型归档、引擎的场景文件、联网玩法的快照，本质是同一件事：带版本的字段清单加数据。本章列出的对策——按名字映射、主版本严格拒绝、旧文件读取回归测试——在实时侧同样成立，而且因为引擎要跨端跨版本，兼容矩阵更长。上一章说的``扩展面档位''在这里也复现：脚本层通常只允许序列化平台认识的类型，原生层可以自定义布局，代价是自己承担兼容矩阵。

\textbf{剖析方法。}两条路都从``一步多久''开始：先测每步的时间分布（装配、通信、求解、输出），再下钻到函数级；把\textbf{通信与同步}当一等公民看，因为并行求解器与多线程引擎的瓶颈常常都在这里，而不是在算术里。差别在目标函数：求解器看扩展效率与单位结果的墙钟，引擎看帧时间的上限与分位。共同的纪律是\textbf{剖析与优化不要同时做}，且任何优化都要留下一个可复算的对照算例——这一条正是本章检查点与导出环节给你的能力。

Part II 的路线按三组展开：\textbf{ch09--ch12} 处理引擎侧的脚本层与原生边界（运行时对象模型、脚本到原生调用、热重载与生命周期）；\textbf{ch13--ch14} 处理实时循环本身（固定步长与变步长的调和、状态同步、序列化与回放）；\textbf{ch15--ch16} 处理资产与工程化（资源装载管线、剖析与性能预算、构建与持续集成）。各章实际主题以目录为准。贯穿这六章的问题与本章同形：\textbf{数据以什么布局进内存、状态以什么格式出内存、时间以什么节奏推进}。

\begin{bestpractice}[title={Part I 收束：四个环节各自的最小可用配置}]
\textbf{网格}：二进制输入、解析一次、重排置换固定并写进日志，字段名与单位在代码里以单一来源定义。
\textbf{并行}：分区可复现、归约顺序固定、用户函数入口做成可重入的纯函数，并为``两核跑一次''配一条守恒闭合检查。
\textbf{检查点}：自描述清单、主版本严格校验、把``右端项会读到的东西''逐项列全并按连续跑与重启跑对拍。
\textbf{后处理}：派生量的定义代码化，快照按物理意义定间隔并在写出时切片，给渲染端的资源单独一条导出路径。
\end{bestpractice}

% ============================================================
% ---- frag_ch09.tex ----
\part{游戏引擎与实时仿真的脚本层}

\chapter{引擎的脚本层次：本机代码、反射、虚拟机与热重载}
\label{ch:ch09-engine-layers}

第一册的前八章处理非实时求解器的扩展层：一次作业可以跑几分钟到几小时，
``改宏、重跑''在节奏上完全可接受。第二部分换场景：实时仿真与游戏引擎。
这里一帧通常只有约 16 毫秒预算，循环一旦上线就不会停，而且每天改动内容的
人多数不是引擎开发者。两条约束叠加，把引擎推向同一个设计：上层接快迭代，
下层担性能。本章先立起后面各章共用的框架——引擎的哪些能力留在本机代码，
哪些能力通过反射暴露给脚本作者，虚拟机坐在链条的哪一段，热重载的边界
到底划在哪里。

\section{为什么引擎必须有两层}
\label{sec:ch09-why-two-layers}

先看生产分工。引擎的性能内核是一小撮 \CPP{} 模块：渲染器、物理、动画求值、
音频 DSP、资源流送。它们对正确性和速度的要求，和第一部分求解器热点几乎同构
——一次多余的指针跳转、一处散乱的内存访问，都会原样反映在帧时间曲线上。
但项目真正量产的内容——关卡、角色、技能、界面逻辑、活动配置——每天都在变，
改动者是策划和玩法程序员，他们的工具链通常是编辑器，不是编译器。让后者为了
改一个掉落率去重编整个引擎，等于把迭代关掉。

于是引擎在本机内核之上嵌一层脚本，把内核的一部分能力包装成可调用的对象。
这条边界可以落到四个具体问题上：谁在改，改完多久生效，错了怎么退，
能不能碰到性能内核。表~\ref{tab:ch09-layer-contrast} 把 \DomCAE{} 里
``宏 + 用户子程序''的分层与 \DomEngine{} 的脚本层并排放在一起。

\begin{table}[ht]
\centering
\caption{\DomCAE{} 用户子程序分层与 \DomEngine{} 脚本层对照}
\label{tab:ch09-layer-contrast}
\begin{tabular}{>{\raggedright\arraybackslash}p{0.15\textwidth}
                >{\raggedright\arraybackslash}p{0.34\textwidth}
                >{\raggedright\arraybackslash}p{0.34\textwidth}}
\toprule
\textbf{问题} & \textbf{\DomCAE{}：宏与用户子程序} & \textbf{\DomEngine{}：脚本与原生内核} \\
\midrule
谁在改 & 求解器用户，多为分析工程师本人，按接口手册写 \CL{} 子程序 & 策划、玩法程序员、技术美术，在编辑器或外部脚本里改 \\
改完多久生效 & 子程序要重编译再链接，宏随作业重执行；都按``下一次运行''计 & 解释层秒级重载；托管层需重启或热补丁；原生层要重编模块 \\
错了怎么退 & 删输入档里的宏、还原子程序源码、重跑基线；失败隔离在一次作业内 & 脚本是资产，随关卡走版本控制，可整体回退；运行期出错可能崩整进程，需要回滚加远程日志 \\
能否碰到性能内核 & 子程序被求解器逐单元热回调，反向影响步长与收敛 & 脚本经反射接口调内核，多数只是``递参数''；只有编译到原生的路径才真正改内核 \\
\bottomrule
\end{tabular}
\end{table}

对照之下能看出两点差异。第一，\DomCAE{} 的边界按``文件''划分：输入档与
子程序源文件互不越界，求解器只认固定的函数签名；\DomEngine{} 的边界按
``类型''划分：一个类里哪些属性可序列化、哪些方法可被脚本调用，逐项声明。
第二，\DomCAE{} 用户即使把子程序写崩，损失通常是一次作业；引擎脚本崩的是
还在跑的玩家进程，所以热重载必须和回滚、隔离一起设计。这两点解释了为什么
引擎要在``能改''与``改不死''之间做远多于 \DomCAE{} 的文章。

\section{四种脚本载体}
\label{sec:ch09-four-carriers}

``脚本层''不是一个东西。按``代码在哪里变成可执行指令''划分，主流引擎
做了四种选择，它们的组合决定了项目的迭代速度和安全边界。

\textbf{解释型内嵌}。引擎里带一个小型解释器——最常见的是 \Lib{Lua}，
也有自研的同类脚本。宿主把 \CL{} 函数注册成脚本环境里的全局符号，
脚本直接调用；边界很薄：参数在脚本栈上逐个取出、换算成 \CL{} 标量，
交给本机函数。清单~\ref{lst:ch09-embed-boundary} 给出这条边界的形状，
它和第一部分``宿主调用用户子程序''的镜像关系一目了然——只是方向反过来，
这次是脚本热路径调用本机内核。

\begin{lstlisting}[language=C++, style=cppstyle,
  caption={解释型内嵌：宿主注册本机函数给 \Lib{Lua}（形状示例，接口以对应语言官方文档为准）},
  label={lst:ch09-embed-boundary}]
// script_host.cpp -- 嵌入解释器的边界形状
extern "C" {
#include "lua.h"       // 参考实现的头文件族
#include "lauxlib.h"
#include "lualib.h"
}

// 把一个本机能力暴露成脚本可调的全局函数
static int l_apply_force(lua_State* L)
{
    double mx = lua_tonumber(L, 1);   // 第 1 个入参
    double fy = lua_tonumber(L, 2);   // 第 2 个入参
    PushToRenderThread(mx, fy);       // 交给性能内核
    lua_pushnumber(L, 0);             // 返回码入栈
    return 1;                         // 栈上有 1 个返回值
}

void InitScripts(lua_State* L)
{
    lua_register(L, "apply_force", l_apply_force);
    luaL_dofile(L, "scripts/character.lua"); // 玩法脚本
}
\end{lstlisting}

\textbf{托管虚拟机}。以 \Lib{Unity} 系列为例：玩法代码用托管语言编写，
先编译成中间语言字节码，编辑器期跑在托管运行时（\Lib{Mono} 一类）上，
打包移动平台时再走 \Tp{IL2CPP} 这条转译路径。特点是开发期与发行期的
执行模型不一致，同一个函数在两边的性能可以差出一个量级；跨语言边界
要做参数序列化（interop），这本身就是可测量的开销。

\textbf{编译到原生}。\Lib{Unreal Engine} 的玩法 \CPP{}、\Lib{Godot}
的 \Tp{GDExtension}、\Tp{IL2CPP} 的产物都属这一类：脚本或脚本形状的
代码最终变成模块级本机产物。热重载能力最弱——轻则重新链接、重则重启，
换来的是执行期没有解释器、调试器能看到真实调用栈。\Tp{GDExtension}
刻意用一套扁平的 \CL{} 接口把 ABI 稳定住，让插件不必和引擎同编译器
同配置，这是``编译到原生''里少见地把 ABI 当卖点的一环。

\textbf{字节码加 JIT}。脚本先编译成自定义字节码，运行时按需解释，
对热点路径即时编译成本机代码。自研引擎的脚本虚拟机与部分商用方案的
托管路径都长成这样；某个具体引擎当前版本有没有 JIT、能编译到什么
粒度，一律以官方文档为准，不要凭印象写进构建脚本。

表~\ref{tab:ch09-carrier-matrix} 把四维对比浓缩在一起。读表建议：
先把项目的``迭代频率''标出来，再倒推该把哪层内容放进哪一列。

\begin{table}[ht]
\centering
\caption{四种脚本载体在热重载、调试、ABI、包体四个维度上的位置}
\label{tab:ch09-carrier-matrix}
\begin{tabular}{>{\raggedright\arraybackslash}p{0.16\textwidth}
                >{\raggedright\arraybackslash}p{0.175\textwidth}
                >{\raggedright\arraybackslash}p{0.175\textwidth}
                >{\raggedright\arraybackslash}p{0.16\textwidth}
                >{\raggedright\arraybackslash}p{0.155\textwidth}}
\toprule
\textbf{载体} & \textbf{热重载能力} & \textbf{调试} & \textbf{ABI 稳定} & \textbf{包体与运行时开销} \\
\midrule
解释型内嵌 & 最强，改文件即生效，可运行中替换模块 & 弱：栈帧跨语言边界常断 & 宿主接口窄，天然稳 & 带解释器，体积开销小，执行最慢 \\
托管虚拟机 & 编辑器期可重载程序集，发行版受限 & 中：托管侧完整，本机侧要符号表 & 依赖运行时版本，边界要 marshalling & 带运行时与元数据，内存开销中 \\
编译到原生 & 弱：多数要重编或重启，函数级热补丁受限 & 最强：真实调用栈与断点 & 看接口设计，扁平 \CL{} 边界最稳 & 包体最省，编译链最长 \\
字节码加 JIT & 中：可换字节码，JIT 边界要重启或失效 & 中：解释段好查，JIT 段行号常对不上 & 字节码格式即 ABI，要版本管理 & 介于托管与原生之间，看 JIT 策略 \\
\bottomrule
\end{tabular}
\end{table}

\section{反射层：引擎为什么要自省}
\label{sec:ch09-reflection}

脚本层能存在，前提不是解释器，而是\emph{反射}：引擎在运行期知道
每个类型有哪些属性、每个属性的名字与类型、该到哪个偏移去读它。
四件事都消费这份信息。编辑器——\DomEngine{} 的 Inspector 或 Details
面板不是为每个类手写表单，而是照反射表自动渲染；序列化——资产写盘
读盘按属性元数据走，才能在不重编代码的前提下增删字段；GC——
回收器必须知道``哪个成员是指向对象的指针''，否则可达性图无从谈起；
脚本绑定——把类的方法自动暴露给脚本，省掉成手写胶水。

\DomEngine{} 的尴尬在于：\CPP{} 的类型系统给不了这些。运行期类型信息
（RTTI）只有类型名和动态转换，没有``这个字段是资产引用、那个字段允许
在编辑器里改''的语义；模板与宏把很多决定留在了编译期；私有成员的
偏移量连编译器自己都不承诺稳定。于是引擎要么付出巨大代价做全语言反射，
要么把反射需求\emph{圈}在一个约定的小范围内——这正是``注解宏 +
代码生成''路线的由来。清单~\ref{lst:ch09-register-pattern} 展示
代码生成器最终要产出的最小形状：一张静态属性表。

\begin{lstlisting}[language=C++, style=cppstyle,
  caption={反射信息的最小形状：属性名、偏移、类型名组成的静态表},
  label={lst:ch09-register-pattern}]
// reflect_pattern.cpp -- 生成代码要落到什么形状
struct PropInfo {
    const char* name;    // 对外可见的属性名
    int         offset;  // 在对象内的字节偏移
    const char* type;    // 反射世界的类型名
};

class UObjectLike {
public:
    // 生成代码负责实现它
    static const PropInfo* Props(int* count);
    virtual ~UObjectLike() = default;
};

class ParticleActor : public UObjectLike {
    float m_mass = 1.f;   // 想让编辑器看见的字段
    float m_life = 2.f;
    // 头文件工具的活：在编译期把
    // 这张表按注解自动写出来
};
\end{lstlisting}

代码生成器（各引擎叫法不同，本书统称``头文件工具''）补的是中间那段：
它解析源码里约定好的宏与注解，产出反射数据、序列化 glue 和注册代码，
再交给常规编译链。代价要认清：构建图里多了一个前置步骤；宏写错会
以晦涩的编译错误或运行期空表的形式回报；反射只覆盖被注解的片段，
没有注解的部分对编辑器、GC、绑定三者同时是隐形的——后文 ch10 的
事故案例正是这条规律的直接推论。

\section{数据导向的一面：帧预算把布局逼向 SoA}
\label{sec:ch09-data-oriented}

反射层服务于``可改''，帧预算服务于``可跑''，两者对内存布局的诉求
正好相反。一个粒子系统里，反射对象意味着虚表、属性表、分散在堆上的
实例；而求解粒子的循环只关心连续的位置、速度与质量。ch01 的
求解器热点已经演示过同一件事：把``对未知数量单元做同一运算''的
字段摆成连续数组，缓存行利用率与自动向量化才有发挥空间。引擎的
实时路径把这条经验制度化，产物就是 ECS 与批处理：同类数据装进
连续容器，系统按批次扫过去，脚本层不再逐对象发消息，而是声明
``这批发给谁''。

\begin{lstlisting}[language=C++, style=cppstyle,
  caption={帧预算下的两种布局：AoS 逐对象交错，SoA 按字段连续},
  label={lst:ch09-soa-storage}]
// particles_storage.cpp -- 同一份数据的两种摆法
// 场景：1024 个粒子，60 fps，一帧预算约 16 ms

struct Particle {              // 被否：AoS，字段交错
    float x, y, z;             // 只改速度时，位置
    float vx, vy, vz;          // 的缓存行被无谓拖进
    float mass;                // 循环
};

struct ParticlesSoA {          // 取用：SoA，按字段连续
    float x[1024], y[1024], z[1024];
    float vx[1024], vy[1024], vz[1024];
    float mass[1024];

    void Integrate(float dt) { // 步长恒定，可向量化
        for (int i = 0; i < 1024; ++i) {
            vx[i] += 0.f;
            x[i] += vx[i] * dt;
            y[i] += vy[i] * dt;
            z[i] += vz[i] * dt;
        }
    }
};
\end{lstlisting}

注意这并不意味着引擎整体抛弃了对象模型。现实形态是两种布局并存：
逻辑、界面、关卡用反射对象图；大规模仿真、粒子、人群用数据容器，
容器里只放``句柄 + 索引''回去查对象图。脚本层通常只被允许摸前者，
后者通过窄接口暴露——这正是下一节``读引擎文档先定位五项''里
``扩展注册方式''要回答的问题。

\begin{compare}[title={反射对象图与数据容器的分工}]
反射对象图按``实例''组织，服务人：可编辑、可序列化、可被 GC 追踪，
单次访问的常数大。数据容器按``字段''组织，服务循环：批量进出、
布局紧凑，但没有自省，改结构就要重编译。判断一段逻辑该落在哪边，
问一句``它每秒被执行的次数是几十还是几万''足够：几十次走对象图，
几万次的热点即使最初写在对象图里，也会在性能修正时被整体搬到
数据容器——搬的过程就是第一部分见过的``热点重写''，只是触发它的
不再是收敛诊断，而是帧率曲线。
\end{compare}

\section{本册 Part II 的路线}
\label{sec:ch09-partii-route}

按本章之后的安排，ch10 到 ch15 依次回答下面几层（具体章题与
顺序以本书目录为准）：ch10 拆 \Lib{Unreal Engine}——头文件工具、
反射对象与模块构建是``注解 + 代码生成''路线最完整的样本；ch11 拆
\Lib{Unity}——托管虚拟机到 AOT 原生的双执行模型；ch12 拆 \Lib{Godot}
——自研脚本语言与 \Tp{GDExtension} 动态库并存；ch13 回到自研引擎视角，
讨论嵌入脚本虚拟机要自己补齐什么；ch14 处理实时仿真的特殊约束——
确定性回放对脚本提出的一致性与版本化要求；ch15 收尾在工具链——
热重载边界怎么划、剖析器怎么归因、打包路径怎么配。

路线展开前，先把读引擎文档的方法立住。引擎文档成百上千页，但任何
一款的脚本层都可以用五个问题钉死：

\begin{guideline}[title={读一款引擎的文档，先定位这五项}]
\begin{enumerate}
\item \textbf{入口语言}——玩法代码官方第一推荐写在哪门语言里，
      次选是什么，官方对两者的性能承诺差多少。
\item \textbf{构建工具}——谁负责从源码到产物：官方自研构建器
      （如 \Lib{Unreal} 系列的 \Tp{UBT}）还是通用 CMake 一类的
      外部构建系统，它决定了编译时长和集成 IDE 的方式。
\item \textbf{扩展注册方式}——新类型、新函数要让引擎看见，是写宏、
      实现接口、写注册文件，还是编辑器里点选；注册失败的症状
      （看不见、不序列化、崩）分别对应什么。
\item \textbf{热重载边界}——改函数体、加成员、改类型层次，三种操作
      分别是否需要重启编辑器；官方热补丁机制覆盖不了的部分由什么
      流程兜底。
\item \textbf{打包路径}——脚本在发行包里以什么形态存在：源文本、
      字节码、本机机器码，资产烘焙成什么；发行版关掉了哪些开发期
      检查，崩溃回来用什么符号还原。
\end{enumerate}
五项都定位到，才谈得上评估一款引擎的迭代成本；缺任何一项，
后面所有章节的对比都无从安放。
\end{guideline}

\begin{bestpractice}[title={给``两层化''项目的三条通用守则}]
\begin{itemize}
\item 先画边界再写代码：把``策划可改''的清单与``必须走构建''的清单
      分开维护在仓库根目录，任何跨界改动都要显式申请。
\item 给脚本层留退路：重载失败、脚本报错时，运行时应能回退到
      上一个已知良好的资产版本，而不是崩整进程——这条在 ch10、
      ch11 会以事故形式反复出现。
\item 性能疑点先问布局再问语言：从对象图搬到数据容器的收益，
      通常大于在脚本层继续抠循环；ch01 的热点分析流程在实时
      路径同样适用。
\end{itemize}
\end{bestpractice}

% ============================================================
% ---- frag_ch10.tex ----
\chapter{Unreal Engine：UHT、UObject 与模块构建}
\label{ch:ch10-unreal}

\Lib{Unreal Engine} 是 ch09 那张载体矩阵里``编译到原生 + 注解代码生成''
一列最完整的样本：玩法代码是真正的 \CPP{}，却又要让编辑器、序列化、GC
和蓝图虚拟机在运行期``看见''它。承担这个翻译工作的是头文件工具
（UnrealHeaderTool，下文简称 UHT）与它要求的一族注解宏。本章以 UE5
系列公开文档为参照，只写有把握的宏名——\Tp{UCLASS}、\Tp{USTRUCT}、
\Tp{UENUM}、\Tp{UPROPERTY}、\Tp{UFUNCTION}、\Tp{GENERATED\_BODY}
这一族——其余接口用形状描述，具体拼写与参数一律以官方文档为准。

\section{一个类型的生命周期：从注解到面板}
\label{sec:ch10-type-lifecycle}

先把链路走通，再谈每一站的费用。你在工程里新增一个玩法类型时，
下面五步按序发生，每一步都有明确的产出：

\begin{enumerate}
\item \textbf{你写带注解的头文件。}类声明套在 \Tp{UCLASS} 宏之后，
      成员用 \Tp{UPROPERTY}、方法用 \Tp{UFUNCTION} 标注；类体内以
      \Tp{GENERATED\_BODY} 起头，为展开的反射成员留出位置。产出：
      一份``合法但不完整''的 \CPP{} 头。
\item \textbf{UHT 在编译器之前跑。}它解析目标涉及的全部模块头文件，
      校验注解语法，然后生成两类东西：每个头对应的
      \Tp{xxx.generated.h}（类内的反射声明、序列化与元数据访问代码），
      以及模块级的注册与序列化 glue，落在 Intermediate 目录。产出：
      反射表与 glue。注解与生成代码由同一次构建绑死，这是``宏写错
      会以晦涩方式回报''的根源。
\item \textbf{常规编译进模块。}\Tp{.generated.h} 必须写在头文件包含
      序列的最后，让编译器先见到真实字段再见到反射展开。生成的
      \Tp{UClass}、\Tp{UProperty}、\Tp{UFunction} 描述随模块二进制
      固化。产出：模块（本机代码加只读反射数据）。
\item \textbf{启动期注册。}无论编辑器还是打包后的游戏，引擎加载
      模块时把反射数据汇入全局类型系统；序列化器、GC、脚本绑定
      都从这张总表出发。产出：可自省的运行期类型层次。
\item \textbf{消费端各取所需。}编辑器按 \Tp{UPROPERTY} 元数据自动
      渲染 Details 面板——策划改的就是这张表指向的字段；资产读写
      走同一条序列化路径；蓝图虚拟机只见被 \Tp{UFUNCTION} 标出的
      子集。产出：``改内容的人''与``写内核的人''共用一份定义。
\end{enumerate}

理解这条链的要点是：反射信息不是运行时从 \CPP{} 类型系统里``问''出
来的——语言自带的 RTTI 提供不了属性语义，这正是上一章反射一节
讨论过的缺口——而是 UHT 按你的注解\emph{另行制造}的一份平行元数据。
注解缺席的地方，编辑器、序列化、GC、蓝图绑定这几个消费端同时失明。

\begin{note}[title={注解族的分工，一句话版}]
\Tp{UCLASS}/{\Bk}\Tp{USTRUCT}/{\Bk}\Tp{UENUM} 声明``这个类型进反射表''；
\Tp{UPROPERTY} 声明``这个成员进序列化与 GC''；\Tp{UFUNCTION} 声明
``这个方法可被蓝图或远程调用机制使用''。三者互相不能替代——尤其
\Tp{USTRUCT} 不等于``自动可编辑''，成员没有 \Tp{UPROPERTY} 一样
不被序列化。细节与规格符全清单以官方文档为准。
\end{note}

\section{构建描述与模块图}
\label{sec:ch10-build-modules}

引擎的构建单元是\textbf{模块}：一个可独立编译链接的产物，对外符号
由模块名命名的导出宏控制（例如工程模块 \Tp{MyGame} 的类会写
\Tp{MYGAME\_API}，该宏由构建系统按包含方向定义为导入或导出）。
每个模块带一份 \Tp{.Build.cs}——用 C\# 编写的构建描述，类继承自
构建系统约定的模块规则基类，声明依赖哪些引擎模块。清单
\ref{lst:ch10-build-cs} 是最小形状。

\begin{lstlisting}[language=C++, style=cppstyle,
  caption={模块构建描述 \Tp{MyGame.Build.cs}（C\# 语法，此处仅示意字段族，枚举拼写以官方文档为准）},
  label={lst:ch10-build-cs}]
// MyGame.Build.cs —— 文件名即模块名加 .Build.cs
using UnrealBuildTool;

public class MyGame : ModuleRules
{
    public MyGame(ReadOnlyTargetRules Target) : base(Target)
    {
        PCHUsage = PCHUsageMode.UseExplicitOrSharedPCHs;

        // 公共头文件里出现过的类型，其模块进公共依赖
        PublicDependencyModuleNames.AddRange(new string[] {
            "Core", "CoreUObject", "Engine", "InputCore"
        });

        // 只在实现文件里用到的模块，收进私有依赖
        PrivateDependencyModuleNames.AddRange(new string[] {
            "Slate", "SlateCore", "Json"
        });

        // 编辑器模块只在编辑器目标里参与链接
        if (Target.Type == TargetRules.TargetType.Editor)
        {
            PrivateDependencyModuleNames.Add("UnrealEd");
        }
    }
}
\end{lstlisting}

模块之上还有\textbf{目标描述}（以 \Tp{.Target.cs} 结尾），它决定最终
可执行文件打包哪些模块、什么配置；模块之间构成依赖图，循环依赖会
被构建器直接拒绝。\textbf{插件}则是可分发的模块捆：一份清单文件加
自己的若干模块加可选的内容资产，别人把目录拷进 \Tp{Plugins} 就能
按 \Tp{.Build.cs} 的声明接入——引擎自身的编辑器、物理、网络功能
大多也以此形态组织，所以``内置还是外购''在构建图上没有区别。

表~\ref{tab:ch10-configs} 给常见编译配置的差别。经验规则：日常
迭代全部在 Development 做，Shipping 只用于验证包体与性能回归，
把它当日常开发配置会丢掉引擎的大半防御。

\begin{table}[ht]
\centering
\caption{编译配置对比（配置名为惯例写法，实际取值以目标描述为准）}
\label{tab:ch10-configs}
\begin{tabular}{>{\raggedright\arraybackslash}p{0.16\textwidth}
                >{\raggedright\arraybackslash}p{0.23\textwidth}
                >{\raggedright\arraybackslash}p{0.23\textwidth}
                >{\raggedright\arraybackslash}p{0.21\textwidth}}
\toprule
\textbf{配置} & \textbf{主要用途} & \textbf{检查与开销} & \textbf{发行形态} \\
\midrule
Debug、DebugGame & 断点级排错，符号全 & 不优化，引擎内部校验最满 & 不外发，只在本机 \\
Development & 编辑器与日常迭代 & 优化部分开启，校验保留，日志可查 & 可出内部包 \\
Shipping & 面向玩家的正式包 & 开发期校验基本移除，体积优先 & 发行；崩溃须靠符号表还原 \\
\bottomrule
\end{tabular}
\end{table}

\section{蓝图与 \CPP{} 的边界}
\label{sec:ch10-blueprint-cpp}

玩法代码不必全部写在 \CPP{}：\Tp{UCLASS(Blueprintable)} 的类可以被
蓝图继承，蓝图是编辑器里的可视化脚本，跑在引擎自带的虚拟机上。
暴露接口靠两个规格符：\Tp{UFUNCTION(Blueprint{\Bk}Callable)} 让方法可被
蓝图节点调用；\Tp{UFUNCTION(Blueprint{\Bk}Native{\Bk}Event)} 让默认实现编译在
本机侧（UHT 会产出带 \Tp{\_{\Bk}Implementation} 后缀的实现函数约定），
蓝图可以整体覆写。代价要算清楚：每次跨界调用，UHT 生成的 exec
stub 要按反射元数据把参数逐个\emph{复制}进虚拟机栈帧——数值按值、
对象与结构体按引用装箱，返回值再抄回来；边界处不可内联，单次的
微秒级开销，放到``每帧每元素''的尺度上就变成肉眼可见的帧时间。

判定该把一段逻辑放哪一层，表~\ref{tab:ch10-placement} 给一个够用
的出发点。

\begin{table}[ht]
\centering
\caption{逻辑放 \CPP{} 还是蓝图的判定表}
\label{tab:ch10-placement}
\begin{tabular}{>{\raggedright\arraybackslash}p{0.30\textwidth}
                >{\raggedright\arraybackslash}p{0.26\textwidth}
                >{\raggedright\arraybackslash}p{0.27\textwidth}}
\toprule
\textbf{问题} & \textbf{倾向 \CPP{}} & \textbf{倾向蓝图} \\
\midrule
每秒执行上万次？ & 是，且考虑搬进数据容器批处理 & 否，或只当触发器 \\
要断点、单测、静态分析？ & 是 & 否 \\
改动者是策划、要求秒级可见？ & 否 & 是 \\
接口要不要跨版本兼容存档？ & 稳定签名写在 \CPP{}，蓝图只消费 & 拼装与演出留蓝图 \\
出错的爆炸半径？ & 崩进程，需重编重启 & 通常只毁一个实例，可回退资产 \\
\bottomrule
\end{tabular}
\end{table}

落地守则：跨边界的接口保持\emph{窄}——传句柄、枚举和少量标量，
不要传大结构体让装箱税翻倍；把``决定''留给 \CPP{}，把``编排''留给
蓝图；同一功能在两边各实现一份是腐化的开始，评审时优先合并。

\section{GC 与指针语义}
\label{sec:ch10-gc-pointers}

引擎对 UObject 采用可达性回收：从一组根（关卡、常驻对象、显式
钉住的引用）出发，\textbf{沿着反射属性表}标记，扫不到的就释放。
于是第一条铁律是：\emph{只有标了 \Tp{UPROPERTY} 的成员才是引用}。
一个指向 UObject 的裸指针，哪怕你天天用它，在 GC 眼里也不存在——
它既不延长对方寿命，也不阻止对方被回收；对象释放后指针变成悬垂
的随机值，读它不一定立刻崩，往往在几帧后以不相干的症状回报。

\begin{lstlisting}[language=C++, style=cppstyle,
  caption={成员声明片段：三种指针语义，三种 GC 命运},
  label={lst:ch10-gc-members}]
// Inventory.h —— 强、弱、软引用与反面教材
UCLASS()
class UInventory : public UActorComponent
{
    GENERATED_BODY()

    // 强引用：进可达性图，延长寿命。
    // 近年的推荐写法（早期以裸指针加 UPROPERTY 为主）
    UPROPERTY()
    TObjectPtr<UItem> Equipped;

    // 弱引用：不延长寿命，专治反向引用成环
    UPROPERTY()
    TWeakObjectPtr<AActor> LastToucher;
    // 用前判活：if (LastToucher.IsValid()) ...

    // 软引用：本质是资产路径，序列化不拖入资源，
    // 使用前按需异步加载（入口以官方文档为准）
    UPROPERTY()
    TSoftObjectPtr<UTexture2D> Icon;

    // 反面教材：裸指针不进属性表，寿命无人担保
    AActor* Tracked = nullptr;   // 随时可能悬垂
};
\end{lstlisting}

三种合法语义各有岗位。强引用指针类型用于``我拥有它''；弱引用用于
``我认识它但不养它''，判活后再解引用；软引用用于``我知道它在资产
库里的位置''，配合异步加载把大资源推迟到真正用到的关卡——开放
世界流送的常规组合就是``关卡里只放软引用，进视野前批量异步请求''，
异步入口的具体类型各版本改动很大，以官方文档为准。裸指针只允许
出现在栈上的瞬态传递，不允许作为成员跨帧存活。

\section{迭代代价：编译时长、分布式编译与打包}
\label{sec:ch10-build-cost}

选了``编译到原生''，就要还它的债：引擎加项目模块的初次全量编译以
小时计，改一个头文件就可能触发成百翻译单元重编。日常迭代靠三件事
撑着：官方的构建工具（增量判断加 UHT 前置）把重编范围压到模块级；
分布式编译与缓存把单机瓶颈摊到整个编译集群——\Lib{Unreal} 线近年
内置了官方加速器（早期以 \Tp{IncreBuild} 之名独立提供，现在归入
\Tp{Unreal Build Accelerator} 一类的官方组件，名称与启用方式以
文档为准），第三方编译缓存（ccache 一类）同样常见；以及 ch09
讲过的热重载机制：函数体内的改动可以近乎实时替换，但\emph{一旦
动到类布局、成员或注解}，反射表和生成代码都要重造，只能回到
全量编译。热补丁能力与
反射变更的这条分界线，是评估这类引擎迭代成本的核心。

打包路径上，内容先经烘焙（cook）把资产转成平台格式，再随所选
目标配置出包；Shipping 包没有编辑器、没有开发期校验，崩溃现场
要靠构建时留存的符号表还原——CI 流水线里``能出包''与``出包后
可诊断''是两个独立的验收项。

\begin{warning}[title={把 UObject 裸指针塞进容器：GC 看不见}]
最常见的事故形态：玩法程序员写了一个 \Tp{std::vector} 装
\Tp{AActor*} 裸指针，缓存``当前存活的敌人''，这个成员\emph{没有}
标 \Tp{UPROPERTY}——它压根进不了反射属性表。敌人死亡被 GC 回收
后，容器里留下的是一串悬垂地址；遍历调用虚函数时，崩的位置离
真正的赋值点隔着几万行。修正方向有二：成员改成带 \Tp{UPROPERTY}
的引擎容器（强引用，注意它同时钉住寿命），或改成弱引用容器并在
访问前判活。根因回到第~\ref{sec:ch10-type-lifecycle}节：反射元数据
是 UHT 按注解\emph{另行制造}的，\CPP{} 语言层面写了成员，不等于
在引擎的世界里声明了引用。
\end{warning}

\begin{bestpractice}[title={在这类工程里控制迭代成本的次序}]
\begin{itemize}
\item 动代码前先问``这是改函数体还是改结构''：前者可走热重载，
      后者按全量编译预算时间，两件事不要混在同一次提交里。
\item 暴露面从窄往宽长：先只加 \Tp{UPROPERTY}，确认序列化与面板
      行为正确，再考虑 \Tp{UFUNCTION} 跨蓝图边界；每多暴露一个
      节点，就多想一次跨界装箱的调用频率。
\item 新增字段前先想旧资产：反射改动意味着已存盘资产的兼容问题，
      引擎的升级通道能兜一部分，兜不住的部分要靠迁移脚本——这条
      成本在第一部分的求解器输入档版本管理里已经付过一遍。
\end{itemize}
\end{bestpractice}

% ============================================================
% ---- frag_ch11.tex ----
\chapter{Unity：原生插件、IL2CPP、Burst 与 Job 系统}
\label{ch:ch11-unity}

\DomEngine{} 侧三个引擎里，Unity 的结构最容易被误解：它不是``一个 \CL/{\Bk}\CPP{} 程序外加脚本胶水''，
而是清清楚楚的两半。原生那一半是编译好的 \CL/{\Bk}\CPP{} 运行时时（渲染、物理、序列化、资源管线），
托管那一半是跑在虚拟器上的 \Lang{C\#} 逻辑层。本章不讲 API 大全，只回答两个问题：
这两半的\textbf{边界在哪里}，以及\textbf{跨过去要付什么}。API 名只列有把握的一族
（\Tp{DllImport}、\Tp{IL2CPP}、\Tp{AOT}、\Tp{Burst}、\Tp{IJobParallelFor}、\Tp{NativeArray}），
其余按调用形状描述，具体签名与参数个数一律以官方文档为准。

\section{托管的一半与原生的一半}
\label{sec:ch11-two-halves}

编辑器进程本身是原生代码，但它承载的脚本域（Assembly/Domain）是托管的：你在
\Tp{Assets/Scripts/} 下写的每一个 \Lang{C\#} 类型都由虚拟器加载、由垃圾回收器管理，
访问引擎对象时通过一层绑定调到原生侧的句柄上。这个划分解释了三件日常体验：

\begin{itemize}
  \item \textbf{改脚本后要等两秒。}编辑器的默认行为是\textbf{域重载}（domain reload）：
        脚本集重新编译后，旧的托管域被整体卸载、新域被重新加载，所有 \Tp{static} 字段归零、
        所有托管对象重建。时间花在卸载、重新 JIT、重新跑一遍 \Tp{InitializeOnLoad} 一类回调上，
        而不是花在磁盘读文件上。
  \item \textbf{播放模式进入慢、状态留不住。}域重载同样发生在进入/退出播放模式时，
        于是静态缓存跨不过一次播放。
  \item \textbf{可以缓解，但有代价。}项目设置里提供``进入播放模式选项''（Enter Play Mode
        Settings），可以选择\textbf{不做}域重载与场景重载，代价是任何依赖``静态字段一开始是零''
        的代码都可能读到上一次播放残留的值。较新的虚拟器（\Lang{.NET} 一侧的可卸载程序集/
        可回收域）提供了更细粒度的做法，Unity 何时、以何种形式采用，以官方文档为准。
\end{itemize}

\begin{note}[title={域重载是原生插件状态的隐形杀手}]
原生插件的 \Tp{static} 全局\textbf{不受}托管域重载影响：域重新加载后，托管侧的句柄对象
已经变成新的、甚至已经消失，而原生侧的堆还活着、还占着上一次的分配。这就是``编辑器里跑第二次
必崩、真机上一切正常''这一类事故的经典成因。规则只有一条：把原生资源的拥有权（owner）
明确写死在托管侧，并用 \Tp{IDisposable} 或最终化器成对释放，绝不依赖进程级全局的初始化状态。
\end{note}

\section{C/C++ 进入 Unity 的唯一入口：本地插件}
\label{sec:ch11-native-plugin}

把已有的 \CL/{\Bk}\CPP{} 求解器、编解码库或厂商 SDK 塞进 Unity，正规路径只有一条：
\textbf{一个 C ABI 导出层 + 托管侧的 \Tp{DllImport} 声明}。C++ 的类、模板与异常不能跨过这个边界，
能过的只有整数、浮点、指针、结构体数组和不透明句柄。清单 \ref{lst:ch11-plugin-c} 是插件侧，
清单 \ref{lst:ch11-plugin-cs} 是托管侧，两边合起来就是完整的``入口形状''。

\begin{lstlisting}[language=C, style=cstyle, caption={本地插件的 C 导出层：不透明句柄加三段式生命周期（create / step / destroy）}, label={lst:ch11-plugin-c}, framexleftmargin=12pt, framexrightmargin=12pt]

// MySolver.c —— 只导出 C ABI，绝不导出 C++ 修饰名
#if defined(_WIN32)
  #define EXPORT __declspec(dllexport)
  #define CALL   __cdecl
#else
  #define EXPORT __attribute__((visibility("default")))
  #define CALL
#endif

typedef struct SolverHandle SolverHandle;   // 不透明句柄：托管侧只见 IntPtr

EXPORT SolverHandle* CALL solver_create(int body_count, int dof);
EXPORT void CALL solver_destroy(SolverHandle* h);

// data 由调用方保证在本次调用期间存活且不被释放
EXPORT int CALL solver_step(SolverHandle* h, float dt,
                            const float* in, float* out, int n);
\end{lstlisting}

\begin{lstlisting}[style=cppstyle, caption={托管侧的声明清单：DllImport 的库名写法决定加载方式}, label={lst:ch11-plugin-cs}, framexleftmargin=12pt, framexrightmargin=12pt]

// NativeSolver.cs —— 托管侧声明（C# 无专属 style，此处复用 cppstyle）
using System;
using System.Runtime.InteropServices;

public static class NativeSolver
{
    // "__Internal"：符号已在宿主可执行文件里（iOS、WebGL、部分编辑器构建）
    // 其它平台写不带后缀/不带路径的库名，由加载器按平台目录解析
    const string LibName = "__Internal";

    [DllImport(LibName, CallingConvention = CallingConvention.Cdecl)]
    public static extern IntPtr solver_create(int bodyCount, int dof);

    [DllImport(LibName, CallingConvention = CallingConvention.Cdecl)]
    public static extern void solver_destroy(IntPtr handle);

    [DllImport(LibName, CallingConvention = CallingConvention.Cdecl)]
    public static extern int solver_step(IntPtr handle, float dt,
                        float[] input, float[] output, int n);
}
\end{lstlisting}

\subsection{平台相关打包与``编辑器能跑真机崩''}
\label{sec:ch11-plugin-packaging}

原生库按平台放在 \Tp{Plugins/{\Bk}\Bk<平台>/} 下，编辑器与真机命中不同的子目录，这就是行为
不一致的根源：加载的是两份不同的二进制、链接的是两份不同的 C 运行库。常见对应关系见
表 \ref{tab:ch11-platform}。

\begin{table}[htbp]
\centering
\caption{本地插件在各目标上的库形态与链接时机}
\label{tab:ch11-platform}
\begin{tabular}{>{\raggedright\arraybackslash}p{0.16\textwidth} >{\raggedright\arraybackslash}p{0.32\textwidth} >{\raggedright\arraybackslash}p{0.32\textwidth}}
\toprule
\textbf{目标} & \textbf{库形态} & \textbf{链接/加载时机} \\
\midrule
编辑器 & 与主机同架构的 \Tp{dll}、\Tp{so} 或 \Tp{dylib}，架构不符则根本加载不到 & 首次调用时延迟加载；与编辑器自带运行库冲突最易出问题 \\
Android & 放在带 ABI 名的目录（\Tp{arm64-v8a} 一类）里的 \Tp{so} & 打包进 APK/AAB，运行时由系统加载器装载 \\
iOS & 静态库（\Tp{a}），或需要额外链接参数的框架 & 构建阶段与 \Tp{IL2CPP} 产物一起链接，故常配 \Tp{\_\_Internal} \\
WebGL & 需配合一个 \Tp{jslib} 声明的中转库 & 没有进程外线程与原生内存的自由领地，多数插件要重写 \\
\bottomrule
\end{tabular}
\end{table}

\section{数据跨越边界的代价}
\label{sec:ch11-marshalling}

调用一次 \Tp{DllImport} 方法的成本远高于一次普通 \Lang{C\#} 调用：参数要\textbf{封送}
（marshalling），字符串要编码转换，数组要么逐元素拷贝要么整体钉住。省钱的关键概念是
\textbf{blittable}：布局在托管与非托管两侧字节一致的类型（整数、浮点、
\Tp{Layout{\Bk}Kind.{\Bk}Sequential} 的结构体数组）可以整块传，其余一律逐个转换。

\begin{lstlisting}[style=cppstyle, caption={两种正确的传数据写法：钉住托管数组，与直接交出 NativeArray 的裸指针}, label={lst:ch11-marshal}, framexleftmargin=12pt, framexrightmargin=12pt]

// A. 托管数组：元素必须 blittable，且调用期间要钉住
public static unsafe void StepPinned(float[] src, float[] dst,
                                     IntPtr h, float dt)
{
    GCHandle pin = GCHandle.Alloc(src, GCHandleType.Pinned);
    try
    {
        float* srcPtr = (float*)pin.AddrOfPinnedObject();
        solver_step(h, dt, srcPtr, dst, src.Length);
    }
    finally { pin.Free(); }   // Free 之后原生侧不得再持有该地址
}

// B. NativeArray：内存在非托管堆上，GC 不会搬动它
//    危险不是"移动"而是"释放"：作业还在跑就 Dispose 即野指针
public static void StepNative(NativeArray<float> src,
                              NativeArray<float> dst,
                              IntPtr h, float dt)
{
    // Unity.Collections 的 unsafe 取值扩展返回一个 IntPtr，
    // 具体名字与拼写以官方文档为准；这里按形状调用
    JobHandle pending = UnsafeStep(dst, src, h, dt);
    pending.Complete();       // 只有回来了，调用方才允许 Dispose
}
\end{lstlisting}

生命周期陷阱的三种典型形态：把 \Tp{NativeArray} 交给原生线程后立即 \Tp{Dispose}；
在作业调度\textbf{期间}修改托管数组的长度（GC 压缩会移动未钉住的对象）；
让原生侧回调持有 \Lang{C\#} 对象的地址。正确的做法只有两种——要么整段数据是
非托管分配、由原生侧拥有；要么托管侧用 \Tp{GCHandle}（\Tp{Alloc} 加 \Tp{Pinned}）
交出句柄并在明确的所有权转移点后 \Tp{Free}。

把边界成本压下来的手段按有效性排序大致是：

\begin{enumerate}
  \item \textbf{减少调用次数}，而不是减少参数个数。一次 \Tp{DllImport} 带 10 万个元素，
        远好过 10 万次各带 1 个元素；接口设计成\textbf{批处理形状}（句柄加缓冲区加长度）
        是最有效的一招。
  \item \textbf{只用 blittable 布局}。结构体加 \Tp{LayoutKind.Sequential}、
        不出现 \Tp{string} 与嵌套引用类型，数组才能整块传。
  \item \textbf{字符串只在建对象时过界}，热路径上传 \Tp{IntPtr} 加长度或干脆用整数枚举。
  \item \textbf{回调反方向也要批量化}：原生侧逐帧回调进托管代码同样要封送，
        改成写进共享缓冲、由托管侧一次性读走。
\end{enumerate}

\section{IL2CPP 与 AOT：iOS 上没有 JIT}
\label{sec:ch11-il2cpp}

\Tp{IL2CPP} 是\textbf{把中间语言先翻译成 \CPP{} 源码、再交给平台工具链编译}的一条后端路径。
它存在的首要理由是 iOS：系统禁止进程写入可执行内存，运行时 JIT 无处容身，只有
\Tp{AOT}（提前编译）产物能上架。翻译带来的连锁约束是：代码在编译期就被固定，
运行期无法生成新类型或新的泛型实例化，未被引用的代码会被裁剪。于是出现一对配套机制：
\textbf{代码裁剪}（managed stripping）按可达性删掉程序集，而反射、
\Tp{UnityEngine.UI} 一类只在运行时按名字用到的类型会被误删，
需要用 \Tp{link.xml} 声明保留，或在类型与成员上打 \Tp{Preserve} 一类的标注。

\begin{lstlisting}[style=yamlstyle, caption={link.xml 的保留声明形状（标签名与层级以官方文档为准）}, label={lst:ch11-linkxml}, framexleftmargin=12pt, framexrightmargin=12pt]

# 只要"运行时按名字用到"，就要在这里显式点名
<link>
  <assembly fullname="MyGame.Logic">
    <type fullname="MyGame.Logic.SolverConfig" preserve="all"/>
  </assembly>
</link>
\end{lstlisting}

\begin{table}[htbp]
\centering
\caption{四种运行环境的能力对比}
\label{tab:ch11-runtime}
\begin{tabular}{>{\raggedright\arraybackslash}p{0.14\textwidth} >{\raggedright\arraybackslash}p{0.17\textwidth} >{\raggedright\arraybackslash}p{0.17\textwidth} >{\raggedright\arraybackslash}p{0.17\textwidth} >{\raggedright\arraybackslash}p{0.17\textwidth}}
\toprule
\textbf{维度} & \textbf{编辑器（虚拟器）} & \textbf{Android} & \textbf{iOS} & \textbf{WebGL} \\
\midrule
执行方式 & 即时编译，改脚本热重载 & 中间语言或 \Tp{IL2CPP} 二选一 & 只能 \Tp{IL2CPP}，\Tp{AOT} & 只能 \Tp{IL2CPP}，产出到 \Tp{wasm} \\
运行期造代码 & 可以 & 视后端 & 不可以 & 不可以 \\
泛型值类型 & 现场实例化 & 视后端 & 编译期须可见 & 编译期须可见 \\
代码裁剪 & 基本不裁 & 可选级别 & 可选级别，误删高发 & 可选级别 \\
原生插件 & 主机架构 \Tp{dll}/{\Bk}\Tp{so} & \Tp{so} 随 ABI 打包 & 静态库进链接 & 需 \Tp{jslib} 中转 \\
多线程 & 无限制 & 无限制 & 无限制 & 受限，按平台规则 \\
\bottomrule
\end{tabular}
\end{table}

实践上的差异只有一条值得反复验证：\textbf{编辑器里跑得通的反射，\Tp{IL2CPP} 目标上不一定活着}。
把裁剪级别开到最高、在真机跑一次冒烟测试，比在编辑器里测一百遍都有效。

\section{Burst 与 Job：什么时候反而不用写插件}
\label{sec:ch11-burst-job}

\Tp{Burst} 是 \Lang{C\#} 侧的编译器后端：把标注过的作业代码编成针对目标指令集优化的
本机代码，能吃上 SIMD。它能编译的前提正是``没有托管引用''——只用值类型、
\Tp{NativeArray}、\Tp{ReadonlyArrayView}（数学库里的只读视图，形状与 \Tp{NativeArray} 相同）
一类的数据。换句话说，Burst 与本地插件解决的是同一个问题（让数学热点跑在原生速度上），
差别在谁拥有那段代码。

\begin{lstlisting}[style=cppstyle, caption={Burst 编译的并行作业：位置数组的显式积分}, label={lst:ch11-burst-job}, framexleftmargin=12pt, framexrightmargin=12pt]

// VerletJob.cs —— float3 来自数学库，全部是值类型
[BurstCompile]
public struct VerletJob : IJobParallelFor
{
    [ReadOnly] public NativeArray<float3> position;
    [ReadOnly] public NativeArray<float3> previous;
    [ReadOnly] public float3 gravity;
    [ReadOnly] public float dtSquare;
    public NativeArray<float3> output;

    public void Execute(int i)
    {
        float3 p = position[i];
        output[i] = p + (p - previous[i]) + gravity * dtSquare;
    }
}

// 调度：第二个实参是线程内的批次粒度，64 是常见起点
var job = new VerletJob { position = pos, previous = prev,
                          gravity = g, dtSquare = dt * dt };
job.output = next;
JobHandle pending = job.Schedule(count, 64);
pending.Complete();     // 依赖链交给调用方时，这里只等自己这一段
\end{lstlisting}

\begin{compare}[title={Burst 作业还是 C/C++ 插件}]
选 Burst 的情形：算法自己写、数据本来就是 \Tp{NativeArray} 上的连续数组、要吃到 Jobs 系统的
自动依赖检查与多核调度，且不愿意维护四套平台的二进制。选本地插件的情形：算法已经是一个
成熟的 \CL/{\Bk}\CPP{} 库（第三方 SDK、公司内已有的求解器）、需要 C 运行时或系统调用、
要用调试器在原生侧单步。混合形态最常见——数据用 \Tp{NativeArray} 组织与复用，
只在真正的外库边界上做一次 \Tp{DllImport}。判断依据不是性能上限，而是
\textbf{谁愿意承担多平台构建与两份生命周期管理的长期成本}。
\end{compare}

\begin{warning}[title={在插件里长期持有托管对象指针}]
最贵的一类事故：把 \Tp{GCHandle} 之外的裸指针（\Tp{ToPointer} 或 \Tp{AddrOfPinnedObject}
的结果）存进原生侧的全局表，然后指望它下次还在。托管堆的压缩、域重载、
播放模式退出都会让它失效，而且失效是静默的——直到某次写入踩坏别的对象。
三条硬规矩：跨边界存活超过一次调用的只有\textbf{不透明句柄}（原生自己 \Tp{malloc} 的东西）；
托管对象需要长期引用时用 \Tp{GCHandle.Alloc(obj)} 这类\textbf{句柄}而非地址，并在明确时点
\Tp{Free}；任何``编辑器里复现、真机不复现''的崩溃，先怀疑指针跨越了上面这两条线。
\end{warning}

% ============================================================
% ---- frag_ch12.tex ----
\chapter{Godot：GDExtension 与 godot-cpp}
\label{ch:ch12-godot-ext}

\DomEngine{} 侧最后一个引擎走的是另一条路线：\Lib{Godot} 是开源的，扩展它不需要逆向、
不需要猜打包规则，而是官方提供一条\textbf{稳定的 ABI 接口}——\Tp{GDExtension}，
配上官方绑定层 \Lib{godot-cpp}。本章的价值不只在 Godot 本身：它是``如何设计一个
能被第三方长期依赖的二进制接口''的公开样本，跟本册 \DomCAE{} 侧的求解器扩展是同一套
工程问题。类名与函数名只写有把握的一族（\Tp{GDExtension}、\Tp{ClassDB}、
\Tp{register\_types}、\Tp{Variant}、\Tp{PackedFloat64Array}、\Tp{.gdextension}），
其余按形状描述，精确拼写与参数个数以官方模板和文档为准。

\section{从 GDScript 到原生：什么时候值得写扩展}
\label{sec:ch12-why-native}

\Lang{GDScript} 跑在引擎自带的虚拟机上：语法为场景树量身定做、与 \Tp{Node}/{\Bk}\Tp{Resource}
的绑定零成本、改完即热重载，代价是数值密集循环明显慢于编译型语言。值得下沉到原生的
判断依据只有三类：\textbf{热点循环}（粒子、刚体求解、寻路、图像核）、\textbf{已有 \CL/{\Bk}\CPP{} 库}
（厂商 SDK、编解码器、公司内求解器）、\textbf{需要系统能力}（共享内存、原始线程、
不受沙箱允许的调用）。其余情况先写 \Lang{GDScript}，用性能剖析确定瓶颈再说。

\Tp{GDNative}（3.x 时代）到 \Tp{GDExtension}（4.x）的变化不是改名，而是接口责任的转移：

\begin{itemize}
  \item \Tp{GDNative} 要求扩展对着\textbf{某一个引擎构建}的接口头文件编译，二进制与
        引擎版本近乎一一绑定，接口换一次就得重编一次，编辑器还要重启。
  \item \Tp{GDExtension} 把接口层做成了\textbf{带兼容性区间的稳定 ABI}：描述文件里声明
        本库适用的引擎版本区间，运行时由引擎提供符号获取函数，扩展按需取自己要的调用。
        同一份库能在同一大版本内的多个构建上加载，热重载与调试器附加也由接口层支持。
\end{itemize}

\begin{table}[htbp]
\centering
\caption{GDNative 与 GDExtension 的接口责任对比}
\label{tab:ch12-abigen}
\begin{tabular}{>{\raggedright\arraybackslash}p{0.16\textwidth} >{\raggedright\arraybackslash}p{0.32\textwidth} >{\raggedright\arraybackslash}p{0.32\textwidth}}
\toprule
\textbf{维度} & \textbf{GDNative（3.x）} & \textbf{GDExtension（4.x）} \\
\midrule
绑定对象 & 特定编辑器构建的头文件 & 运行时符号获取加版本区间声明 \\
版本兼容 & 近乎一对一，升级即重编 & 大版本内一份库可用，区间写在描述文件 \\
加载方式 & 常需重启编辑器 & 支持重载，开发期改完再看 \\
接口面 & 直接暴露的函数表 & 按层级注册的初始化回调 \\
\bottomrule
\end{tabular}
\end{table}

\section{注册链条：描述文件、入口符号与 register types}
\label{sec:ch12-registration}

一个原生扩展进入引擎要过三道门。第一道是 \Tp{.gdextension} \textbf{描述文件}，
它告诉引擎``入口符号叫什么、每个平台的库文件在哪、兼容哪些引擎版本''（清单
\ref{lst:ch12-gdext}）。第二道是入口符号：一个 \CL{} 链接约定的导出函数，
收到引擎的函数表与自身回调表，把两侧对接起来。第三道是 \Tp{register\_types}：
把自定义类登记进 \Tp{ClassDB}，绑定方法、属性与信号。

\begin{lstlisting}[style=yamlstyle, caption={particle solver.gdextension 的描述形状（键名以官方模板生成的为准）}, label={lst:ch12-gdext}, framexleftmargin=12pt, framexrightmargin=12pt]

# 三段：兼容性/入口、各平台库路径、可选依赖
config_version=5
compatibility_minimum="4.1"
compatibility_maximum="4.x"
entry_symbol="library_entry"

[libraries]
# 键名由 平台 . 构建 . 架构 三段拼成，值相对于本文件所在目录
windows.editor.x86_64="bin/solver.windows.editor.x86_64.dll"
windows.release.x86_64="bin/solver.windows.release.x86_64.dll"
linux.editor.x86_64="bin/solver.linux.editor.x86_64.so"
macos.editor.universal="bin/solver.macos.editor.universal.dylib"
\end{lstlisting}

\Tp{initialize} 与 \Tp{deinitialize} 只负责把生命周期接上，具体登记动作在
\Tp{register\_{\Bk}types.cpp} 里；登记还分\textbf{层级}（核心、场景、编辑器），
类要出现在场景树与编辑器面板上就得注册在正确的层级。清单 \ref{lst:ch12-register}
是最小骨架。

\begin{lstlisting}[language=C++, style=cppstyle, caption={register types 的注册骨架：入口对接与 ClassDB 登记}, label={lst:ch12-register}, framexleftmargin=12pt, framexrightmargin=12pt]

// register_types.cpp —— godot-cpp 模板给出同名文件，此处为形状骨架
#include <godot_cpp/core/class_db.hpp>
#include "particle_solver.hpp"

using namespace godot;

// 初始化按层级分段调用：要进场景树的类注册在场景层级
void initialize_solver_module(ModuleInitializationLevel p_level)
{
    if (p_level != MODULE_INITIALIZATION_LEVEL_SCENE) { return; }
    ClassDB::register_class<ParticleSolver>();   // 编辑器与脚本都能看见
}

void uninitialize_solver_module(ModuleInitializationLevel p_level)
{
    if (p_level != MODULE_INITIALIZATION_LEVEL_SCENE) { return; }
    // 类由引擎注销；这里只清自己分配的全局状态
}

extern "C"
{
    // 描述文件的 entry_symbol 指向这里。下面这个函数的宏名、参数个数
    // 与返回类型请以官方模板为准，此处只表达"三个实参对接两侧"的形状
    GDE_EXPORT GDEExtensionBool solver_library_entry(
        GDEInterfaceGetProcAddress p_get_proc_address,
        const GDELibraryCallbacks p_library,
        GDEObjectPtr p_token)
    {
        return install_entry(p_get_proc_address, p_library, p_token,
                             initialize_solver_module,
                             uninitialize_solver_module);
    }
}
\end{lstlisting}

\section{类型桥：String、Array、Variant 与 Inspector}
\label{sec:ch12-type-bridge}

跨接口的值分两等。一类是\textbf{引擎值类型}的 C++ 表示：\Tp{String}、\Tp{StringName}、
\Tp{Array}、\Tp{Dictionary}、\Tp{Variant}，以及打包数组如 \Tp{PackedFloat64Array}——
它们是引用计数的包装，内部数据由引擎的内存分配器管理，\textbf{不要}按 \CL{} 结构体
去 memcpy 或跨线程裸传。一类是反射暴露出去的成员：\Tp{ClassDB} 里登记的属性会
自动出现在编辑器 Inspector 与 \Lang{GDScript} 的自动补全中，方法用
\Tp{D\_METHOD} 加名字与参数名绑定，信号用信号描述信息登记。

\begin{lstlisting}[language=C++, style=cppstyle, caption={把一个类暴露给脚本与 Inspector：方法、属性、信号三件事}, label={lst:ch12-class}, framexleftmargin=12pt, framexrightmargin=12pt]

// particle_solver.hpp / .cpp —— 反射宏与 bind_methods 是固定套路
class ParticleSolver : public RefCounted {
    GDCLASS(ParticleSolver, RefCounted)   // 反射样板，模板里都有

    double stiffness_ = 1.0;
    PackedFloat64Array positions_;        // 与脚本共享的打包数组

  protected:
    static void bind_methods();           // 由 GDCLASS 宏约定调用点

  public:
    void set_stiffness(double p_v) { stiffness_ = p_v; }
    double get_stiffness() const { return stiffness_; }
    void step(double p_delta);
    PackedFloat64Array get_positions() const { return positions_; }
};

void ParticleSolver::bind_methods()
{
    ClassDB::bind_method(D_METHOD("set_stiffness", "value"),
                         &ParticleSolver::set_stiffness);
    ClassDB::bind_method(D_METHOD("get_stiffness"),
                         &ParticleSolver::get_stiffness);
    // ADD_PROPERTY 的实参是属性信息加读写方法名（形状，细节看模板）
    ADD_PROPERTY(PropertyInfo(Variant::FLOAT, "stiffness"),
                 "set_stiffness", "get_stiffness");
    ADD_SIGNAL(MethodInfo("settled"));
}
\end{lstlisting}

暴露之后脚本侧能做什么是关键：属性可以写在场景文件里、能在 Inspector 上拖、
能被 \Tp{set} 与 \Tp{get} 反射访问；方法直接调用；信号用 \Tp{connect} 连到
\Lang{GDScript} 函数。基类选择决定身份——继承 \Tp{Node} 就进场景树、
继承 \Tp{Resource} 就变成可存盘的资产、继承 \Tp{RefCounted} 就交给引用计数。
序列化与 Inspector 是同一件事的两面：\textbf{只有被登记为属性的数据才会进
\Tp{.tscn}/{\Bk}\Tp{.tres}}，私有成员不进去；因此属性名一旦发布就是数据格式的一部分，
改名会让旧场景里的值读不回来（引擎有属性重命名机制，但需要显式声明，以文档为准）。

从 \Tp{Variant} 里取值是另一个容易踩坑的地方。它是引擎的\textbf{万能值容器}，
C++ 侧的读法有两层：一层是构造/转换（拿不到就报错或回退默认值），
一层是显式判别类型再取值；数值在 \Tp{Variant} 里以引擎内部宽度存储，
不要假设它保留你放进去的那个 \CPP{} 类型。写接口时的建议顺序是：
\textbf{能用具体类型（\Tp{int64\_t}、\Tp{double}、\Tp{String}）就不要用 \Tp{Variant}，
能用打包数组就不要用 \Tp{Array}}——前者一次转换、后者一次内存拷贝，
而 \Tp{Variant} 套 \Tp{Array} 套 \Tp{Dictionary} 的读法会把转换成本摊到每个元素上。

\section{构建与分发：多平台库与 zip 目录约定}
\label{sec:ch12-build}

官方 \Lib{godot-cpp} 模板默认用 SCons 描述构建，用 CMake 同样可行（清单
\ref{lst:ch12-cmake}）；要产出的是一组按``平台.构建.架构''命名的库，路径与描述文件
里的键一一对应。这里有一个容易被忽略的分岔：\textbf{编辑器模板与导出模板是两种构建}。
在编辑器里跑扩展需要 editor 构建（带编辑器侧的类与回调）；玩家包用引擎导出模板
（export templates）打包，扩展必须再出一份 release 构建，否则导出后类直接加载失败。

\begin{lstlisting}[language=CMakeLD, style=cmakestyle, caption={CMake 的最小构建形状：静态绑定层加共享扩展库}, label={lst:ch12-cmake}, framexleftmargin=12pt, framexrightmargin=12pt]

# 目标名与库文件后缀要和 .gdextension 里的键对得上
cmake_minimum_required(VERSION 3.20)
project(particle_solver CXX)
set(CMAKE_CXX_STANDARD 17)

add_subdirectory(godot-cpp)        # 绑定层，编一次可复用
add_library(solver SHARED
    src/register_types.cpp
    src/particle_solver.cpp)
target_link_libraries(solver PRIVATE godot_cpp)   # 链接目标名以模板为准
set_target_properties(solver PROPERTIES
    PREFIX ""                      # 去掉 lib 前缀，与描述文件路径一致
    OUTPUT_NAME "solver.linux.editor.x86_64")
\end{lstlisting}

\begin{lstlisting}[style=shellstyle, caption={发布包目录约定：一个 zip 里 addons 下带描述文件与多平台库}, label={lst:ch12-package}, framexleftmargin=12pt, framexrightmargin=12pt]

# 用户把 addons/ 整目录解到项目根，Godot 自动识别插件
D=dist/addons/particle_solver
mkdir -p $D/bin/linux $D/bin/windows $D/bin/macos
cp build/linux.x86_64.so        $D/bin/linux/
cp build/windows.editor.x86.dll $D/bin/windows/
cp build/macos.universal.dylib  $D/bin/macos/
cp particle_solver.gdextension  $D/
cp -r gdscript/                 $D/          # 脚本侧包装与示例场景
(cd dist && zip -r particle_solver.zip addons)
\end{lstlisting}

两套构建体系各有取舍，选择依据是团队已有的习惯：

\begin{itemize}
  \item \textbf{SCons}（\Lib{godot-cpp} 模板默认）：一条命令按平台自动出文件名与
        编译开关，交叉编译与 release/editor 两种 target 已经写好，改动少。
  \item \textbf{CMake}：接进已有的求解器工程与 CI 更顺，能用现成的工具链文件与
        预编译依赖，代价是输出文件名、符号可见性与运行库设置要自己逐平台对齐。
  \item 共同点：绑定层只编一次当静态库用，扩展库保持\textbf{单一入口符号}；
        发布前用同版本引擎的命令行模式加载一次，确认注册日志没有报错。
\end{itemize}

\begin{note}[title={分发时最常漏的一件事}]
只交了 editor 构建的库、描述文件里 release 键指向不存在的文件、
或者库路径按绝对路径写——三者都会让``编辑器里一切正常，导出后类变成未知类型''。
发布前在空项目里做一次导出冒烟，比在开发项目里点一百次可靠。
\end{note}

\section{与 CAE 侧的共通处：小接口加反射元数据}
\label{sec:ch12-cae-crossover}

把本章和本册 \DomCAE{} 侧那一族的扩展模式并排看，结构是一样的：

\begin{itemize}
  \item \textbf{参数走字典/属性，不走代码。}\Tp{Variant} 加 \Tp{ClassDB} 的属性反射，
        让编辑器与脚本按名字读写求解器参数；\DomCAE{} 侧的对应物是``参数字典或
        关键字表加一份由求解器解释的配置文件''。两侧都是\textbf{数据描述驱动}，
        代码只负责实现。
  \item \textbf{实现要编进宿主或按其规则装载。}自定义资源进 Godot 需要扩展库与
        引擎同版本；自定义求解器进 \DomCAE{} 宿主要编译进可执行文件或走固定的
        插件 ABI。两者的版本矩阵都是分发成本的大头。
  \item \textbf{接口面决定长期成本。}\Tp{GDExtension} 之所以把 ABI 层做成
        带兼容性区间的形式，就是把``第三方可依赖的面积''压到最小；
        本册前面每一章的边界层设计都是同一个方向。
\end{itemize}

\begin{bestpractice}[title={扩展接口越小越好}]
把这条当作验收标准：一次发布只暴露\textbf{一个入口符号}、\textbf{一组必要的类与方法}、
\textbf{一份可枚举的属性表}。具体的取舍是——能不暴露全局函数就不暴露（走类的静态方法）；
参数用引擎值类型或打包数组，别自定义一套跨边界的结构体布局；数据密集的用
\Tp{PackedFloat64Array} 一类连续数组一次性交出，而不是逐元素调用；
版本的进出都写在描述文件里，不要靠``猜引擎版本号''。接口每多一个符号，
就多一份要跨平台、跨版本维护的兼容性债务——这一条对 \DomEngine{} 和 \DomCAE{} 都成立。
\end{bestpractice}

% ============================================================
% ---- frag_ch13.tex ----
\chapter{物理与中间件：PhysX、Havok、Bullet 与 Jolt 的 C++ 接口}
\label{ch:ch13-physics}

实时物理是引擎里最像``求解器''的一块：输入是一堆刚体、关节与接触，输出是一帧内
被近似满足的运动方程。它和本册第一部分讲的离线求解器共享同一套词汇（时间步、
迭代、收敛、稳定性），却多了一条硬约束：一帧只有 16.7 ms 或 33.3 ms，物理只是
被切出来的一小块。所以这一章是两部分的接缝——把时间步与收敛接到帧预算上，
再接到中间件的 C++ 接口形态上。

\begin{guideline}[title={本章口径}]
\begin{itemize}
  \item 只写有把握的类型名与调用形状：PhysX 的 \Tp{PxScene}、\Tp{PxMaterial}、
        \Tp{Px{\Bk}Rigid{\Bk}Actor} 一族，Bullet 的 \Tp{bt{\Bk}Discrete{\Bk}Dynamics{\Bk}World}，
        Jolt 的 \Tp{Physics{\Bk}System}。
  \item Havok 只讨论授权与集成形态，不给类名与函数签名。
  \item 参数个数、枚举拼写与许可条款都随版本演进，落地前以官方文档为准。
\end{itemize}
\end{guideline}

\section{三个不可调和项：正确性、稳定性与帧预算}
\label{sec:ch13-triangle}

离线求解器的目标函数是``在给定残差容差内把问题解对''，运行时长不进目标函数；
实时物理的目标函数是``在固定墙上时间里给出一个不被玩家看出破绽的答案''。
三条要求两两冲突：

\begin{itemize}
  \item \textbf{正确性}：接触不穿透、约束不被违反、动量与能量不发生非物理漂移；
        靠更小的时间步、更多接触点、更严格的求解获得。
  \item \textbf{稳定性}：不发散、不抖动、堆叠不``呼吸''、对参数不敏感；
        靠更多迭代、软约束参数与合理的 mass ratio 获得。
  \item \textbf{帧预算}：物理耗时必须有上界，且上界近似与实体数量无关；它要求把
        前两项的手段（步长、迭代次数、接触对数）从``直到收敛''改成``固定上限''。
\end{itemize}

一条粗糙但有用的账：记 $S$ 为本帧 substep 数、$B$ 为宽相耗时、$I$ 为每个
substep 的求解迭代次数、$N$ 为窄相耗时、$Z$ 为一次迭代扫全部约束的耗时，则

\[
  T \approx S \times (B + I \times (N + Z))
\]

帧预算守住的是 $T$ 的上界，能调的只有 $S$、$I$ 与接触规模，三种手段的性价比
完全不同：加大 $I$ 是最便宜的稳定性手段（宽相与窄相不做两遍，只多扫几次约束），
但收益不线性，从 4 到 8 有意义，从 16 到 32 只是把已经看得过去的堆叠再压一压；
加大 $S$（把固定步长切小）是正确性手段，快速物体穿薄墙、弹簧失稳基本都靠它，
代价按 $1/\Tp{dt}$ 线性放大，还要付``同一帧里世界被推进两次''的逻辑副作用；
减接触规模（睡眠、island 切分、碰撞分层、把不参与物理的几何剔出场景）
是唯一同时降 $B$、$N$、$Z$ 且不牺牲精度的手段。

关键概念差异在这里：离线求解器有残差，可以``解到 $10^{-6}$ 就停''，实时求解器没有
可上报的全局残差，只做固定次数的近似扫描。所以实时物理交付的不是``收敛的解''，
而是``一个在 $I$ 次扫描后可行度足够高的速度场与接触冲量''，剩下的违反量靠软约束
参数、位置修正与人工限幅去掩盖。把它当成``迭代到收敛为止''的求解器来调，
是唯一一定会超预算的思路。

\subsection{固定步长与累加器}
\label{sec:ch13-fixed-step}

渲染帧率是浮动的，物理步长不是。所有中间件的步进入口都是同一个形状：一次调用
推进一个固定的 $\Tp{dt}$（常见 1/60 或 1/120 s），一帧内调用几次由你自己决定。
典型做法是累加器：把本帧真实时间加进去，够一个步长就步进一次，并给本帧的
步进次数设上限。上限没设好会出``死亡螺旋''：物理超时 $\rightarrow$ 帧时间变长
$\rightarrow$ 追更需要更多 substep $\rightarrow$ 更超时。追赶上限、以及放弃追赶
时怎么处理插值残差（渲染侧需要一个 $[0,1)$ 的插值因子），
清单 \ref{lst:ch13-tick-budget} 里给全了。

\section{中间件作为``只暴露 C++ 接口的库''}
\label{sec:ch13-middleware-shape}

把物理中间件当第三方 C++ 库看，它的使用形状高度一致，四家都是同一段谱：

\begin{enumerate}
  \item \textbf{先给运行时}：多数 SDK 要求你提供分配器、日志与断言回调，
        或从它自带的默认实现里选一个。这一步决定后面所有内存落在哪个堆上，
        也决定主机平台能不能过审。
  \item \textbf{再造库实例与描述符}：PhysX 是 Foundation、描述符、Physics 实例
        这样一层层下来；Bullet 与 Jolt 直接把部件当构造参数传。
  \item \textbf{组装世界}：碰撞配置、派发器、宽相、约束求解器四个部件进世界
        构造函数，四者的生命周期都必须长于世界。
  \item \textbf{材质与形状}：这两类是进程级可复用对象，不该每个 actor 新建；
        然后造 actor 并加入世界，之后才允许步进，销毁顺序严格反着来。
\end{enumerate}

所有权（谁负责释放）是三家的主要风格分歧，也是新手最常见的内存错误来源：

\begin{itemize}
  \item Bullet 是手工内存风格：形状与刚体你自己 \Tp{new}，\Fn{addRigidBody}
        之后世界只借用不接管内存，必须先从世界移除再删除对象；
        部件在栈上先声明、世界后声明，析构顺序才正确。
  \item PhysX 是引用计数风格：\Tp{PxMaterial}、\Tp{PxShape}、\Tp{Px{\Bk}Rigid{\Bk}Actor}
        这一族带 add{\Bk}Ref/{\Bk}release 语义，\Tp{PxScene} 接管 actor 之后仍要按计数释放；
        材质与形状可被多个 actor 引用，这正是省内存的正道。
  \item Jolt 走句柄风格：\Tp{Physics{\Bk}System} 持有 body 存储，接口层
        （\Tp{BodyInterface} 一类）用 body ID 而不是裸指针交流，
        于是可串行化、可在任务系统里安全传递；代价是每次取数据
        都要过一遍带锁的 ID 解析。
\end{itemize}

回调是第三块共同纪律。接触、睡眠/唤醒、形状过滤这三类回调都在求解线程上、
关键路径里被高频调用，一帧几千到几万次是正常量级。因此回调体必须短小、
无阻塞、不改变世界：不分配内存、不加会等待的锁、不做 IO、不调用脚本、
不增删 actor（这些只能在步进之外做）。正确姿势是只往预分配的定长缓冲里
记一条事件，主循环在步进结束后消费。

清单 \ref{lst:ch13-bullet-world} 给一个能直接编的最小骨架（Bullet 形状；
换 PhysX 时 ``四个部件 $\rightarrow$ 世界 $\rightarrow$ actor'' 的顺序完全一样）：

\begin{lstlisting}[language=C++, style=cppstyle, caption={建世界、加两个刚体、步进、读回变换（Bullet 形状，可编译骨架）}, label={lst:ch13-bullet-world}]
#include "btBulletDynamicsCommon.h"

int main()
{
    /* 1) 四个部件先声明，世界后声明：析构顺序才正确 */
    btDefaultCollisionConfiguration cfg;
    btCollisionDispatcher dispatcher(&cfg);
    btDbvtBroadphase broadphase;                  /* 宽相：AABB 树 */
    btSequentialImpulseConstraintSolver solver;   /* 顺序冲量求解器 */
    btDiscreteDynamicsWorld world(&dispatcher, &broadphase,
                                 &solver, &cfg);
    world.setGravity(btVector3(0.0, -9.81, 0.0));

    /* 2) 形状进程级复用，不要每帧新建 */
    btBoxShape* box = new btBoxShape(btVector3(0.5f, 0.5f, 0.5f));
    btBoxShape* groundShape =
        new btBoxShape(btVector3(20.0f, 1.0f, 20.0f));

    /* 3) 动态体：motion state 是变换的读写通道 */
    btDefaultMotionState* ms = new btDefaultMotionState(
        btTransform(btQuaternion(0.0f, 0.0f, 0.0f, 1.0f),
                    btVector3(0.0f, 5.0f, 0.0f)));
    btVector3 inertia;
    box->calculateLocalInertia(1.0f, inertia);
    btRigidBody::btRigidBodyConstructionInfo ci(
        1.0f, ms, box, inertia);          /* 质量 1 kg */
    ci.m_friction    = 0.6f;
    ci.m_restitution = 0.2f;
    btRigidBody* body = new btRigidBody(ci);
    body->setUserPointer(reinterpret_cast<void*>(0x1001));
    world.addRigidBody(body);

    /* 4) 静态地面：质量为 0，不需要 motion state */
    btRigidBody::btRigidBodyConstructionInfo gi(
        0.0f, nullptr, groundShape, btVector3(0.0f, 0.0f, 0.0f));
    btRigidBody* ground = new btRigidBody(gi);
    ground->setWorldTransform(btTransform(btQuaternion(
        0.0f, 0.0f, 0.0f, 1.0f), btVector3(0.0f, -1.0f, 0.0f)));
    world.addRigidBody(ground);

    /* 5) 步进：真实时间、本帧最大 substep 数、固定步长 */
    for (int frame = 0; frame < 180; ++frame)
        world.stepSimulation(1.0f / 60.0f, 2, 1.0f / 60.0f);

    /* 6) 读回变换：渲染侧每帧从这里取值 */
    btTransform now;
    ms->getWorldTrans(now);
    printf("y = %f\n", now.getOrigin().getY());

    /* 7) 销毁：先出世界，再还内存，顺序严格反着来 */
    world.removeRigidBody(body);
    world.removeRigidBody(ground);
    delete body;          /* 不接管 ms，也不接管形状 */
    delete ground;
    delete ms;
    delete box;
    delete groundShape;
    return 0;
}
\end{lstlisting}

几处要盯住的地方：\Fn{stepSimulation} 的三个实参是``真实经过时间、本帧允许的最大
substep、固定步长''，传法不对会出现``慢动作''或``一帧跳两步''；\Fn{setUserPointer}
存的整型 tag 是回调里找回游戏对象的可靠通道（不要存指向会被搬迁的容器的指针）；
\Tp{btDefaultMotionState} 是物理与渲染之间那层``谁是权威''的开关——引擎自己管
变换时要用 \Fn{setWorldTransform} 加 \Fn{activate} 强制写入，否则求解器会把它
拉回去。``删刚体不等于删 motion state 与形状''这一点各家都成立，
只是 PhysX 用 release 而不是 \Tp{delete}。

清单 \ref{lst:ch13-tick-budget} 是与 SDK 无关的驱动层，这段代码是帧预算
真正落地的地方：

\begin{lstlisting}[language=C++, style=cppstyle, caption={固定步长驱动：累加器、追赶上限与放弃策略}, label={lst:ch13-tick-budget}]
class PhysicsTick
{
public:
    static constexpr float kFixedDt  = 1.0f / 60.0f;
    static constexpr int   kMaxSteps = 3;         /* 帧预算的硬上界 */
    static constexpr float kMaxLag   = 0.10f;     /* 超过就放弃追赶 */

    /* worldSeconds 是本帧真实时间，来自帧计时而不是系统时钟差 */
    void Advance(float worldSeconds)
    {
        if (worldSeconds > kMaxLag)               /* 异常帧先夹紧 */
            worldSeconds = kMaxLag;
        accumulator_ += worldSeconds;
        int steps = 0;
        while (accumulator_ >= kFixedDt && steps < kMaxSteps) {
            m_beforeStep = Capture();             /* 回放与网络用 */
            TickOneStep(kFixedDt);                /* 中间件的步进 */
            DrainContactEvents();                 /* 步进后再消费事件 */
            accumulator_ -= kFixedDt;
            ++steps;
        }
        if (accumulator_ > kFixedDt) {
            /* 没追完：宁可让物理落后，也不能让这一帧无限长 */
            accumulator_ = 0.0f;
            ++m_dropped;
        }
        m_alpha = accumulator_ / kFixedDt;    /* 渲染插值因子，恒在 [0,1) */
    }

private:
    float accumulator_ = 0.0f;
    float m_alpha      = 0.0f;
    int   m_dropped    = 0;
};
\end{lstlisting}

\begin{note}[title={事件消费放在步进之后，而不是回调当场处理}]
接触事件``当场处理''听着省事，实际是把脚本调用、内存分配、世界修改这三件
最不该出现在求解路径上的事串了进去，而且处理顺序和迭代次数绑定，同一算例
两次跑出不一致的逻辑结果。把回调降级为``只写环形缓冲''、消费放在整帧步进
之后，顺序才可控、可统计，也才能在一帧内合并重复事件。
\end{note}

\section{四个库的取向差别}
\label{sec:ch13-libraries}

四个产品在精度与可控性、可修改度（源码可得性）、平台与主机认证、许可模型
这些维度上取向不同。表 \ref{tab:ch13-libs} 是本册口径的概括，许可一栏只写
``口径''不写条款细节——宽松许可也会随大版本改历史，签约前务必读当前版本的
LICENSE 与官方公告。

\begin{table}[htbp]
\centering
\footnotesize
\caption{四家实时物理中间件的接口、构建、许可与可改造程度（口径概括）}
\label{tab:ch13-libs}
\begin{tabular}{>{\raggedright\arraybackslash}p{0.10\textwidth}
                >{\raggedright\arraybackslash}p{0.13\textwidth}
                >{\raggedright\arraybackslash}p{0.18\textwidth}
                >{\raggedright\arraybackslash}p{0.21\textwidth}
                >{\raggedright\arraybackslash}p{0.18\textwidth}}
\toprule
名称 & 接口语言 & 构建与交付方式 & 许可口径 & 可改造程度 \\
\midrule
PhysX & \CPP{} 头文件为主，上层再包脚本 & 官方源码仓库可 CMake 构建，也发预编译包；主机版走单独通道 & 新版本为 BSD 类宽松开源；早期主线版本历史上曾是专有授权 & 高：可 fork、可换分配器与序列化，能改到求解器内部 \\
Havok & \CPP{} 大表面接口，配套序列化资产 & 签约后交付预编译 SDK，构建期与发布期要绑定许可信息；主机版另有平台认证 & 商业授权，按项目或席位签约，条款不公开 & 低：只有接口与数据格式，无源码；扩展靠官方提供的自定义约束 \\
Bullet & \CPP{}，可自封 \CL{} 边界 & 源码 + CMake，可选 GPU 计算部件 & zlib 类宽松许可 & 高：碰撞配置、宽相、求解器都可整体替换 \\
Jolt & \CPP{}，句柄式接口 & 源码 + CMake，接口层薄，容易并进自己的构建 & MIT 类宽松许可 & 高：宽相、任务系统、body 存储都可换；设计目标里就含跨平台确定性 \\
\bottomrule
\end{tabular}
\end{table}

读表时抓三条。\textbf{精度与可控性}：开源的三家可以把求解器换成自己的实现，
这是``可控''的真正含义，不只是``能调试''；Havok 的可控性体现在参数表面很大、
序列化管线很成熟，而不是能改代码。\textbf{平台与认证}：主机平台关心的是
有没有已认证的二进制、是否满足内存与固件调用约束、能不能提供认证所需的
确定性证据——商用中间件在这条路上的积累比开源版本长，这是很多团队仍选
商业授权的真实理由，和``谁更准''关系不大。\textbf{许可模型}：宽松许可
（BSD、MIT、zlib 一族）意味着可以自由修改与再分发，但不意味着可以拿别人的
品牌名宣传；商业授权意味着有支持渠道与认证包，也意味着升级节奏不由你决定。

\section{与离线仿真的概念对照}
\label{sec:ch13-vs-offline}

第一部分讲过的概念在这里都能找到``变形后的对应物''，但对应关系不是``同名''而是
``同一问题的不同预算答案''。表 \ref{tab:ch13-mapping} 的读法：左列``有没有被判据
终止''，右列``只有次数上限''。

\begin{table}[htbp]
\centering
\caption{离线求解概念与实时物理概念的映射（左：可判敛；右：只能计数）}
\label{tab:ch13-mapping}
\begin{tabular}{>{\raggedright\arraybackslash}p{0.30\textwidth}
                >{\raggedright\arraybackslash}p{0.48\textwidth}}
\toprule
离线求解里的概念 & 实时物理里的对应做法 \\
\midrule
一致切线刚度、整体雅可比组装 & 每个约束的 effective mass（形如 $J M^{-1} J^{T}$）在每次迭代里就地重算，不做全局组装，也不进内存池 \\
残差判敛、``迭代到容差为止'' & 固定 \Tp{iteration} 次数的顺序冲量扫描；没有全局残差可读，判据退化为``画面上看不出破绽'' \\
预条件（Jacobi、ICC、多重网格） & warm start：拿上一帧的冲量当本次迭代的初值；软约束参数（ERP/CFM 或按频率给的刚度阻尼）起同样的``把问题变温和''的作用 \\
大规模稀疏线性系统的直接法 & 无对应能力。实时不做直接法，只迭代有限次；``解得好''变成``初值给得好 + 迭代顺序排得好'' \\
隐式积分的无条件稳定、CFL 条件 & 固定步长 + substep 上限 + 速度上限与最大分离速度（depenetration velocity）限幅 \\
松弛因子、阻尼迭代 & 迭代次数、求解顺序与 island 划分；GS 排序质量在这里就是``预条件质量'' \\
同一算例两次跑出同一结果 & 同构建、同硬件、同线程数下才可复现；跨平台、跨 SIMD 宽度与并行 island 顺序变化都会破坏回放一致性 \\
\bottomrule
\end{tabular}
\end{table}

其中 ``warm start 相当于预条件''这一条最值得展开：离线算例里预条件矩阵要离线
构造、在线复用，而实时物理里上一帧的冲量场就是免费的``近似解''。所以物理对
``帧间连续性''极其敏感——一帧改一次质量、把 body 传送过去、或者切步长，都会
让初值失效，症状是``参数没改却突然抖''。

\section{数值与调试}
\label{sec:ch13-numerics}

\textbf{宽相与窄相}是两个必须分开看性能与正确性的阶段。宽相只回答``这两个
AABB 可能碰吗''，产物是候选对；窄相才生成流形与接触点，回答``以多深、在哪些点
接触''。宽相成本靠分层与休眠来压，窄相成本靠形状复杂度与接触点数上限来压。
把``宽相耗时高''归因于``接触算得慢''是常见误判：先看候选对数量。

\textbf{穿透与抖动}是两个不同病因。穿透多半是速度与厚度之比超限：一个物体一步
之内的位移大于碰撞体厚度，离散步进就会漏过（tunneling），解法是加大 $S$
（切小步长）、限制最大线速度、给薄物体单独开连续碰撞检测。抖动多半是求解器
给不出满足全部约束的解：深堆叠、mass ratio 悬殊、一个 body 被多个关节同时
钉住；解法是减少约束、把质量比压进一个量级、加大迭代或让静止物体休眠。

\textbf{offset 一族参数}管的是``什么时候开始算接触''与``稳定下来停在哪''。
PhysX 的形状侧有一对 contact offset 与 rest offset：前者是预测距离，接触对在这
个距离上就被生成（回调也在这里触发），后者是最终保持的间隙且必须小于前者；
Bullet 与 Jolt 用类似的``碰撞预测距离/margin''概念，margin 越大宽相候选对越多。
调大它们是让接触事件提前触发的常用手段，代价直接落在宽相耗时上；各家参数名
与挂载位置不同，以官方文档为准。

\textbf{确定性与跨平台}是这一层的``回放问题''。浮点加法不满足结合律，SIMD 宽度、
FMA 指令、线程数与 island 遍历顺序任一改变都可能改变结果位型，于是同一份输入
在 x86 与 ARM、单线程与多线程下跑出不同轨迹，录像回放与回归测试全部失效。
可行的口径有三条：（1）要确定性就把步长、迭代数、线程数与遍历顺序钉死，逻辑里
优先用句柄 ID 而不是指针；（2）网络同步优先``服务端权威 + 客户端只预测画面''，
把不一致留在看不见的一侧；（3）物理回归测试做成``同一构建同一硬件''的位级比较，
跨平台只比统计量（穿透率、耗时分布）。部分开源实现把跨平台确定性写成设计目标，
是否成立要看它自己文档里的条件，不要假设。

\begin{warning}[title={在物理回调里做内存分配或 IO}]
最贵的一类事故不是求解慢，而是回调把求解路径变成了串行路径：
\begin{itemize}
  \item \textbf{分配}：回调里 \Tp{new}、构造 \Tp{std::string}、往 \Tp{vector}
        里 push 都可能触发堆锁与重分配；求解是多线程的，一次堆锁能把整帧时间
        放大数倍，而且只在接触最多的那一帧出现。
  \item \textbf{IO 与日志}：一次接触一行写进文本文件，几万个接触对就把磁盘变成
        瓶颈；即使有缓冲，字符串格式化也仍在求解线程上做。
  \item \textbf{改世界}：增删 actor、改形状、改碰撞层等于在迭代中途重排存储
        结构，多数 SDK 直接定义为未定义行为，症状是偶发崩溃或整帧错位。
  \item \textbf{进脚本}：穿一次脚本边界（下一章的主题）不但贵，还可能触发脚本侧
        的 GC、异常与再入。
\end{itemize}
正确形状是：回调体内只做``读 tag + 写预分配的定长环形缓冲 + 返回''，容量满时丢弃
并计数，缓冲大小按``一帧最大接触数''而不是平均值预留；主循环在整帧步进结束后
统一消费、合并，再分发给逻辑与脚本。
\end{warning}

\subsection*{本章小结}
实时物理的全部工程判断都可以回到同一个式子：$T \approx S \times (B + I \times (N + Z))$。
正确性与稳定性只能靠 $S$ 与 $I$ 换，帧预算只能靠减少接触规模换，于是
``离线思维''（迭代到收敛）必须换成``实时思维''（固定次数近似 + warm start +
限幅）。接口层面四家都是同一段谱：给运行时、组装世界部件、复用材质与形状、
加 body 步进；分歧集中在所有权模型（手工 \Tp{delete}、引用计数 release、
句柄 ID）与许可/认证形态。下一章把镜头移到另一条边界：脚本层怎么跨过 C++ 接口。

% ============================================================
% ---- frag_ch14.tex ----
\chapter{脚本绑定层：Lua、sol2、AngelScript 与反射代价}
\label{ch:ch14-binding}

脚本层的价值是``改数据不改二进制''，代价是每次从脚本碰一次 C++ 都要过一道边界。
绑定层就是这道边界上的全部代码：它把 C++ 的类、函数与对象，翻译成脚本语言里
一个能调用、能索引、能比较的东西。这一章讲三件事——边界上的调用约定、
对象生命周期归谁、以及每跨一次界的账单。

\begin{guideline}[title={本章口径}]
\begin{itemize}
  \item 只写有把握的类型与函数名：\Lib{Lua} 的 \Tp{lua\_State} 与 \Tp{luaL\_*}
        一族、\Lib{sol2} 的 \Tp{sol::state} 与 \Tp{new\_usertype}、
        \Lib{AngelScript} 的注册接口概念。
  \item Lua 的具体 API 在 5.1、5.2、5.4 之间有增删（例如注册函数的接口），
        落地时以所用版本的 \Tp{lauxlib.h} 与官方手册为准。
  \item 性能一栏只给量级与测法，不给具体机器上的绝对数字。
\end{itemize}
\end{guideline}

\section{绑定层为什么存在：跨界与生命周期}
\label{sec:ch14-why}

脚本自己也能算账，但引擎里真正值钱的东西都在 C++ 侧：场景图、渲染设备、
物理世界、资源句柄。绑定层存在的理由就一条：让脚本能碰到这些对象，
同时把``谁负责它们还活着''这件事定清楚。生命周期归属有三种基本形态，
三者的事故形态完全不同。

\begin{table}[htbp]
\centering
\footnotesize
\caption{脚本对象生命周期的三种归属模型}
\label{tab:ch14-lifetime}
\begin{tabular}{>{\raggedright\arraybackslash}p{0.20\textwidth}
                >{\raggedright\arraybackslash}p{0.24\textwidth}
                >{\raggedright\arraybackslash}p{0.34\textwidth}}
\toprule
归属模型 & 脚本侧拿到的句柄是什么 & 典型事故与对策 \\
\midrule
C++ 拥有，脚本借用 & 裸指针或整型 ID，被包进 userdata & 引擎侧销毁对象后脚本仍调用它，成为悬垂引用。对策：ID 加代际号（generation），或让句柄在对象析构时被主动置空 \\
脚本拥有，C++ 引用 & 带 \Tp{\_\_gc} 元方法的 userdata，内含一份对象拷贝 & C++ 侧存了指向 userdata 内部的指针，脚本表被清空后 GC 何时跑不确定，引用先发失效。对策：跨界只传值或传 ID，不传内部指针 \\
共同拥有，计数管理 & 引用计数句柄（智能指针包一层） & 计数环与 GC 互相牵住造成泄漏，或退出顺序错乱导致析构访问已释放的设备。对策：明确退出阶段，先禁脚本再拆引擎 \\
\bottomrule
\end{tabular}
\end{table}

按``绑定代码从哪来''分，路线只有三档，代价与自由度也三档分明：

\begin{itemize}
  \item \textbf{手写 C API 绑定}：一切自己写，边界上做什么、拷贝什么、
        什么时候报错全在自己手里。可读性与可控性最好，工程量与漏写检查的风险最大。
  \item \textbf{模板库或生成器}（\Lib{sol2}、代码生成工具、\Lib{pybind11}）：
        用 C++ 模板或外部描述把上面那圈样板生成出来。开发速度最快，
        代价落在编译时间与可读的报错信息上。
  \item \textbf{引擎自带反射}：类型信息由引擎的反射表提供（宏或代码生成物化），
        脚本按名字访问对象。编辑期与调试期最灵活，
        代价是每次调用都要在运行时查表与打包参数。
\end{itemize}

\section{Lua 与 C 的栈式契约}
\label{sec:ch14-lua-stack}

Lua 的 C 接口不是``函数指针加参数列表''，而是一棵\textbf{虚拟栈}：调用时实参
已经在栈上，被调函数从栈上取，返回值压回栈顶，再用一个整数告诉 Lua``我压了几个
返回值''。栈在 C 侧是 \Tp{lua\_State*}（Lua 里叫 thread），一个 \Tp{lua\_State}
一条栈；协程是另一条 \Tp{lua\_State}，与主线程共享全局表但有自己的栈。索引规则
要背下来：绝对索引从 1 开始数到 \Fn{lua\_gettop}，负索引从 -1（栈顶）往回数，
\Tp{lua\_upvalueindex}(k) 让一个 C 函数带几个隐藏槽位。

清单 \ref{lst:ch14-lua-cfunction} 是最小完整的一段：一个 userdata 型对象加
两个 C 函数。每行右侧注明它对栈做了什么。

\begin{lstlisting}[language=C, style=cstyle, caption={\Lib{Lua} C 函数：压栈、检查、返回值的完整契约}, label={lst:ch14-lua-cfunction}]
/* 对象在 Lua 里的样子：一块托管内存 + 一张元表 */
static const char* kVecMeta = "Demo.Vector3"; /* 元表的注册名 */

/* Lua 侧写法：local v = demo.new(1, 2, 3) */
static int vec_new(lua_State* L)
{
    double x = luaL_checknumber(L, 1); /* 读索引 1，不弹栈 */
    double y = luaL_checknumber(L, 2);
    double z = luaL_checknumber(L, 3);

    double* v = (double*)lua_newuserdata(L, 3 * sizeof(double));
                                         /* 压入一块内存：栈 .. u */
    v[0] = x; v[1] = y; v[2] = z;
    luaL_setmetatable(L, kVecMeta);  /* 关联元表，栈不变 */
    return 1;                 /* 栈顶那 1 个值就是返回值 */
}

/* 用法：demo.dot(a, b)。进来时栈里已经是 u u */
static int vec_dot(lua_State* L)
{
    double* a = (double*)luaL_checkudata(L, 1, kVecMeta);
    double* b = (double*)luaL_checkudata(L, 2, kVecMeta);
    double s = a[0] * b[0] + a[1] * b[1] + a[2] * b[2];
    lua_pushnumber(L, s);      /* 栈：u u n */
    return 1;
}

static const luaL_Reg kFuncs[] = { /* 名字到函数指针的表 */
    { "new", vec_new },
    { "dot", vec_dot },
    { NULL,  NULL }            /* 空项结尾，长度靠它算 */
};

/* require("demo") 的装载入口：返回模块表 */
int luaopen_demo(lua_State* L)
{
    luaL_newmetatable(L, kVecMeta); /* 建元表并压栈：栈 m */
    lua_pushvalue(L, -1);           /* 复制栈顶：栈 m m */
    lua_setfield(L, -2, "__index"); /* m.__index = m */
    luaL_setfuncs(L, kFuncs, 0);    /* C 函数挂进 m */
    return 1;                       /* m 就是模块表 */
}
\end{lstlisting}

几处约定要单独说：\Tp{luaL\_checknumber} 与 \Tp{luaL\_checkudata} 这类
\Tp{luaL\_check\_*} 是``要么拿到合法值、要么以 Lua 错误退出''，它们不会让
你的 C 函数返回不合法的值，但会用 longjmp 离开当前栈帧——所以
\textbf{在被 Lua 调用的 C 函数里，不要依赖局部 C++ 对象的析构来清理资源}：
要么把资源登记给 \Tp{\_\_gc}，要么由 C++ 侧用 \Fn{lua\_pcall} 反向进入脚本，
并在边界上保证异常被接住。表示对象有两条路：
\textbf{light userdata} 只是一个裸指针（无自己的内存块、不能挂元表的
finalizer、比较就是比地址），适合传``宿主拥有、随时可能失效''的句柄；
\textbf{full userdata} 是一段由 Lua GC 管理的内存，可以挂元表、可以带
\Tp{\_\_gc}，适合``脚本持有所有权''的对象。要把 C++ 对象放进 full userdata，
常见做法是放一个指针并配一个 placement 构造的包装体，\Tp{\_\_gc} 里再显式
调析构，因为 GC 不会替你调 C++ 的析构函数。

\begin{warning}[title={栈索引在循环里漂移}]
最典型的栈错误是``边压边数''：在循环里 \Fn{lua\_pushvalue}、\Fn{lua\_rawgeti}
而不弹回，索引 -1 指的已经不是你以为的那个值，症状是取到 nil 或取到隔壁的字段。
纪律有三条：（1）每个辅助函数开头注释``进来时栈上有什么''，
结尾注释``出去时应该有什么''；（2）临时值用完立刻 \Fn{lua\_pop}，
数量一次给够；（3）需要长期保存的东西不要留在栈上，存进注册表
（\Fn{luaL\_ref} 给一个整数句柄）或 upvalue 槽里。Lua 那条伪栈有元素数上限，
超出时 \Fn{lua\_checkstack} 直接报错，递归展开的绑定代码尤其容易踩。
\end{warning}

\section{C++ 侧的模板库：sol2 自动化了什么}
\label{sec:ch14-sol2}

\Lib{sol2}（以及同类的 \Tp{sol::state}/{\Bk}\Tp{sol::table} 封装）做的事情，是把
清单 \ref{lst:ch14-lua-cfunction} 里那一圈样板变成一行注册：类型怎么在
userdata 里落地、成员指针怎么转成 C 函数、参数怎么逐个检查与转换、
C++ 异常怎么变成 Lua 错误、多个返回值怎么打包。

\begin{lstlisting}[language=C++, style=cppstyle, caption={\Lib{sol2} 注册一个类型与一个工厂函数}, label={lst:ch14-sol2-bind}]
#include <sol/sol.hpp>
#include <vector>

struct Enemy
{
    float hp = 100.0f;
    float TakeDamage(float amount) { hp -= amount; return hp; }
};

std::vector<Enemy> g_pool;      /* 所有权留在 C++ 这一侧 */

void BindEnemy(sol::state& lua)
{
    /* 名字、构造方式、成员与方法：一次写全 */
    lua.new_usertype<Enemy>("Enemy",
        sol::constructors<Enemy()>(),
        "hp",          &Enemy::hp,
        "take_damage", &Enemy::TakeDamage);

    /* C++ 异常被翻译成 Lua error；返回指针意味着脚本只借用 */
    lua.set_function("spawn", [](float hp) -> Enemy* {
        g_pool.push_back(Enemy{hp});
        return &g_pool.back();  /* 扩容后旧指针会失效 */
    });

    lua.script("local e = spawn(50.0)\n"
               "print(e.hp, e:take_damage(20.0))");
}

void Init()
{
    sol::state lua;
    lua.open_libraries(sol::lib::base);  /* 只开需要的库 */
    BindEnemy(lua);
}
\end{lstlisting}

代价是实打实的四条：

\begin{itemize}
  \item \textbf{编译时间}：每个注册点都要展开一大套模板，绑定的 .cpp
        很容易成为全工程最慢的翻译单元。缓解办法是把注册集中到一个
        或少数几个 .cpp，头文件里只留声明。
  \item \textbf{报错信息}：签名写错（成员指针类型不符、忘了默认构造）会报出
        几十行模板实例化轨迹。定位方法是把报错里的第一个实例化点
        对回注册语句，而不是从最后一行往上读。
  \item \textbf{隐式拷贝}：默认按值把对象塞进 userdata，意味着每次从脚本
        取一个对象可能触发一次拷贝构造；带 \Tp{std::string} 或容器的类型
        在热点里会被这个拷爆。要借用就显式返回指针或引用型包装。
  \item \textbf{异常翻译}：把 C++ 异常变成 Lua 错误很方便，但代价是你的
        析构与 RAII 依赖栈展开，跨语言边界抛异常的路径要单独测。
\end{itemize}

\section{其它路线：编译期绑定与运行时反射}
\label{sec:ch14-other-routes}

\Lib{AngelScript} 走的是``先注册、再编译''的路线：绑定信息在初始化时一次性
交给脚本引擎（类型、方法、全局函数各一条注册调用），脚本源码编译成它自己的
字节码后才执行。因为类型在编译期就定下来了，调用点不需要按名字查表、
参数不需要装箱，边界上的动作可以提前算好；它的脚本函数也更容易被本机执行
或 JIT 一类的加速方案接管（这类方案与版本、平台相关，以官方文档为准）。
代价是绑定代码是显式的一整套注册，写起来比模板库更机械。

\Lib{pybind11} 一类的 \PY{} 绑定主要不在游戏运行路径上，而在工具与自动化侧：
离线批处理、资产流水线、回归测试脚本、把求解器或物理世界当库来调。
它换来的是完整的 \PY{} 生态（数值库、报表、CI），代价是 GIL、
类型转换与容器拷贝，以及``生产环境里没有解释器''这条硬边界。

第三条路线是``通过反射访问引擎对象''：引擎把类型信息物化成反射表
（宏或代码生成），脚本或编辑器按名字取成员、调方法、显示属性面板。
它的灵活性最高，成本结构也最特殊：每次访问都要过一次
``名字 $\rightarrow$ 成员/函数 $\rightarrow$ 参数打包 $\rightarrow$ 调用''的
运行时流程。编辑期把结果缓存下来很便宜，运行期每帧几千次就变成负担。

\begin{compare}[title={把一次调用的成本拆开看}]
\begin{itemize}
  \item \textbf{查名字}：全局表、元表 \Tp{\_\_index}、反射字段查表。
        这是脚本侧最大的一项，通常比实际调用本身贵。
  \item \textbf{打包参数}：数字要转成 \Tp{lua\_Number}，字符串要过编码，
        容器要逐元素转换（装箱）；返回值同一遍程。
  \item \textbf{对象解析}：userdata 取指针、ID 查池、可能带一次锁或一次
        有效性检查。
  \item \textbf{真正的 C++ 调用}：往往是这几项里最便宜的一个。
\end{itemize}
所以``把函数改快''通常没用，``少调用几次''才有用。
\end{compare}

\section{边界上的性能账}
\label{sec:ch14-cost}

量级感（同一台机器、Release 构建、自己测才算数）大致是表
\ref{tab:ch14-cost} 这样。注意这些数字的真正用途是让你判断
``每帧调多少次''，不是让你去比较哪个绑定库更快。

\begin{table}[htbp]
\centering
\footnotesize
\caption{一次跨界调用的成本量级与主要来源（口径表，须自测核对）}
\label{tab:ch14-cost}
\begin{tabular}{>{\raggedright\arraybackslash}p{0.22\textwidth}
                >{\raggedright\arraybackslash}p{0.14\textwidth}
                >{\raggedright\arraybackslash}p{0.42\textwidth}}
\toprule
调用形态 & 量级 & 成本主要花在哪 \\
\midrule
栈上取两三个数、直接算完返回的手写 C 函数 & 十纳秒级 & 入栈与 \Tp{luaL\_check\_*}；实际算术几乎不占 \\
元表方法（冒号语法）调用一个成员函数 & 数十纳秒级 & 多一次 \Tp{\_\_index} 查表与 self 参数打包 \\
模板库注册的类型方法，带引用返回 & 数十到数百纳秒 & 参数逐个转换、返回值打包、可能一次拷贝构造 \\
按名字经反射表调用引擎方法 & 微秒级 & 名字查表、参数装箱为 variant、生命周期检查 \\
同一次调用里再进物理或渲染子系统 & 由被调用方决定 & 边界开销被淹没，此时应优化调用次数与算法 \\
\bottomrule
\end{tabular}
\end{table}

测法比数字更重要。三条做法：（1）\textbf{微基准}——固定次数（$10^{6}$ 到
$10^{7}$ 次）循环调用，结果写进一个不会被优化掉的变量，分别测冷启动与热循环，
并把``空函数''的对照组一起测出来；（2）\textbf{计数而不是猜}——在每个绑定入口
累加计数器，帧末打印每类调用的次数与总耗时，热点脚本的问题九成是
``次数''而不是``单次太慢''；（3）\textbf{在真实帧里测}——微基准给的相对
比例，到了带缓存、带分支的调用序列里会变，最终以帧 profiler 的数字为准。

清单 \ref{lst:ch14-batch} 是这条账唯一有效的还法：把循环留在原生一侧。

\begin{lstlisting}[language=C, style=cstyle, caption={批量接口：一次跨界处理一整批参数，循环留在 C 这一侧}, label={lst:ch14-batch}]
/* 参数：索引 1 是 Lua 表，索引 2 是条数， bodies 走 upvalue */
static int batch_apply_impulse(lua_State* L)
{
    luaL_checktype(L, 1, LUA_TTABLE);            /* 表必须是表 */
    const lua_Integer n = luaL_checkinteger(L, 2);
    RigidBody* bodies = (RigidBody*)
        lua_touserdata(L, lua_upvalueindex(1));  /* 对象池指针 */

    for (lua_Integer i = 1; i <= n; ++i) {
        lua_rawgeti(L, 1, i);          /* 压入第 i 条：栈 .. t */
        if (!lua_istable(L, -1))
            return luaL_error(L, "entry %d not a table", (int)i);
        lua_getfield(L, -1, "ix");     /* .. t ix */
        lua_getfield(L, -1, "iz");     /* .. t ix iz */
        bodies[i - 1].vx += (float)lua_tonumber(L, -2);
        bodies[i - 1].vz += (float)lua_tonumber(L, -1);
        lua_pop(L, 3);                 /* 弹出 ix、iz 与那条表 */
    }
    lua_pushinteger(L, n);             /* 返回处理条数 */
    return 1;
}
\end{lstlisting}

\begin{bestpractice}[title={把循环留在原生一侧，只暴露批量接口}]
\begin{itemize}
  \item 需要遍历大量对象的操作（给一千个单位加冲量、刷新一批材质）
        写成\textbf{一次调用带一批数据}，而不是让脚本写 for 循环。
        跨界次数从 $N$ 降到 1，脚本侧的查表与装箱也一起消失。
  \item 批量参数用扁平数组或定长记录，不用每元素一张表：表要逐字段
        \Tp{getfield}，数组只要按下标取值。
  \item 热点接口提供``一次拿多个值''的读接口（位置、速度、状态一起返回），
        避免脚本为了三个字段跨界三次。
  \item 对象句柄在脚本侧存成 ID 或 userdata，不要每帧用名字去引擎里找对象；
        名字查表是反射路线最大的单项开销。
  \item 稳定后再放开粒度：先用计数与帧 profiler 找出次数最多的边界，
        只把那几处改成批量，别为了统一风格把所有接口都做成聚合入口。
\end{itemize}
\end{bestpractice}

\subsection*{本章小结}
绑定层的全部设计问题只有两个：调用约定（\Lib{Lua} 是栈式契约，
\Lib{AngelScript} 与反射路线是查表与打包契约）与生命周期归属
（C++ 拥有、脚本拥有、共同拥有三种，事故形态各异）。三档实现路线的取舍
是``样板自己写''换可控、``模板库生成''换开发速度、``引擎反射''换编辑期灵活。
性能结论只有一句：单次跨界是十纳秒到微秒级，真正决定帧时间的是次数，
所以优化顺序永远是先减次数、再改批量、最后才看单次实现。

% ============================================================
% ---- frag_ch15.tex ----
\chapter{实时仿真与硬件在环：确定性时间步、HIL 与网络同步}
\label{ch:ch15-hil}

前两部分讲的分别是两端：\DomCAE{} 批处理仿真不关心墙钟，算多久都可以；
\DomEngine{} 的帧循环最在乎 16 ms 一格的节奏。中间还有一大片实时仿真：
仿真器必须和真实设备或真人玩家在墙钟里闭合回路，又不能像游戏那样把掉帧当作无伤大雅。
本章按三个问题展开：时间循环对外承诺什么、承诺如何落成代码、事后如何复现，
并对照 Part I 的检查点机制。

\section{实时性的三种承诺}
\label{sec:ch15-promises}

``实时''不是一种性质，而是三种互不蕴含的承诺，混用它们是这一层最常见的概念事故：

\begin{itemize}
  \item \textbf{硬实时}：错过截止期本身就是失败，无论结果算得多对。
        一个 HIL 台架没能在被测 ECU 的下一个采样窗口前给出转速，
        ECU 会判传感器故障并进入保护逻辑——实验已经作废，哪怕仿真器随后算出来了。
  \item \textbf{软实时}：偶尔迟到是可接受的退化，只要平均跟得上、退化有限。
        游戏帧循环是标准例子：掉几帧没人受伤，代价只是手感短暂变差。
  \item \textbf{确定性}：同一输入序列产生同一轨迹，与墙钟无关。
        它不要求快，只要求可复现。
\end{itemize}

三种承诺在实现手段上交叠但互相独立：硬实时系统未必确定（抢占让步与步的边界移动），
确定性程序未必实时。\DomCAE{} 批处理通常只要第三条里最弱的一档——科学上可重复即可；
\DomEngine{} 单机游戏要第二条加第三条的弱化形式；
HIL 要前两条，还要第三条保证回归测试与认证送审可信。

\begin{table}[htbp]
\centering
\caption{三种承诺的失败含义、度量方式与典型需求方}
\label{tab:ch15-promise}
\begin{tabular}{>{\raggedright\arraybackslash}p{0.12\textwidth} >{\raggedright\arraybackslash}p{0.22\textwidth} >{\raggedright\arraybackslash}p{0.20\textwidth} >{\raggedright\arraybackslash}p{0.18\textwidth}}
\toprule
承诺 & 违反意味着什么 & 用什么指标度量 & 谁必须较真 \\
\midrule
硬实时 & 整个任务失败，结果迟到即作废：被测件或真实被控对象直接受到物理影响 & 最坏执行时间与截止期之比、截止期 miss 次数是否恒为零 & HIL 台架、被控对象在线的半实物实验 \\
软实时 & 退化而非失败：卡顿、音画不同步、输入响应延迟变大，但世界仍然自洽 & 帧时间分布的 P99、每秒长帧次数、连续掉帧的最长长度 & 游戏帧循环、音视频管线 \\
确定性 & 复现失败：回归测试结果不可信，事故现场无法回放定位 & 同输入两次运行的状态哈希是否逐位一致、回放首次分歧发生在第几步 & 联机对局校验、反作弊、认证送审 \\
\bottomrule
\end{tabular}
\end{table}

为什么游戏对确定性的要求是``弱化''的：竞技向实现要求``同一输入脚本得到同一对局''，
用于回放、观战与反作弊；但渲染、网络抖动、资源加载顺序这些外围不参与一致性判定。
工程上于是把``必须确定的部分''（游戏逻辑、物理步进）与``允许发散的展示层''切开，这个切法正是 \S\ref{sec:ch15-net} 网络同步的前提。

\section{时间步的形状：定步长与变步长}
\label{sec:ch15-timestep}

\DomEngine{} 侧的主流形状是\textbf{固定步长加累加器}：真实时间先攒进一个累加器，
每攒够一个物理步长就推进一次仿真；渲染时刻的状态用相邻两个物理状态插值得到。
优点是仿真内核完全不知道帧率的存在：物理函数每次收到同样的 \Tp{dt}，
缓存友好、调试可预测，而且改显示器刷新率不会改轨迹。

\begin{lstlisting}[language=C++, style=cppstyle, caption={定步长累加器循环：最大子步与追赶上限缺一不可}, label={lst:ch15-fixed-step}]
// 帧循环只负责攒时间；物理永远以整步推进，两个上限各防各的病
constexpr double kFixedDt     = 1.0 / 120.0;  // 物理步长，与产品协商后钉死
constexpr int    kMaxSubsteps = 8;            // 单帧最多补算几步
constexpr double kMaxLag      = 0.25;         // 追赶上限（秒），超出直接丢弃

class FixedStepLoop {
public:
    void tick(double wallNow) {                 // 每帧一次，wallNow 必须取单调时钟
        double frameTime = wallNow - lastWall_;
        lastWall_ = wallNow;
        if (frameTime > kMaxLag)                // 断点恢复、系统睡眠都会灌进巨大的
            frameTime = kMaxLag;                // frameTime：宁可虚拟钟落后于墙钟
        accumulator_ += frameTime;

        int sub = 0;
        while (accumulator_ >= kFixedDt && sub < kMaxSubsteps) {
            sampleInputs();                      // 输入按步对齐，一步只取一份
            physicsStep(kFixedDt);               // 纯函数：只依赖本步状态与输入
            accumulator_ -= kFixedDt;
            simTime_  += kFixedDt;
            ++step_; ++sub;
        }
        if (sub == kMaxSubsteps && accumulator_ >= kFixedDt) {
            accumulator_ = 0.0;                  // 物理本身跟不上墙钟：清账降级
            ++spiralCount_;                      // 该计数非零说明步进成本超标
        }
    }
    double alpha() const { return accumulator_ / kFixedDt; }  // 渲染插值因子
private:
    double accumulator_ = 0.0, simTime_ = 0.0, lastWall_ = 0.0;
    int    step_ = 0, spiralCount_ = 0;
};
\end{lstlisting}

清单 \ref{lst:ch15-fixed-step} 里两个上限针对两种病：\Tp{kMaxSubsteps} 防
``死亡螺旋''（一步超时导致下帧补更多步、下帧更超时）；\Tp{kMaxLag} 防一次性灌入的
巨大帧时间，触发时让仿真时间落后现实时间，由产品决定此时是降物理频率还是接受落后。
渲染状态要落在两个物理状态之间，所以插值是定步长方案的固定配套：

\begin{lstlisting}[style=cppstyle, caption={相邻两个物理快照的插值：帧可见状态只由插值产生}, label={lst:ch15-interp}]
struct Snap { float p[3]; float q[4]; };   // 位置与朝向
Snap prevSnap_{}, curSnap{};               // 每个物理步结束后同步更新

void renderView(double alpha) {            // alpha 即上一个清单里累加器的余量比例
    float p[3];
    for (int i = 0; i < 3; ++i)
        p[i] = prevSnap_.p[i] + float(alpha * double(curSnap.p[i] - prevSnap_.p[i]));
    // 朝向用球面插值；线性插值再归一化会在角速度大时系统性变慢
    drawBody(p, slerpApprox(prevSnap_.q, curSnap.q, alpha));
}
\end{lstlisting}

对帧时刻状态做\textbf{外推}（沿速度把物体往前猜）只在两种情况下合理：
目标由别的机器驱动、本端只能猜（\S\ref{sec:ch15-net} 的远端实体），
或者显示绝不能落后于输入。本地物理里外推意味着渲染可能进入
``下一个步长才判定存在''的空间——穿墙轨迹与被拉回的橡皮筋都是它的症状。

\DomCAE{} 侧的求解器普遍是\textbf{变步长加误差控制}：步长不是约定值，
而由局部截断误差估计反推——误差超容差就砍半步重试，连续几步富余就放大；
stiff 系统还会在隐式求解与事件定位（过零检测）之间切换步长。
两个域对``时间''的分歧基本都源自这两种形状。

\begin{compare}[title={定步长与变步长的分工与代价}]
\begin{itemize}
  \item 谁定义时间：定步长由仿真自己钉死；变步长由误差容差与事件位置决定。
  \item 可复现性：定步长天然逐位可复现；变步长还要固定容差、最大步长与事件顺序。
  \item 对外接口：定步长容易接实时循环与总线节拍；变步长需要 \Tp{advance\_until} 这种按外部时刻推进的接口，且每步耗时不均。
  \item 撕裂风险：定步长来自渲染采样到不一致的状态对；变步长来自两个子系统各自推进到不同时刻后交换数据。
\end{itemize}
\end{compare}

所谓\textbf{物理撕裂}，指同一时刻存在两份权威状态：渲染、其他子系统或客户端
看到的不是同一次步进的产物。典型来源：多个子系统跑在不同频率（骨骼每帧、
布料每三步）而耦合点上的力与位置来自不同时刻；插值用的状态对时序没对齐；
回滚或热重载后只恢复了一半状态。修法统一：每个可对外见的实体固定权威步长，
跨频率耦合走清单 \ref{lst:ch15-interp} 的快照对，而不是各取``最近的值''。

\section{硬件在环：被测件才是时钟主人}
\label{sec:ch15-hil-loop}

硬件在环（\Tp{HIL}）把闭环里的一环换成真硬件：被测的是真实 ECU 或飞控计算机，
仿真侧提供它周围``整个世界''的模型——转速、轮速、母线电压、总线报文。
实时承诺由此从``尽量快''升级为硬性契约：模型必须恰好在被测件的控制周期内给出响应。
100 Hz 的环路意味着 10 ms 内跑完全部模型并完成一次总线交换；超了，
被测件看到的世界就发生时间跳变，轻则报故障码，重则进保护逻辑把台架停掉。

接口层常见两类总线：车载以 CAN 为主，仲裁速率约 1 Mbit 级、一帧数据段 8 字节，
报文调度自带优先级与排队延迟，较新平台往车载以太网（含 TSN 一类实时增强）迁移；
航空侧则多见 ARINC 429、离散量与模拟量采集。无论哪种，延迟预算都要拆开逐项算，按最坏情形而非平均值：

\begin{itemize}
  \item 模型执行：本步全部模型的总执行时间，用 WCET（最坏执行时间）口径。
  \item 采集与激励：把上一周期收到的报文解析进状态、把本周期输出写回硬件；
        协议带时间戳时还要对齐到统一步节拍。
  \item 传输与排队：报文在控制器与收发器队列里的滞留，帧长除以速率是硬下限。
  \item 台架固有延迟：仿真主机中断响应与板卡 DMA 周期，商用台架（\Lib{dSPACE} SCALEXIO、\Lib{NI} VeriStand、Speedgoat 一类）的规格书会给出，以所购配置为准。
\end{itemize}

\begin{note}[title={延迟预算表要留一条明确的余量线}]
硬实时的判据是 ``最坏执行时间加传输上限不超过截止期''，不是 ``平均执行 3 ms''。
缓存缺失、总线仲裁风暴、后台维护任务都造成长尾：预算表每一项写分布的上端，
整表留一到两成余量，逼近余量线即扩容或砍模型，而不是祈祷分布不发作。
\end{note}

为了让最坏执行时间可算、可送审，HIL 模型通常由模型编译器（Simulink Coder、
TARGETBOX 一类）生成\textbf{定长代码}：每个采样点一次调用、静态分配、无递归、
无动态内存、循环次数编译期有界。清单 \ref{lst:ch15-hil-model} 给出这种代码的形状。

\begin{lstlisting}[language=C, style=cstyle, caption={HIL 模型步函数的形状：静态存储、单精度、定步长}, label={lst:ch15-hil-model}]
/* 生成代码的形态示意：一个采样点一次调用，接口与缓冲全部静态分配 */
#include <math.h>
#include <stdint.h>

#define HIL_DT     0.001f          /* 1 ms 步节拍，与被测件控制周期对齐 */
#define HIL_NIN    4u              /* 输入通道：油门、转速、电压、温度 */
#define HIL_NOUT   2u              /* 输出通道：转速、电压 */

typedef struct {
    float x1;                       /* 一阶状态：转速惯性 */
    float x2;                       /* 一阶状态：母线电容电压 */
} PlantState;

static PlantState g_st;             /* 禁止 malloc：地址与执行时间都有界 */

void hil_init(void) { g_st.x1 = 800.0f; g_st.x2 = 12.0f; }

void hil_step(const float in[HIL_NIN], float out[HIL_NOUT])
{
    const float torque = in[0] * 40.0f;                 /* 油门到力矩，线性近似 */
    const float load   = 0.02f * g_st.x1 * g_st.x1;     /* 风阻类负载项 */

    g_st.x1 += HIL_DT * (torque - load) * 0.1f;         /* 显式欧拉：定长无分支 */
    if (g_st.x1 < 0.0f) g_st.x1 = 0.0f;

    g_st.x2 += HIL_DT * (in[2] - g_st.x2) * 0.5f;       /* 电压一阶跟踪 */

    out[0] = g_st.x1;                                   /* 发送任务负责打成帧 */
    out[1] = g_st.x2;                                   /* 步函数内禁止发送与打印 */
}
\end{lstlisting}

四条代码约束各自指向可验证性：\Tp{float32} 优先，因为被测 ECU 的 MCU 普遍只有
单精度 FPU，双精度要么软件模拟拖慢节拍、要么精度白白超出被测件，同精度还让
``仿真与被测件看到的数值一致''在送审时可解释，范围问题用定标（量程加偏移）解决；
禁动态分配，因为堆分配失败、碎片与延迟都是截止期的敌人，静态存储让内存占用
成为链接期常量；定长无分支循环，因为 WCET 工具按指令路径计数，
数据相关的长循环直接破坏上界；输出交给发送任务，因为步函数里做总线发送或打印，
等于把不可控的驱动延迟塞进闭环节拍。

\section{确定性与可复现}
\label{sec:ch15-determinism}

同一输入得到同一轨迹，听起来是``反正代码一样''，实际会被四类因素打破：

\begin{itemize}
  \item \textbf{浮点结合律失效}：\Tp{(a+b)+c} 与 \Tp{a+(b+c)} 逐位不同。
        优化换了求和顺序、向量宽度变了归约分组，结果就跟不上：128 位、256 位、
        512 位的水平归约加法树形状不同，同一份代码在两台 CPU 上就能差出末位。
  \item \textbf{指令选择}：融合乘加（FMA）与``分别乘、舍入、再相加''结果不同；
        GCC、Clang 默认允许 contraction，MSVC 默认保守，跨编译器构建先统一这一项。
  \item \textbf{并行顺序}：OpenMP 归约按线程数分组，任务窃取让分组跨运行变化；
        容器迭代序、线程本地缓存的写入序也会渗进结果。
  \item \textbf{非仿真输入}：墙钟读数、未初始化填充、地址相关的哈希种子、
        异步回调到达顺序——每一项都必须先在步进边界被正规化，否则无从复现。
\end{itemize}

可复现的工程手段是\textbf{记录-回放}：把``仿真外来的信息''全部记成每步一帧（输入、事件数、随机消耗数、状态哈希），帧很小；重放时只消费帧与种子。

\begin{lstlisting}[style=cppstyle, caption={记录端：输入正规化进帧，全仿真只允许一条随机流}, label={lst:ch15-trace}]
struct TraceFrame {
    uint64_t stepId;
    int32_t  inputs[16];        // 总线或手柄输入，定点化后记录，不记录原始时间戳
    uint32_t eventCount;        // 本步注入的事件条数（事件体另进一条有序队列）
    uint32_t rngDraws;          // 本步从唯一种子流取用了几个随机数
    uint64_t stateHash;         // 本步末的全局状态哈希，粗粒度即可
};

class Trace {
    std::mt19937_64 rng_;       // 全仿真仅此一条随机流：禁止再碰 device_random
    std::FILE*      fp_;
public:
    void begin(uint64_t seed) { rng_.seed(seed); /* 帧头：seed、源码版本、编译开关 */ }
    uint64_t nextRandom() { return rng_(); }
    void commit(uint64_t id, const int32_t (&in)[16], uint32_t draws) {
        TraceFrame f{id, {}, eventCount_, draws, hashWorldState()};
        std::memcpy(f.inputs, in, sizeof f.inputs);
        std::fwrite(&f, sizeof f, 1, fp_);
    }
};
\end{lstlisting}

\begin{lstlisting}[style=cppstyle, caption={回放端：只吃帧文件与种子，第一条分歧步就是事故现场}, label={lst:ch15-replay}]
bool replay(const char* tracePath, World& w) {
    std::FILE* fp = std::fopen(tracePath, "rb");
    if (!fp) return false;
    TraceHeader h{};  std::fread(&h, sizeof h, 1, fp);
    w.beginFromSeed(h.seed);                       // 先恢复种子流，再开始步进
    for (uint64_t step = 0;; ++step) {
        TraceFrame f{};
        if (std::fread(&f, sizeof f, 1, fp) != 1) break;   // 帧读完：正常结束
        w.injectInputs(f.inputs);                  // 输入正规化是外部信息的唯一入口
        w.stepFixedDt();                           // 定步长；步长与 dt 记在帧头里
        if (w.hashState() != f.stateHash) { reportDivergence(step); return false; }
    }
    return w.rngDrawTotal() == h.rngDrawsTotal;    // 随机消耗数也要逐级对上
}
\end{lstlisting}

哈希用粗粒度就行（对全部实体状态定点化后做一次 FNV 或 xxHash），目的不是防伪，
而是把逐步比较变成常数时间。这套机制与 Part I 的检查点形状像、目的相反：
检查点保存``已经算出来的结果''为了少算，记录帧保存``算出它所需的外部信息''为了重算可信。

\begin{table}[htbp]
\centering
\caption{检查点与记录帧、联机快照帧的三方对照}
\label{tab:ch15-checkpoint}
\begin{tabular}{>{\raggedright\arraybackslash}p{0.16\textwidth} >{\raggedright\arraybackslash}p{0.20\textwidth} >{\raggedright\arraybackslash}p{0.22\textwidth} >{\raggedright\arraybackslash}p{0.18\textwidth}}
\toprule
维度 & 批处理仿真检查点（Part I） & 实时仿真记录帧 & 联机对局帧与快照 \\
\midrule
目的 & 断点续算与故障恢复，节省的是整批机时 & 验证确定性，把事故精确定位到单一步 & 纠错回滚、观战与断线重连 \\
保存什么 & 全场变量加求解器内部态：矩阵缓存、历史变量、迭代号与随机态 & 输入正规化帧加种子与源码版本，状态本身不落地 & 周期性权威状态，加一段时间窗内的完整输入日志 \\
恢复代价 & I/O 可达整个内存足迹，恢复后还要重建求解器的各种缓存 & 单帧很小，但要从头重放整段，用一段计算换可信度 & 回滚只重放最近几百毫秒，前提是输入历史保持完整 \\
\bottomrule
\end{tabular}
\end{table}

三者还能叠着用：HIL 台架既打检查点又录输入帧，但检查点里不含``世界为什么变成这样''的因果信息，别指望靠它调试非确定性。

\section{网络同步：把时间卖给远端}
\label{sec:ch15-net}

多人游戏里每个客户端各有一条时间线，网络同步的本质是让多条时间线对
``现在世界什么样''达成一致。基本盘是\textbf{客户端-服务器权威}：
服务器（或专用权威进程）持有唯一可信状态，客户端只做预测与呈现。在此之上是三种手段：

\begin{itemize}
  \item \textbf{快照与内插}：服务器周期广播完整或增量快照，客户端在相邻两份
        快照之间插值渲染；代价是每个远端实体恒定显示一到两个周期之前的过去。
  \item \textbf{外推}：不等快照、沿当前速度往前猜，延迟感最低，但猜错要纠正，
        表现为回拉与抖动。
  \item \textbf{回滚（rollback）}：本地输入立即执行并渲染，延迟到达的输入进来时
        把仿真时间倒回、用修正后的输入历史重放到当前。它要求逻辑层确定
        （\S\ref{sec:ch15-promises} 的那个切分），并为每个实体保留输入历史环。
\end{itemize}

延迟补偿是服务器侧的镜像操作：判定命中时把世界回卷到开火者按下扳机的那一刻，
用他当时看到的几何做检测（Valve 系引擎在高 tick 率服务器上的标准做法）。
它公平地补偿高延迟玩家，代价是低延迟玩家可能``明明先打中却死了''——两种失真不可同时消除，选哪种是产品题不是代码题。

\begin{table}[htbp]
\centering
\caption{四种同步手段在帧预算、状态量与带宽上的权衡}
\label{tab:ch15-nettrade}
\begin{tabular}{>{\raggedright\arraybackslash}p{0.13\textwidth} >{\raggedright\arraybackslash}p{0.16\textwidth} >{\raggedright\arraybackslash}p{0.15\textwidth} >{\raggedright\arraybackslash}p{0.13\textwidth} >{\raggedright\arraybackslash}p{0.16\textwidth}}
\toprule
手段 & 帧预算成本 & 需要保存的状态量 & 带宽 & 典型适用 \\
\midrule
全量快照广播 & 序列化整个权威态，每帧成本随实体数线性增长 & 每实体一份可随时序列化的完整权威状态 & 最高，与实体规模同阶 & 实体少的回合制与按局结算类对局 \\
增量快照 & 脏标记维护与新旧状态逐帧比较的成本 & 前一帧基线加差异链，客户端丢包后要重同步基线 & 中，取决于变化率而非总规模 & 大多数 MMO、赛车与开放世界 \\
命令输入加本地确定模拟 & 回滚重放的步数直接吃掉每帧预算 & 输入历史环加可回退的状态栈，两者都要定期裁剪 & 最低：只传输入，不传状态 & 格斗、高速竞技射击与音游联机 \\
状态同步加外推 & 几乎为零，只付显示层的插值与纠正 & 只需当前瞬时位置与速度向量，无历史 & 中，纠正包频率按误差浮动 & 不要求可预测性的非关键远端对象 \\
\bottomrule
\end{tabular}
\end{table}

表中``命令输入''一列的低带宽有一个前置条件：模拟必须确定。这正是多数团队翻车之处：

\begin{warning}[title={把非确定性的物理直接当权威状态复制给每个客户端}]
最常见的错误设计是：服务器跑一套``通用''物理（变步长、依赖线程数或向量宽度
归约的求解器），把每个刚体的状态逐帧广播，同时客户端也跑同一套物理做本地预测。
这个设计两头落空：客户端本地模拟与服务器``同名不同值''，预测每帧都被纠正，
表现为持续回拉；全量广播又吃满带宽；回滚因不敢重放而形同虚设。
出路只有两条，选定其一再开工：把参与预测的物理改造成确定性内核
（定步长、定点或严格约束的浮点、固定归约顺序，参照 \S\ref{sec:ch15-determinism}），
或者干脆不让客户端模拟它——纯状态同步，接受固定的显示延迟。
\end{warning}

\subsection*{本章小结}
实时这一层的三块基石：把时间步形状选对（\S\ref{sec:ch15-timestep}），
把延迟预算按最坏情形算（\S\ref{sec:ch15-hil-loop}），把确定性做成可验证的工程（\S\ref{sec:ch15-determinism}）。
网络同步只是把这三块放上时间线错位的信道。

\begin{bestpractice}[title={进入实时联调前钉死四件事}]
\begin{enumerate}
  \item 权威时钟归属：谁定义``现在''——被测件、服务器还是本地累加器，白纸黑字。
  \item 帧边界：哪些信息必须正规化到步进边界（输入、事件、随机数），禁止绕开。
  \item 一致性判据：状态哈希的粒度与比较频率，先于任何优化定下来。
  \item 退化行为：超截止期、掉线、回滚过深时各层的降级顺序，写成一张表。
\end{enumerate}
这四条在批处理仿真里全都不存在，进 HIL 与联机的前两周会全部找上门。
\end{bestpractice}

% ============================================================
% ---- frag_ch16.tex ----
\part{共用技术、选型与工程实践}

\chapter{数据交换管线：USD、glTF、VTK 与 CGNS}
\label{ch:ch16-dataexchange}

进入第三部分：从这里开始讲两个域\emph{共用}的技术层。第一个问题是
``把几何、场和场景搬到另一个程序里''——\DomCAE{} 把结果交给后处理与交付方，
\DomEngine{} 把美术产出搬进运行时，两边最终都落到同一族格式上：
glTF、USD、VTK 家族、CGNS。本章讲这四者的分工、\CL/{\Bk}\CPP{} 侧的读写路径，
以及交换数据里必须显式写下的约定。

\section{两类数据：场景与场}
\label{sec:ch16-two-kinds}

交换的对象其实只有两类，先分清楚再谈格式。

\textbf{可视化导向的场景}回答``看起来是什么''：变换层级（父子节点、
世界到本地的矩阵）、网格（顶点缓冲加索引缓冲）、材质（PBR 参数与贴图引用）、
动画（关键帧曲线与蒙皮权重）、相机与灯光。它面向渲染器的访问模式：
按绘制批次遍历、按材质分组、索引指向顶点池，体量以 MB 计，
发布前一次性烘好，之后基本只读。

\textbf{仿真导向的场}回答``算出来是什么''：拓扑（节点、单元及其连接关系）、
单元类型（四面体、六面体、棱柱……不同阶次）、变量（挂在节点、单元中心
还是面上的多分量数组）、时间步序列，以及每个量的单位与区域语义。
它面向分析访问模式：切片、等值面、沿程提取、逐分量统计，体量以 GB 计，
每个计算步都在追加写。

一套格式很难同时伺候两者：场景格式假设``少量大网格加材质引用''，
场格式假设``海量小数组加属性元数据''；场景要的是 GPU 友好的交错布局，
场要的是按变量连续存放以便切片；场景的演化半径是美术流程，
场的演化半径是求解器迭代。硬把一类塞进另一类，得到的通常是
``两头都要在导入器里写转换层''的中间态——这正是下面四个格式各占一角的原因。

\section{四种格式的分工}
\label{sec:ch16-formats}

\textbf{glTF}（Khronos 组的运行时场景交换格式）把场景写成一个 JSON 清单
加二进制缓冲（或合并成单文件 glb）：访问器直接描述缓冲里每段的偏移、
分量数与数值类型，导入侧几乎不需要解析就知道怎么把字节交给顶点缓冲。
材质走金属度-粗糙度模型，几何压缩（Draco、meshopt）走扩展位。
它是``交接给运行时''这一段的通用货币。
\textbf{USD}（Pixar 开源、现归 AOUSD 联盟演进的 OpenUSD）解决的是
``大型场景的分层组装''：一切皆 prim 与属性，场景由许多层（layer）叠加而成，
下层来自引用、覆盖来自上层，变体（variant）让同一资产携带多套参数。
文本 usda 便于评审与手改，二进制 usdc 便于加载，usdz 把资源打包成单文件。
它不是为运行时效率设计的，而是为协作编辑与镜头合成设计的。

\textbf{VTK 家族}是可视化管线里的事实标准容器：数据集类型分多块
（polydata、结构网格、非结构网格、图像数据、多重块），网格与场同仓，
每个数组带名字、分量数与元数据数组。XML 系列（vtp/vtu/vts/vti/vtm）
是当前主力，legacy 格式仍在大量存量文件里。ParaView 直接以它为输入。

\textbf{CGNS} 是 CFD 的数据推荐规程（AIAA 标准）：以``基-区-网格-场''
的树状结构组织结构化与非结构化网格，场的位置用网格位置属性区分
顶点中心、单元中心与面中心，区域名与族名靠命名约定承载边界语义，
每个量可以显式附上单位注记。存储后端有 HDF5 与经典 Fortran 二进制两种。

\begin{table}[htbp]
\centering
\caption{四种格式的定位对照}
\label{tab:ch16-formats}
\begin{tabular}{>{\raggedright\arraybackslash}p{0.08\textwidth} >{\raggedright\arraybackslash}p{0.14\textwidth} >{\raggedright\arraybackslash}p{0.16\textwidth} >{\raggedright\arraybackslash}p{0.11\textwidth} >{\raggedright\arraybackslash}p{0.14\textwidth} >{\raggedright\arraybackslash}p{0.15\textwidth}}
\toprule
格式 & 面向对象 & 承载内容 & 文本或二进制 & 谁常用它 & 跨语言方式 \\
\midrule
glTF & 运行时场景交换 & 层级、网格、PBR 材质、动画 & JSON 或 glb 二进制 & 引擎导入器、Web 查看器 & C++ 库加各语言绑定 \\
USD & 分层组装与合成 & prim 层级、属性、变体、时间码 & usda 文本或 usdc 二进制 & 影视管线、大型场景编辑 & C++ 库为主，Python 脚本层完整 \\
VTK 系 & 可视化管线 & 网格拓扑加多分量场数组 & XML 压缩为主，存量为 legacy 文本 & 后处理工具、科学计算 & C++ 核心，Python 包装同源 \\
CGNS & CFD 交换规程 & 结构化与非结构网格、场、单位 & HDF5 或经典二进制 & 气动、燃烧求解器生态 & C 与 Fortran 双接口库 \\
\bottomrule
\end{tabular}
\end{table}

选型上有个速判法：交付给``人看''的运行时场景走 glTF；交付给``多方编辑''
的镜头与大型场景走 USD；交付给``分析与再计算''的体场走 VTK 系或 CGNS。
四个并不互斥——同一资产常常以 USD 编辑、以 glTF 出运行时、以 VTK 系出分析。

\section{C/C++ 侧的读写路径}
\label{sec:ch16-cpp-path}

库依赖：glTF 有轻量单头文件库与全功能导入器两类选择；OpenUSD 自带
C++ 库与工具链；VTK 系建议直接用 VTK 库或独立解析 XML 结构；
CGNS 用官方参考库的中间层——其调用形状是``先定位后读取''：
把游标移动到树中某节点，再按类型调相应的读函数取数组，具体函数名与参数
以官方文档为准。许可口径都要以各源包自带的 LICENSE 为准再进构建清单，
OpenUSD 与 VTK、CGNS 参考库均为宽松许可，但静态链接与再分发条款仍要逐包核对。

真正的工程问题不是``能不能读到''，而是``读到我这里要搬几次''。
外部格式的磁盘布局与内部布局之间，典型有三次拷贝：磁盘到系统页缓存、
页缓存到解析缓冲、解析缓冲到最终数组。压缩数据（Draco、meshopt、XML 的 appended 编码）
还要多出解码一次。值得争取的是最后一步：解析完成后直接落进 SoA 目标数组，
而不是先堆一个临时对象图再转存。清单 \ref{lst:ch16-mesh-layout} 与
清单 \ref{lst:ch16-loader} 给一个自包含的最小闭环：读 VTK legacy 文本网格，
转成内部 SoA，算单元体积派生量，再按清单 \ref{lst:ch16-dumper} 写回成
ParaView 可直接打开的带派生场的文件。

\begin{lstlisting}[language=C++, style=cppstyle, caption={内部布局与单元体积：四面体公式加六面体 Kuhn 剖分}, label={lst:ch16-mesh-layout}]
// mesh.h —— 内部布局：每个数组独立连续，热点循环不跨结构体跳转
#include <cmath>
#include <cstdint>
#include <vector>

struct MeshSoA {
    std::vector<float>    px, py, pz;   // 节点坐标，SoA 三数组
    std::vector<uint32_t> cellType;     // VTK 单元类型号（10 四面体，5 六面体）
    std::vector<uint32_t> start, len;   // 单元 c 的索引在 idx 的 [start[c], +len[c])
    std::vector<int32_t>  idx;          // 展平连接表，不含 VTK 的行首计数
    std::vector<float>    vol;          // 派生量：单元体积，计算后写回
};

static inline float tetVol(const MeshSoA& m, int32_t a, int32_t b, int32_t c, int32_t d)
{   // 三条棱向量混合积除以六；先取差再相乘，避免坐标同数量级时的灾难性抵消
    double ax = m.px[b]-m.px[a], ay = m.py[b]-m.py[a], az = m.pz[b]-m.pz[a];
    double bx = m.px[c]-m.px[a], by = m.py[c]-m.py[a], bz = m.pz[c]-m.pz[a];
    double cx = m.px[d]-m.px[a], cy = m.py[d]-m.py[a], cz = m.pz[d]-m.pz[a];
    return float(std::fabs(ax*(by*cz - bz*cy) - ay*(bx*cz - bz*cx) + az*(bx*cy - by*cx)) / 6.0);
}

static const int32_t kHexToTet[6][4] = {   // 共享体对角线 0-6 的六四面体剖分
    {0,1,2,6}, {0,1,5,6}, {0,4,5,6}, {0,3,7,6}, {0,4,7,6}, {0,3,2,6},
};

static inline float hexVol(const MeshSoA& m, const int32_t* v)
{
    float s = 0.0f;
    for (int t = 0; t < 6; ++t)
        s += tetVol(m, v[kHexToTet[t][0]], v[kHexToTet[t][1]],
                    v[kHexToTet[t][2]], v[kHexToTet[t][3]]);
    return s;
}
\end{lstlisting}

\begin{lstlisting}[language=C++, style=cppstyle, caption={读取 VTK legacy ASCII 非结构网格到内部布局（解析边读边落 SoA，无中间对象图）}, label={lst:ch16-loader}]
// read_legacy.cpp —— 只为演示路径：真实项目请交给成熟解析库
#include "mesh.h"
#include <cstdio>
#include <cstring>

bool loadLegacyVtk(const char* path, MeshSoA& m)
{
    std::FILE* fp = std::fopen(path, "rt");
    if (!fp) return false;
    char line[256], kw[64];
    for (int i = 0; i < 3; ++i)                       // 版本行、标题行、ASCII 行
        if (!std::fgets(line, sizeof line, fp)) { std::fclose(fp); return false; }
    long n = 0, nCells = 0, nTok = 0;
    if (std::fscanf(fp, "%63s %63s", kw, line) != 2   // DATASET UNSTRUCTURED_GRID
        || std::strcmp(kw, "DATASET") != 0) { std::fclose(fp); return false; }
    if (std::fscanf(fp, "%63s %ld %63s", kw, &n, line) != 3
        || std::strcmp(kw, "POINTS") != 0) { std::fclose(fp); return false; }
    m.px.resize(n); m.py.resize(n); m.pz.resize(n);
    for (long i = 0; i < n; ++i)
        if (std::fscanf(fp, "%f %f %f", &m.px[i], &m.py[i], &m.pz[i]) != 3)
            { std::fclose(fp); return false; }
    if (std::fscanf(fp, "%63s %ld %ld", kw, &nCells, &nTok) != 3
        || std::strcmp(kw, "CELLS") != 0) { std::fclose(fp); return false; }
    m.start.resize(nCells); m.len.resize(nCells);      // nTok = 索引数 + 每行计数
    m.idx.reserve(nTok - nCells);                      //   即展平后纯索引的总数
    for (long c = 0; c < nCells; ++c) {
        int cnt = 0;
        std::fscanf(fp, "%d", &cnt);                   // VTK 行首的节点数
        m.start[c] = (uint32_t)m.idx.size(); m.len[c] = (uint32_t)cnt;
        for (int k = 0; k < cnt; ++k) { int v = 0; std::fscanf(fp, "%d", &v);
                                       m.idx.push_back(v); }
    }
    if (std::fscanf(fp, "%63s %ld", kw, &n) != 2
        || std::strcmp(kw, "CELL_TYPES") != 0) { std::fclose(fp); return false; }
    m.cellType.resize(n);
    for (long c = 0; c < n; ++c)
        if (std::fscanf(fp, "%u", &m.cellType[c]) != 1) { std::fclose(fp); return false; }
    std::fclose(fp);
    return true;                                       // 到这里没有构建任何临时对象图
}
\end{lstlisting}

\begin{lstlisting}[language=C++, style=cppstyle, caption={写出带派生场的 legacy 文件：体积作为单元中心标量随网格同行输出}, label={lst:ch16-dumper}]
// write_legacy.cpp —— 派生量与场数据一起交付，后处理端零改造可见
#include "mesh.h"
#include <cstdio>

void computeVolumes(MeshSoA& m)                       // 两种单元各走各的公式
{
    m.vol.assign(m.cellType.size(), 0.0f);
    for (size_t c = 0; c < m.cellType.size(); ++c) {
        const int32_t* v = &m.idx[m.start[c]];
        if      (m.cellType[c] == 10) m.vol[c] = tetVol(m, v[0], v[1], v[2], v[3]);
        else if (m.cellType[c] ==  5) m.vol[c] = hexVol(m, v);
    }
}

bool dumpLegacyVtk(const char* path, const MeshSoA& m)
{
    std::FILE* fp = std::fopen(path, "wt");
    if (!fp) return false;
    std::fprintf(fp, "# vtk DataFile Version 3.0\nvolume demo\nASCII\n"
                      "DATASET UNSTRUCTURED_GRID\n");
    std::fprintf(fp, "POINTS %zu float\n", m.px.size());
    for (size_t i = 0; i < m.px.size(); ++i)
        std::fprintf(fp, "%g %g %g\n", m.px[i], m.py[i], m.pz[i]);
    std::fprintf(fp, "CELLS %zu %zu\n", m.len.size(), m.idx.size() + m.len.size());
    for (size_t c = 0; c < m.len.size(); ++c) {
        std::fprintf(fp, "%u", m.len[c]);
        for (uint32_t k = 0; k < m.len[c]; ++k)
            std::fprintf(fp, " %d", m.idx[m.start[c] + k]);
        std::fprintf(fp, "\n");
    }
    std::fprintf(fp, "CELL_TYPES %zu\n", m.len.size());
    for (size_t c = 0; c < m.cellType.size(); ++c) std::fprintf(fp, "%u\n", m.cellType[c]);
    std::fprintf(fp, "CELL_DATA %zu\n", m.vol.size());        // 派生场：单元中心
    std::fprintf(fp, "SCALARS volume float 1\nLOOKUP_TABLE default\n");
    for (size_t c = 0; c < m.vol.size(); ++c) std::fprintf(fp, "%g\n", m.vol[c]);
    std::fclose(fp);
    return true;
}
\end{lstlisting}

大规模读写靠两条纪律。分块：一次不要搬运整个域——VTK 的 XML 系列支持
按块读取与多重块清单，CGNS 的多区结构天然就是分块单元，HDF5 后端
还能按数据集分块压缩；按区或按块流式处理，内存峰值与文件大小解耦。
并行：写侧多进程各写自己的分片文件、由主文件汇总引用（VTK 的并列
分片系列即为此设计），读侧把分片分给线程池；共享单文件句柄的并发写在
大多数后端上并不安全，别省那一步聚合。

\begin{lstlisting}[language=ShellLD, style=shellstyle, caption={格式间的批处理转换示例：文本与二进制互转、打包、并行读取校验}, label={lst:ch16-convert}]
# USD：文本 usda 与二进制 usdc 互转，目录打包成单文件 usdz
usdcat scene.usda -o scene.usdc
usdzip -o package.usdz usd/ textures/

# glTF：单文件 glb 展开为 JSON 加分离缓冲（便于进版本控制）
gltf-transform expand model.glb model.gltf

# VTK 系：用 ParaView 无头端把分片系列重建成单文件并降采样
pvbatch rebuild.py case_*.vtu merged.vtu

# CGNS：用参考库自带工具做结构校验（函数级检查按官方文档为准）
cgnscheck aero_case.cgns
\end{lstlisting}

\begin{note}[title={派生量放在哪一边算}]
同一个体积、应变或涡量，可以在求解器里顺手算、在交换前离线算、
也可以在查看器表达式里现算。判据是消费方数量与数据体量：多方复用就落地成
场数组随文件走；只有单一查看方、且原始场已带全，就留给查看方的计算器，
别让每个交换文件里都养一套含义漂移的派生副本。
\end{note}

\section{与引擎资产管线及后处理的交界}
\label{sec:ch16-pipeline}

\DomEngine{} 侧的导入器把 glTF/USD 搬进来，接着是烘焙（bake）：展 UV、
烘光照贴图、生成 LOD、网格体重组织、贴图上压缩块（BC 系列、ASTC 一类按平台选）。
交换格式与运行时格式分离的原因可以列成一张表，但根因只有一条：
编辑时优化的是``人能改''，运行时优化的是``机器读得快''，两者的数据结构相反。

\begin{table}[htbp]
\centering
\caption{编辑时与运行时两种形态的分工}
\label{tab:ch16-pipeline}
\begin{tabular}{>{\raggedright\arraybackslash}p{0.16\textwidth} >{\raggedright\arraybackslash}p{0.32\textwidth} >{\raggedright\arraybackslash}p{0.32\textwidth}}
\toprule
维度 & 编辑时形态（USD、源 glTF、DCC 场景） & 运行时形态（cooked 资产、glb、引擎容器） \\
\midrule
内容 & 全量元数据、层与变体、未展 UV 的高模 & 烘焙后的低模、图集、光照贴图、LOD 链 \\
访问 & 随机读写、局部修改、频繁增删属性 & 顺序流式加载、内存映射、基本只读 \\
体积 & 大，含大量版本期才有意义的信息 & 小，按平台压缩并砍掉不可见信息 \\
演化 & 跟美术与镜头流程，天天变 & 跟版本发布，按批重新烘焙 \\
\bottomrule
\end{tabular}
\end{table}

与 Part I 后处理的交界是同一份场数据的两个出口：一路给 ParaView
一类分析工具（VTK 系或 CGNS，切片、等值面、沿程线都是查看器侧现算）；
一路给引擎着色器（把体场烘成 3D 纹理或顶点属性，颜色映射在着色器里做，
数据侧只保证数值与布局对齐）。值得避免的反模式是``给每个出口单独再算一遍''：
派生量在源端定一次口径，两边引用同一份结果，
否则同一物理量会在两个查看器里显示成两个数。

\begin{compare}[title={交换格式与运行时格式的四步转换检查单}]
\begin{itemize}
  \item 导入：保留层级与命名语义，坐标系与单位先按 \S\ref{sec:ch16-units} 校准。
  \item 烘焙：不可逆步骤集中在这一步，保留源文件以便重烤。
  \item 压缩：按目标平台选择，压缩率与解码成本要在加载预算里实测。
  \item 校验：运行时只信任带校验和的容器，导入失败要能定位到源文件对象。
\end{itemize}
\end{compare}

\section{元数据与单位：交换文件的第五列}
\label{sec:ch16-units}

绝大多数``数据没错但画面全歪''的事故来自三件事没有显式写下来。
\textbf{单位制}：glTF 约定米、USD 用米每单位属性承载长度基准、
CGNS 可以给每个量附单位注记，而 VTK 的数据文件基本不带单位字段——
约定全靠文档，于是必须把单位写进属性名或随文件发一页说明书。
CFD 侧还有无量纲化的传统：同一数组在不同文件里可能一个是物理值一个
是参考量比值，读的人无从分辨。

\textbf{坐标系}：手性（左右手）、上轴（Y 还是 Z）、前向轴与深度范围，
glTF 与多数 DCC 工具是右手 Y 上，USD 两种约定都常见、靠属性声明，
渲染坐标到世界坐标再叠加锚点与原点偏移。转换永远发生在导入边界，
绝不要让它下沉进运行时数据里。

\textbf{时间基}：求解器时间（物理秒，重启后可能带偏移）、帧时间
（tick 序号乘固定步长，即 ch15 里累加器循环所钉的那个 \Tp{dt}）、
USD 动画的时间码、查看器动画的快照序号，四者换算关系要落在文件里
或说明书里，不能只在写文件的人脑子里。

\begin{guideline}[title={交换格式里必须显式写下的五件事}]
\begin{enumerate}
  \item 长度单位与换算因子（以米为基准时的米每单位值），无量纲量单独声明。
  \item 坐标系约定：手性、上轴、前向轴、原点锚在哪里。
  \item 时间基：物理秒还是步序号，起点与参考偏移，采样间隔是否均匀。
  \item 场的位置与语义：顶点中心、单元中心还是面中心，区域名代表什么边界。
  \item 数值布局：单双精度、字节序、行列主序，压缩时附原始校验和。
\end{enumerate}
写不下这五条的格式不是格式的错，是\emph{写文件的人}把上下文当成了私有知识。
\end{guideline}

\subsection*{常见陷阱}
\begin{warning}[title={把格式选择当成数据契约}]
选了 USD 不等于场景可组装，选了 VTK 不等于场量对齐：格式只规定容器，
不规定``这里的压力是静压还是总压''``这段动画曲线的时间基准''。
跨团队交换前先签一份数据契约（上述五条的实例化版本），
再谈格式与库；顺序颠倒时，第一次联调会同时暴露格式与语义两层问题，
谁也修不动。另一个高频坑是拿 legacy 文本格式当长期归档：它没有压缩、
没有分块随机读，GB 级场数据会在 I/O 上付出数倍代价。
\end{warning}

本章的清单 \ref{lst:ch16-mesh-layout}、\ref{lst:ch16-loader}、
\ref{lst:ch16-dumper} 是刻意自包含的：不依赖任何交换库，
为的是把``外部文件到内部布局再到派生场''这条路径本身摆清楚——
换成熟解析库时，变的只是第一跳，后面每一跳的纪律原样保留。

% ============================================================
% ---- frag_ch17.tex ----
\chapter{性能与工具链：剖析、GPU 并行、构建与许可}
\label{ch:ch17-perf-toolchain}

前面各章讲``把用户逻辑塞进哪一层''，本章讲``塞进去之后怎么把它跑快、跑稳、跑成能交付的东西''。两个域的算法几乎没有共同点，工程手段却高度重合：都要先测量再动手，都要处理符号与构建类型，都要把依赖的许可口径写进交付清单。真正的分岔只在时间约束的形式上：\DomCAE{} 卖墙钟时间与可复现性，\DomEngine{} 卖每帧的预算。

工具口径先钉住：本章只写稳定的工具族与用法形状——CPU 侧的 \Lib{perf}、\Lib{VTune}、\Lib{Tracy}，GPU 侧的 \Lib{RenderDoc}、\Lib{nsight-compute}，异构编译侧的 \Lib{nvcc}、\Lang{HIP}、\Lang{SYCL}、\Lang{Vulkan compute} 四条路径，并行侧的 \Lang{OpenMP} 与 \Lib{TBB} 一族。厂商版本号一律不写死，命令行选项与界面路径都按版本演进，参数以所用版本的官方文档为准。

\section{两种性能问题：吞吐型与帧预算型}
\label{sec:ch17-two-metrics}

\textbf{吞吐型问题}的目标只有一个：把一次算例的墙钟时间压下去，同时不改结果。它的成本结构是分层可加的——网格读入、初始化、每步的装配与线性求解、后处理与写出；测量单位是``一次算例''与``一次迭代''，判据是阶段占比与并行扩展效率。因为允许跑几十分钟到几小时，采样可以拉长、算例可以重跑、剖析器自身开销能摊薄到忽略不计。

\textbf{帧预算型问题}的目标是分布而不是均值：一帧必须在预算内做完（60 帧对应每帧约 16.6 ms，90 帧约 11.1 ms），超了就是肉眼可见的卡顿。它的成本结构是并发重叠的——逻辑、物理步进、动画、渲染提交、GPU 执行、驱动与显存管理跑在不同线程甚至不同处理器上；测量单位是``一帧''，判据是帧时间的分位数与尖刺频率。平均帧率达标而 99 分位超时，体验照样是``时不时卡一下''。

指标口径不同，工具选择就分岔。吞吐型先要回答``哪个函数在热点里''，采样式 CPU 剖析直接给答案；帧预算型先要回答``这一帧里哪一段占掉了预算''，要靠带时间轴的区间剖析与一帧的 GPU 录制。用错工具的两种典型症状：拿火焰图去看卡顿（采样把一次尖刺抹进平均值），拿逐阶段计时器去看向量化不足（时间戳分辨率看不见循环内部的访存行为）。

\begin{table}[htbp]
\centering
\caption{症状 $\rightarrow$ 先看哪个指标 $\rightarrow$ 先动什么工具}
\label{tab:ch17-symptom}
\begin{tabular}{>{\raggedright\arraybackslash}p{0.20\textwidth}
                >{\raggedright\arraybackslash}p{0.29\textwidth}
                >{\raggedright\arraybackslash}p{0.30\textwidth}}
\toprule
\textbf{症状} & \textbf{先看的指标} & \textbf{先用的工具或手段} \\
\midrule
\DomCAE{}：同算例比上一版慢 & 阶段占比（装配、求解、IO）、每迭代时间 & 分阶段计时日志，加采样式剖析（\Lib{perf}、\Lib{VTune}） \\
\DomCAE{}：加核不加速 & 扩展效率、同步与通信占比、每线程工作量 & 线程数或 rank 数扫描，加剖析器的线程视图 \\
\DomCAE{}：单核就慢 & 每周期指令数、缓存未命中率、向量化是否发生 & \Lib{perf} 的微架构计数，加编译器向量化报告 \\
\DomEngine{}：均值够而偶发卡顿 & 帧时间 99 分位、每帧分配字节数 & 区间式逐帧剖析（\Lib{Tracy} 一类）加分配统计 \\
\DomEngine{}：CPU 帧时间正常而画面慢 & GPU 帧时间、draw call 数、带宽 & \Lib{RenderDoc} 一类的帧捕获 \\
两侧共有：只有打包后出问题 & 构建类型与符号完整性、依赖解析顺序 & 带调试信息的优化构建复现（\S\ref{sec:ch17-profiling}） \\
\bottomrule
\end{tabular}
\end{table}

两域共用的一条纪律是：动手之前先留可重跑的基线——\DomCAE{} 侧是固定输入、固定线程数、固定停止准则的一次运行及其监测量，\DomEngine{} 侧是固定种子、固定输入序列的一段帧轨迹。没有基线时，``变快了''与``结果变了''会同时出现，而你分不清后者的代价能不能接受；基线怎么存、怎么比，见第 \ref{ch:ch16-dataexchange} 章与 \S\ref{sec:ch17-build-license}。

\section{CPU 侧的公共手段：布局、向量化与确定性的代价}
\label{sec:ch17-cpu}

不论求解器还是引擎，CPU 热点的成因分布很集中，因此手段只有三条主线，外加一条关于可复现性的取舍。

\textbf{其一，数据布局。} 内层循环遍历的多个量要连续存放：AoS（结构体数组，每个元素含全部字段）在只读其中两个字段时把缓存行浪费在无关字段上，SoA（每字段一个连续数组）让内层沿连续地址推进。并行时另一件事同样重要：给每个线程私有的累加缓冲并按缓存行大小（常见 64 字节）对齐或填充，否则两个线程写同一条缓存行的不同字段会互相无效化对方的副本——归约看着没有竞争却退化成串行。\DomEngine{} 的任务系统里与 \DomCAE{} 的 \Lang{OpenMP} 归约里是同一个坑。

\textbf{其二，向量化。} 让编译器把内层循环摊成 SIMD 的前提是可检查的：内层没有控制流分叉（分支折成掩码，或把不合格元素提前剔除成紧凑数组）、没有指针别名（\Tp{restrict} 存在的理由，也是多数求解器把热点写在 \CL{} 而非 \CPP{} 里的原因）、内层没有除法与库函数调用（除法换成预先算好的倒数，超越函数换成表或多项式近似，但要连同末位差异记账）。有没有真的向量化不要靠猜，读编译器的向量化报告，与剖析结果对照着看。

\textbf{其三，内存访问剖析。} 把``慢''归因清楚才能改对地方：每周期指令数偏低而缓存未命中率高，是布局与遍历顺序问题；未命中率正常而算术时间仍长，是算法或分支预测问题；两者都正常而墙钟仍长，多半在同步、通信或 IO，不在核里。\Lib{perf} 与 \Lib{VTune} 的微架构计数就是给这一步判定用的。

\textbf{其四，并行归约的确定性代价。} \Lang{OpenMP} 的归约、\Lib{TBB} 的并行归约、MPI 的规约类调用都不承诺加法顺序：线程数或 rank 数一变，部分和的分组就变，结果在末位上不同；末位差异还可能改变判敛路径，表现为迭代次数变化。三种口径按需要选一：（1）固定顺序归约——每线程、每 rank 的部分和按索引排序后串行合并，代价是一次同步与一点时间；（2）高精度累加（成对或补偿求和）把噪声压到判据之下，但换不来逐位一致；（3）承认非确定，把回归基线从``逐位相同''改成``范数容差加监测量趋势''，与第 \ref{ch:ch08-solver-engineering} 章的判据口径连着。实时侧取向相反：跨平台一致的手感比末位可复现更值钱，但``允许的差异''要写成明确判据，见第 \ref{ch:ch13-physics}、\ref{ch:ch15-hil} 章。

清单 \ref{lst:ch17-aos-soa} 把同一段逻辑写两遍：左边是难向量化的形状，右边是它的可向量化改写。这段逻辑是``把所有面上的加权通量求和''，在求解器里是残差与装配核的最小抽象，在引擎里对应粒子系统对邻域贡献的累加。

\begin{lstlisting}[language=C++, style=cppstyle,
  caption={同一段通量累加核：AoS 带分支版与 SoA 可向量化版},
  label={lst:ch17-aos-soa}]
// 版本 A：结构体数组 + 内层分支 + 内层除法。读着顺，向量化不上
struct FaceA { double cx, cy, cz, area, ux, uy, uz, cond; };

double fluxSumA(const std::vector<FaceA>& f)
{
    double s = 0.0;
    for (const auto& x : f) {
        if (x.cond <= 0.0) continue;         // 分叉：向量化第一障碍
        s += x.area * (x.ux * x.cx + x.uy * x.cy
             + x.uz * x.cz) / x.cond;        // 内层除法：换不掉就摊不开
    }
    return s;
}

// 版本 B：同一逻辑的 SoA 写法。分支折进数据，除法换成预置倒数
struct FaceB {
    std::vector<double> cx, cy, cz, ux, uy, uz;
    std::vector<double> area, weight, invCond;   // invCond：非法面预置 0
};

double fluxSumB(const FaceB& f, std::vector<double>& acc)
{
    const int n = static_cast<int>(f.area.size());
    for (int i = 0; i < n; ++i)              // 段 1：只有连续乘加
        acc[i] = f.area[i] * f.weight[i] * (f.ux[i] * f.cx[i]
                 + f.uy[i] * f.cy[i] + f.uz[i] * f.cz[i]);
    double s = 0.0;
    for (int i = 0; i < n; ++i) s += acc[i] * f.invCond[i];  // 段 2：归约
    return s;
}
\end{lstlisting}

两边换来的是三件事：内层没有分叉（\Tp{continue} 变成 \Tp{invCond} 里的零）、除法不再出现在循环体里、两次遍历各自都是连续地址上的乘加与一个归约。代价写在明面上：字段拆开后要自己维护一致性（多一个 \Tp{acc} 缓冲、多一份失效逻辑），倒数换除法会改末位——这个改动必须与基线容差一起决定，不能在调优的兴奋里偷偷做。\Tp{acc} 这里按单缓冲写，真并行时应展开成每线程一段、按缓存行对齐、最后一次合并，正是本节第一条要防的伪共享。

\section{GPU 计算：四条路的取舍与上卡之前的三件事}
\label{sec:ch17-gpu}

把热点搬上 GPU 有四种入口，算法能力差别不大，差别在可控性与交付面：

\begin{table}[htbp]
\centering
\caption{四条异构计算路径的可移植性、代价与适用场景}
\label{tab:ch17-gpu-routes}
\begin{tabular}{>{\raggedright\arraybackslash}p{0.12\textwidth}
                >{\raggedright\arraybackslash}p{0.17\textwidth}
                >{\raggedright\arraybackslash}p{0.23\textwidth}
                >{\raggedright\arraybackslash}p{0.26\textwidth}}
\toprule
\textbf{路径} & \textbf{可移植性} & \textbf{控制权与代价} & \textbf{什么时候选它} \\
\midrule
\Lang{CUDA}（\Lib{nvcc}） & 绑定单一厂商的 GPU & 最细的调度与内存控制，内核分析器与数学库生态最完整 & 已有 \Lang{CUDA} 资产；要榨内核性能；要用现成高性能库 \\
\Lang{HIP} & 源码形状贴近 \Lang{CUDA}，可落到第二家平台 & 迁移层薄，运行时接口仍有平台绑定 & 把既有 \Lang{CUDA} 求解器带进另一家硬件，不想重写算法 \\
\Lang{SYCL} & 标准驱动的单源异构，同一份 \CPP{} 换后端 & 抽象层吃掉一部分控制权，成熟度按后端差别明显 & 交付面要跨多家硬件；团队以 \CPP{} 为主；不追极致调优 \\
\Lang{Vulkan compute} & 跟着图形 API 走，与渲染共用队列与内存模型 & 显式内存与同步全自己写，开发量最大 & 引擎已有 \Lang{Vulkan} 栈；计算与渲染强耦合（粒子、后处理、GPU 驱动的步进） \\
\bottomrule
\end{tabular}
\end{table}

上卡之后成本结构换了主角，三笔账必须先算。\textbf{启动开销}：每次内核启动有固定的微秒级代价，单次工作量只有几千次的循环在 CPU 上是纳秒级、在 GPU 上被启动开销淹没，对策是算子融合与批量提交，而不是把离线求解器里``一算子一内核''的习惯照搬过去。\textbf{数据搬运}：主机与设备之间的拷贝受链路与页管理限制，可分页内存的拷贝还会被驱动串行化；对策是数据常驻、用锁页内存与异步传输让搬运与计算重叠、把必须回传的量压到一个标量。\textbf{同步与判敛}：把中间结果读回主机再决定是否继续，等于每次插入一次全同步——残差判敛留在设备上、只回传一个标志位，这是上卡时的设计要求而非后期优化。

在此之上，把求解器热点搬上 GPU 之前该做的三件事，顺序不能颠倒：

\begin{enumerate}
\item \textbf{先量清楚它是不是真的算得慢}。用 \S\ref{sec:ch17-two-metrics} 的分阶段计时确认时间在这个核里，再看算术强度（每次访存做多少浮点操作）。热点若在稀疏随机访存、在 IO、在集合通信，上卡多半不会赢。
\item \textbf{先把形状改成设备友好}。索引规则化（分块、按线程映射）、SoA 或交错布局、邻域读收敛到共享内存、避免一个线程写别人正在读的槽位。第 \ref{ch:ch06-openfoam}、\ref{ch:ch08-solver-engineering} 章的离散格式在这一步差别很大：格式越规则，搬上去越顺。
\item \textbf{先建对拍再换实现}。留一份 CPU 参考实现，容差按第 \ref{ch:ch16-dataexchange} 章的基线口径定死，先对单个内核、再对单个时间步、最后对整段历史。缺这一步时，GPU 版本的末位差异会被误读成算法错误，或者反过来。
\end{enumerate}

清单 \ref{lst:ch17-gpu-kernel} 给一个通用骨架：分块加共享内存的重叠边界，主机侧只启动与首尾搬运。形状对 \Lang{CUDA} 与 \Lang{HIP} 都成立（后者多数运行时接口以 \Tp{hip} 前缀对应同类调用）；\Lang{SYCL} 与 \Lang{Vulkan compute} 的同义写法换成``队列提交加缓冲区或无别名的设备指针''，分块与边界处理的思想一致。

\begin{lstlisting}[language=C++, style=cppstyle,
  caption={迭代型内核骨架：分块、共享边界、乒乓缓冲与首尾搬运},
  label={lst:ch17-gpu-kernel}]
// 形状示意：一次 Jacobi 扫掠。具体 API 名以官方文档为准
#define TILE 32                        // 块边长：与设备调度规模匹配

__global__ void jacobiStep(const double* __restrict__ in,
                           double* __restrict__ out,
                           const double* __restrict__ rhs,
                           int nx, int ny)
{
    __shared__ double tile[TILE + 2][TILE + 2];   // 一圈 halo 存邻居
    const int gx = blockIdx.x * TILE + threadIdx.x;
    const int gy = blockIdx.y * TILE + threadIdx.y;
    if (gx >= nx || gy >= ny) return;             // 块外线程退出
    const int tx = threadIdx.x + 1, ty = threadIdx.y + 1;
    tile[ty][tx] = in[gy * nx + gx];              // 沿列读，可合并
    __syncthreads();                              // 等整块数据到位
    const double s = tile[ty-1][tx] + tile[ty+1][tx]
                   + tile[ty][tx-1] + tile[ty][tx+1];
    out[gy * nx + gx] = 0.25 * (s + rhs[gy * nx + gx]);
}

// 主机侧：只在头尾搬运，推进留在设备里
void sweepOnDevice(double* hu, const double* hr, int nx, int ny, int sweeps)
{
    const size_t bytes = sizeof(double) * nx * ny;
    double *src, *dst, *dr;
    cudaMalloc(&src, bytes); cudaMalloc(&dst, bytes); cudaMalloc(&dr, bytes);
    cudaMemcpy(src, hu, bytes, cudaMemcpyHostToDevice);   // 一次上行
    cudaMemcpy(dr,  hr, bytes, cudaMemcpyHostToDevice);
    const dim3 blk(TILE, TILE), grd((nx+TILE-1)/TILE, (ny+TILE-1)/TILE);
    for (int k = 0; k < sweeps; ++k) {
        jacobiStep<<<grd, blk>>>(src, dst, dr, nx, ny);
        double* t = src; src = dst; dst = t;    // 乒乓，免每步清零
    }
    cudaMemcpy(hu, src, bytes, cudaMemcpyDeviceToHost);   // 一次下行
    cudaFree(src); cudaFree(dst); cudaFree(dr);
}
\end{lstlisting}

两处取舍值得说明。\Tp{tile} 多出一圈 halo 是为了把邻域读变成共享内存读，块尺寸与 halo 的比决定``有效算术除以搬运量''，块一小固定开销就抬头；更进一步是把整个时间推进留在设备上，只在判敛或写出时回传，否则每次回传都是一个同步点。写出尤其要在设备上先压缩或抽取切片，再按第 \ref{ch:ch16-dataexchange} 章的格式路由传回——
整场搬运是大规模算例里最常见的隐性成本。

\section{剖析、火焰图、帧捕获与逐算子计时}
\label{sec:ch17-profiling}

采集方式只有两种，回答的问题不同。\textbf{采样式}按固定间隔中断并记录调用栈，开销低、给统计意义上的热点；火焰图的条宽就是它的折叠结果，即同一调用路径的样本占比。读火焰图要记住三件事：条宽是占比不是绝对时长、递归会被折叠成一条、它看不见``等待''（阻塞在锁或 IO 上而没被调度的线程不贡献样本）。\textbf{区间式}在代码里显式插入成对标记（zone、scope、label），开销按标记次数计费，但给出时间轴、并发关系与边界，因此是帧预算问题的唯一解：要知道尖刺出现在哪一帧的哪一段，只能靠区间。

\begin{compare}[title={帧捕获与逐算子计时：边界不同，手段就不通用}]
\DomEngine{} 需要\textbf{帧捕获}，因为一帧是天然的事务边界：把一帧的全部绘制调用、状态切换与资源绑定录下来，就能离线逐个回放、单独看某一次调用改了什么。它的前提是该域本来就在一个命令队列上提交工作；仿真的算例没有``帧''这个边界，捕获无从对齐。

\DomCAE{} 需要\textbf{逐算子计时}，因为一次算例的时间是阶段累加的：读网格、初始化、装配、线性求解、更新、写出，每段一个计时器，按迭代或时间步汇总落盘。引擎里同样可以记每帧的系统级耗时，但拿它去找``哪一次绘制贵''无能为力，那必须回到帧捕获与内核级分析。
\end{compare}

符号与构建类型决定这些手段能不能读出人话。发布构建通常全优化并剥离符号，栈展开信息一缺，火焰图上只剩地址或一段模板名；纯调试构建又慢得不能用来测性能。工程上稳妥的口径是\textbf{带调试信息的优化构建}（\Tp{CMAKE\_BUILD\_TYPE} 取 RelWithDebInfo 一类，源码构建里对应``开发期带调试信息的优化配置''），并保留帧指针或等价的展开数据，让采样器能还原调用链。另外三件事持续影响可读性：内联让时间归属偏到调用者、链接期优化让函数边界模糊、模板与重载需要事后还原成可读名。引擎侧还多一层：脚本或可视化脚本层的标记若没接进剖析器，时间轴上只剩运行时函数名，看不见内容层谁在驱动。

清单 \ref{lst:ch17-profile-run} 给一段采集的形状，参数一律写占位符：换算法换机器时只有路径与规模在变，命令结构不变。

\begin{lstlisting}[language=ShellLD, style=shellstyle,
  caption={采集清单：采样、微架构计数、内核指标与帧级记录},
  label={lst:ch17-profile-run}]
# 1) 吞吐侧：采样式热点。先用最小可复现算例，别拿生产算例试
perf record -g -o sim.data -- ./solver --case <case> --threads <T>
perf report -i sim.data --stdio > hotspots.txt

# 2) 同一算例再看微架构计数：分辨"算得慢"与"取得慢"
perf stat -e task-clock,cycles,instructions,cache-misses -- ./solver --case <case>

# 3) 并行效率：固定输入扫一遍线程数，把每阶段占比落盘
for T in 1 2 4 8 16; do
  ./solver --case <case> --threads $T --stage-timing stage_$T.csv
done

# 4) 编译侧：让编译器自述向量化结果（开关名按所用编译器文档选）
cc -O2 -march=<arch> -fopt-info-vec -c src/kernel.cpp -o /dev/null

# 5) GPU 计算：内核级指标，抓少数几个启动就够，全量会慢到不可用
ncu --set full --launch-skip <K> --launch-count 1 -o gpu ./solver --case <case>

# 6) 帧预算侧：区间剖析要程序内埋点，采集端与查看端同版本
#    用 Tracy 的捕获端录制一段帧（参数按其文档），存成可回放的轨迹文件
# 7) 帧捕获：抓一帧，逐个回放到可疑的绘制调用；异步提交要开记录选项
\end{lstlisting}

第 3 段的循环是两域都用得上的最小测量法：一张``阶段占比随线程数变化''的表，比任何单点剖析都能更快说明该动哪里。区间工具的插桩宏要收在自己的构建开关里，交付构建关掉，否则埋点本身进了帧预算；写出的字节数与打开文件次数也建议记进阶段日志，IO 常在榜首。

\section{构建与许可：一套 CMake 服务两种产物}
\label{sec:ch17-build-license}

\DomCAE{} 侧的产物通常是可复用的求解器核心库（供人链接、供人调用），\DomEngine{} 侧的产物通常是运行时插件（供人装载）。一个仓库同时出这两种产物时，构建要点集中在四处。\textbf{目标形态}：核心库编成位置无关的目标（否则进不了插件），插件用模块目标而不是普通共享库，跨边界只暴露 \CL{} 形状的入口（第 \ref{ch:ch01-boundary} 章那份合同），默认隐藏符号可见性、显式导出少数字符串名。\textbf{ABI 口径}：链接进同一进程的东西必须在位数、编译器运行库选择（Windows 上调试与非调试运行库不可混）、标准库调试宏、结构体对齐默认值、浮点模型上一致；升级时要重编用户子程序与插件，原因就在这节（第 \ref{ch:ch03-fluent-udf}、\ref{ch:ch05-abaqus-usub} 章各有实例）。\textbf{搜索路径}：Linux 与 macOS 在安装阶段写相对加载路径，让插件带着依赖一起搬；Windows 没有等价机制，依赖必须与可执行文件或插件同目录，或显式进搜索路径。引擎打包会把库搬进包内目录，于是任何绝对路径的构建残留都会在真机上失效——``编辑器正常、真机崩''要第一时间查这条（第 \ref{ch:ch10-unreal}、\ref{ch:ch11-unity} 章的打包口径）。\textbf{配置面}：工具链文件与配置预设把``编译器、架构、选项''钉成一份可提交的文件，让本地与 CI 同口径；后端差异（是否启用某家 GPU）用选项开关，而不是散落的宏。

\begin{lstlisting}[language=CMakeLD, style=cmakestyle,
  caption={一份构建脚本同时出求解器库与引擎插件（形状示意）},
  label={lst:ch17-cmake}]
cmake_minimum_required(VERSION 3.20)     # 下限按团队口径，别各处自写
project(sim_ext LANGUAGES CXX C)
option(USE_GPU_BACKEND "启用异构计算后端" OFF)

add_library(simcore STATIC src/csr.cpp src/solve.cpp)
target_compile_features(simcore PRIVATE cxx_std_17)
set_target_properties(simcore PROPERTIES POSITION_INDEPENDENT_CODE ON)
if(USE_GPU_BACKEND)                       # 后端差异只在这里出现一次
  target_sources(simcore PRIVATE src/jacobi_kernel.cu)
endif()
install(TARGETS simcore EXPORT simcoreTargets ARCHIVE DESTINATION lib)
install(EXPORT simcoreTargets NAMESPACE simcore::
        DESTINATION lib/cmake/simcore)   # 求解器侧：下游 find_package 可用

add_library(simplugin MODULE src/plugin_main.cpp)   # 引擎侧：C 形状入口
set_target_properties(simplugin PROPERTIES CXX_VISIBILITY_PRESET hidden)
if(APPLE)
  set_target_properties(simplugin PROPERTIES INSTALL_RPATH "@loader_path")
elseif(UNIX)
  set_target_properties(simplugin PROPERTIES INSTALL_RPATH "$ORIGIN")
endif()
# Windows 无等价机制：依赖 DLL 必须与插件同目录或进 PATH
\end{lstlisting}

许可口径是交付的一部分，不是法务的附录。按对商用发行的影响，依赖大致分三层：

\begin{table}[htbp]
\centering
\caption{依赖许可的三层口径与发行时的动作}
\label{tab:ch17-license}
\begin{tabular}{>{\raggedright\arraybackslash}p{0.16\textwidth}
                >{\raggedright\arraybackslash}p{0.31\textwidth}
                >{\raggedright\arraybackslash}p{0.31\textwidth}}
\toprule
\textbf{层级} & \textbf{典型口径} & \textbf{发行时的动作} \\
\midrule
宽松族 & BSD、MIT 一类允许闭源再分发；Apache 2.0 同为宽松但附带专利授权与声明义务 & 随制品附许可证文本与声明清单，保留版权与 NOTICE 类内容 \\
弱 copyleft & LGPL 一类允许被使用，但要求用户能替换该库；静态链接把问题带到可重链接上 & 优先动态装载；静态链接时提供可重链接的物件文件或等价替换手段 \\
强 copyleft & GPL 一类覆盖到衍生作品并影响发行方式，AGPL 把义务延伸到网络服务 & 在架构上划成独立进程或独立交付件；能不能用先过许可判定再写代码 \\
\bottomrule
\end{tabular}
\end{table}

两点常见误判值得点名。第一，``用户写的扩展代码只是因为调用了接口''通常不等于接口所有者的作品，但这个判断取决于许可文本与调用方式，属于先确认再承诺的条款，不该由工程师替法务下结论。第二，插件与打包让许可义务在交付时才成形：开发期随便用的依赖，进了发行包就可能要求提供源码获取途径或替换说明，所以许可扫描要挂在打包步骤上而不是发布之后。

CI 侧三条纪律与两域都相关：\textbf{缓存键要包含影响 ABI 与数值的输入}（编译器标识、关键选项、平台），否则``CI 过了本地不过''多半是被旧制品骗了；\textbf{制品仓库区分源码制品与二进制制品}，记录构建配置与依赖版本，让``复现一次历史算例''成为可能；\textbf{回归基线分层存}（冒烟、数值容差比对、少量全量快照），别让每次提交都跑全量。数值基线的比较口径与快照怎么存，见第 \ref{ch:ch16-dataexchange} 章。

\begin{guideline}[title={离线与实时共用的工程纪律：可执行清单}]
\begin{itemize}
\item 动手前四问：时间花在哪个阶段、单位是均值还是分位、改了结果会不会变、怎么知道我改坏了。答不上就先加计时与基线，不要先改代码。
\item 每个热点改动留两版共存，用一个开关切换、用同一份输入对拍；对拍判据写成范数容差加监测点，不写``看起来一样''。
\item 剖析构建一律用带调试信息的优化配置；交付构建关掉区间埋点，把插桩宏收进一个编译开关里。
\item GPU 只在形状合适时上：先量算术强度，再改索引形状，最后才换后端；启动、搬运与内核时间要分开记。
\item 跨进程边界的合同只用 \CL{} 类型与版本号；插件隐藏符号，依赖随产物走相对路径。
\item 许可清单与构建元数据随制品发布：许可证、编译器与选项口径、基线来源，缺一项就把``可复现''变成口头承诺。
\end{itemize}
\end{guideline}

\begin{warning}[title={把剖析结果当结论的三种方式}]
\begin{itemize}
\item \textbf{只信一条火焰图}：采样式看不见等待与顺序，把``被调度到的时间''当成``全部时间''。并行慢与 IO 慢在火焰图上长得一样，必须补阶段计时与计数。
\item \textbf{用未剥离的构建测性能}：符号齐了，代码也变了。要么接受带调试信息的优化构建的少量差异，要么把剖析与基准分两次跑，绝不混着报数字。
\item \textbf{在优化里偷偷改数值语义}：倒数换除法、折分支、开快速数学选项，都会在末位甚至判敛路径上留痕。这类改动要走显式评审并同步更新基线容差；把``快了 8\%''和``迭代次数从 41 变 39''写进同一张单子的，才算合格的优化提交。
\end{itemize}
\end{warning}

\subsection*{本章小结}
性能工作在两侧是同一套纪律、两套指标：\DomCAE{} 用阶段占比与扩展效率回答``墙钟去哪了''，\DomEngine{} 用帧分位与帧捕获回答``这一帧去哪了''。CPU 侧的布局、向量化与内存访问剖析两边共用，只有确定性代价需要各表其词——离线要末位可复现，实时要手感一致。GPU 不是另一种语言的问题而是另一种形状的问题：先量、再改形状、再建对拍。构建与许可决定这些东西能不能交付，ABI、搜索路径与许可证清单要和性能数字一起进制品仓库。第 \ref{ch:ch18-decision} 章把这些判据压成决策树与排障手册。

% ============================================================
% ---- frag_ch18.tex ----
\chapter{选型决策树与排障手册}
\label{ch:ch18-decision}

这一章是全书的收束工具，不是感想集。它假定你已经读过前面各章，只做三件事：把``该走哪条路''
压成能照着回答的判据，把``出了问题从哪查起''压成一张可对照的表，把前面各章的口径汇成一页式
对照。写法上尽量用``如果……就走……''的形，而不是``某某更好''的形——两域的产品迭代都很快，
能长期成立的判据只有任务形状、时间约束与交付面这三样。

\section{第一层：问题属于哪一类}
\label{sec:ch18-classify}

第一层要分出三类任务：\textbf{批处理仿真}（离线算完再评价）、\textbf{实时仿真}（物理与逻辑在
硬预算内推进）、\textbf{交互渲染}（以画面与输入为主，物理是配角）。判据是三问，按顺序问，
第一问就能把大多数情况分开。

\textbf{第一问：这份结果要拿去签字吗？} 要——结论会进设计评审、强度评定、合规材料，那么属于
批处理仿真：它接受长墙钟时间，要求收敛研究、网格无关性验证与可复现记录（第
\ref{ch:ch01-boundary} 章的分界，工程细节见第 \ref{ch:ch08-solver-engineering} 章）。不要，
只要看着对、手感对、能交互，就落到 \DomEngine{} 一侧（第 \ref{ch:ch09-engine-layers} 章）。

\textbf{第二问：时间约束是总量上限还是硬分布上限？} 只受排队截止时间约束，允许一次算几小时
——批处理侧。每帧必须做完，超了立刻可见——实时侧（预算与降级策略在第 \ref{ch:ch15-hil}
章，物理侧的取舍在第 \ref{ch:ch13-physics} 章）。这一问解决``看着像仿真但其实是渲染''的情况：
把物理解析精度降一档来换帧率，是实时域的正常动作；在批处理域里等价动作通常叫``模型简化''，
要走显式的误差评估。

\textbf{第三问：用户逻辑被调用的频次与可见性？} 每个积分点、每个面、每帧都调——逻辑必须靠近
内核（编译型扩展或原生插件）。只在建模期、工况切换期执行一次——脚本层更好维护（第
\ref{ch:ch02-ansys-layers} 章的层划分）。要访问内核内部数据结构，也是靠近内核的那一侧。

\begin{table}[htbp]
\centering
\small
\caption{三类任务的分界判据与对应章}
\label{tab:ch18-three-kinds}
\begin{tabular}{>{\raggedright\arraybackslash}p{0.21\textwidth}
                >{\raggedright\arraybackslash}p{0.28\textwidth}
                >{\raggedright\arraybackslash}p{0.27\textwidth}}
\toprule
\textbf{任务类型} & \textbf{判据（何时判给它）} & \textbf{主要相关章} \\
\midrule
批处理仿真 & 精度与可复现优先；单次求解可跑分钟到小时；结果要签字归档 &
第 \ref{ch:ch01-boundary}、\ref{ch:ch08-solver-engineering}、\ref{ch:ch16-dataexchange} 章 \\
实时仿真 & 每帧硬预算；允许降级换稳定性；状态要可回放、可同步 &
第 \ref{ch:ch09-engine-layers}、\ref{ch:ch13-physics}、\ref{ch:ch15-hil} 章 \\
交互渲染 & 主角是画面与输入延迟；物理与仿真只是场景元素 &
第 \ref{ch:ch10-unreal}、\ref{ch:ch11-unity}、\ref{ch:ch12-godot-ext} 章 \\
\bottomrule
\end{tabular}
\end{table}

\begin{compare}[title={如果……就走……：第一层的六条分支}]
\begin{itemize}
\item 如果一次算例跑几十分钟到几天、结果要提交评审：走 \DomCAE{} 批处理路线，入口先从脚本层
      与批处理参数开始（第 \ref{ch:ch02-ansys-layers} 章），热点逻辑再考虑编译型扩展。
\item 如果结果每帧都参与决策（角色、载具、可破坏场景）：走 \DomEngine{} 路线，先定分层与帧
      预算（第 \ref{ch:ch09-engine-layers}、\ref{ch:ch15-hil} 章），再选物理解析层。
\item 如果要离线高保真算一遍、线上实时表现另一遍：两条路线都要，但共享一个离散格式或降阶模型，
      用两个构建目标区分（第 \ref{ch:ch14-binding} 章的绑定层是常见交接点）。
\item 如果只是把既有工况批量跑完、不改内核：留在脚本宏层，不要引入编译扩展——扩展层的维护
      成本在这里换不回收益。
\item 如果要改本构、源项或单元方程：必须走编译型用户函数（第 \ref{ch:ch03-fluent-udf}、
      \ref{ch:ch05-abaqus-usub} 章），因为只有它落在积分点频次上。
\item 如果要在不重启进程的前提下换实现、或把第三方库带进项目：走运行时插件与反射（第
      \ref{ch:ch12-godot-ext}、\ref{ch:ch10-unreal} 章），代价是自己维护稳定 ABI。
\end{itemize}
\end{compare}

\section{第二层：在选定域内选扩展入口}
\label{sec:ch18-entry-choice}

第二层的候选只有四个：脚本宏、编译型用户函数、运行时插件、源码级修改。判据是五项——迭代速度
需求、要访问的内部量、部署面、许可与合规、团队语言栈。下表按``多数情况下''给值，个别产品有
特例，以第 \ref{ch:ch01-boundary} 章的三形态口径为准。

\begin{table}[htbp]
\centering
\footnotesize
\setlength{\tabcolsep}{4pt}
\caption{扩展入口的四选一：五项判据对照}
\label{tab:ch18-entry-matrix}
\begin{tabular}{>{\raggedright\arraybackslash}p{0.13\textwidth}
                >{\raggedright\arraybackslash}p{0.13\textwidth}
                >{\raggedright\arraybackslash}p{0.13\textwidth}
                >{\raggedright\arraybackslash}p{0.14\textwidth}
                >{\raggedright\arraybackslash}p{0.14\textwidth}
                >{\raggedright\arraybackslash}p{0.14\textwidth}}
\toprule
\textbf{入口} & \textbf{迭代速度} & \textbf{可访问内部量} & \textbf{部署面} &
\textbf{许可与合规} & \textbf{语言栈} \\
\midrule
脚本宏 & 命令与工况级；每步内部才重算 & 只见命令面与已导出变量 & 随算例或场景走，无需工具链 &
一般不触发额外义务 & \PY{}、领域宏语言 \\
编译型用户函数 & 每积分点、每面、每步 & 场量、材料槽、雅可比 & 要匹配产品自带工具链与位数 &
注意链接的运行时与第三方库许可 & \CL{} 为主 \\
运行时插件 & 每帧或每次调用 & 反射元数据、子系统句柄 & 打包清单与加载顺序都要管 &
发行时要列许可证清单 & \CPP{} 为主 \\
源码级 & 与内核同速，可改一切 & 全部 & 自己维护分支与升级 & 上游许可决定发行方式 &
\CPP{}、\Lang{Fortran} \\
\bottomrule
\end{tabular}
\end{table}

四条最短路径值得单独记住。其一，任务能拆成粗粒度步骤就留在脚本层；一旦脚本层出现每帧几千次
的循环，就该下沉（第 \ref{ch:ch11-unity} 章的托管层开销、第 \ref{ch:ch02-ansys-layers} 章的
参数化循环是同一件事的两种表现）。其二，要碰本构、源项、单元方程，只有编译型入口，且要连同
工具链与重编纪律一起接受（第 \ref{ch:ch03-fluent-udf}、\ref{ch:ch05-abaqus-usub} 章）。其三，
要热重载与不重启换实现，就走插件；同时接受 ABI、静态初始化顺序与符号可见性这三笔长期债务（第
\ref{ch:ch12-godot-ext}、\ref{ch:ch10-unreal} 章）。其四，要跨语言把内核能力暴露给脚本，先定绑定
层的类型合同，别逐项手抄（第 \ref{ch:ch14-binding} 章）。

\section{第三层：数据交换与发行路径}
\label{sec:ch18-exchange-routing}

第 \ref{ch:ch16-dataexchange} 章按格式本身给了四者的分工，这里按``给谁看''与``进不进回归''把它们
收拢成场景。\textbf{USD}：多方编辑、大型场景组装与合成走它，分层与变体正是它的价值，也是它不
该直接当运行时交付物的原因。\textbf{glTF}：交付给运行时与对外查看的默认出口，几何与材质一次
烘好、之后只读，属于``发布路径''而不是``编辑路径''。\textbf{VTK 家族}：分析侧的主力容器，网格
与场同仓、数组自带名字与元数据，后处理工具直接吃它；体量大，按切片与关键截面输出而不是整场
搬运。\textbf{CGNS}：CFD 生态内的网格与场交换规程，带单位注记、靠命名约定承载边界语义，适合
跨工具复算。四类之外还有一类必须显式存在：文本类（CSV、JSON 一类）——它不在那四种格式里，却
是回归基线的主力载体，因为可 diff、体积小、脚本能直接读；代价是位数与单位要写进格式约定，
否则``回归差异''其实是打印差异。

发行路径的判定顺序是：先定``基线存什么''再定``过程存什么''。基线取``文本类监测量加若干关键
截面''，因为全量场快照对浮点顺序敏感且体积不可持续；过程数据取 VTK 系或 CGNS，并只在判敛失败
时保留一次全量；交付物按受众分流——给引擎与查看器的走 glTF（编辑面留在 USD），给分析与再
计算的走 VTK 系或 CGNS，无论哪一路都配一份文本摘要与单位说明。比较口径分三层：冒烟层只判
``是否算完、是否出现非有限值''；数值层用范数容差加监测点容差；快照层用少量全量比对。逐位相同
只允许用于确定性构建（单线程或固定顺序归约，见第 \ref{ch:ch08-solver-engineering} 章与第
\ref{ch:ch17-perf-toolchain} 章的归约讨论）。

\begin{lstlisting}[style=yamlstyle,
  caption={结果输出与回归路由清单（示意：键名由项目自定，不属任何产品）},
  label={lst:ch18-routing}]
# outputs: 每个产物声明三件事——给谁看、走哪种格式、放哪里留多久
outputs:
  - role: baseline          # 回归基线与评审表格：可 diff、体积小
    format: csv_or_json
    fields: [monitor_mass, monitor_drag, probe_max_T]
    sink: artifact_store
    keep: 90d
  - role: editing           # 多方编辑与大型场景组装
    format: usd
    variant_sets: true      # 同一资产携带多套参数
    sink: shared_fs
    keep: until_release
  - role: runtime           # 交付给引擎与查看器：一次性烘好、之后只读
    format: gltf
    bake: mesh_material     # 几何与材质在发布前落地
    sink: package
  - role: analysis          # 后处理与再计算：切片优先，别整场搬
    format: vtk_family
    slice_only: true
    sink: shared_fs
    keep: 7d
  - role: exchange_cfd      # 跨工具的网格与场：单位注记要显式写
    format: cgns
    units: explicit
    sink: archive
regression:
  layers: [smoke, numeric, snapshot]   # 三层判据，见上面的发行路径段
  numeric:
    metric: relative_l2     # 逐位比较只对确定性构建开放
    rel_tol: 1.0e-8
    probes: absolute
    abs_tol: 1.0e-10
  on_fail: keep_last_full_field        # 只在失败时留一次全量
\end{lstlisting}

\section{排障手册：症状、根因、定位手段}
\label{sec:ch18-troubleshoot}

用法约定：先按症状定位行，再照``第一步手段''做，不要跳步。表里的章号是解释原理的地方，本章只给
判定动作；凡涉及具体产品选项名，一律以官方文档为准。

\begingroup\small
\begin{longtable}{>{\raggedright\arraybackslash}p{0.22\textwidth}
                 >{\raggedright\arraybackslash}p{0.21\textwidth}
                 >{\raggedright\arraybackslash}p{0.24\textwidth}
                 >{\raggedright\arraybackslash}p{0.09\textwidth}}
\caption{排障手册：症状、最可能根因与第一步定位手段}\label{tab:ch18-troubleshoot}\\
\toprule
\textbf{症状} & \textbf{最可能的根因} & \textbf{第一步定位手段} & \textbf{相关章} \\
\midrule
\endfirsthead
\toprule
\textbf{症状} & \textbf{最可能的根因} & \textbf{第一步定位手段} & \textbf{相关章} \\
\midrule
\endhead
\midrule
\endfoot
\bottomrule
\endlastfoot
用户函数结果只在并行下与串行不同 & 逐点钩子里做了全局汇总，或用静态变量缓存，每个 rank 各一份 & 先跑单 rank 复现，再把汇总移到每迭代一次的钩子里 & \ref{ch:ch03-fluent-udf}、\ref{ch:ch08-solver-engineering} \\
编译加载的用户子程序报未定义符号 & 工具链、位数或运行时库与产品构建不一致 & 比对导出符号表；核对位数、运行库与调试宏 & \ref{ch:ch03-fluent-udf}、\ref{ch:ch05-abaqus-usub} \\
迭代次数随线程数变化 & 并行归约顺序未定，残差末位差异改变判敛路径 & 用固定顺序归约或串行基准复跑一次做对照 & \ref{ch:ch06-openfoam}、\ref{ch:ch17-perf-toolchain} \\
同一算例在不同机器上残差曲线不同 & 快速数学选项、底层线代库的分派差异 & 关掉快速数学、固定线程数与库分派再跑 & \ref{ch:ch08-solver-engineering}、\ref{ch:ch17-perf-toolchain} \\
大规模结果取回太慢 & 全量场逐步输出，网络拷贝成为主项 & 先量写出耗时占比，再改切片抽取与设备侧压缩 & \ref{ch:ch16-dataexchange}、\ref{ch:ch02-ansys-layers} \\
引擎插件在真机崩溃而编辑器正常 & 编辑器是非优化或带符号的宿主构建，真机做了剥离与打包裁剪 & 用带调试信息的发布构建复现，查打包日志里该资源是否被裁 & \ref{ch:ch09-engine-layers}、\ref{ch:ch10-unreal} \\
打包后插件加载失败、找不到依赖库 & 加载路径解析差异，或依赖没进包内目录 & 核对包内目录布局；用系统工具列依赖再比对 & \ref{ch:ch10-unreal}、\ref{ch:ch17-perf-toolchain} \\
热重载后旧的反射数据仍在 & 注册表卸载路径没清理，或序列化数据带旧字段静默注入 & 打印注册条目数与结构体哈希，确认新旧不共存 & \ref{ch:ch12-godot-ext}、\ref{ch:ch09-engine-layers} \\
物理回放跨平台不一致 & 浮点环境、解析库版本、固定步长与插值口径不同 & 关无关分支，固定种子跑同一输入序列比轨迹 & \ref{ch:ch13-physics}、\ref{ch:ch15-hil} \\
平均帧率够但偶发卡顿 & 托管层回收突发、同步资源加载、流送线程抢占 & 帧时间轴上看尖刺落在哪个 zone 与哪次分配 & \ref{ch:ch11-unity}、\ref{ch:ch17-perf-toolchain} \\
CPU 帧时间正常而画面仍慢 & GPU 侧提交堆积、绘制调用数量或带宽不足 & 抓一帧逐个回放，区分提交量与占用率 & \ref{ch:ch10-unreal}、\ref{ch:ch17-perf-toolchain} \\
GPU 内核比预期慢很多 & 启动与搬运主导，访存不合并，块内分支发散 & 分开记搬运时间与内核时间，再看内核级指标 & \ref{ch:ch17-perf-toolchain} \\
剖析图上只剩地址或模板乱码 & 构建剥离符号、缺栈展开信息、未启用调试信息 & 换带调试信息的优化构建，保留帧指针重采 & \ref{ch:ch17-perf-toolchain} \\
脚本层循环把整个流程拖慢 & 每次调用跨进程往返，循环留在脚本层 & 统计调用次数；合并成批量调用或下沉原生层 & \ref{ch:ch01-boundary}、\ref{ch:ch02-ansys-layers} \\
CI 通过而本地不过（或相反） & 缓存命中旧制品；并行度、随机种子或容差不一致 & 清缓存重跑一次；比对制品哈希与构建元数据 & \ref{ch:ch16-dataexchange}、\ref{ch:ch17-perf-toolchain} \\
\end{longtable}
\endgroup

表里前五行集中在 \DomCAE{}，中间六行集中在 \DomEngine{}，最后四行两侧共有——这个分布本身就
是结论：\DomCAE{} 的故障多与``数值口径与工具链''有关，\DomEngine{} 的故障多与``构建与资源管线''
有关。跨行通用的起手式是二分：把输入规模、线程数、开关、提交历史各砍一半，缩小可疑集，再去
读对应章的机制解释。清单 \ref{lst:ch18-triage} 给两个最常用的排除脚本形状。

\begin{lstlisting}[language=ShellLD, style=shellstyle,
  caption={两个起手式：并行相关差异的复现，与加载期符号问题的核对},
  label={lst:ch18-triage}]
# A) 症状随并行配置变化时：同一输入跑三种配置，只比监测量文件
for T in 1 4 8; do
  ./solver --case <case> --threads $T --out run_$T.csv
done
python3 cmp.py run_1.csv run_4.csv run_8.csv --rel-tol 1e-10
# 差异随 T 出现又消失 -> 看归约顺序；恒定存在 -> 看工具链与库分派

# B) 插件或扩展加载失败：先看符号，再看路径
nm -D --defined-only build/libsimplugin.so > plugin.syms
grep -c sim_plugin_desc plugin.syms        # 入口符号在不在
ldd build/libsimplugin.so | grep "not found"   # 依赖解析齐不齐
# Windows 侧对应动作：用工具列导出表与依赖清单，比同名文件的位置
\end{lstlisting}

\section{一页式对照表与收束}
\label{sec:ch18-onepage}

表 \ref{tab:ch18-onepage} 把各章口径汇在一页：入口语言指``用户主要写什么''，注册方式指``实现
怎么被产品发现''，热重载边界指``改完是否免重启、免重编''，打包路径指``交付时它跟谁走''。
``部分''一栏的具体条件见对应章；版本与界面路径以官方文档为准。

\begingroup\footnotesize
\begin{longtable}{>{\raggedright\arraybackslash}p{0.14\textwidth}
                 >{\raggedright\arraybackslash}p{0.15\textwidth}
                 >{\raggedright\arraybackslash}p{0.19\textwidth}
                 >{\raggedright\arraybackslash}p{0.18\textwidth}
                 >{\raggedright\arraybackslash}p{0.14\textwidth}}
\caption{产品与引擎的一页式对照：入口、注册、热重载与打包}\label{tab:ch18-onepage}\\
\toprule
\textbf{产品或引擎} & \textbf{入口语言} & \textbf{扩展注册方式} & \textbf{热重载边界} &
\textbf{打包路径} \\
\midrule
\endfirsthead
\toprule
\textbf{产品或引擎} & \textbf{入口语言} & \textbf{扩展注册方式} & \textbf{热重载边界} &
\textbf{打包路径} \\
\midrule
\endhead
\midrule
\endfoot
\bottomrule
\endlastfoot
\Lib{Fluent} UDF & \CL{}，解释或编译加载 & \Tp{DEFINE\_*} 宏名即注册 & 部分：改源码需重载重编 & \Tp{libudf} 源树随算例归档（\ref{ch:ch03-fluent-udf}） \\
\Lib{Mechanical} 系 & 脚本宏加原生扩展 & 工程树内对象与加载项注册 & 部分：宏可重跑，原生要重载 & 工程文件与扩展随包（\ref{ch:ch04-mechanical-ext}） \\
\Lib{APDL} & 产品宏语言 & 命令流即接口 & 无：重跑整条流水 & 宏与参数文件随算例（\ref{ch:ch02-ansys-layers}） \\
\Lib{Abaqus} 子程序 & \CL、\CPP、\Lang{Fortran} & 约定名查找加补丁注入 & 无：重编重跑 & 输入、子程序源码、编译选项（\ref{ch:ch05-abaqus-usub}） \\
\Lib{OpenFOAM} 及同类 & \CPP{} 动态库 & 运行字典选类型加库名，链接即注册 & 好：重编库后重跑 & 源码或 \Tp{so} 加库路径（\ref{ch:ch06-openfoam}） \\
\Lib{Unreal} & 脚本层加 \CPP & 反射元数据生成注册 & 部分：脚本层可热改 & cook 清单决定内容进不进包（\ref{ch:ch10-unreal}） \\
\Lib{Unity} & 托管语言加原生插件 & 属性与生命周期回调 & 好：编辑器内重载，原生域重载 & 插件目录进构建目标（\ref{ch:ch11-unity}） \\
\Lib{Godot} & 场景脚本加原生模块 & 类注册表绑定场景节点 & 好：编辑器重启注册表 & 模块编进引擎或导出共享库（\ref{ch:ch12-godot-ext}） \\
\end{longtable}
\endgroup

最后把三层的判据合成一句可执行的顺序，这就是本章真正想留给读者的东西：任何一次接入，先在
最小闭环上把四步走通——构建产出、扩展注册生效、一次求解或一帧步进跑完、结果按
\S\ref{sec:ch18-exchange-routing} 的口径落到可比的位置。四步通了之后再谈规模与优化，否则
优化的对象是不可信的，而不可信的``快了 30\%''比慢更贵。清单 \ref{lst:ch18-minloop} 给这段
最小闭环的形状。

\begin{lstlisting}[language=ShellLD, style=shellstyle,
  caption={端到端最小闭环脚手架：构建、注册生效、跑通一步、与基线比},
  label={lst:ch18-minloop}]
# minloop.sh —— 四步都要有明确产物，缺一步就不要开始优化
set -euo pipefail
cmake --build <build_dir> -j <N>                  # 1) 构建产出
./solver --case <tiny_case> --check-registration  # 2) 扩展被找到并调用一次
./solver --case <smoke_case> --threads <T> --out run.csv   # 3) 跑通一个算例
python3 cmp.py --run run.csv --base <baseline.csv> --rel-tol 1e-8   # 4) 可比
# 引擎侧同形：构建目标 -> 插件加载成功日志 -> 固定种子步进 600 帧 -> 轨迹比对
./headless_replay --seed <S> --frames 600 --out trace.json
python3 cmp_trace.py --trace trace.json --base <baseline.json>
\end{lstlisting}

\begin{bestpractice}[title={先做端到端最小闭环，再谈优化}]
\begin{itemize}
\item 一次接入只做一条通路：一个最小扩展、一个最小算例或场景、一个基线文件。它们都通过之后
      才加规模，因为``规模问题''和``接口没接通''混在一起时无法归因。
\item 判据在写代码前定：数值用容差、实时用帧预算分位、交付用清单齐全度。事后定判据会被结果
      牵着走。
\item 两侧共同的第一步是``把阶段边界暴露出来''：\DomCAE{} 侧是逐阶段计时与可复现配置，
      \DomEngine{} 侧是帧边界与系统级 zone。有了边界，工具才有可对齐的单位。
\item 每引入一个``提速手段''（快速数学选项、解释改编译、CPU 改 GPU、脚本下沉原生），都要在
      最小闭环上重跑一次并留下差异记录；没有记录的提速不算完成。
\item 最后：选型不是选最强的组合，而是选能让最小闭环最快跑通、并且能在交付时把许可与元数据
      说清楚的组合。前三层的判据都是为了这个目的服务。
\end{itemize}
\end{bestpractice}

% ============================================================
% ---- frag_apA.tex ----
\appendix
\chapter{术语与缩写对照}
\label{ch:apA-glossary}

本附录把全书正文反复使用的术语收进一张跨页长表。它不是行业词典，而是\textbf{本书口径}的对照表：
同一个缩写在不同产品里常有不同的实现边界，正文写作时只能取其中一种说法，这里把那种说法固定下来，
免得读者拿着表里的定义去对官方手册时发现措辞不一致。

表按大致顺序排列：先\DomCAE{} 侧的产品与求解概念，再并行与数据布局，然后是本侧的引擎概念、
实时与网络、数据交换格式，最后是接口与发行层面的词。检索时可以只看词条列，那一列刻意压得很短。

\begin{note}[title={同一概念在两侧叫法不同}]
三组最容易混的词：求解器侧的``时间步''与引擎侧的``逻辑帧''是同一个循环变量的两种用法，
但前者可以几十秒一步、后者必须落在帧预算内；求解器侧的``检查点''与引擎侧的``存档''
在数据形状上接近，可检查点要求能改变续算的物理量都写进去，存档往往只留玩法相关的子集；
两侧的``热重载''都不是语言层的特性，而是宿主进程把新库换进旧调用表的能力。
\end{note}

\footnotesize
\begin{longtable}{>{\raggedright\arraybackslash}p{0.17\textwidth}>{\raggedright\arraybackslash}p{0.15\textwidth}>{\raggedright\arraybackslash}p{0.49\textwidth}}
\caption{\DomCAE{} 与 \DomEngine{} 术语与缩写对照（本书口径）}\label{tab:apA-glossary}\\
\toprule
词条 & 英文或缩写 & 一句话定义（本书口径） \\
\midrule
\endfirsthead
\toprule
词条 & 英文或缩写 & 一句话定义（见表~\ref{tab:apA-glossary}） \\
\midrule
\endhead
\endfoot
\bottomrule
\endlastfoot
计算机辅助工程 & CAE & 结构、流体、热、电磁等工程仿真产品的统称，本书特指商用求解器一侧。 \\
计算流体力学 & CFD & 以离散化的 Navier-Stokes 方程族为核心的求解类别，本书主要讨论有限体积法一支。 \\
有限元分析 & FEA & 用分片多项式单元离散偏微分方程的数值方法，也是结构类产品的主体。 \\
参数化设计语言 & APDL & 命令行式的批处理脚本语言，一行一条命令，负责建模、求解与后处理的串接。 \\
用户定义函数 & UDF & 流体侧用 \CL{} 写的扩展，编译成共享对象后由求解器在固定钩子处调用。 \\
用户材料子程序 & UMAT & 隐式动力学里描述应力更新的本构入口，由积分点循环反复调用。 \\
显式材料子程序 & VUMAT & 显式动力学对应的本构入口，一次处理一批积分点，对向量化更友好。 \\
用户子程序 & user routine & 求解器留给用户的 Fortran 挂钩族，名称与参数表由产品固定。 \\
后处理 & post-processing & 求解结束后把结果读出来做显示、报表与二次分析的一段流程。 \\
S-函数 & S-Function & 仿真建模环境里的模块级扩展，既是一个函数也是一个可连线的块。 \\
求解器 & solver & 把离散方程推进到下一步或下一个平衡态的那段核心代码。 \\
装配 & assembly & 把每个单元或控制体的局部贡献累加进全局矩阵与右端项的过程。 \\
稀疏矩阵 & sparse matrix & 只存非零元素及其位置的矩阵，工程仿真的规模全靠它撑着。 \\
区域分解 & domain decomp. & 把网格切成若干子域分给不同进程，子域之间靠边界数据互相耦合。 \\
晕单元 & halo cell & 为让子域能独立算梯度而额外复制过来的邻接层单元。 \\
鬼单元 & ghost cell & 与晕单元同类的叫法，强调这些单元只读、不参与本域装配。 \\
检查点 & checkpoint & 让算例能够中断后续算的中间状态快照，含场量、历史变量与时间信息。 \\
结构数组 & SoA & 同一字段的值连续存放的布局，SIMD 与缓存友好的默认选择。 \\
数组结构 & AoS & 一个对象的全部字段连续存放的布局，读写直观但易拖累带宽。 \\
单指令多数据 & SIMD & 一条指令处理一组数据的硬件能力，用户子程序性能常取决于它。 \\
消息传递接口 & MPI & 多进程之间交换数据的标准接口，也是并行求解器的默认通信层。 \\
热重载 & hot reload & 不重启宿主进程就把扩展库的新版本换进调用表的能力。 \\
反射 & reflection & 运行时按名字查询类型、属性与方法的能力，编辑器生态的地基。 \\
代码生成器 & codegen & 由声明或元数据产出样板源码的工具，用来省掉手写绑定。 \\
头文件工具 & UHT & 扫描反射标注并生成元数据与样板代码的前端步骤。 \\
对象基类 & UObject & 引擎对象体系的公共基类，带名字、outer 与序列化能力。 \\
垃圾回收 & GC & 由运行时判定对象不可达并统一回收的机制，不同于引用计数。 \\
域重载 & DRC & 编辑器在被重载类的实例上保存并重放状态的一套流程。 \\
转译后端 & IL2CPP & 把中间语言脚本转成 \CPP{} 再本地编译的脚本后端。 \\
向量化编译器 & Burst & 把受约束的作业代码编成向量化机器码的编译器。 \\
作业系统 & Job System & 把任务切片调度到多核并跟踪依赖关系的运行时。 \\
脚本扩展 & GDExtension & 用原生 \CPP{} 给编辑器注册新类型、由脚本按名使用的机制。 \\
类注册表 & ClassDB & 运行时按名字登记与查找类型的中心表。 \\
动态值 & Variant & 能携带任意受支持类型并可序列化的值容器。 \\
实体组件系统 & ECS & 以实体 id 加组件数据加系统遍历取代深继承树的数据组织法。 \\
帧预算 & frame budget & 一帧内允许各类工作占用的时间上限，通常以毫秒计。 \\
固定步长 & fixed step & 逻辑推进采用恒定时间增量、与渲染帧率解耦的做法。 \\
硬件在环 & HIL & 把真实控制器接进仿真回路、由模型充当被控对象的验证形态。 \\
确定性 & determinism & 同一输入与同一种子必然得到同一结果的性质。 \\
记录与回放 & record replay & 把输入序列存下来以便逐帧重现过程的取数方式。 \\
回滚式同步 & rollback & 先本地预测、状态分歧时退回到旧状态重放的同步策略。 \\
宽相检测 & broad phase & 用层次包围盒之类的粗筛快速排除不可能碰撞的对。 \\
窄相检测 & narrow phase & 对宽相留下的候选对精确求交点、法向与深度的阶段。 \\
传输格式 & glTF & 面向运行时的三维资产交换格式，强调可直接加载的布局。 \\
通用场景描述 & USD & 以分层与组合组织大规模场景与几何的描述体系。 \\
可视化工具包 & VTK & 科学可视化与几何处理的库族，常用于结果呈现。 \\
计算流网格标准 & CGNS & 流体网格与解字段的交换规范，配一套按名访问的库。 \\
二进制接口 & ABI & 编译产物之间在符号、调用约定与内存布局上的隐式契约。 \\
C 接口层 & C ABI & 只依赖 \CL{} 符号与结构布局的最小公共 ABI，跨语言边界的常客。 \\
封送 & marshalling & 在两套内存布局之间搬运并转换参数的过程。 \\
可位拷贝 & blittable & 托管与原生两侧内存表示一致、可直接搬过去的类型。 \\
许可传染 & copyleft 传播 & 上游许可要求衍生分发物继承同样条款的性质。 \\
运行时库路径 & RPATH & 写进可执行文件里、告诉动态链接器去哪里找库的路径字段。 \\
\end{longtable}
\normalsize

\section{几条边界约定}

\begin{itemize}
\item ``扩展''在本表里只指被宿主在运行时或链接时调用的用户代码，不包括改宿主源码的补丁。
\item ``插件''指有独立打包形态的扩展集合；一个插件里可以有多个扩展入口，也可以只是资源。
\item ``模块''指一个编译与链接单元，跟发行形态无关；同一模块可以只编进宿主，也可以发成插件。
\item 本表把``并行''与``分布式''分开：前者默认同一台机器上的多线程与 SIMD，后者默认跨进程的 MPI。
\item ``序列化''只表示把内存对象写成可持久或可传输的字节，不含跨语言表示转换，后者是封送。
\item ``元数据''指登记出来的数据，``反射''指运行时按名字查询这些数据的通道；两者经常同现但不同义。
\item 本表不为任何产品维护兼容矩阵；具体名称与参数一律回到该产品当时的官方文档确认。
\end{itemize}

\section{本书严格区分的词对}

这些词在行业里常被当同义词用，本书为了把代价说清楚，必须按下面的口径分开。
混用不会报错，只会让讨论对不上号，所以这一节的优先级高于表里的定义。

\begin{itemize}
\item 单元与控制体：有限元法的离散基元带节点与形函数，有限体积法的离散基元带通量面；
  两者的用户入口能读到的量不一样，本书从不互换这两个词。
\item 载荷步、时间步与帧：载荷步是求解流程里的一次完整求解设置，
  时间步是同一个载荷步内部的推进单位，帧是引擎的渲染与逻辑推进单位；三者时长互不相关。
\item 残差与容差：残差是方程不满足的程度，容差是判定``可以停''的门限；
  把``容差调小''说成``残差变小''是新手最常见的口误。
\item 热重载与域重载：前者换的是库，后者换的是类型定义并要在实例上重放状态；
  可重放的范围由宿主决定，不由用户决定。
\item 确定性与位级一致：确定性只要求同输入同结果，位级一致额外要求浮点运算顺序固定；
  跨机器、跨线程数的``可重现''通常只能做到前者。
\item 检查点与结果文件：检查点为续算服务，字段选择服从数值可恢复性；
  结果文件为呈现服务，字段选择服从可读与可裁剪。
\item 宽相与窄相：两相的成本结构与控制参数完全不同，混着谈``碰撞性能''没法得出结论。
\item 热路径与热重载：前者说的是被反复调用的代码，后者说的是不换进程替换库的动作；
  两个词只差一个字，指的却是完全不同的两件事，本书一律写全称。
\item 帧率与推进频率：帧率属于渲染，推进频率属于逻辑时间轴；
  两者数值相同只是巧合，把``掉帧''说成``算得慢''会掩盖真正的原因。
\end{itemize}

\begin{warning}[title={一个名字的两种所指}]
本书最容易出事的是几个跨域同名词。\Tp{field} 在求解器侧指``场量''，在引擎侧指类的成员变量；
\Tp{state} 在结构侧用户子程序里指历史变量，在引擎里指对象的可序列化状态；
\Tp{cell} 在流体侧指控制体，在实体组件布局里指组件槽位；
\Tp{assembly} 在求解器侧指组装全局矩阵，在构建语境里指编译产物单元。
读到这些词时先确认所在域，再确认它在附录~\ref{ch:apA-glossary} 里取的是哪个义项。
\end{warning}

\section{易混词速记}

\begin{itemize}
\item UDF 与用户子程序：前者是流体侧按钩子调用的 \CL{} 扩展，后者多指结构侧的 Fortran 挂钩，
  两者在编译时机、可用变量与错误处理上完全不同，不能指望同一套写法两边通吃。
\item 检查点与序列化结果文件：检查点服务于``继续算''，字段选择服从数值可恢复性；
  结果文件服务于``给人看''，字段选择服从可读与可裁剪。同一份数据很少两头都合适。
\item 反射与动态值：反射回答``这个类型有什么''，动态值回答``这个格子当前装的是什么''。
  引擎的脚本边界通常同时需要两者，缺一个就只剩硬编码。
\item 确定性与固定步长：固定步长只是确定性的必要条件之一，浮点归约顺序、
  容器遍历顺序与随机数流分配都能把它破坏掉。
\item 宽相与窄相：调物理性能时先确认瓶颈在哪一相，两相的可优化手段几乎不重叠。
\end{itemize}

\begin{guideline}[title={读这张表的正确方式}]
把它当成写作与检索的对照物，而不是事实来源。表里的一句话定义是为了让正文里某个词不再含糊，
一旦要落笔写具体入口的签名、参数个数或默认值，请以官方文档为准；
本书刻意不给小版本号，就是这个原因。
\end{guideline}

% ============================================================
% ---- frag_apB.tex ----
\chapter{扩展入口与钩子速查}
\label{ch:apB-cheatsheet}

这一张表是全书的``入口索引''：当你已经知道要改的是哪一侧、哪一个产品，只想快速确认
``这门扩展用什么语言写、什么时候被编出来、什么时候被调起来、能不能看到内部量、最后以什么形态发出去''，
就查这里。表的六列正好对应设计一个扩展前必须回答的六个问题，缺一个就没法估工。

用表时有三条限制要记住。第一，``可扩展语言''一列只写本书讨论过的入口，不是产品的全部扩展能力清单。
第二，``运行时机''决定了一次调用能摊到多少计算量，这一列比语言列更重要。
第三，``发行形态''决定你要不要碰编译链、要不要管 \Tp{RPATH}、要不要跟许可条款较劲。
名字、宏与参数一律回官方文档核对；本书刻意不给小版本号，所以表里也不会出现``某年版才支持''这类限定。

\begin{compare}[title={三族入口的差别其实只有一处}]
\begin{itemize}
\item 链接期挂钩（结构侧用户子程序）：登记动作发生在编译链接时，代价是每次都要重编求解器，
  好处是调用点极快、参数布局固定、可以向量化。
\item 装载期挂钩（流体侧 UDF、引擎侧扩展库）：宿主进程起来后把库换进调用表，
  代价是首调用有一次建表开销，好处是能不换宿主地迭代。
\item 运行期登记（脚本与绑定层）：把名字塞进可查找的表，随时可改，
  代价是每次跨边界的调用都要封送，性能取决于你一次搬多少个元素。
\end{itemize}
\end{compare}

\scriptsize
\begin{longtable}{>{\raggedright\arraybackslash}p{0.11\textwidth}>{\raggedright\arraybackslash}p{0.12\textwidth}>{\raggedright\arraybackslash}p{0.14\textwidth}>{\raggedright\arraybackslash}p{0.14\textwidth}>{\raggedright\arraybackslash}p{0.15\textwidth}>{\raggedright\arraybackslash}p{0.14\textwidth}}
\caption{跨产品扩展入口速查：语言、编译时机、运行时机、内部量可达性与发行形态}\label{tab:apB-hooks}\\
\toprule
产品或引擎 & 可扩展语言 & 编译或注册时机 & 运行时机 & 可达的内部量 & 发行形态 \\
\midrule
\endfirsthead
\toprule
产品或引擎 & 可扩展语言 & 编译或注册时机 & 运行时机 & 可达的内部量 & 发行形态 \\
\midrule
\endhead
\endfoot
\bottomrule
\endlastfoot
Ansys Mechanical & APDL、\CL & 命令流在执行时逐条解释；编译型扩展先编成库再挂 & 建模、求解前后与每个载荷步的边界 & 命令流层看得到全局模型；扩展层只见宿主暴露的量 & 命令流文件随算例走，扩展另发动态库 \\
Ansys Fluent（UDF） & \CL & 编译期编成库，或由求解器在载入时解释并建钩子表 & 每个迭代、每个时间步的钩子处 & 可读写线程上的单元、面与场量 & 求解器进程加载 \Tp{.so} 或 \Tp{.dll} \\
Abaqus 用户子程序 & Fortran、\CPP & 链接期把子程序编进求解器可执行 & 每个增量步、每个积分点 & 可读写应力、状态变量与材料雅可比 & 静态或动态链接进求解器 \\
LS-DYNA & Fortran、\CPP & 与求解器一起重新编译可执行 & 每个时间步由关键字指定的位置 & 只见关键字与约定参数暴露的量 & 重新编译的求解器可执行 \\
COMSOL & \CPP、内置函数 & 建模阶段登记；编译型扩展需配套开发包 & 算子求值与后处理表达式处 & 有限，走公开 API 与变量名 & 模型文件加动态库 \\
Simulink（S-Function） & \CL、\CPP & 建模期登记块，编译成块库或即时编译产物 & 每个采样周期，跟随所在块的速率 & 只有端口、参数与自身的工作向量 & 块库或即时编译产物 \\
OpenFOAM & \CPP & 独立编译成动态库，运行期按字典里的类型名加载 & 由类型名与构造时机决定调用点 & 可直接访问场、网格与求解器对象 & 动态库加算例目录 \\
Unreal Engine & \CPP、蓝图 & 构建期由元数据工具生成登记码再参与编译 & 对象生命周期、每帧回调、回收与序列化点 & 全量，经反射与对象体系访问 & 模块或插件包 \\
Unity & 托管脚本、转译后端 & 编译期产出程序集，发布时可转成本地代码 & 组件生命周期与固定步长的物理循环 & 托管层可见量受运行时边界限制 & 程序集加包 \\
Godot & \CPP（GDExtension） & 装载时由扩展描述文件登记入口符号 & 每帧处理与固定步长处理两个回调 & 经类注册表与动态值边界访问 & 动态库加脚本资源 \\
物理库（PhysX、Bullet、Jolt） & \CPP & 与调用方一起编译，运行期没有注册步骤 & 物理步的宽相与窄相两个阶段 & 完全可达，但生命周期归调用方管 & 静态库或动态库 \\
绑定层（Lua、sol2） & \CPP 加脚本 & 启动时把函数与类型按名字注册进表 & 脚本被调用的任意时刻 & 只暴露注册过的那些量 & 解释器加绑定库 \\
\end{longtable}
\normalsize

\section{流体侧 \CL{} 钩子族的最小骨架}

Fluent 的 UDF，以及各家流体产品里``用 \CL{} 写一个函数然后由求解器调''的入口，形状高度一致：
入口是普通 \CL{} 函数，第一个实参常常是代表``当前区域或当前线程''的不透明指针，
其余是索引与进出变量；注册靠一个宏把函数名与钩子类型登记进求解器的表；
一个钩子会被调用很多次，所以里面不该有分配、不该有阻塞调用、不该有磁盘写。
下面只给调用形状，具体名与参数以官方文档为准。

真正要在这一族里做对的事情只有三件：把取值与写值都走宿主给的接口而不是自己算地址；
把每步都要用的常数放在只跑一次的入口里；把索引合法性检查放在最前面，
因为一个越界访问会把整个并行作业的全部进程一起带走。

\begin{lstlisting}[language=C, style=cstyle, caption={\CL{} 钩子族的最小形状：入口函数、注册宏与一次遍历}, label={lst:apB-udf-shape}]
/* C 钩子的形状。名字、宏与参数个数以官方文档为准。 */
typedef struct RegionOpaque RegionData;   /* 不透明句柄 */

#define HOOK_PROFILE(name) \
    double name(RegionData *rg, int index)

/* 每个面上的索引都会调一次：只算，不分配，不打日志 */
HOOK_PROFILE(inlet_speed)
{
    double area = 0.0;      /* 面积由宿主接口取，此处占位 */
    double flux = 0.0;      /* 写回宿主的出口变量 */
    if (rg == 0 || index < 0) {
        return 0.0;         /* 非法入参先返回，别越界 */
    }
    area = 1.0;             /* 换成官方的取值调用 */
    flux = 2.0 * area;      /* 换成官方的写值调用 */
    return flux;
}
\end{lstlisting}

\section{结构侧 Fortran 子程序族的最小骨架}

Abaqus、LS-DYNA 这一类的用户子程序是另一路传统：入口名和参数表由产品定死，
用户不注册、不改名，靠链接期把自己塞进求解器；调用发生在每个增量步的每个积分点或每个接触对，
所以单次成本极小、次数极大。这一族里决定性能的是能不能让编译器向量化，
而向量化又回到数组按 AoS 还是 SoA 摆放、循环边界是否可知。

还有一点常被忽略：这类入口里不能改全局状态以外的东西，
显式与隐式两条路径给的参数含义不同，同一份本构代码换入口就要重查参数含义。
下面是形状示意，不含真实参数表，具体名与参数以官方文档为准。

\begin{lstlisting}[style=fortranstyle, caption={Fortran 用户子程序族的最小形状：固定名字、批量积分点、纯计算}, label={lst:apB-fortran-shape}]
! 形状示意：子程序名与参数表由产品规定，以官方文档为准。
subroutine constitutive_shape(npt, dtime, stress, statev)
  implicit none
  integer, parameter :: ndim = 6
  integer, intent(in) :: npt
  double precision, intent(in) :: dtime
  double precision, intent(inout) :: stress(ndim, npt)
  double precision, intent(inout) :: statev(npt)
  integer :: i
  ! 一次处理一批积分点，便于编译器把内层循环向量化
  do i = 1, npt
    statev(i) = statev(i) + dtime
    stress(1, i) = stress(1, i) * 0.9d0
  end do
  return
end subroutine constitutive_shape
\end{lstlisting}

\section{引擎侧 \CPP{} 类型注册族的最小骨架}

Unreal 的模块与反射、Unity 的组件、Godot 的 GDExtension、OpenFOAM 的运行时选择表，
共同点是先``把类型按名字交给宿主''，再由宿主在固定时刻调你的构造函数或虚函数。
差别只在登记发生的位置：构建期生成登记码、装载期读元数据、还是运行期查一张表。
骨架只体现这三步——声明类型、登记、实现被回调的时机。名字与宏以官方文档为准。

这一族的坑在于回调时机而不是回调内容：每帧回调与固定步长回调看起来都叫 update，
但前者次数随帧率浮动，后者次数随时间轴锁定，把带物理意义的推进写进前者，
就会得到``帧率高时跑得慢''的怪结果。OpenFOAM 一侧对应的差别是构造时机：
对象由字典决定类型名，读表与分配都发生在构造里。

\begin{lstlisting}[style=cppstyle, caption={引擎侧类型注册的三步形状：声明、登记、被回调}, label={lst:apB-engine-shape}]
// 三步形状：继承宿主基类、把类型交出去、实现宿主调的回调。
// 真实基类名、宏名与回调签名以对应引擎的官方文档为准。
class DemoMover
{
public:
    // 1) 登记：宿主装载时调用，把类型名交给类注册表
    static void RegisterType(class Registry &reg);
    // 2) 每帧回调：delta 随帧率浮动，别放推进逻辑
    void OnFrame(double delta);
    // 3) 固定步长回调：步长恒定，物理与推进放这里
    void OnFixedStep(double fixedDelta);
};

void DemoMover::RegisterType(Registry &reg)
{
    reg.Declare("DemoMover");   // 只是形状，不是真实 API
}
\end{lstlisting}

\section{脚本与绑定层的最小骨架}

最后一族是``脚本层调用原生实现''：Lua 加 sol2、Godot 的脚本、以及产品自带的 Python 客户端。
共同形状是把原生函数按名字塞进一个可查找的表，脚本侧按名字取。
跨边界的数据必须是 blittable 的，否则每次调用都要封送；这一族的成本在边界上，
不在于被调函数本身，所以批量传数据永远优于逐元素传。

脚本层还有一处独有能力：它是唯一常能改完就生效的入口，
因此适合放在开发循环的最外层当调度器，把稳定的数值内核留在编译型入口里。
名字、装饰器与注册接口以官方文档为准。

\begin{lstlisting}[language=Python, style=pythonstyle, caption={脚本绑定层的最小形状：注册、按名调用、批量搬运}, label={lst:apB-binding-shape}]
# 绑定层形状：原生侧登记名字，脚本侧按名字调用。
# 注册接口与装饰器的真实拼写以对应产品的官方文档为准。


def register_binding(registry):
    registry.define("run_step")        # 名字对上即可
    registry.attach("run_step", step)


def advance_loop(nsteps, dt_ms):
    for i in range(nsteps):
        call("run_step", dt_ms)        # 逐步跨边界，代价在封送
    return nsteps * dt_ms


def advance_batch(steps, dt_ms):
    call("run_step_batch", steps, dt_ms)
    return steps * dt_ms               # 一次搬运，永远更好
\end{lstlisting}

\begin{bestpractice}[title={选入口之前先量一次调用频率}]
把速查表的``运行时机''列乘上``一次调用允许多少时间''，再决定用哪一族入口。
每个积分点调一次的入口，写多少优化都不如少分配一次；
每帧调一次的入口，真正的敌人往往是跨语言边界的封送；
建模阶段只跑一次的入口，可以放心用解释型脚本换开发速度。
反过来也一样：先确认调用频率，再决定要不要为了``能热重载''而接受一层动态库的封送成本。
\end{bestpractice}

% ============================================================
% ---- frag_apC.tex ----
\chapter{版本口径、许可与参考路线}
\label{ch:apC-versions}

前两章附录解决``这个词是什么意思''和``这个入口长什么样''，这一章解决三件落地的事：
本书为什么不锁小版本号、把扩展随产品发出去之前要过哪三道许可判断、
以及把求解器侧的库和引擎侧的插件放进同一个构建目录时该怎么组织目标。
最后给三条阅读路线，按你原本的身份选一条走。

\section{本书的版本口径}

全书统一采用``系列加年代''的版本口径：写产品时说到系列名、接口形态与大致年代，
不写死小版本号，也不写``某年版才支持某特性''这类断言。这么定有三个理由。

第一，这两侧产品的接口演进节奏差别极大。求解器侧的用户子程序与钩子参数表十年不动的情况很常见，
而引擎侧的反射、打包与脚本后端几乎年年改名；把两者写成同一种``版本号断言''，
一半会立刻过期。第二，本书的重点是形状与代价，形状比版本号稳定得多，
而具体可用性只有你手上那份安装才说了算。第三，一旦写了小版本号，读者会把它当成兼容承诺，
而本书没有任何测试覆盖，给不出这种承诺。

所以约定是：凡涉及``能不能用、叫什么名、参数几个''，回到下表对应的官方文档类别现场确认；
凡涉及``这一层大概怎么工作、代价在哪''，以本书为准。

\scriptsize
\begin{longtable}{>{\raggedright\arraybackslash}p{0.13\textwidth}>{\raggedright\arraybackslash}p{0.24\textwidth}>{\raggedright\arraybackslash}p{0.25\textwidth}>{\raggedright\arraybackslash}p{0.19\textwidth}}
\caption{域与官方文档类别的对应：该去哪一份文档确认哪一类事实}\label{tab:apC-docmap}\\
\toprule
域 & 需要现场确认的事 & 该看的文档类别 & 典型触发时机 \\
\midrule
\endfirsthead
\toprule
域 & 需要现场确认的事 & 该看的文档类别 & 典型触发时机 \\
\midrule
\endhead
\endfoot
\bottomrule
\endlastfoot
数值内核 & 可用的单元类型、湍流模型、并行规模上限 & 求解器手册的数值方法与限制章节 & 算例可行性评估阶段 \\
用户子程序与钩子 & 入口名、参数表与允许读写的量 & 参考手册里的用户子程序或 UDF 章 & 动手写扩展之前 \\
脚本与命令流层 & 命令名、属性拼写与返回值形态 & 脚本语言参考加版本发布说明 & 自动化脚本迁移时 \\
引擎对象与反射 & 注册宏、生命周期回调与序列化边界 & 引擎的本地语言程序手册 & 新模块搭骨架时 \\
引擎资产与打包 & 序列化字段、插件描述与依赖声明 & 资产管理与打包文档 & 出第一个包之前 \\
物理与仿真库 & 步进接口、约束求解器选项 & 库自带的 API 文档与变更日志 & 换库或换大版本时 \\
数据交换格式 & 支持的字段、扩展与版本协商 & 格式的规范文本与 schema & 实现读写器时 \\
并行与运行时 & 进程拓扑、启动器与环境变量 & 安装与高性能计算文档 & 首次上集群时 \\
许可与发行 & 能否随产品分发、需要附哪些声明 & 上游许可文件与常见问题 & 打包发布之前 \\
构建工具链 & 生成器、语言标准与开发包版本 & 构建工具自己的官方文档 & 写连续集成脚本时 \\
\end{longtable}
\normalsize

\begin{note}[title={把版本口径当成写作纪律而不是免责声明}]
读者拿本书去对官方文档时，希望看到的不是``本书可能过时''，
而是明确知道哪一类事实归本书讲、哪一类事实必须去现场确认。
上表就是这条边界的可执行版本：左边两列是本书敢讲的，右边两列是本书不敢讲的。
\end{note}

\section{许可与发行的三步判断}

把扩展随产品分发时，许可问题的复杂度全部来自一件事：你的分发物里包含了别人的东西。
宽松许可（BSD、MIT、Apache-2.0 这一族）通常只要求你保留版权声明与许可文本，
允许把成果闭源发出去；著佐权（copyleft，GPL、LGPL、AGPL 这一族）则要求衍生分发物继承条款，
差别就在于``你还能不能不公开源码''。三者内部差异不小，
例如 LGPL 允许动态链接使用而不要求开源宿主，而 AGPL 把网络提供服务也算成分发。

判断只做三步，顺序不能颠倒。

\begin{enumerate}
\item 读上游许可原文。只看三件事：是否要求源码开放、是否要求署名与许可文本随附、
  有没有把``修改''与``链接''分开规定的条款。同时留意例外条款、专利条款与双许可选项。
\item 看链接方式。静态链接、动态链接、把源码抄进自己的树、只按名字调用另一个进程，
  这四种形态在很多许可文本里的待遇完全不同。别凭工程直觉，按许可自己的用词去对。
\item 看分发物。你在客户机器上留下哪些东西：可执行文件、动态库、头文件、脚本、数据表，
  还是只有一个安装包。还要看谁分发给谁、是否连带产品一起发、是否需要用户提供源码获取途径。
\end{enumerate}

\begin{warning}[title={把 copyleft 依赖静态链进闭源插件}]
最常见的一类风险形态是：扩展里静态链接了一个 GPL 类依赖，
产物只有一个插件动态库，看起来``没有分发别人的源码''。
静态链接通常会把上游条款带到整个产物上，而插件形态在不少许可文本里还会被单独讨论。
本书不下法律结论，只给出判断方法：先确认上游许可对``链接''与``衍生作品''的原话怎么说，
再确认你的发布渠道到底交付了哪些文件，最后确认有没有更宽松的替代上游或者双许可路径。
三步里任何一步没做过，就不要把结论当成可行。拿不准时找法务或上游确认，
比自己按工程直觉推断便宜得多。
\end{warning}

\begin{compare}[title={两种常见规避形态的代价差别}]
把 copyleft 依赖换成实现同一功能的宽松许可库，代价是一次功能对齐与精度验证，
好处是发布渠道不再依赖条款解释。把依赖做成独立进程、只通过文件套接口或本地套接口通信，
代价是数据搬运与封送的开销，好处是分发物边界清晰。
两者的取舍取决于你要保护的是发行进度还是运行性能，不要指望同时得到。
\end{compare}

\section{把两侧制品放进同一个构建}

跨产品搬数据的项目里，最省心的仓库形态是：一套 CMake 目录同时产出求解器侧的库与引擎侧的插件，
两侧共用的数据结构单独成一个目标，谁都不许反过来依赖对方。
下面这份清单给出目标形状，选项名可按项目调整，不含任何厂商私有工具文件。

\begin{lstlisting}[style=cmakestyle, caption={两侧制品共存的最小目标结构：求解器侧库、引擎侧插件与共用类型}, label={lst:apC-build}]
# 一个目录同时产出求解器侧的钩子库与引擎侧的插件。
# 目标是形状清晰，选项可按项目调整。
cmake_minimum_required(VERSION 3.20)
project(bridge C CXX Fortran)

# 共用的内存布局：必须两边逐字节一致
add_library(bridge_types STATIC types.cpp)
target_include_directories(bridge_types
  PUBLIC include)

# 求解器侧：用户子程序与钩子，编成可被宿主链接的库
add_library(cae_hooks SHARED udf_shape.c constitutive.f90)
target_link_libraries(cae_hooks PRIVATE bridge_types)
target_compile_definitions(cae_hooks
  PRIVATE HOOK_ABI_V1)

# 引擎侧：类型注册与回调入口，编成模块
add_library(engine_plugin MODULE engine_shape.cpp)
set_target_properties(engine_plugin PROPERTIES
  PREFIX "" OUTPUT_NAME demo_plugin)
target_link_libraries(engine_plugin
  PRIVATE bridge_types)

# 发行时把随附动态库的位置写进运行时搜索路径
set_target_properties(engine_plugin PROPERTIES
  INSTALL_RPATH "$ORIGIN/../lib")

install(TARGETS cae_hooks engine_plugin
  LIBRARY DESTINATION lib
  ARCHIVE DESTINATION lib)
\end{lstlisting}

\begin{bestpractice}[title={共用类型是唯一的双向耦合点}]
求解器侧与引擎侧不要互相链接，只共用 \Tp{bridge\_types} 这一个目标：
它只包含数据布局与常量，不含任何一侧的运行时调用。
这样任何一侧换版本，另一侧只需要重新核对偏移与对齐，
不会出现``改了引擎的一个函数名导致求解器链接失败''这种跨域事故。
\end{bestpractice}

\section{三类读者的阅读路线}

三条路线都假设你已经能读懂 \CL{} 与 \CPP{} 的基本语法。
方括号里的标签是本书对应的章节或附录，按需跳读即可，不必按顺序通读。

\textbf{路线一：做仿真的读者}，目标是自己把本构、边界或后处理挂进求解器。

\begin{enumerate}
\item 先读术语与入口两本附录，把``子程序、钩子、脚本层''三个词的范围钉住 [附录 A、B]。
\item 读 \DomCAE{} 侧的分册章节：离散化、装配与求解器推进的形状 [正文前半]。
\item 挑一个入口族，照最小骨架写出空实现，只打印调用次数，确认调用时机 [附录 B]。
\item 补并行与数据布局：区域分解、晕单元与 SoA 对性能的影响 [正文中段]。
\item 最后读检查点、结果格式与数据交换的章节，把算例产出接给下游 [正文后半]。
\end{enumerate}

\textbf{路线二：做引擎的读者}，目标是把外部物理或数值能力接进自己的循环。

\begin{enumerate}
\item 读术语附录里引擎那一组：反射、代码生成、序列化与回收 [附录 A]。
\item 读扩展入口速查表，选定你要用的那一族入口 [附录 B]。
\item 按帧预算与固定步长两节确定``推进逻辑该挂在哪个回调'' [正文引擎部分]。
\item 读插件与模块的打包形态、依赖声明与热重载边界 [正文引擎部分]。
\item 最后读许可与构建两节，确认你的包能不能带着第三方库发出去 [本章]。
\end{enumerate}

\textbf{路线三：要在两侧之间搬数据的读者}，目标是一条稳定的导入导出管线。

\begin{enumerate}
\item 先读交换格式一节：网格、场量与元数据各有什么规范 [正文数据部分]。
\item 读术语里的 ABI、C 接口层、封送与 blittable 四条 [附录 A]。
\item 用速查表确认两侧各自的入口能不能同时被你的管线驱动 [附录 B]。
\item 抄本章的构建目标结构，把两侧制品装进同一次构建 [本章]。
\item 读确定性、记录与回放相关章节，决定管线要不要保证可重现 [正文数据部分]。
\item 收尾读许可三步法，把第三方格式库的声明文件放进安装包 [本章]。
\end{enumerate}

\begin{guideline}[title={三条路线共用的落地顺序}]
不管从哪条路线进来，落地顺序都是同一套：先用空实现确认调用时机，
再用最小数据确认边界形状，然后才谈性能与并行，最后才谈发行与许可。
把顺序倒过来做的项目，通常会把大量时间花在``接口本来就不允许''的事情上。
\end{guideline}

\end{document}
